\section{Metatheory}
\label{sec:metatheory-full}

\subsection{Standard Lemmas}
\label{sec:lemmas}

\begin{oxlemma}{Canonical Forms}{lemma:canonical-forms}
  If
  \oxtypjudge{\oxglobalctx}{\oxtvarctx}{\oxkontctx}{\oxvarctx}{\oxvalue}{\oxtype}{\oxvarctx}
  then

  \begin{enumerate}
  \item if $\oxtype = \oxtbool$, then $\oxvalue = \oxtrue$ or $\oxvalue =
    \oxfalse$.
  \item if $\oxtype = \oxtnum$, then $\oxvalue = \oxnum$.
  \item if $\oxtype = \oxtunit$, then $\oxvalue = \oxunit$.
  \item if $\oxtype = \oxtref{\oxprov}{\oxmuta}{\oxsitype}$, then $
    \oxvalue$ is of the form $\oxptr{\oxreferent}$.
  \item if $\oxtype = \oxtref{\oxprov}{\oxmuta}{\oxtslice{\oxsitype}}$, then $
    \oxvalue$ is of the form $\oxsliceptr{\oxreferent}{i}{j}$.
  \item if $\oxtype = \oxtfunext{
      \overline{\oxenvvar} \oxcomma \overline{\oxabsprov} \oxcomma
      \overline{\oxtvar}
    }{\oxsitype_1 \oxdotsc \oxsitype_n}{\oxsitype_r}{}{
      \overline{\oxabsprov_1 : \oxabsprov_2}
    }$, then $\oxvalue$ is of the form $\oxfnname$.
  \item if $\oxtype = \oxtclosure{
        \oxsitype_1 \oxdotsc \oxsitype_n
      }{\oxsitype_r}{\oxframe}$, then $\oxvalue$ is of the form
      $\oxclosureval{\oxstore}{
        \oxascribe{\oxid_1}{\oxsitype_1} \oxdotsc
        \oxascribe{\oxid_n}{\oxsitype_n}
      }{\oxsitype_r}{\oxexpr}$.
  \item if $\oxtype = \oxtarr{\oxtype^\prime}{\oxnum}$, then $\oxvalue$ is of
    the form $\oxarr{\oxvalue_1 \oxdotsc \oxvalue_n}$.
  \item if $\oxtype = \oxtslice{\oxtype^\prime}$, then $\oxvalue$ is of
    the form $\oxsliceval{\oxvalue_1 \oxdotsc \oxvalue_n}$.
  \item if $\oxtype = \oxtprod{\oxtype_1 \oxdotsc \oxtype_n}$, then $\oxvalue$
    is of the form $\oxprod{\oxvalue_1 \oxdotsc \oxvalue_n}$.
  \item if $\oxtype = \oxtsum{\oxtype_1}{\oxtype_2}$, then $\oxvalue$ is of
    either the form $\oxinl{\oxtype_1}{\oxtype_2}{\oxvalue^\prime}$ or
    $\oxinr{\oxtype_1}{\oxtype_2}{\oxvalue^\prime}$.
  \end{enumerate}
\end{oxlemma}

\begin{proof}
  By inspection of the grammar of values and typing rules.
\end{proof}

\begin{oxlemma}{Preservation of Types under Substitution}{lemma:substitution}
  \hfill
  \begin{enumerate}
  \item If $\oxtypjudge{\oxglobalctx}{ \oxextendctx{\oxtvarctx}{
        \oxtvarctxentry{\oxtvar}{\oxktype} }
    }{\oxkontctx}{\oxvarctx}{\oxexpr}{\oxtype}{\oxvarctx^\prime}$ and
    $\oxtypevalidity{\oxglobalctx}{\oxtvarctx}{\oxvarctx}{\oxtype^\prime}$, then
    $\oxtypjudge{\oxglobalctx}{\oxtvarctx}{\oxkontctx}{\oxvarctx}{
      \oxsubst{\oxexpr}{\oxtype^\prime}{\oxtvar}
    }{\oxsubst{\oxtype}{\oxtype^\prime}{\oxtvar}}
    {\oxsubst{\oxvarctx^\prime}{\oxtype^\prime}{\oxtvar}}$
  \item If $\oxtypjudge{\oxglobalctx}{ \oxextendctx{\oxtvarctx}{
        \oxtvarctxentry{\oxabsprov}{\oxkprov} }
    }{\oxkontctx}{\oxvarctx}{\oxexpr}{\oxtype}{\oxvarctx^\prime}$ and
    $\oxrgnvalidity{\oxtvarctx}{\oxvarctx}{\oxprov}$, then
    $\oxtypjudge{\oxglobalctx}{\oxtvarctx}{\oxkontctx} {\oxvarctx}{
      \oxsubst{\oxexpr}{\oxprov}{\oxabsprov}
    }{\oxsubst{\oxtype}{\oxprov}{\oxabsprov}}
    {\oxsubst{\oxvarctx^\prime}{\oxprov}{\oxabsprov}}$
  \item If $\oxtypjudge{\oxglobalctx}{ \oxextendctx{\oxtvarctx}{
        \oxtvarctxentry{\oxenvvar}{\oxkenv} }
    }{\oxkontctx}{\oxvarctx}{\oxexpr}{\oxtype}{\oxvarctx^\prime}$ and
    $\oxenvwellformed{\oxglobalctx}{\oxtvarctx}{\oxvarctx}{\oxenv}$, then
    $\oxtypjudge{\oxglobalctx}{\oxtvarctx}{\oxkontctx} {\oxvarctx}{
      \oxsubst{\oxexpr}{\oxenv}{\oxenvvar}
    }{\oxsubst{\oxtype}{\oxenv}{\oxenvvar}}
    {\oxsubst{\oxvarctx^\prime}{\oxenv}{\oxenvvar}}$
  \end{enumerate}
\end{oxlemma}

\begin{proof}
  By induction on the typing derivation.
\end{proof}

\subsection{Referent Lemmas}
\label{sec:ref-lemmas}

\begin{oxlemma}{Well-Formed References Evaluate to Well-Typed Values}
  {lemma:wf-refs-evalref}
  If $\oxreferentvalidity{\oxglobalctx}{\oxvarctx}{\oxreferent}{\oxxitype}$ and
  $\oxstorevalidity{\oxglobalctx}{\oxvarctx}{\oxstore}$, then
  $\oxevalreferent{\oxstore}{\oxreferent}{\oxvaluectx}{\oxvalue}$.
\end{oxlemma}

\begin{proof}
  We proceed by induction on $\oxreferentvalidity{\oxglobalctx}{\oxvarctx}
  {\oxreferent}{\oxxitype}$. There are six cases: \oxname{WF-RefId},
  \oxname{WF-RefProj}, \oxname{WF-RefIndexArray}, \oxname{WF-RefIndexSlice},
  \oxname{WF-RefSliceArray}, and \oxname{WF-RefSliceSlice}. Each of these cases
  has a corresponding evaluation rule:

  \vspace{1em}
  \noindent \framebox[\textwidth]{
    \figuresize
    \begin{mathpar}
      \WFRefId \and \ERId
    \end{mathpar}
  }
  \vspace{1em}

  For the base case, we consider the frame of $\oxvarctx$ which contains
  $\oxid$. By inversion of \oxname{WF-StackFrame} for the portion of the
  derivation $\oxstorevalidity{\oxglobalctx}{\oxvarctx}{\oxstore}$ pertaining to
  that frame, we have $\forall \oxid \in \oxdomain{\oxstackframe}. \;
  \oxtypjudge{\oxglobalctx}{\oxemptyctx}{\oxemptyctx}{\oxnewframe{\oxvarctx}{\oxframe}}
  {(\oxnewframe{\oxstore}{\oxstackframe})(\oxid)}{
    (\oxnewframe{\oxvarctx}{\oxframe})(\oxid)
  }{\oxnewframe{\oxvarctx}{\oxframe}}$. Focusing on our particular $\oxid$, we
  have both that $\oxlookup{\oxstore}{\oxid}{\oxvalue}$ and that
  $\oxtypjudge{\oxglobalctx}{\oxemptyctx}{\oxemptyctx}{\oxnewframe{\oxvarctx}{\oxframe}}
  {\oxvalue}{(\oxnewframe{\oxvarctx}{\oxframe})(\oxid)}{
    \oxnewframe{\oxvarctx}{\oxframe}
  }$, finishing the case. The remaining cases follow:

  \vspace{1em}

  \noindent \framebox[\textwidth]{
    \figuresize
    \begin{mathpar}
      \WFRefProj \and \ERProj
    \end{mathpar}
  }

  \vspace{1em}

  \noindent \framebox[\textwidth]{
    \figuresize
    \begin{mathpar}
      \WFRefIndexArray \and \ERIndexArray
    \end{mathpar}
  }

  \vspace{1em}

  \noindent \framebox[\textwidth]{
    \figuresize
    \begin{mathpar}
      \WFRefIndexSlice \and \ERIndexSlice
    \end{mathpar}
  }

  \vspace{1em}

  \noindent \framebox[\textwidth]{
    \figuresize
    \begin{mathpar}
      \WFRefSliceArray \and \ERSliceArray
    \end{mathpar}
  }

  \vspace{1em}

  \noindent \framebox[\textwidth]{
    \figuresize
    \begin{mathpar}
      \WFRefSliceSlice \and \ERSliceSlice
    \end{mathpar}
  }

  The proof for each case is identical: apply the induction hypothesis and then
  \Lemma{lemma:canonical-forms} and then the evaluation rule on the right. For
  the well-typed portion, apply inversion on the typing rule for the appropriate
  value.
\end{proof}

\begin{oxlemma}{Place Expressions Reduce}{lemma:place-exprs-reduce}
  If $\oxcomputetynoprov{\oxtvarctx}{\oxvarctx}{\oxmuta}{\oxplaceexpr}{\oxxitype}$
  and $\oxstorevalidity{\oxglobalctx}{\oxvarctx}{\oxstore}$,
  then $\oxnorm{\oxstore}{\oxplaceexpr}{\oxreferent}{\oxvaluectx}{\oxvalue}$ and
  $\oxtypjudge{\oxglobalctx}{\oxtvarctx}{\oxkontctx}{\oxvarctx}{\oxvalue}{\oxxitype}{\oxvarctx}$.
\end{oxlemma}

\begin{proof} 
  We proceed by induction on
  $\oxcomputetynoprov{\oxtvarctx}{\oxvarctx}{\oxmuta}{\oxplaceexpr}{\oxxitype}$.
  There are three cases: \oxname{TC-Var}, \oxname{TC-Proj}, and \oxname{TC-Deref}.

  \vspace{1em}
  \noindent \framebox[\textwidth]{
    \figuresize
    \begin{mathpar}
      \TCVar \and \PId
    \end{mathpar}
  }
  \vspace{1em}
  
  For \oxname{TC-Var}, we consider the piece of the derivation for
  $\oxstorevalidity{\oxglobalctx}{\oxvarctx}{\oxstore}$ (from our premise)
  for the frame containing $\oxid$. By inversion on \oxname{WF-StackFrame},
  we have $\forall \oxid \in \oxdomain{\oxstackframe}. \;
  \oxtypjudge{\oxglobalctx}{\oxemptyctx}{\oxemptyctx}{\oxnewframe{\oxvarctx}{\oxframe}}
  {(\oxnewframe{\oxstore}{\oxstackframe})(\oxid)}{
    (\oxnewframe{\oxvarctx}{\oxframe})(\oxid)
  }{\oxnewframe{\oxvarctx}{\oxframe}}$. This immediately gives us that
  $\oxlookup{\oxstore}{\oxid}{\oxvalue}$ and that
  $\oxtypjudge{\oxglobalctx}{\oxtvarctx}{\oxkontctx}{\oxvarctx}{\oxvalue}
  {\oxxitype}{\oxvarctx}$. To construct our premise for \oxname{P-Referent},
  we apply \oxname{ER-Id} to $\oxlookup{\oxstore}{\oxid}{\oxvalue}$.

  \vspace{1em}
  \noindent \framebox[\textwidth]{
    \figuresize
    \begin{mathpar}
      \TCProj \and \PProj
    \end{mathpar}
  }
  \vspace{1em}
  
  For \oxname{TC-Proj}, we apply our induction hypothesis to
  $\oxcomputety{\oxtvarctx}{\oxvarctx}{\oxmuta}{
    \oxplaceexpr
  }{
    \oxtprod{\oxsitype_1 \oxdotsc \oxsitype_i \oxdotsc \oxsitype_n}
  }{\oxset{\overline{\oxprov_p}}}$ from the premise of \oxname{TC-Proj} and get
  $\oxnorm{\oxstore}{\oxplaceexpr}{\oxreferent}{\oxvaluectx}{\oxvalue}$ and
  $\oxtypjudge{\oxglobalctx}{\oxtvarctx}{\oxkontctx}{\oxvarctx}{\oxvalue}
  {\oxtprod{\oxsitype_1 \oxdotsc \oxsitype_i \oxdotsc \oxsitype_n}}{\oxvarctx}$.
  Then, by \Lemma{lemma:canonical-forms}, we know that $\oxvalue$ must be of the
  form $\oxprod{\oxvalue_1 \oxdotsc \oxvalue_i \oxdotsc \oxvalue_n}$. We can use
  this and the definition of $\oxnorm{\oxstore}{\oxplaceexpr}{\oxreferent}
  {\oxvaluectx}{\oxvalue}$ to get
  $\oxnorminner{\oxstore}{\oxplaceexprctx}{\oxid}{\oxreferent}{\oxvaluectx}{
    \oxprod{\oxvalue_1 \oxdotsc \oxvalue_i \oxdotsc \oxvalue_n}
  }$ (where $\oxplaceexprctx[\oxid] = \oxplaceexpr$). This is precisely the
  premise of \oxname{P-Proj} and thus we can use that. We also have by inversion
  of \oxname{T-Tuple} for
  $\oxtypjudge{\oxglobalctx}{\oxtvarctx}{\oxkontctx}{\oxvarctx}{\oxvalue}
  {\oxtprod{\oxsitype_1 \oxdotsc \oxsitype_i \oxdotsc \oxsitype_n}}{\oxvarctx}$
  that $\oxtypjudge{\oxglobalctx}{\oxtvarctx}{\oxkontctx}{\oxvarctx}{\oxvalue_i}
  {\oxsitype_i}{\oxvarctx}$.

  \vspace{1em}
  \noindent \framebox[\textwidth]{
    \figuresize
    \begin{mathpar}
      \TCDeref
    \end{mathpar}
  }
  \vspace{1em}

  For \oxname{TC-Deref}, we apply our induction hypothesis to 
  $\oxcomputety{\oxtvarctx}{\oxvarctx}{\oxmuta}{
    \oxplaceexpr
  }{
    \oxtref{\oxprov}{\oxmuta^\prime}{\oxxitype}
  }{\oxset{\overline{\oxprov_p}}}$ to get
  $\oxnorm{\oxstore}{\oxplaceexpr}{\oxreferent}{\oxvaluectx}{\oxvalue}$ and
  $\oxtypjudge{\oxglobalctx}{\oxtvarctx}{\oxkontctx}{\oxvarctx}{\oxvalue}
  {\oxtref{\oxprov}{\oxmuta^\prime}{\oxxitype}}{\oxvarctx}$. Then, by
  \Lemma{lemma:canonical-forms}, we know that $\oxvalue$ must of the form
  $\oxptr{\oxreferent}$. We now have five subcases to consider depending on
  whether $\oxreferent$ is of $\oxplace$, $\oxindex{\oxreferent_3}{i}$, or
  $\oxslice{\oxreferent}{i}{j}$, and for the latter two, whether $\oxxitype$ is
  $\oxtarr{\oxsitype}{\oxnum}$ or $\oxtslice{\oxsitype}$.

  \vspace{1em}
  \noindent \framebox[\textwidth]{
    \figuresize
    \begin{mathpar}
      \PDerefPtr \and \PDerefIndexPtrArray \and \PDerefIndexPtrSlice \and
      \PDerefSlicePtrArray \and \PDerefSlicePtrSlice
    \end{mathpar}
  }
  \vspace{1em}

  In all these cases, we know structurally that $\oxplaceexprctx = \oxhole$
  since \oxname{TC-Deref} has no context outside of the dereference. So, for
  each of them, we need to be able to show $\oxnorm{\oxhole}{\oxreferent}{
    \oxreferent^\prime
  }{\oxvaluectx}{\oxvalue^\prime}$. Inversion on \oxname{T-Pointer} gives us
  $\oxreferentvalidity{\oxglobalctx}{\oxvarctx}{\oxreferent}{\oxxitype}$. We can
  then apply \Lemma{lemma:wf-refs-evalref} to get
  $\oxevalreferent{\oxglobalctx}{\oxvarctx}{\oxvaluectx}{\oxvalue}$. Then, we
  can apply \oxname{P-Referent} to this to produce the derivation we need to
  apply the appropriate rule. For \oxname{P-DerefIndexPtrArray} and
  \oxname{P-DerefSlicePtrArray}, we apply \Lemma{lemma:canonical-forms}
  to get that the value is an array. For \oxname{P-DerefIndexPtrSlice} and
  \oxname{P-DerefSlicePtrSlice}, we apply \Lemma{lemma:canonical-forms} to get
  that the value is a slice value.
\end{proof}

\begin{oxlemma}{Reduced Place Expressions Produce Valid Referents}
  {lemma:norm-for-valid-referents}
  If $\oxstorevalidity{\oxglobalctx}{\oxvarctx}{\oxstore}$
  and $\oxnorm{\oxstore}{\oxplaceexpr}{\oxreferentctx[\oxplace]}
  {\oxvaluectx}{\oxvalue}$,
  then
  $\oxreferentvalidity{\oxglobalctx}{\oxvarctx}{\oxreferentctx[\oxplace]}{\oxxitype}$.
\end{oxlemma}

\begin{proof}
  We start by rewriting 
  $\oxnorm{\oxstore}{\oxplaceexpr}{\oxreferentctx[\oxplace]}
  {\oxvaluectx}{\oxvalue}$ with its definition to get
  $\oxnorminner{\oxstore}{\oxplaceexprctx}{\oxid}{
    \oxreferentctx[\oxplace]
  }{\oxvaluectx}{\oxvalue}$ where $\oxplaceexpr = \oxplaceexprctx[\oxid]$. We
  then proceed by induction by cases (note this means our induction hypothesis
  is really about the rewritten form).

  \vspace{1em}
  \noindent \framebox[\textwidth]{
    \figuresize
    \begin{mathpar}
      \PId \and \WFRefId
    \end{mathpar}
  }
  \vspace{1em}

  \oxname{P-Referent} only applies if the context is $\oxhole$ which is only the
  case if our original place expression was $\oxid$. We can rewrite with this
  knowledge to see that we really have
  $\oxevalreferent{\oxstore}{\oxid}{\_}{\oxvalue}$ in our premise. Inversion on
  \oxname{ER-Id} gives us $\oxlookup{\oxstore}{\oxvalue}$ Then, we consider the
  frame of $\oxvarctx$ which contains $\oxid$. By inversion of
  \oxname{WF-StackFrame} for the portion of the derivation
  $\oxstorevalidity{\oxglobalctx}{\oxvarctx}{\oxstore}$ pertaining to that
  frame, we have $\forall \oxid \in \oxdomain{\oxstackframe}. \;
  \oxtypjudge{\oxglobalctx}{\oxemptyctx}{\oxemptyctx}{\oxnewframe{\oxvarctx}{\oxframe}}
  {(\oxnewframe{\oxstore}{\oxstackframe})(\oxid)}{
    (\oxnewframe{\oxvarctx}{\oxframe})(\oxid)
  }{\oxnewframe{\oxvarctx}{\oxframe}}$. Focusing on our particular $\oxid$, we
  have both that $\oxlookup{\oxvarctx}{\oxid}{\oxvalue}$. We can then apply
  \oxname{WF-RefId}.

  \vspace{1em}
  \noindent \framebox[\textwidth]{
    \figuresize
    \begin{mathpar}
      \PProj \and \WFRefProj
    \end{mathpar}
  }
  \vspace{1em}

  Applying the induction hypothesis to
  $\oxnorminner{\oxstore}{\oxplaceexprctx}{\oxreferent_1}{\oxreferent_2}{\_}{
    \oxprod{\oxvalue_0 \oxdotsc \oxvalue_i \oxdotsc \oxvalue_n}
  }$ gives us $\oxreferentvalidity{\oxglobalctx}{\oxvarctx}{\oxreferent_2}{
    \oxtprod{\oxsitype_0 \oxdotsc \oxsitype_i \oxdotsc \oxsitype_n}
  }$. We can then apply \oxname{WF-RefProjection}.
  
  \vspace{1em}
  \noindent \framebox[\textwidth]{
    \figuresize
    \begin{mathpar}
      \PDerefPtr
    \end{mathpar}
  }
  \vspace{1em}

  Applying the induction hypothesis to 
  $\oxnorminner{\oxstore}{\oxplaceexprctx}{\oxplace}{\oxreferent_2}{\_}{\oxvalue}$
  gives us $\oxreferentvalidity{\oxglobalctx}{\oxvarctx}{\oxreferent_2}{
    \oxxitype
  }$.

  \vspace{1em}
  \noindent \framebox[\textwidth]{
    \figuresize
    \begin{mathpar}
      \PDerefIndexPtrArray \and \WFRefIndexArray
    \end{mathpar}
  }
  \vspace{1em}

  Applying the induction hypothesis to
  $\oxnorminner{\oxstore}{\oxplaceexprctx}{\oxreferent_2}{\oxreferent_3}{\_}{
    \oxarr{\oxvalue_0 \oxdotsc \oxvalue_i \oxdotsc \oxvalue_n}
  }$ gives us $\oxreferentvalidity{\oxglobalctx}{\oxvarctx}{\oxreferent_3}{
    \oxtarr{\oxsitype}{\oxnum}
  }$. Then, we can apply \oxname{WF-RefIndexArray} to get
  $\oxreferentvalidity{\oxglobalctx}{\oxvarctx}{\oxindex{\oxreferent_3}{i}}{
    \oxsitype
  }$.

  \vspace{1em}
  \noindent \framebox[\textwidth]{
    \figuresize
    \begin{mathpar}
      \PDerefIndexPtrSlice \and \WFRefIndexSlice
    \end{mathpar}
  }
  \vspace{1em}

  Applying the induction hypothesis to
  $\oxnorminner{\oxstore}{\oxplaceexprctx}{\oxreferent_2}{\oxreferent_3}{\_}{
    \oxsliceval{\oxvalue_0 \oxdotsc \oxvalue_i \oxdotsc \oxvalue_n}
  }$ gives us $\oxreferentvalidity{\oxglobalctx}{\oxvarctx}{\oxreferent_3}{
    \oxtslice{\oxsitype}
  }$. Then, we can apply \oxname{WF-RefIndexSlice} to get
  $\oxreferentvalidity{\oxglobalctx}{\oxvarctx}{\oxindex{\oxreferent_3}{i}}{
    \oxsitype
  }$.

  \vspace{1em}
  \noindent \framebox[\textwidth]{
    \figuresize
    \begin{mathpar}
      \PDerefSlicePtrArray \and \WFRefSliceArray
    \end{mathpar}
  }
  \vspace{1em}

  Applying the induction hypothesis to
  $\oxnorminner{\oxstore}{\oxplaceexprctx}{\oxreferent_2}{\oxreferent_3}{\_}{
    \oxarr{
      \oxvalue_0 \oxdotsc \oxvalue_i \oxdotsc \oxvalue_j \oxdotsc \oxvalue_n
    }
  }$ gives us $\oxreferentvalidity{\oxglobalctx}{\oxvarctx}{\oxreferent_3}{
    \oxtarr{\oxsitype}{\oxnum}
  }$. Then, we can apply \oxname{WF-RefSliceArray} to get
  $\oxreferentvalidity{\oxglobalctx}{\oxvarctx}{\oxslice{\oxreferent_3}{i}{j}}{
    \oxsitype
  }$.

  \vspace{1em}
  \noindent \framebox[\textwidth]{
    \figuresize
    \begin{mathpar}
      \PDerefSlicePtrSlice \and \WFRefSliceSlice
    \end{mathpar}
  }
  \vspace{1em}

  Applying the induction hypothesis to
  $\oxnorminner{\oxstore}{\oxplaceexprctx}{\oxreferent_2}{\oxreferent_3}{\_}{
    \oxsliceval{
      \oxvalue_0 \oxdotsc \oxvalue_i \oxdotsc \oxvalue_j \oxdotsc \oxvalue_n
    }
  }$ gives us $\oxreferentvalidity{\oxglobalctx}{\oxvarctx}{\oxreferent_3}{
    \oxtslice{\oxsitype}
  }$. Then, we can apply \oxname{WF-RefSliceSlice} to get
  $\oxreferentvalidity{\oxglobalctx}{\oxvarctx}{\oxslice{\oxreferent_3}{i}{j}}{
    \oxsitype
  }$.
\end{proof}

\begin{oxlemma}{Reduced Place Expressions Have Roots in Loan Sets}
  {lemma:norm-place-prefix-in-loans}
  If $\oxstorevalidity{\oxglobalctx}{\oxvarctx}{\oxstore}$,
  $\oxnorm{\oxstore}{\oxplaceexpr}{\oxreferentctx[\oxplace]}
  {\oxvaluectx}{\oxvalue}$, and
  $\oxmusafety{\oxemptyctx}{\oxkontctx}{\oxvarctx}{\oxmuta}{\oxplaceexpr}{
    \oxset{\overline{\oxloan}}
  }$, then $\oxreferent = \oxreferentctx[\oxplace]$ and
  $\oxloanpkg{\oxmuta}{\oxplace} \in \oxset{\overline{\oxloan}}$.
\end{oxlemma}

\begin{proof}
  We proceed by induction on $\oxmusafety{\oxemptyctx}{\oxkontctx}{\oxvarctx}{\oxmuta}{
    \oxplaceexpr
  }{
    \oxset{\overline{\oxloan}}
  }$. There are ordinarily three cases: \oxname{O-SafePlace}, \oxname{O-Deref},
  and \oxname{O-DerefAbs}. However, \oxname{O-DerefAbs} requires the type
  variable context to contain entries, and thus can be immediately discharged by
  contradiction. This leaves us with only \oxname{O-SafePlace} and
  \oxname{O-Deref}.

  \vspace{1em}
  \noindent \framebox[\textwidth]{
    \figuresize
    \begin{mathpar}
      \OSafePlace
    \end{mathpar}
  }
  \vspace{1em}

  \oxname{O-SafePlace} tells us that our $\oxplaceexpr$ is in fact a place
  $\oxplace$ meaning that it does not contain any dereferences. As such, we know
  that $\oxnorm{\oxstore}{\oxplaceexpr}{\oxreferentctx[\oxplace]}
  {\oxvaluectx}{\oxvalue}$ must have been derived using a combination of
  \oxname{P-Referent} and \oxname{P-Proj} corresponding to the structure of
  $\oxplace$. The resulting referent in such a case is precisely $\oxplace$
  (meaning $\oxreferentctx = \oxhole$), which we know is in the output
  immediately from the definition of \oxname{O-SafePlace}.
  
  \vspace{1em}
  \noindent \framebox[\textwidth]{
    \figuresize
    \begin{mathpar}
      \ODeref
    \end{mathpar}
  }
  \vspace{1em}

  In the premise of \oxname{O-Deref}, we have a number of ownership safety
  derivations corresponding to each of the loans for the pointer being
  dereferenced. Since we know we have a dereference, we know that we must have
  derived $\oxnorm{\oxstore}{\oxplaceexpr}{\oxreferentctx[\oxplace]}
  {\oxvaluectx}{\oxvalue}$
  using one of the five dereference rules at the appropriate point
  (\oxname{P-DerefPtr}, \oxname{P-DerefIndexPtrArray},
  \oxname{P-DerefIndexPtrSlice}, \oxname{P-DerefSlicePtrArray}, and
  \oxname{P-DerefSlicePtrSlice}). Each of which share a common premise (at least
  when sufficiently generalized):
  $\oxnorminner{\oxstore}{\oxplaceexprctx}{\oxreferent_2}
  {\oxreferent_3}{\oxvaluectx}{\oxvalue}$. Here, $\oxreferent_2$ corresponds to
  the referent of the pointer we are dereferencing. As such, we know that one of
  the derivations of ownership safety corresponds to that particular referent.
  So, we can apply our induction hypothesis and get that
  $\oxloanpkg{\oxmuta}{\oxplace} \in \oxset{
    \overline{\oxloanpkg{\oxmuta}{\oxplaceexpr^\prime_i}} }$ for the appropriate
  ownership safety derivation numbered i. The final output is the union of all
  of these sets, and thus we can generalize to $\oxloanpkg{\oxmuta}{\oxplace}
  \in \oxset{ \overline{\oxloanpkg{\oxmuta}{\oxplaceexpr^\prime_1}} \oxdotsc
    \overline{\oxloanpkg{\oxmuta}{\oxplaceexpr^\prime_n}} \oxdotsc
    \oxloanpkg{\oxmuta}{\oxplaceexprctx[\oxderef{\oxplace}]}}$.
\end{proof}

\subsection{Preservation under Region Rewriting Lemmas}
\label{sec:rewriting-lemmas}

\begin{oxlemma}{Ownership Safety is Preserved under Region Rewriting}
  {lemma:subtyping-ownership-safety} If $\oxmusafetyinner{\oxemptyctx}{\oxkontctx}{\oxvarctx}
  {\oxmuta}{\overline{\oxplace_e}}{\oxplaceexpr}{\oxset{\overline{\oxloan^\prime}}}$ and
  $\oxtunify{\oxemptyctx}{\oxkontctx}{\oxvarctx}{\oxtype_1}{\oxtype_2}{\oxvarctx^\prime}$ and
  $\oxset{\overline{\oxplace_e}} \subseteq \oxset{\overline{\oxplace_e^\prime}}$ then
  $\oxmusafetyinner{\oxemptyctx}{\oxkontctx}{\oxvarctx^\prime}
  {\oxmuta}{\overline{\oxplace_e}}{\oxplaceexpr}{\oxset{\overline{\oxloan^{\prime\prime}}}}$.
\end{oxlemma}
\begin{proof}
  We proceed by induction on the region rewriting judgement. We note that if
  $\oxrewritemode$ is $=$, then by inspection of the outlives judgement, $\oxvarctx
  = \oxvarctx^\prime$, so the proof follows immediately from the premise. So
  consider when $\oxrewritemode$ is $+$. The only case that doesn't follow
  immediately by induction and application of premises is \oxname{RR-Reference},
  and in this case the only interesting part of the proof is the outlives
  constraint.

  Proceeding by induction on the outlives constraint, the only interesting case
  is \\ \oxname{OL-CombineConcrete}.

  \vspace{1em}
  \noindent \framebox[\textwidth]{ \figuresize
    \begin{mathpar}
      \UCombineLocalProvs
    \end{mathpar}
  } \vspace{1em}

  We want to show that $\oxmusafetyinner{\oxtvarctx}{\oxkontctx}
  {\oxvarctx[\oxcprov_2 \mapsto \oxset{\overline{\oxloan}}]}
  {\oxmuta}{\overline{\oxplace_e}}{\oxplaceexpr}{\oxset{\overline{\oxloan^{\prime\prime}}}}$.
  Proceed by induction on the ownership safety judgement in the premise.

  \vspace{1em}
  \noindent \framebox[\textwidth]{ \figuresize
    \begin{mathpar}
      \OSafePlace
    \end{mathpar}
  } \vspace{1em}

  Let $\oxcprov^\prime$ be an arbitrary region. There are two cases to prove,
  depending which part of the disjunction is true for the premise
  $\oxmusafetyinner{\oxtvarctx}{\oxkontctx}{\oxvarctx}{\oxmuta}{\overline{\oxplace_e}}{\oxplace}
  {\oxset{\overline{\oxloan^\prime}}}$.

  If the first part was true, then we need to show that $\forall
  \oxloanpkg{\oxmuta^\prime}{\oxplaceexprctx[\oxplace^\prime]} \in
  \oxvarctx[\oxcprov_2 \mapsto
    \oxset{\overline{\oxloan^{\prime\prime}}}](\oxcprov^\prime).\ (\oxmuta =
  \oxmut \vee \oxmuta^\prime = \oxmut) \implies
  \oxrelevant{\oxplace^\prime}{\oxplace}$. This is only interesting when
  $\oxcprov^\prime = \oxcprov_2$. Using the fact that $\oxvarctx[\oxcprov_2
    \mapsto \oxset{\oxloans}](\oxcprov_2) = \oxvarctx(\oxcprov_1) \cup
  \oxvarctx(\oxcprov_2)$, we need to show $\forall
  \oxloanpkg{\oxmuta^\prime}{\oxplaceexprctx[\oxplace^\prime]} \in
  \oxvarctx(\oxcprov_1) \cup \oxvarctx(\oxcprov_2).\ (\oxmuta = \oxmut \vee
  \oxmuta^\prime = \oxmut) \implies \oxrelevant{\oxplace^\prime}{\oxplace}$.
  This is immediate if we can show that $\oxcprov_1$ and $\oxcprov_2$ are not
  excluded. This is immediate from the region not reborrowed judgement. For
  $\oxcprov_1$ or $\oxcprov_2$ to be excluded, for each reference
  $\oxplace^{\prime\prime}$ that has $\oxcprov_1$ or $\oxcprov_2$, there would
  have to be a loan of the form $\oxplaceexprctx[\oxderef
    \oxplace^{\prime\prime}]$, but such loans are precisely what the region not
  reborrowed judgement excludes.

  If the second part was true, then we can prove the second part immediately
  from the hypothesis, because the types of references are unchanged,
  $\oxkontctx$ is unchanged, and the exclusion list can only grow.

  \vspace{1em}
  \noindent \framebox[\textwidth]{ \figuresize
    \begin{mathpar}
      \ODeref
    \end{mathpar}
  } \vspace{1em}

  Firstly, note that the exclusion list will be equal if $\oxcprov \neq
  \oxcprov_2$, and will be potentially larger if $\oxcprov = \oxcprov_2$.
  Therefore we can immediately apply our induction hypothesis to get ownership
  safety for $\oxplaceexprctx[\oxplaceexpr_i]$ under
  $\oxvarctx[\oxcprov_2 \mapsto \oxloans]$.

  For the rest of the case, apply identical reasoning to that in the
  \oxname{O-SafePlace} case.

  \vspace{1em}
  \noindent \framebox[\textwidth]{ \figuresize
    \begin{mathpar}
      \ODerefAbs
    \end{mathpar}
  } \vspace{1em}

  Since $\oxtvarctx = \oxemptyctx$, there are no valid reference types that have
  an abstract region, meaning the first hypothesis is a contradiction.

\end{proof}

\begin{oxlemma}{Type Computation is Preserved under Region Rewriting}
  {lemma:subtyping-type-computation} If
  $\oxcomputetynoprov{\oxemptyctx}{\oxnewframe{\oxvarctx}{\oxframe}}{\oxmuta}{\oxplaceexpr}{\oxsitype}$
  and
  $\oxtunify{\oxemptyctx}{\oxkontctx}{\oxvarctx}{\oxtype_1}{\oxtype_2}{\oxvarctx^\prime}$
  then
  $\oxcomputetynoprov{\oxemptyctx}{\oxnewframe{\oxvarctx^\prime}{\oxframe}}{\oxmuta}{\oxplaceexpr}{\oxsitype}$.
\end{oxlemma}

\begin{proof}
  The proof is immediate by inspection of the type computation judgement,
  because the only things considered in the judgement are the types of places in
  $\oxvarctx$, which cannot change through the region rewriting judgement (in
  other words, $\oxdomain{\oxvarctx} = \oxdomain{\oxvarctx^\prime}$).
\end{proof}

\begin{oxlemma}{Outlives is Preserved under Region Rewriting}
  {lemma:subtyping-outlives} If $\oxrunify{\oxemptyctx}{\oxkontctx}{\oxnewframe{\oxvarctx}{\oxframe}}{\oxcprov_1}{\oxcprov_2}{\oxnewframe{\oxvarctx_o}{\oxframe_o}}$ and
  $\oxtunify{\oxemptyctx}{\oxkontctx}{\oxvarctx}{\oxtype_1}{\oxtype_2}{\oxvarctx^\prime}$ then
  $\oxrunify{\oxemptyctx}{\oxkontctx}{\oxnewframe{\oxvarctx^\prime}{\oxframe}}{\oxcprov_1}{\oxcprov_2}{\oxnewframe{\oxvarctx_o^\prime}{\oxframe_o^\prime}}$.
\end{oxlemma}
\begin{proof}
  We proceed by induction on the outlives judgement. The only interesting cases
  are \oxname{OL-CombineConcrete} and \oxname{OL-CheckConcrete}. In both cases,
  the only non immediate premise is the region not reborrowed judgement. Proceed
  by induction over the region rewriting hypothesis, and in the interesting case
  \oxname{RR-Reference}, proceed by induction over the outlives judgement. In
  this case, the only interesting case is when the loan sets in $\oxvarctx$
  potentially change, \oxname{OL-CombineConcrete}. But note that no new loans
  are generated, only loans are copied into other sets. For this reason, the
  region not reborrowed judgements we already have are sufficient, because these
  loans that are now potentially in two loan sets were already found to not
  contain any problematic reborrows.
\end{proof}

\begin{oxlemma}{Region Rewriting is Preserved under Region Rewriting}
  {lemma:subtyping-subtyping} If
  $\oxtunify{\oxemptyctx}{\oxkontctx}{\oxnewframe{\oxvarctx}{\oxframe}}{\oxtype_1}{\oxtype_2}{\oxnewframe{\oxvarctx_o}{\oxframe_o}}$ and
  $\oxtunify{\oxemptyctx}{\oxkontctx}{\oxvarctx}{\oxtype_1^\prime}{\oxtype_2^\prime}{\oxvarctx^\prime}$ then
  $\oxtunify{\oxemptyctx}{\oxkontctx}{\oxnewframe{\oxvarctx^\prime}{\oxframe}}{\oxtype_1}{\oxtype_2}{\oxnewframe{\oxvarctx_o^\prime}{\oxframe_o^\prime}}$.
\end{oxlemma}
\begin{proof}
  Proceed by induction over the the region rewriting judgement. The only
  interesting case is \oxname{RR-Reference}, for which we just apply
  \Lemma{lemma:subtyping-outlives}.
\end{proof}

\begin{oxlemma}{Region Rewriting is Preserved by Garbage Collecting Loans}
  {lemma:subtyping-gc} If
  $\oxtunify{\oxemptyctx}{\oxkontctx}{\oxvarctx}{\oxtype_1^\prime}{\oxtype_2^\prime}{\oxvarctx^\prime}$
  then
  $\oxtunify{\oxemptyctx}{\oxkontctx}{\oxgcloans{\oxkontctx}{\oxvarctx}}{\oxtype_1^\prime}{\oxtype_2^\prime}{\oxvarctx^{\prime\prime}}$.
\end{oxlemma}
\begin{proof}
  We proceed by induction over the region rewriting judgement, in which the only
  interesting case is \oxname{RR-Reference}. We then proceed by induction over
  the outlives relation, in which the only interesting cases are
  \oxname{OL-CombineConcrete} and \oxname{OL-CheckConcrete}. The only
  interesting part of the judgement is the region not reborrowed, and this is
  immediate because garbage collection will only potentially remove some loans.
\end{proof}

\begin{oxlemma}{Closure Body Typing is Preserved under Region Rewriting}
  {lemma:subtyping-closure-typing} If
  $\oxtypjudge{\oxglobalctx}{\oxemptyctx}{\oxkontctx}{\oxnewframe{\oxvarctx}{\oxframe}}
  {\oxexpr}{\oxtype}{\oxnewframe{\oxvarctx_o}{\oxframe_o}}$ and
  $\oxtunify{\oxemptyctx}{\oxkontctx}{\oxvarctx}{\oxtype_1}{\oxtype_2}{\oxvarctx^\prime}$
  then
  $\oxtypjudge{\oxglobalctx}{\oxemptyctx}{\oxkontctx}{\oxnewframe{\oxvarctx^\prime}{\oxframe}}
  {\oxexpr}{\oxtype}{\oxnewframe{\oxvarctx_o^\prime}{\oxframe_o^\prime}}$ and
  $\oxtunify{\oxemptyctx}{\oxkontctx}{\oxvarctx_o^\prime}{\oxtype_1}{\oxtype_2}{\oxvarctx_o^{\prime\prime}}$.
\end{oxlemma}
\begin{proof}
  Proceed by induction over the typing derivation for $\oxexpr$.

  The \oxname{T-Abort}, \oxname{T-Function}, \oxname{T-Unit}, \oxname{T-u32},
  \oxname{T-True}, and \oxname{T-False} cases follow immediately.

  The \oxname{T-LetRegion}, \oxname{T-While}, \oxname{T-Closure},
  \oxname{T-Tuple}, \oxname{T-Array}, \oxname{T-Slice}, \oxname{T-Drop},
  \oxname{T-Left}, and \oxname{T-Right} cases all follow immediately from the
  induction hypothesis.

  The \oxname{T-Seq} case follows from the induction hypothesis and
  \Lemma{lemma:subtyping-gc}.

  The \oxname{T-Branch}, \oxname{T-Let}, and \oxname{T-Match} cases follow from
  the induction hypothesis, \Lemma{lemma:subtyping-subtyping}, and
  \Lemma{lemma:subtyping-gc}. Note the reborrow restriction follows immediately
  from the fact that rewriting can at most union together loan sets, which means
  the overall loans considered for the region not reborrowed judgement are the
  same in the context after rewriting.

  The \oxname{T-Move}, \oxname{T-Copy}, \oxname{T-Borrow},
  \oxname{T-BorrowIndex}, \oxname{T-BorrowSlice}, \oxname{T-IndexCopy},
  \oxname{T-ForArray}, and \oxname{T-ForSlice} cases follow from the induction
  hypothesis, \Lemma{lemma:subtyping-ownership-safety}, and
  \Lemma{lemma:subtyping-type-computation}.

  The \oxname{T-AppFunction} and \oxname{T-AppClosure} cases follow from the
  induction hypothesis, \Lemma{lemma:subtyping-outlives}, and the fact that
  context, region, and type well formedness aren't affected by changes in the
  loan sets.
\end{proof}

\begin{oxlemma}{Value Typing is Preserved under Region Rewriting}
  {lemma:subtyping-value-typing} If
  $\oxtypjudge{\oxglobalctx}{\oxemptyctx}{\oxkontctx}{\oxvarctx}
  {\oxvalue}{\oxtype}{\oxvarctx}$ and
  $\oxtunify{\oxemptyctx}{\oxkontctx}{\oxvarctx}{\oxtype_1}{\oxtype_2}{\oxvarctx^\prime}$
  then $\oxtypjudge{\oxglobalctx}{\oxemptyctx}{\oxkontctx}{\oxvarctx^\prime}
  {\oxvalue}{\oxtype}{\oxvarctx^\prime}$.
\end{oxlemma}

\begin{proof}
  We proceed by induction on the value typing.

  \vspace{1em}
  \noindent \framebox[\textwidth]{ \figuresize
    \begin{mathpar}
      \inferrule[T-Pointer]{
        %% premises
        \oxreferentvalidity{\oxglobalctx}{\oxvarctx}{\oxreferentctx[\oxplace]}{\oxxitype} \\
        \oxloanpkg{\oxmuta}{\oxplace} \in \oxvarctx(\oxcprov) \\
      }{
        %% conclusions
        \oxtypjudge{\oxglobalctx}{\oxemptyctx}{\oxkontctx}{\oxvarctx}{
          \oxptr{\oxreferentctx[\oxplace]} }{
          \oxtref{\oxcprov}{\oxmuta}{\oxxitype} }{\oxvarctx} }
    \end{mathpar}
  } \vspace{1em}

  The \oxname{T-Pointer} case is immediate, because by inspection of the
  referent well formedness, there is no reliance on loan sets, and the loan is
  preserved since by inspection of the rewriting judgement, the loan sets either
  stay the same or potentially grow.
  
  \vspace{1em}
  \noindent \framebox[\textwidth]{ \figuresize
    \begin{mathpar}
      \TClosureVal
    \end{mathpar}
  } \vspace{1em}

  First, we invert the stack frame typing hypothesis to get that $\forall \oxid
  \in \oxdomain{\oxstackframe}.\;
  \oxtypjudge{\oxglobalctx}{\oxemptyctx}{\oxkontctx}{\oxnewframe{\oxvarctx}{\oxframe_c}}{\oxstackframe(\oxid)}{\oxframe_c(\oxid)}{\oxnewframe{\oxvarctx}{\oxframe_c}}$.
  We can apply the induction hypothesis to each of these statements, and apply
  \oxname{WF-Frame} to get
  $\oxsubstorevalidity{\oxglobalctx}{\oxvarctx^\prime}{\oxstackframe_c}{\oxframe_c}$.

  For the typing of the body, we can apply \Lemma{lemma:subtyping-closure-typing}.
\end{proof}

\begin{oxlemma}{Stack Well-Formedness is Preserved under Region Rewriting}
  {lemma:stack-validity-region-rewriting} If
  $\oxstorevalidity{\oxglobalctx}{\oxvarctx}{\oxstore}$ and
  $\oxtunify{\oxemptyctx}{\oxkontctx}{\oxvarctx}{\oxtype_1}{\oxtype_2}{\oxvarctx^\prime}$
  then $\oxstorevalidity{\oxglobalctx}{\oxvarctx^\prime}{\oxstore}$.
\end{oxlemma}

\begin{proof}
  We proceed by induction on the stack typing derivation.
  
  \vspace{1em}
  \noindent \framebox[\textwidth]{ \figuresize
    \begin{mathpar}
      \VStack \and \VStackEmpty
    \end{mathpar}
  } \vspace{1em}

  The \oxname{WF-StackEmpty} case is immediate. In the \oxname{WF-StackFrame}
  case, we get the well formedness in the premise from our induction hypothesis.
  What's left to show is that for all of the values $\oxvalue$ in the stack
  frame, they remain well typed in $\oxvarctx^\prime$. This follows from
  applying \Lemma{lemma:subtyping-value-typing}.
\end{proof}

\subsection{Preservation under Drops and Garbage Collection Lemmas}
\label{sec:environment-relation-lemmas}

\begin{oxlemma}{Values Change Environments in Limited Ways}{lemma:values-dont-change}
  If \oxtypjudge{\oxglobalctx}{\oxtvarctx}{\oxkontctx}{\oxvarctx}{\oxvalue}
  {\oxtype}{\oxvarctx^\prime}, then
  $\oxsubtypectx{\oxglobalctx}{\oxtvarctx}{\oxvarctx}{\oxvarctx^\prime}$.
\end{oxlemma}

\begin{proof}
  We proceed by induction on the structure of the typing derivation. Since we
  assume that the expression being typed is a value, we need only consider the
  cases that can be used to type a value.

  For many cases, the output environments are precisely the input environments,
  and thus this holds immediately. These cases are \oxname{T-Unit},
  \oxname{T-u32}, \oxname{T-True}, \oxname{T-False}, \oxname{T-Pointer},
  \oxname{T-Function}, \oxname{T-ClosureValue}, and \oxname{T-Dead}.

  For \oxname{T-Tuple}, \oxname{T-Array}, \oxname{T-Left}, and \oxname{T-Right},
  knowing that we have a value means that all of the subterms are themselves
  values, and thus we can apply our induction hypothesis to them in sequence
  (relying on the transitivity of $\oxsubtype$ for stack typings).

  This leaves us with one remaining case: \oxname{T-Drop}.

  \vspace{1em}
  \noindent \framebox[\textwidth]{ \figuresize
    \begin{mathpar}
      \TDrop
    \end{mathpar}
  } \vspace{1em}

  For \oxname{T-Drop}, we apply our induction hypothesis to
  $\oxtypjudge{\oxglobalctx}{\oxtvarctx}{\oxkontctx}{
    \oxtupdate{\oxvarctx}{\oxplace}{\oxsidtype_\oxplace}
  }{\oxexpr}{\oxsxtype}{\oxvarctx_f}$ which tells us that
  $\oxsubtypectx{\oxglobalctx}{\oxtvarctx}{
    \oxtupdate{\oxvarctx}{\oxplace}{\oxsidtype_\oxplace} }{\oxvarctx_f}$. Then,
  by \oxname{R-Env}, we have that
  $\oxsubtypectx{\oxglobalctx}{\oxtvarctx}{\oxvarctx}
  {\oxtupdate{\oxvarctx}{\oxplace}{\oxsidtype_\oxplace}}$. Then, by
  transitivity, we have
  $\oxsubtypectx{\oxglobalctx}{\oxtvarctx}{\oxvarctx}{\oxvarctx_f}$.
\end{proof}


\begin{oxlemma}{Type Computation is Preserved in Related Environments}
  {lemma:tc-related-contexts} If $\oxsubtypectx{\oxglobalctx}{\oxtvarctx}
  {\oxvarctx}{\oxvarctx^\prime}$ and
  $\oxcomputety{\oxtvarctx}{\oxvarctx}{\oxmuta}{\oxplaceexprctx[\oxplace]}{\oxtype}{\oxset{\overline{\oxprov}}}$
  and $\oxvarctx(\oxplace) = \oxvarctx^\prime(\oxplace)$, then
  $\oxcomputety{\oxtvarctx}{\oxvarctx^\prime}{\oxmuta}{\oxplaceexprctx[\oxplace]}{\oxtype}{\oxset{\overline{\oxprov}}}$.
\end{oxlemma}
\begin{proof}
  We proceed by induction on the type computation derivation. \oxname{TC-Var}
  follows immediately by the same type hypothesis, and \oxname{TC-Proj} follows
  from applying the induction hypothesis. All that is left is \oxname{TC-Deref}.

  \vspace{1em}
  \noindent \framebox[\textwidth]{ \figuresize
    \begin{mathpar}
      \TCDeref
    \end{mathpar}
  } \vspace{1em}

  First, we can apply the induction hypothesis to get the type computation for
  $\oxplaceexpr$. Then, all that's left is to show the outlives constraint, but
  this is immediate because $\oxtvarctx$ is unchanged and both $\oxvarctx$ and
  $\oxvarctx^\prime$ have the exact same domains.

\end{proof}

\begin{oxlemma}{Ownership Safety Preserved in Related Environments}
  {lemma:ownership-safety-related-envs} If
  $\oxmusafetyinner{\oxtvarctx}{\oxkontctx}{\oxvarctx}{\oxmuta}{\overline{\oxplace_e}}{\oxplaceexpr}
  {\oxset{\overline{\oxloan}}}$ and
  $\oxsubtypectx{\oxglobalctx}{\oxtvarctx}{\oxvarctx}{\oxvarctx^\prime}$ and
  $\oxcomputetynoprov{\oxtvarctx}{\oxvarctx^\prime}{\oxmuta}{\oxplaceexpr}{\oxxitype}$
  and $\oxplaceexpr = \oxplaceexprctx[\oxplace_\oxplaceexpr]$ and
  $\oxvarctx(\oxplace_\oxplaceexpr) = \oxvarctx^\prime(\oxplace_\oxplaceexpr)$,
  then
  $\oxmusafetyinner{\oxtvarctx}{\oxkontctx}{\oxvarctx^\prime}{\oxmuta}{\overline{\oxplace_e}}{\oxplaceexpr}
  {\oxset{\overline{\oxloan}}}$.
\end{oxlemma}

\begin{proof}
  We proceed by induction on the $\omega$-safety derivation, for which there are
  three cases to consider.

  \vspace{1em}
  \noindent \framebox[\textwidth]{ \figuresize
    \begin{mathpar}
      \OSafePlace
    \end{mathpar}
  } \vspace{1em}

  We'd like to show that \oxname{O-SafePlace} can be applied with context
  $\oxvarctx^\prime$. First, note that for any $\oxcprov^\prime$, if the right
  side of the or is true for $\oxvarctx$ with $\overline{\oxplace}$ then it will
  be true for $\oxvarctx^\prime$ with $\overline{\oxplace}$. That is, if all of
  the pointers with region $\oxcprov^\prime$ in $\oxvarctx$ are in the exclusion
  list $\overline{\oxplace}$, then all of the pointers with region
  $\oxcprov^\prime$ in $\oxvarctx^\prime$ are also in the exclusion list
  $\overline{\oxplace}$. Note that $\oxkontctx$ is unchanged between the two.
  Therefore, the only cases we need to consider are where $\oxcprov^\prime$
  occurs in pointers in $\oxvarctx$ and $\oxvarctx^\prime$ that do not occur in
  $\overline{\oxplace}$.

  Since the only allowed change to loan sets is emptying, and an emptied loan
  set has the left side of the disjunction as vacuously true, and if the loan
  set is the same we have the condition from the ownership safety in the
  premise, we are done.

  \vspace{1em}
  \noindent \framebox[\textwidth]{ \figuresize
    \begin{mathpar}
      \ODeref
    \end{mathpar}
  } \vspace{1em}

  Firstly, we have that $\oxvarctx(\oxplace_i) = \oxvarctx^\prime(\oxplace_i)$,
  because $\oxvarctx^\prime(\oxplace_i)$ must be an initialized type by the type
  computation premise, and the only changes in types between $\oxvarctx$ and
  $\oxvarctx^\prime$ allowed by the environment relation is dropping some types
  to uninitialized.

  Second, note that $\oxvarctx^\prime(\oxcprov) = \oxvarctx(\oxcprov)$ since
  $\oxvarctx^\prime(\oxplace)$ being a reference with region $\oxcprov$ means we
  can't empty the loan set. So we proceed by applying the induction hypothesis
  for all $n$ loans, noting that the type computation requirement follows from
  the well formedness of $\oxvarctx^\prime$.

  Finally, we have to show the statement about no conflicting loans, but here
  the argument is identical to that in the \oxname{O-SafePlace} case. If the
  loan set is empty then we're done, otherwise we just use the ownership safety
  premise.

  \vspace{1em}
  \noindent \framebox[\textwidth]{ \figuresize
    \begin{mathpar}
      \ODerefAbs
    \end{mathpar}
  } \vspace{1em}

  This case proceeds similarly to the \oxname{O-Deref} case, but with an added
  application of \Lemma{lemma:tc-related-contexts} to get the type computation,
  and no application of any induction hypothesis.
\end{proof}

\begin{oxlemma}{Types Are Well Formed in Related Environments}
  {lemma:types-well-formed} If $\oxsubtypectx{\oxglobalctx}{\oxtvarctx}
  {\oxvarctx}{\oxvarctx^\prime}$ and $\oxtypevalidity{\oxglobalctx}{\oxtvarctx}
  {\oxvarctx}{\oxxitype}$ and $\forall \oxcprov$ that occur in $\oxxitype$,
  $\oxvarctx(\oxcprov) = \oxvarctx^\prime(\oxcprov)$, then
  $\oxtypevalidity{\oxglobalctx}{\oxtvarctx} {\oxvarctx^\prime}{\oxxitype}$.
\end{oxlemma}

\begin{proof}
  We proceed by induction on the type well formedness derivation. The only case
  that doesn't follow directly from induction and the fact that $\oxtvarctx$ and
  $\oxkontctx$ are unchanged between the two related environments is
  \oxname{WF-Ref}.

  \vspace{1em}
  \noindent \framebox[\textwidth]{ \figuresize
    \begin{mathpar}
      \VRef
    \end{mathpar}
  } \vspace{1em}

  Firstly we apply our induction hypothesis to get that
  $\oxtypevalidity{\oxglobalctx}{\oxtvarctx}
  {\oxvarctx^\prime}{\oxxitype_\oxplaceexpr}$. What's left to show is the loan
  set condition on $\oxcprov$. If $\oxvarctx^\prime(\oxcprov) = \emptyset$, then
  we're done. Otherwise, we just need that the type computation still holds,
  which we get from \Lemma{lemma:tc-related-contexts}. We know the places in
  these place expressions all have the same type in $\oxvarctx$ and
  $\oxvarctx^\prime$ because between these two contexts the only changes allowed
  that could cause problems here are dropping one of these places, but then
  $\oxvarctx^\prime$ would not be well formed since there would be an invalid
  loan.

\end{proof}

\begin{oxlemma}{Related Environments Remain Well-Formed}
  {lemma:related-contexts-well-formed} If
  $\oxsubtypectx{\oxglobalctx}{\oxtvarctx} {\oxvarctx}{\oxvarctx^\prime}$ and
  $\oxctxswellformed{\oxglobalctx}{\oxtvarctx}
  {\oxnewframe{\oxvarctx}{\oxframe_{c}}}{\oxkontctx}$ then
  $\oxctxswellformed{\oxglobalctx}{\oxtvarctx}
  {\oxnewframe{\oxvarctx^\prime}{\oxframe_{c}}}{\oxkontctx}$.
\end{oxlemma}

\begin{proof}
  From the well formedness of $\oxnewframe{\oxvarctx}{\oxframe_c}$, we know that
  the places and disjointness conditions both hold. We also know that the occurs
  in restriction holds, because we at most have the same alive types. By
  \Lemma{lemma:types-well-formed}, noting that
  $\oxsubtypectx{\oxglobalctx}{\oxtvarctx}{\oxnewframe{\oxvarctx}{\oxframe_c}}{\oxnewframe{\oxvarctx^\prime}{\oxframe_c}}$
  is immediate, we know that the types remain well formed in the environment. We
  also have the well formedness of $\oxvarctx^\prime$ as a premise of the
  related environments judgement. All that's left to show is the loan set
  condition. But for this all we have to show is that each place computes to
  some type, which follows from \Lemma{lemma:tc-related-contexts}. We know the
  types of the places in each place expression remain the same because the only
  allowed changes between $\oxvarctx$ and $\oxvarctx^\prime$ are that places can
  be dropped and loan sets emptied, but if one such place was dropped, then
  $\oxvarctx^\prime$ would have not been well formed.
\end{proof}

\begin{oxlemma}{Related Input Environments Produce Similar Output Environments}
  {lemma:related-input-output-differences} If:
  \begin{itemize}
  \item
    $\oxtypjudge{\oxglobalctx}{\oxtvarctx}{\oxkontctx}{\oxvarctx_1}{\oxexpr_1}{\oxtype_1}{\oxvarctx_2}$
  \item
    $\oxtypjudge{\oxglobalctx}{\oxtvarctx}{\oxkontctx}{\oxvarctx_1}{\oxexpr_2}{\oxtype_2}{\oxvarctx_3}$
  \item
    $\oxsubtypectx{\oxglobalctx}{\oxtvarctx}{\oxvarctx_1}{\oxvarctx_1^\prime}$
  \item
    $\oxtypjudge{\oxglobalctx}{\oxtvarctx}{\oxkontctx}{\oxvarctx_1^\prime}{\oxexpr_1}{\oxtype_1}{\oxvarctx_2^\prime}$
  \item
    $\oxtypjudge{\oxglobalctx}{\oxtvarctx}{\oxkontctx}{\oxvarctx_1^\prime}{\oxexpr_2}{\oxtype_2}{\oxvarctx_3^\prime}$
  \item
    $\oxsubtypectx{\oxglobalctx}{\oxtvarctx}{\oxvarctx_2}{\oxvarctx_2^\prime}$
  \item
    $\oxsubtypectx{\oxglobalctx}{\oxtvarctx}{\oxvarctx_3}{\oxvarctx_3^\prime}$
  \item $\forall \oxid \in \oxdomain{\oxvarctx_2}$, $\oxvarctx_2(\oxid) =
    \oxvarctx_3(\oxid)$ and $\oxvarctx_2^\prime(\oxid) =
    \oxvarctx_3^\prime(\oxid)$
  \item $\forall \oxcprov$ that occur in $\oxexpr_1$ or $\oxexpr_2$ or
    $\oxtype_1$ or $\oxtype_2$, $\oxvarctx_1(\oxcprov) =
    \oxvarctx_1^\prime(\oxcprov)$
  \end{itemize}

  then $\forall \oxcprov \in \oxdomain{\oxvarctx_1}$, if
  $\oxvarctx_2^\prime(\oxcprov) = \emptyset$ and $\oxvarctx_3^\prime(\oxcprov)
  \neq \emptyset$, then $\oxvarctx_2(\oxcprov) = \emptyset$, and if
  $\oxvarctx_3^\prime(\oxcprov) = \emptyset$ and $\oxvarctx_2^\prime(\oxcprov)
  \neq \emptyset$, then $\oxvarctx_3(\oxcprov) = \emptyset$.
\end{oxlemma}
\begin{proof}
  The proofs for both statements in the conclusion follow identically, so
  without loss of generality it suffices to show that if
  $\oxvarctx_2^\prime(\oxcprov) = \emptyset$ and $\oxvarctx_3^\prime(\oxcprov)
  \neq \emptyset$, then $\oxvarctx_2(\oxcprov) = \emptyset$. Note there are two
  cases to consider: that the loan set was empty all along, or that the loan set
  was at some point non empty, but then got garbage collected.

  First, at some point between $\oxvarctx_1^\prime$ and $\oxvarctx_2^\prime$,
  $\oxcprov$ mapped to a non empty set of loans but then was garbage collected.
  In this case, $\oxvarctx_2^\prime$ must not contain any references that
  contain $\oxcprov$, since otherwise it would have been invalid to garbage
  collect $\oxcprov$. But then since $\oxvarctx_2^\prime$ and
  $\oxvarctx_3^\prime$ agree on types, it must be the case that it was also
  garbage collected in $\oxvarctx_3^\prime$, which is a contradiction with the
  fact that $\oxvarctx_3^\prime(\oxcprov)$ is non empty, so this case is
  impossible.

  Second, at each step of the derivation between $\oxvarctx_1^\prime$ and
  $\oxvarctx_2^\prime$, $\oxcprov$ mapped to empty. If $\oxvarctx_1(\oxcprov)$
  also was empty, then this means $\oxvarctx_2(\oxcprov)$ is also empty, and
  we're done. Otherwise, $\oxcprov$ was garbage collected between $\oxvarctx_1$
  and $\oxvarctx_1^\prime$. But then $\oxcprov$ must be free in $\oxexpr_2$ for
  loans to have been added between $\oxvarctx_1^\prime$ and
  $\oxvarctx_3^\prime$, which means the loan set could not have been emptied
  between $\oxvarctx_1$ and $\oxvarctx_1^\prime$, which is a contradiction.
\end{proof}

\begin{oxlemma}{Outlives Preserves Related Environments}
  {lemma:outlives-related} If
  $\oxrunify{\oxtvarctx}{\oxkontctx}{\oxvarctx}{\oxprov_1}{\oxprov_2}{\oxvarctx_o}$, and
  $\oxsubtypectx{\oxglobalctx}{\oxtvarctx}{\oxvarctx}{\oxvarctx^\prime}$ and
  $\oxctxswellformed{\oxglobalctx}{\oxtvarctx}{\oxvarctx_o}{\oxkontctx}$ and
  $\oxvarctx(\oxprov_1) = \oxvarctx^\prime(\oxprov_1)$ and $\oxvarctx(\oxprov_2)
  = \oxvarctx^\prime(\oxprov_2)$, then
  $\oxrunify{\oxtvarctx}{\oxkontctx}{\oxvarctx^\prime}{\oxprov_1}{\oxprov_2}{\oxvarctx_o^\prime}$,
  and
  $\oxsubtypectx{\oxglobalctx}{\oxtvarctx}{\oxvarctx_o}{\oxvarctx_o^\prime}$.
  and $\oxvarctx_o(\oxprov_1) = \oxvarctx_o^\prime(\oxprov_1)$ and
  $\oxvarctx_o(\oxprov_2) = \oxvarctx_o^\prime(\oxprov_2)$
\end{oxlemma}

\begin{proof}
  Proceed by induction on the outlives derivation. \oxname{OL-Refl},
  \oxname{OL-Trans}, \oxname{OL-AbstractConcrete}, and \oxname{OL-BothAbstract}
  are immediate.

  \oxname{OL-ConcreteAbstract} follows from additionally applying
  \Lemma{lemma:tc-related-contexts}. The condition on the place having the same
  type follows from the fact that $\oxplaceexpr$ is a loan and
  $\oxvarctx^\prime(\oxcprov)$ is not emptied, so we could not have dropped the
  place.

  \oxname{OL-CheckConcrete} is immediate, because the occurs before condition is
  unaffected since the domains are equal, and the region not reborrowed
  judgement is unaffected by adding loans that are already in other loan sets.

  This leaves two cases which proceed similarly: \oxname{OL-CombineConcrete} and
  \oxname{OL-CombineConcreteUnrestricted}.
  
  \vspace{1em}
  \noindent \framebox[\textwidth]{ \figuresize
    \begin{mathpar}
      \UCombineLocalProvs
    \end{mathpar}
  } \vspace{1em}

  Since $\oxvarctx^\prime(\oxcprov_1) = \oxvarctx(\oxcprov_1)$ and
  $\oxvarctx^\prime(\oxcprov_2) = \oxvarctx(\oxcprov_2)$,
  $\oxvarctx^\prime(\oxcprov_1) \cup \oxvarctx^\prime(\oxcprov_2) =
  \oxvarctx(\oxcprov_1) \cup \oxvarctx^\prime(\oxcprov_2)$. The region not
  reborrowed judgement is unaffected by adding loans that are already in other
  loan sets, so those conditions are also still true. The rest of the conditions
  are immediate: the equality on $\oxcprov_1$ and $\oxcprov_2$'s loan sets, the
  closure restriction since types are at most the same, and well formedness.
\end{proof}

\begin{oxlemma}{Related Environments Preserved by Region Rewriting}
  {lemma:unified-types-well-formed} If
  $\oxtunify{\oxtvarctx}{\oxkontctx}{\oxvarctx}{\oxsitype_1}{\oxsitype_2}{\oxvarctx_o}$, and
  $\oxsubtypectx{\oxglobalctx}{\oxtvarctx}{\oxvarctx}{\oxvarctx^\prime}$ and
  $\oxctxswellformed{\oxglobalctx}{\oxtvarctx}{\oxvarctx_o}{\oxkontctx}$ and $\forall
  \oxcprov$ that occur in $\oxsitype_1$ or $\oxsitype_2$, $\oxvarctx(\oxcprov) =
  \oxvarctx^\prime(\oxcprov)$, then
  $\oxtunify{\oxtvarctx}{\oxkontctx}{\oxvarctx^\prime}{\oxsitype_1}{\oxsitype_2}{\oxvarctx_o^\prime}$,
  and
  $\oxsubtypectx{\oxglobalctx}{\oxtvarctx}{\oxvarctx_o}{\oxvarctx_o^\prime}$,
  and $\forall \oxcprov$ that occur in $\oxsitype_1$ or $\oxsitype_2$,
  $\oxvarctx_o(\oxcprov) = \oxvarctx_o^\prime(\oxcprov)$.
\end{oxlemma}

\begin{proof}
  Proceed by induction on the region rewriting derivation. The only interesting
  case is \oxname{RR-Reference}, which proceeds by
  \Lemma{lemma:outlives-related} in addition to applying the induction
  hypothesis.
\end{proof}

\begin{oxlemma}{Expression Typing Preserved in Related Environments}
  {lemma:expression-typing-related-envs} Let $\oxexpr$ be a surface expression
  as defined on page 1. If
  $\oxtypjudge{\oxglobalctx}{\oxtvarctx}{\oxkontctx}{\oxnewframe{\oxvarctx}{\oxframe}}
  {\oxexpr}{\oxtype}{\oxnewframe{\oxvarctx_o}{\oxframe_o}}$ and
  $\oxsubtypectx{\oxglobalctx}{\oxtvarctx}{\oxnewframe{\oxvarctx}{\oxframe}}
  {\oxnewframe{\oxvarctx^\prime}{\oxframe}}$ and $\oxfvars{\oxexpr} =
  \overline{\oxid_f} \subseteq \oxdomain{\oxframe}|_\oxid$ and $\forall \oxcprov
  \in \oxfprovs{\oxexpr}.$ $\oxcprov \in \oxdomain{\oxframe}$, and $\forall
  \oxcprov$ that occur a type in $\overline{\oxframe(\oxid_f)}$,
  $\oxvarctx(\oxcprov) = \oxvarctx^\prime(\oxcprov)$ then
  $\oxtypjudge{\oxglobalctx}{\oxtvarctx}{\oxkontctx}{\oxnewframe{\oxvarctx^\prime}{\oxframe}}
  {\oxexpr}{\oxtype}{\oxnewframe{\oxvarctx_o^\prime}{\oxframe_o}}$ and
  $\oxsubtypectx{\oxglobalctx}{\oxtvarctx}{\oxnewframe{\oxvarctx_o}{\oxframe_o}}
  {\oxnewframe{\oxvarctx_o^\prime}{\oxframe_o}}$ and $\forall \oxcprov$ that
  occur a type in $\overline{\oxframe(\oxid_f)}$, $\oxvarctx_o(\oxcprov) =
  \oxvarctx_o^\prime(\oxcprov)$.
\end{oxlemma}
\begin{proof}
  Proceed by induction on the typing derivation for $e$. In the cases of
  \oxname{T-Abort}, \oxname{T-Function}, \oxname{T-Unit}, \oxname{T-u32},
  \oxname{T-True}, and \oxname{T-False}, the results are immediate.

  In the cases of \oxname{T-LetRegion}, \oxname{T-While}, \oxname{T-ForArray},
  \oxname{T-ForSlice}, \oxname{T-Closure}, \oxname{T-Left}, and
  \oxname{T-Right}, they all follow immediately from induction hypotheses.

  For each of the following cases, the convention is that the statement in the
  box is our assumption, and we want to prove the same statement with
  $\oxvarctx^\prime$ replaced for each $\oxvarctx$.

  \vspace{1em}
  \noindent \framebox[\textwidth]{ \inferrule[T-Tuple]{
      %% premises
      \forall i \in \oxset{1 \oxdots n}. \;
      \oxtypjudge{\oxglobalctx}{\oxtvarctx}{
        \oxextendctx{\oxkontctx}{
          \oxkontctxentry{\oxsitype_1} \oxdotsc \oxkontctxentry{\oxsitype_{i-1}}
        }
      }{
        \oxnewframe{\oxvarctx_{i-1}}{\oxframe_{i-1}}
      }{ \oxexpr_i
      }{\oxsitype_i}{\oxnewframe{\oxvarctx_i}{\oxframe_i}} 
    }{
      %% conclusions
      \oxtypjudge{\oxglobalctx}{\oxtvarctx}{\oxkontctx}{\oxnewframe{\oxvarctx_0}{\oxframe_0}}{
        \oxprod{\oxexpr_1 \oxdotsc \oxexpr_n} }{\oxtprod{\oxsitype_1 \oxdotsc
          \oxsitype_n}}{ \oxnewframe{\oxvarctx_n}{\oxframe_n} } } }

  We have $n$ induction hypotheses, each giving us the properties for input
  context $\oxnewframe{\oxvarctx^\prime_{i-1}}{\oxframe^\prime_{i-1}}$ and output context
  $\oxnewframe{\oxvarctx^\prime_i}{\oxframe^\prime_i}$.

  Given these resulting $n$ typing judgements, we get from applying
  \oxname{T-Tuple} that
  $\oxtypjudge{\oxglobalctx}{\oxtvarctx}{\oxkontctx}{\oxnewframe{\oxvarctx^\prime_0}{\oxframe_0}}{
    \oxprod{\oxexpr_1 \oxdotsc \oxexpr_n} }{\oxtprod{\oxsitype_1 \oxdotsc
      \oxsitype_n}}{ \oxnewframe{\oxvarctx^\prime_o}{\oxframe^\prime_o} }$, as
  well as the related environments judgement
  $\oxsubtypectx{\oxglobalctx}{\oxtvarctx}{\oxnewframe{\oxvarctx_n}{\oxframe_n}}{\oxnewframe{\oxvarctx_n^\prime}{\oxframe_n^\prime}}$.

  The cases for \oxname{T-Array} and \oxname{T-Slice} proceed identically. This
  reasoning is also used in the \oxname{T-App} case.


  \vspace{1em}
  \noindent \framebox[\textwidth]{ \figuresize
    \begin{mathpar}
      \inferrule[T-Branch]{
        %% premises
        \oxtypjudge{\oxglobalctx}{\oxtvarctx}{\oxkontctx}{\oxnewframe{\oxvarctx}{\oxframe}}{
          \oxexpr_1
        }{\oxtbool}{\oxnewframe{\oxvarctx_1}{\oxframe_1}} \\
        \oxtypjudge{\oxglobalctx}{\oxtvarctx}{\oxkontctx}{\oxnewframe{\oxvarctx_1}{\oxframe_1}}{
          \oxexpr_2
        }{\oxsitype_2}{\oxnewframe{\oxvarctx_2}{\oxframe_2}} \\\\
        \oxtypjudge{\oxglobalctx}{\oxtvarctx}{\oxkontctx}{\oxnewframe{\oxvarctx_1}{\oxframe_1}}{
          \oxexpr_3
        }{\oxsitype_3}{\oxnewframe{\oxvarctx_3}{\oxframe_3}} \\
        \oxsitype = \oxsitype_2 \vee \oxsitype = \oxsitype_3 \\\\
        \oxtunify{\oxtvarctx}{\oxkontctx}{\oxnewframe{\oxvarctx_2}{\oxframe_2}}
        {\oxsitype_2}{\oxsitype}{\oxnewframe{\oxvarctx_{2s}}{\oxframe_{2s}}} \\
        \oxtunify{\oxtvarctx}{\oxkontctx}{\oxnewframe{\oxvarctx_3}{\oxframe_3}}
        {\oxsitype_3}{\oxsitype}{\oxnewframe{\oxvarctx_{3s}}{\oxframe_{3s}}} \\
        \oxnewframe{\oxvarctx_{2s}}{\oxframe_{2s}} \oxintersect \oxnewframe{\oxvarctx_{3s}}{\oxframe_{3s}} = \oxnewframe{\oxvarctx_o}{\oxframe_o} \\
      }{
        %% conclusions
        \oxtypjudge{\oxglobalctx}{\oxtvarctx}{\oxkontctx}{\oxnewframe{\oxvarctx}{\oxframe}}{
          \oxbranch{\oxexpr_1}{\oxexpr_2}{\oxexpr_3}
        }{\oxsitype}{\oxnewframe{\oxvarctx_o}{\oxframe_o}} }
    \end{mathpar}
  } \vspace{1em}

  By our induction hypothesis we get that
  $\oxtypjudge{\oxglobalctx}{\oxtvarctx}{\oxkontctx}{\oxnewframe{\oxvarctx^\prime}{\oxframe}}{\oxexpr_1}{\oxtbool}{\oxnewframe{\oxvarctx_1^\prime}{\oxframe_1}}$,
  and
  $\oxtypjudge{\oxglobalctx}{\oxtvarctx}{\oxkontctx}{\oxnewframe{\oxvarctx_1^\prime}{\oxframe_1}}{\oxexpr_2}{\oxtbool}{\oxnewframe{\oxvarctx_2^\prime}{\oxframe_2}}$,
  and
  $\oxtypjudge{\oxglobalctx}{\oxtvarctx}{\oxkontctx}{\oxnewframe{\oxvarctx_1^\prime}{\oxframe_1}}{\oxexpr_3}{\oxtbool}{\oxnewframe{\oxvarctx_3^\prime}{\oxframe_3}}$
  with
  $\oxsubtypectx{\oxglobalctx}{\oxtvarctx}{\oxnewframe{\oxvarctx_2}{\oxframe_2}}{\oxnewframe{\oxvarctx_2^\prime}{\oxframe_2}}$
  and
  $\oxsubtypectx{\oxglobalctx}{\oxtvarctx}{\oxnewframe{\oxvarctx_3}{\oxframe_3}}{\oxnewframe{\oxvarctx_3^\prime}{\oxframe_3}}$.

  Next we want to show that
  $\oxtunify{\oxtvarctx}{\oxkontctx}{\oxnewframe{\oxvarctx_2^\prime}{\oxframe_2}}{\oxsitype_2}{\oxsitype}{\oxnewframe{\oxvarctx_{2s}^\prime}{\oxframe_{2s}}}$,
  $\oxsubtypectx{\oxglobalctx}{\oxtvarctx}{\oxnewframe{\oxvarctx_{2s}}{\oxframe_{2s}}}{\oxnewframe{\oxvarctx_{2s}^\prime}{\oxframe_{2s}}}$,
  $\oxtunify{\oxtvarctx}{\oxkontctx}{\oxnewframe{\oxvarctx_3^\prime}{\oxframe_3}}{\oxsitype_3}{\oxsitype}{\oxnewframe{\oxvarctx_{3s}^\prime}{\oxframe_{3s}}}$,
  and
  $\oxsubtypectx{\oxglobalctx}{\oxtvarctx}{\oxnewframe{\oxvarctx_{3s}}{\oxframe_{3s}}}{\oxnewframe{\oxvarctx_{3s}^\prime}{\oxframe_{3s}}}$,
  which all follow from applying \Lemma{lemma:unified-types-well-formed}. To do
  this lemma application, we just need to show that for all $\oxcprov$ in
  $\oxsitype_1$, $\oxsitype_2$ and $\oxsitype$, $\oxvarctx(\oxcprov) =
  \oxvarctx^\prime(\oxcprov)$ which follows from the premise.

  Finally, we need to show that
  $\oxsubtypectx{\oxglobalctx}{\oxtvarctx}{\oxnewframe{\oxvarctx_o}{\oxframe_o}}{\oxnewframe{\oxvarctx_o^\prime}{\oxframe_o}}$.
  The well formedness condition on $\oxvarctx_o^\prime$ follows immediately
  since all types are the same as in $\oxvarctx_2^\prime$ and
  $\oxvarctx_3^\prime$ and the loan sets are just unioned, meaning reference
  types remain valid and we can compute types for all loans.

  The equal or empty condition follows from the fact that $\oxvarctx_2^\prime$
  and $\oxvarctx_3^\prime$ both agree on types by
  \Lemma{lemma:related-input-output-differences}, which means they drop exactly
  the same entries. For any regions emptied, either the same regions are
  emptied, or the region was emptied in the corresponding smaller context
  $\oxvarctx_2$ or $\oxvarctx_3$. Otherwise the loan sets are untouched.

  Finally, both of these are preserved when adding on the same frame, so we're
  done.

  \vspace{1em}
  \noindent \framebox[\textwidth]{ \figuresize
    \begin{mathpar}
      \inferrule[T-Match]{
        %% premises
        \oxtypjudge{\oxglobalctx}{\oxtvarctx}{\oxkontctx}{\oxnewframe{\oxvarctx}{\oxframe}}{
          \oxexpr
        }{\oxtsum{\oxsitype_l}{\oxsitype_r}}{\oxnewframe{\oxvarctx_1}{\oxframe_1}} \\
        \forall \oxcprov \in \oxfprovs{\oxtsum{\oxsitype_l}{\oxsitype_r}}. \;
        \oxnotreborrowed{\oxnewframe{\oxvarctx_1}{\oxframe_1}}{\oxcprov} \\\\
        \oxtypjudge{\oxglobalctx}{\oxtvarctx}{\oxkontctx}
        {\oxextendctx{\oxnewframe{\oxvarctx_1}{\oxframe_1}}{
            \oxvarctxentry{\oxid_1}{\oxsitype_l} }}{ \oxexpr_1
        }{\oxsitype_2}{\oxextendctx{\oxnewframe{\oxvarctx_2}{\oxframe_2}}{
            \oxvarctxentry{\oxid_1}{\oxsdtype_l}
          }} \\\\
        \oxtypjudge{\oxglobalctx}{\oxtvarctx}{\oxkontctx}{\oxextendctx{\oxnewframe{\oxvarctx_1}{\oxframe_1}}{
            \oxvarctxentry{\oxid_2}{\oxsitype_r} }}{ \oxexpr_2
        }{\oxsitype_3}{\oxextendctx{\oxnewframe{\oxvarctx_3}{\oxframe_3}}{
            \oxvarctxentry{\oxid_2}{\oxsdtype_r}
          }} \\
        \oxsitype = \oxsitype_2 \vee \oxsitype = \oxsitype_3 \\\\
        \oxtunify{\oxtvarctx}{\oxkontctx}{\oxnewframe{\oxvarctx_2}{\oxframe_2}}
        {\oxsitype_2}{\oxsitype}{\oxnewframe{\oxvarctx_{2s}}{\oxframe_{2s}}} \\
        \oxtunify{\oxtvarctx}{\oxkontctx}{\oxnewframe{\oxvarctx_3}{\oxframe_3}}
        {\oxsitype_3}{\oxsitype}{\oxnewframe{\oxvarctx_{3s}}{\oxframe_{3s}}} \\
        \oxnewframe{\oxvarctx_{2s}}{\oxframe_{2s}} \oxintersect \oxnewframe{\oxvarctx_{3s}}{\oxframe_{3s}} = \oxnewframe{\oxvarctx_o}{\oxframe_o} \\
      }{
        %% conclusions
        \oxtypjudge{\oxglobalctx}{\oxtvarctx}{\oxkontctx}{\oxnewframe{\oxvarctx}{\oxframe}}{
          \oxmatch{\oxexpr}{\oxid_1}{\oxexpr_1}{\oxid_2}{\oxexpr_2}
        }{\oxsitype}{\oxnewframe{\oxvarctx_o}{\oxframe_o}} }
    \end{mathpar}
  } \vspace{1em}

  This case follows almost identically to the \oxname{T-Branch} case above. The
  only structural difference is that the expression typing judgements for
  $\oxexpr_1$ and $\oxexpr_2$ have $\oxid_1$ and $\oxid_2$ respectively in their
  environments, but we know we can remove $\oxid_1$ and $\oxid_2$ from each side
  and keep the contexts well formed, since nothing that comes before $\oxid_1$
  or $\oxid_2$ can refer to it, and we know that $\oxsitype_1$ and $\oxsitype_2$
  cannot in any way refer to $\oxid_1$ or $\oxid_2$ because we have from the
  region rewriting judgements in the premises that the types are well formed in
  $\oxnewframe{\oxvarctx_2}{\oxframe_2}$ and
  $\oxnewframe{\oxvarctx_3}{\oxframe_3}$ respectively. We also need to show that the region not reborrowed judgement still holds, but this is immediate because at most $\oxvarctx_1^\prime$ has the same types as $\oxvarctx_1$. 
  
  \vspace{1em}
  \noindent \framebox[\textwidth]{ \figuresize
    \begin{mathpar}
      \inferrule[T-Let]{
        %% premises
        \oxtypjudge{\oxglobalctx}{\oxtvarctx}{\oxkontctx}{\oxnewframe{\oxvarctx}{\oxframe}}{
          \oxexpr_1
        }{\oxsitype_1}{\oxnewframe{\oxvarctx_1}{\oxframe_1}} \\
        \oxtunify{\oxtvarctx}{\oxkontctx}{\oxnewframe{\oxvarctx_1}{\oxframe_1}}
                 {\oxsitype_1}{\oxsitype_a}{\oxnewframe{\oxvarctx_{1s}}{\oxframe_{1s}}} \\\\
        \forall \oxcprov \in \oxfprovs{\oxsitype_a}. \;
        \oxnotreborrowed{\oxnewframe{\oxvarctx_{1s}}{\oxframe_{1s}}}{\oxcprov} \\\\
        \oxtypjudge{\oxglobalctx}{\oxtvarctx}{\oxkontctx}{
          \oxgcloans{\oxkontctx}{\oxextendctx{\oxnewframe{\oxvarctx_{1s}}{\oxframe_{1s}}}{\oxvarctxentry{\oxid}{\oxsitype_a}}}
        }{ \oxexpr_2 }{\oxsitype_2}{
          \oxextendctx{\oxnewframe{\oxvarctx_2}{\oxframe_2}}{
            \oxvarctxentry{\oxid}{\oxsdtype} }
        } \\
      }{
        %% conclusions
        \oxtypjudge{\oxglobalctx}{\oxtvarctx}{\oxkontctx}{\oxnewframe{\oxvarctx}{\oxframe}}{
          \oxlet{\oxid}{\oxsitype_a}{\oxexpr_1}{\oxexpr_2} }{\oxsitype_2}{
          \oxnewframe{\oxvarctx_2}{\oxframe_2} } }
    \end{mathpar}
  } \vspace{1em}

  Firstly, we apply our induction hypothesis to get that $\oxexpr_1$ is well
  typed with input environment $\oxnewframe{\oxvarctx^\prime}{\oxframe}$ and
  output environment $\oxnewframe{\oxvarctx_1^\prime}{\oxframe_1}$ with
  $\oxsubtypectx{\oxglobalctx}{\oxtvarctx}{\oxnewframe{\oxvarctx_1}{\oxframe_1}}{\oxnewframe{\oxvarctx_1^\prime}{\oxframe_1}}$.
  Then, we apply \Lemma{lemma:unified-types-well-formed} to get
  $\oxsubtypectx{\oxglobalctx}{\oxtvarctx}{\oxnewframe{\oxvarctx_{1s}}{\oxframe_{1s}}}{\oxnewframe{\oxvarctx_{1s}^\prime}{\oxframe_{1s}}}$.
  In order to apply this lemma we need to know that for any $\oxcprov$ that
  occur in $\oxsitype_1$ or $\oxsitype_a$,
  $\oxnewframe{\oxvarctx_1}{\oxframe_1}(\oxcprov) =
  \oxnewframe{\oxvarctx_1^\prime}{\oxframe_1}(\oxcprov)$, which we have as a
  conclusion from the previous application of the induction hypothesis.

  The region not reborrowed judgement holds immediately, because at most
  $\oxvarctx_1^\prime$ has the same types as $\oxvarctx_1$.

  To apply our induction hypothesis on $\oxexpr_2$ and continue the case, we need that \\
  $\oxsubtypectx{\oxglobalctx}{\oxtvarctx}{\oxgcloans{\oxkontctx}{\oxnewframe{\oxvarctx_{1s}}{\oxframe_{1s}},
      \oxid :
      \oxsitype_a}}{\oxgcloans{\oxkontctx}{\oxnewframe{\oxvarctx_{1s}^\prime}{\oxframe_{1s}},
      \oxid : \oxsitype_a}}$. But this is immediate by definition since gcloans
  can only empty loan sets for regions for which there are no types that contain
  them, which is allowed by \oxname{R-Env}.

  Our final obligation to apply the induction hypothesis is that for any
  $\oxcprov$ that occurs in a type in $\oxframe_{1s}$ but is not in
  $\oxframe_{1s}$, we need that
  $\oxgcloans{\oxkontctx}{\oxnewframe{\oxvarctx_{1s}}{\oxframe_{1s}}}(\oxcprov)
  =
  \oxgcloans{\oxkontctx}{\oxnewframe{\oxvarctx_{1s}^\prime}{\oxframe_{1s}}}(\oxcprov)$.
  We already have that $\oxnewframe{\oxvarctx_{1s}}{\oxframe_{1s}}(\oxcprov) =
  \oxnewframe{\oxvarctx_{1s}^\prime}{\oxframe_{1s}}(\oxcprov)$, so we just need
  to know that $\exists \oxplace : \oxtype \in \oxvarctx_{1s}$, where $\oxcprov$
  occurs in $\oxtype$, and $\oxnewframe{\oxvarctx_{1s}}{\oxframe_{1s}}(\oxplace)
  = \oxnewframe{\oxvarctx_{1s}}{\oxframe_{1s}}^\prime(\oxplace)$. But we said
  that $\oxcprov$ is contained in a type in $\oxframe_{1s}$, so the place for
  that type is one such place, so we cannot empty the loan set.

  \vspace{1em}
  \noindent \framebox[\textwidth]{ \figuresize
    \begin{mathpar}
      \inferrule[T-Seq]{
        %% premises
        \oxtypjudge{\oxglobalctx}{\oxtvarctx}{\oxkontctx}{\oxnewframe{\oxvarctx}{\oxframe}}{
          \oxexpr_1
        }{\oxsitype_1}{\oxnewframe{\oxvarctx}{\oxframe}_1} \\\\
        \oxtypjudge{\oxglobalctx}{\oxtvarctx}{\oxkontctx}{\oxgcloans{\oxkontctx}{\oxnewframe{\oxvarctx_1}{\oxframe_1}}}{
          \oxexpr_2
        }{\oxsitype_2}{\oxnewframe{\oxvarctx_2}{\oxframe_2}} \\
      }{
        %% conclusions
        \oxtypjudge{\oxglobalctx}{\oxtvarctx}{\oxkontctx}{\oxnewframe{\oxvarctx}{\oxframe}}{
          \oxseq{\oxexpr_1}{\oxexpr_2}
        }{\oxsitype_2}{\oxnewframe{\oxvarctx_2}{\oxframe_2}} }
    \end{mathpar}
  } \vspace{1em}

  Firstly, we apply our induction hypothesis to get that $\oxexpr_1$ is well
  typed with input environment $\oxnewframe{\oxvarctx^\prime}{\oxframe}$ and
  output environment $\oxnewframe{\oxvarctx_1^\prime}{\oxframe_1}$, with
  $\oxsubtypectx{\oxglobalctx}{\oxtvarctx}{\oxnewframe{\oxvarctx_1}{\oxframe_1}}{\oxnewframe{\oxvarctx_1^\prime}{\oxframe_1}}$.
  We need to know that
  $\oxsubtypectx{\oxglobalctx}{\oxtvarctx}{\oxgcloans{\oxkontctx}{\oxnewframe{\oxvarctx_1}{\oxframe_1}}}{\oxgcloans{\oxkontctx}{\oxnewframe{\oxvarctx_1^\prime}{\oxframe_1}}}$
  before we can apply our induction hypothesis to finish the proof. But this
  fact is trivial by the definitions, since \textrm{gc-loans} can only empty
  regions that are not in initialized types in the context, which is allowed in
  \oxname{R-Env}.

  Our final obligation to apply the induction hypothesis is that for any
  $\oxcprov$ that occurs in a type in $\oxframe_{1}$ but is not in
  $\oxframe_{1}$, we need that
  $\oxgcloans{\oxkontctx}{\oxnewframe{\oxvarctx_{1}}{\oxframe_{1}}}(\oxcprov) =
  \oxgcloans{\oxkontctx}{\oxnewframe{\oxvarctx_{1}^\prime}{\oxframe_{1}}}(\oxcprov)$.
  We already have that $\oxnewframe{\oxvarctx_{1}}{\oxframe_{1}}(\oxcprov) =
  \oxnewframe{\oxvarctx_{1}^\prime}{\oxframe_{1}}(\oxcprov)$, so we just need to
  know that $\exists \oxplace : \oxtype \in \oxvarctx_{1}$, where $\oxcprov$
  occurs in $\oxtype$, and $\oxnewframe{\oxvarctx_{1}}{\oxframe_1}(\oxplace) =
  \oxnewframe{\oxvarctx_{1}^\prime}{\oxframe_1}(\oxplace)$. But since $\oxcprov$
  occurs in a type in $\oxframe_1$, the place that maps to that type is such a
  place.
  
  \vspace{1em}
  \noindent \framebox[\textwidth]{ \figuresize
    \begin{mathpar}
      \inferrule[T-Drop]{
        %% premises
        \oxlookup{\oxvarctx}{\oxplace}{\oxsitype_\oxplace} \\
        \oxtypjudge{\oxglobalctx}{\oxtvarctx}{\oxkontctx}{
          (\oxnewframe{\oxvarctx}{\oxframe})[\oxplace \mapsto
          \oxsidtype_\oxplace]
        }{\oxexpr}{\oxsxtype}{\oxnewframe{\oxvarctx_o}{\oxframe_o}} \\
      }{
        %% conclusions
        \oxtypjudge{\oxglobalctx}{\oxtvarctx}{\oxkontctx}{\oxnewframe{\oxvarctx}{\oxframe}}{
          \oxexpr }{\oxsxtype}{\oxnewframe{\oxvarctx_o}{\oxframe_o}} }
    \end{mathpar}
  } \vspace{1em}

  In order to apply our induction hypothesis and finish the case, we only need
  to show that
  $\oxsubtypectx{\oxglobalctx}{\oxtvarctx}{(\oxnewframe{\oxvarctx}{\oxframe})[\oxplace
    \mapsto
    \oxsidtype_\oxplace]}{(\oxnewframe{\oxvarctx^\prime}{\oxframe})[\oxplace
    \mapsto \oxsidtype_\oxplace]}$, which is immediate by the definition of
  related contexts. Note that
  $(\oxnewframe{\oxvarctx^\prime}{\oxframe})[\oxplace \mapsto
  \oxsidtype_\oxplace]$ is well formed because
  $(\oxnewframe{\oxvarctx}{\oxframe})[\oxplace \mapsto \oxsidtype_\oxplace]$ is
  well formed. There cannot be any loans to $\oxplace$ because in the
  $(\oxnewframe{\oxvarctx^\prime}{\oxframe})[\oxplace \mapsto
  \oxsidtype_\oxplace]$ because those loans would be there in
  $(\oxnewframe{\oxvarctx}{\oxframe})[\oxplace \mapsto \oxsidtype_\oxplace]$.

  \vspace{1em}
  \noindent \framebox[\textwidth]{ \figuresize
    \begin{mathpar}
      \inferrule[T-App]{
        %% premises
        \overline{\oxenvwellformed{\oxglobalctx}{\oxtvarctx}{\oxnewframe{\oxvarctx}{\oxframe}}{\oxenv}} \\
        \overline{\oxrgnvalidity{\oxtvarctx}{\oxnewframe{\oxvarctx}{\oxframe}}{\oxprov}} \\
        \overline{\oxtypevalidity{\oxglobalctx}{\oxtvarctx}{\oxnewframe{\oxvarctx}{\oxframe}}{\oxsitype}} \\
        \delta = \oxsubstmany{
          \oxsubstmany{
            \oxsubstmany{\cdot}{\oxenv}{\oxenvvar}
          }{\oxprov}{\oxabsprov}
        }{\oxsitype}{\oxtvar} \\\\
        \oxtypjudge{\oxglobalctx}{\oxtvarctx}{\oxkontctx}{\oxnewframe{\oxvarctx}{\oxframe}}{
          \oxnoseqexpr_f }{ \oxtfunext{ \overline{\oxenvvar} \oxcomma
            \overline{\oxabsprov} \oxcomma \overline{\oxtvar} }{ \oxsitype_1
            \oxdotsc \oxsitype_n
          }{\oxsitype_f}{\oxenv_c}{\overline{\oxabsprov_1: \oxabsprov_2}}
        }{\oxnewframe{\oxvarctx_0}{\oxframe_0}} \\\\
        \forall i \in \oxset{1 \oxdots n}. \;
        \oxtypjudge{\oxglobalctx}{\oxtvarctx}{
          \oxextendctx{\oxkontctx}{
            \delta(\oxsitype_1) \oxdots \delta(\oxsitype_{i-1})
          }
        }{\oxnewframe{\oxvarctx_{i-1}}{\oxframe_{i-1}}}{
          \oxnoseqexpr_i }{\delta(\oxsitype_i)}{\oxnewframe{\oxvarctx_i}{\oxframe_i}} \\
        \forall i \in \oxset{ 1 \oxdots n }. \;
    \forall \oxcprov \in \oxfprovs{\oxsitype_i}. \;
    \oxnotreborrowed{\oxnewframe{\oxvarctx_n}{\oxframe_n}}{\oxcprov} \\
        \oxrunifymany{\oxtvarctx}{\oxkontctx}{\oxnewframe{\oxvarctx_n}{\oxframe_n}}{
          \oxsubstmany{\oxabsprov_2}{\oxprov}{\oxabsprov} }{
          \oxsubstmany{\oxabsprov_1}{\oxprov}{\oxabsprov}
        }{\oxnewframe{\oxvarctx_b}{\oxframe_b}} \\
      }{
        %% conclusions
        \oxtypjudge{\oxglobalctx}{\oxtvarctx}{\oxkontctx}{\oxnewframe{\oxvarctx}{\oxframe}}{
          \oxapp{\oxnoseqexpr_f}{ \overline{\oxenv} \oxcomma \overline{\oxprov}
            \oxcomma \overline{\oxsitype} }{ \oxnoseqexpr_1 \oxdotsc
            \oxnoseqexpr_n } }{ \oxsubstmany{ \oxsubstmany{
              \oxsubstmany{\oxsitype_f}{\oxenv}{\oxenvvar}
            }{\oxprov}{\oxabsprov} }{\oxsitype}{\oxtvar}
        }{\oxnewframe{\oxvarctx_b}{\oxframe_b}} }
    \end{mathpar}
  } \vspace{1em}

  In the case of \oxname{T-App}, we firstly must prove the well formedness
  properties:
  \begin{itemize}
  \item
    $\oxenvwellformed{\oxglobalctx}{\oxtvarctx}{\oxnewframe{\oxvarctx^\prime}{\oxframe}}{\oxenv}$.
    Since $\oxtvarctx$ is unchanged, \oxname{WF-Env} is the only interesting
    case.
    
    \vspace{1em}
    \noindent \hspace{-2.5em} \framebox[\textwidth]{ \figuresize
      \begin{mathpar}
        \inferrule[WF-Env]{
          %% premises
          \oxvarctxwellformed{\oxglobalctx}{\oxtvarctx}{
            \oxnewframe{\oxnewframe{\oxvarctx}{\oxframe}}{\oxframe_c} } }{
          %% conclusions
          \oxenvwellformed{\oxglobalctx}{\oxtvarctx}{\oxnewframe{\oxvarctx}{\oxframe}}{
            \oxframe_c } }
      \end{mathpar}
    } \vspace{1em}

    Let $\oxenv = \oxframe_e$. We want to show that
    $\oxctxswellformed{\oxglobalctx}{\oxtvarctx}{\oxnewframe{\oxnewframe{\oxvarctx^\prime}{\oxframe}}{\oxframe_c}}{\oxkontctx}$
    given
    $\oxctxswellformed{\oxglobalctx}{\oxtvarctx}{\oxnewframe{\oxnewframe{\oxvarctx}{\oxframe}}{\oxframe_c}}{\oxkontctx}$,
    which is immediate from \Lemma{lemma:related-contexts-well-formed}.
  \item $\oxrgnvalidity{\oxtvarctx}
    {\oxnewframe{\oxvarctx^\prime}{\oxframe}}{\oxprov}$, which is immediate from
    the premises since related loan environments have the same domains and
    $\oxtvarctx$ is the same.
  \item $\oxtypevalidity{\oxglobalctx}{\oxtvarctx}
    {\oxnewframe{\oxvarctx^\prime}{\oxframe}}{\oxsitype}$, which is immediate
    from \Lemma{lemma:types-well-formed}. We just need that for the regions that
    occur in the type, their loan sets are unchanged, but we get that from the
    premise, because the function argument is either: locally defined, in which
    case it can only use and produce types accessible in the context; an
    argument, in which case its arguments are also part of the argument type; or
    a global function, in which case these types do not contain any non abstract
    regions which are replaced with concrete regions all in $\oxframe$.
  \end{itemize}

  For the rest of the application case, we can apply our induction hypothesis on
  the function and the arguments, additionally applying the substitution lemma,
  \Lemma{lemma:substitution}, where needed. The region not reborrowed condition
  is true by the fact that at most the types between $\oxvarctx_n$ and
  $\oxvarctx_n^\prime$ are the same. The last part about outlives follows from
  \Lemma{lemma:outlives-related}, where we have the condition on the loan sets
  from the conclusion of the application of the induction hypothesis.

  \vspace{1em}
  \noindent \framebox[\textwidth]{ \figuresize
    \begin{mathpar}
      \inferrule[T-AppClosure]{
        %% premises
        \oxtypjudge{\oxglobalctx}{\oxtvarctx}{\oxkontctx}{\oxnewframe{\oxvarctx}{\oxframe}}{
          \oxexpr_f
        }{
          \oxtfun{}{
            \oxsitype_1 \oxdotsc \oxsitype_n
          }{\oxsitype_f}{\oxenv_c}
        }{\oxnewframe{\oxvarctx_{0s}}{\oxframe_{0s}}} \\\\
        \forall i \in \oxset{1 \oxdots n}. \;
        \oxtypjudge{\oxglobalctx}{\oxtvarctx}{
          \oxextendctx{\oxkontctx}{
            \oxsitype_1 \oxdots \oxsitype_{i-1}
          }
        }{\oxnewframe{\oxvarctx_{i-1s}}{\oxframe_{i-1s}}}{
          \oxexpr_i
        }{
          \oxsitype_{i^\prime}
        }{\oxnewframe{\oxvarctx_i}{\oxframe_i}} \\
        \oxtunify[\oxcombineunrest]{\oxtvarctx}{\oxkontctx}{\oxnewframe{\oxvarctx_i}{\oxframe_i}}
                 {\oxsitype_{i^\prime}}{\oxsitype_i}{\oxnewframe{\oxvarctx_{is}}{\oxframe_{is}}} \\\\
                 \forall i \in \oxset{ 1 \oxdots n }. \;
                 \forall \oxcprov \in \oxfprovs{\oxsitype_i}. \;
                 \oxnotreborrowed{\oxnewframe{\oxvarctx_{ns}}{\oxframe_{ns}}}{\oxcprov} \\
      }{
        %% conclusions
        \oxtypjudge{\oxglobalctx}{\oxtvarctx}{\oxkontctx}{\oxnewframe{\oxvarctx}{\oxframe}}{
          \oxapply{\oxexpr_f}{
            \oxexpr_1 \oxdotsc \oxexpr_n
          }
        }{
          \oxsitype_f
        }{\oxnewframe{\oxvarctx_{ns}}{\oxframe_{ns}}}
      }

    \end{mathpar}
  }

  In the case of \oxname{T-AppClosure}, we follow a very similar procedure to
  \oxname{T-AppFunction}, but with an empty substituion, and the addition that
  we need to apply \Lemma{lemma:unified-types-well-formed} to handle the
  rewriting.

  In the cases of \oxname{T-Move}, \oxname{T-Copy}, \oxname{T-Borrow},
  \oxname{T-BorrowIndex}, \oxname{T-BorrowSlice}, \oxname{IndexCopy}, they all
  follow from the induction hypothesis and additionally applying
  \Lemma{lemma:ownership-safety-related-envs} and
  \Lemma{lemma:tc-related-contexts}. Note we get the place having the right type
  requirement for \oxname{T-Move} from the fact that the place must be in
  $\oxframe$ since it is a free variable.

  The remaining cases of \oxname{T-Assign} and \oxname{T-AssignDeref} proceed
  similarly. Firstly, we apply the induction hypothesis on the expression, then
  \Lemma{lemma:ownership-safety-related-envs} and
  \Lemma{lemma:tc-related-contexts}, and finally we get well formedness and
  relatedness on the output environment by applying
  \Lemma{lemma:unified-types-well-formed}. Note we get the place having the same
  type requirement for type computation from the fact that the place must be in
  $\oxframe$ since it is a free variable.

\end{proof}

\begin{oxlemma}{Referent Well Formedness Preserved in Related Environments}
  {lemma:referent-related-envs} If
  $\oxsubtypectx{\oxglobalctx}{\oxtvarctx}{\oxvarctx}{\oxvarctx^\prime}$ and
  $\oxreferentvalidity{\oxglobalctx}{\oxvarctx}{\oxreferentctx[\oxplace]}{\oxxitype}$
  and $\oxvarctx(\oxplace) = \oxvarctx^\prime(\oxplace)$, then
  $\oxreferentvalidity{\oxglobalctx}{\oxvarctx^\prime}{\oxreferentctx[\oxplace]}{\oxxitype}$.
\end{oxlemma}
\begin{proof}
  Proceed by induction on the referent validity derivation. The only case that
  doesn't follow immediately from premises and the induction hypothesis in
  \oxname{WF-RefId}, which follows from the equal types premise.
\end{proof}

\begin{oxlemma}{Value Typing Preserved in Related Environments}
  {lemma:value-typing-related-envs} If
  $\oxsubtypectx{\oxglobalctx}{\oxemptyctx}{\oxvarctx} {\oxvarctx^\prime}$,
  and $\oxtypjudge{\oxglobalctx}{\oxemptyctx}{\oxemptyctx}{\oxvarctx}{\oxvalue}
  {\oxvarctx(x)}{\oxvarctx}$, then
  $\oxtypjudge{\oxglobalctx}{\oxemptyctx}{\oxemptyctx}{\oxvarctx^\prime}
  {\oxvalue}{\oxvarctx^\prime(x)}{\oxvarctx^\prime}$.
\end{oxlemma}

\begin{proof}
  Proceed by simultaneous induction on the typing derivation and the stack frame
  well formedness.
  Since we know the expression is already a value, we restrict ourselves
  only to those cases that type values: \oxname{T-Unit}, \oxname{T-u32},
  \oxname{T-True}, \oxname{T-False}, \oxname{T-Tuple}, \oxname{T-Array},
  \oxname{T-Left}, \oxname{T-Right}, \oxname{T-Dead}, \oxname{T-Pointer}, and
  \oxname{T-ClosureValue}.

  For \oxname{T-Unit}, \oxname{T-u32}, \oxname{T-Dead}, \oxname{T-True}, and
  \oxname{T-False}, this holds trivially. For \oxname{T-Tuple},
  \oxname{T-Array}, \oxname{T-Left}, and \oxname{T-Right}, this holds directly
  by repeated application of our induction hypothesis. This leaves us with
  four cases.

  \vspace{1em}
  \noindent \framebox[\textwidth]{ \figuresize
    \begin{mathpar}
      \TClosureVal
    \end{mathpar}
  } \vspace{1em}

  For the \oxname{T-ClosureValue} case, firstly we want to show
  $\oxsubstorevalidity{\oxglobalctx}{\oxvarctx^\prime}{\oxstackframe}{\oxframe_c}$.
  This follows immediatedly from applying the induction hypothesis for each value.

  Then to finish the closure case, it suffices to show\\
  $\oxtypjudge{\oxglobalctx}{\oxemptyctx}{\oxemptyctx}{\oxnewframe{\oxvarctx^\prime}{\oxframe_c}}
  {\oxexpr}{\oxsitype_r}{\oxnewframe{\oxvarctx_o^\prime}{\oxframe}}$, which
  follows immediately from \Lemma{lemma:expression-typing-related-envs}.
  

  \vspace{1em}
  \noindent \framebox[\textwidth]{ \figuresize
    \begin{mathpar}
      \TPointer
    \end{mathpar}
  } \vspace{1em}
  
  If $x$ was dropped, then $\oxvarctx^\prime(x) = \oxvarctx(x)^\dagger$. Then
  the proof follows immediately from \oxname{T-Dead}.
  
  If $x$ was not dropped, then $\oxvarctx(x) = \oxvarctx^\prime(x)$. All that
  is left to show is that that the referent remains well formed, and the loan
  $\oxloanpkg{\oxmuta}{\oxplace}$ is in $\oxvarctx^\prime(\oxcprov)$. The
  first condition follows from \Lemma{lemma:referent-related-envs}. The second
  condition is immediate because the only potential changes allowed in the
  related environment to loan sets is emptying the loan sets of regions if
  there's no references with the region in their type, and this particular
  reference is a reference with the region, so emptying the loan set is ruled
  out.
\end{proof}
  
\begin{oxlemma}{Value Typing Fixed on Output Environments}
  {lemma:value-typing-output} If
  $\oxtypjudge{\oxglobalctx}{\oxtvarctx}{\oxkontctx}{\oxvarctx}{\oxvalue}
  {\oxtype}{\oxvarctx^\prime}$, then
  $\oxtypjudge{\oxglobalctx}{\oxtvarctx}{\oxkontctx}{\oxvarctx^\prime}
  {\oxvalue}{\oxtype}{\oxvarctx^\prime}$.
\end{oxlemma}

\begin{proof}
  Immediate by induction on the typing derivation. The only non immediate case
  is \oxname{T-Pointer}, where we also need to apply
  \Lemma{lemma:referent-related-envs}.
\end{proof}

\begin{oxlemma}{Stack Validity is Preserved in Related Environments}
  {lemma:stack-validity-related-envs} If
  $\oxstorevalidity{\oxglobalctx}{\oxvarctx}{\oxstore}$ and
  $\oxsubtypectx{\oxglobalctx}{\oxemptyctx}{\oxvarctx}{\oxvarctx^\prime}$, then
  $\oxstorevalidity{\oxglobalctx}{\oxvarctx^\prime}{\oxstore}$.
\end{oxlemma}

\begin{proof}
  We proceed by induction over the well typedness of the store.

  \vspace{1em}
  \noindent \framebox[\textwidth]{ \figuresize
    \begin{mathpar}
      \VStack\and \VStackEmpty
    \end{mathpar}
  } \vspace{1em}

  The interesting case is when the stack is non empty. Then we have that
  $\oxstorevalidity{\oxglobalctx}{\oxvarctx^\prime}{\oxstore}$ and want to show
  that
  $\oxstorevalidity{\oxglobalctx}{\oxnewframe{\oxvarctx^\prime}{\oxframe}}{\oxnewframe{\oxstore}{\oxstackframe}}$.
  The requirement on the domain is immediate since related environments have the
  same domains. What's left to show is that the values in the store remain well
  typed under the new environment. This follows from repeated applications of
  \Lemma{lemma:value-typing-related-envs}
\end{proof}

\subsection{Preservation When Popping a Stack Frame Lemmas}
\label{sec:pop-lemmas}

\begin{oxlemma}{Stack Validity is Preserved When Popping A Stack Frame}
  {lemma:stack-envminus} If
  $\oxstorevalidity{\oxglobalctx}{\oxnewframe{\oxvarctx}{\oxframe}}{\oxnewframe{\oxstore}{\oxstackframe}}$,
  then $\oxstorevalidity{\oxglobalctx} {\oxvarctx}{\oxstore}$.
\end{oxlemma}

\begin{proof}
  Immediate by inversion on \oxname{WF-StackFrame} which gives us
  $\oxstorevalidity{\oxglobalctx}{\oxvarctx}{\oxstore}$.
\end{proof}

\subsection{Preservation under Well Typed Extension Lemmas}
\label{sec:extension-lemmas}

\begin{oxlemma}{Ownership Safety is Preserved under Well-Typed Extensions}
  {lemma:ownership-well-typed-extension} If
  \begin{enumerate}
  \item
    $\oxtypevalidity{\oxglobalctx}{\oxemptyctx}{\oxnewframe{\oxvarctx}{\oxframe}}{\oxsitype_\oxid}$
  \item
    $\oxmusafetyinner{\oxemptyctx}{\oxkontctx}{\oxnewframe{\oxvarctx}{\oxframe}}{\oxmuta}{\overline{\oxplace}}{\oxplaceexpr}{\oxset{\oxloans}}$
  \item $\oxrootof{\oxplaceexpr} \in \oxdomain{\oxframe}$
  \item $\forall \oxcprov \in \oxfprovs{\oxsitype_\oxid}. \;
    \oxnotreborrowed{\oxvarctx}{\oxcprov}$
  \item $\forall \oxcprov \in \oxfprovs{\oxsitype_\oxid}. \; \forall \oxplace_i
    \in \oxset{\overline{\oxplace}}. \; \oxvarctx(\oxplace_i) =
    \oxtref{\oxcprov^{\prime\prime}}{\oxmuta}{\oxxitype} \implies
    \oxcprov^{\prime\prime} \neq \oxcprov$
  \end{enumerate}
  then
  $\oxmusafetyinner{\oxemptyctx}{\oxkontctx}{\oxnewframe{\oxextendctx{\oxvarctx}{\oxvarctxentry{\oxid}{\oxsitype_\oxid}}}{\oxframe}}{\oxmuta}{\overline{\oxplace}}{\oxplaceexpr}{\oxset{\oxloans}}$.
\end{oxlemma}

\begin{proof}
  We proceed by induction on the ownership safety derivation.

  \vspace{1em}
  \noindent \framebox[\textwidth]{ \figuresize
    \begin{mathpar}
      \OSafePlace
    \end{mathpar}
  } \vspace{1em}

  We'd like to apply \oxname{O-SafePlace} to show
  $\oxmusafetyinner{\oxemptyctx}{\oxkontctx}{ \oxnewframe{
      \oxextendctx{\oxvarctx}{\oxvarctxentry{\oxid}{\oxsitype_\oxid}}
    }{\oxframe} }{\oxmuta}{ \overline{\oxplace_e} }{ \oxplace }{
    \oxset{\oxloanpkg{\oxmuta}{\oxplace}} }$. Let
  $\oxloanctxentry{\oxcprov^\prime}{\oxset{\oxloans}} \in \oxvarctx$. Note that
  necessarily $\oxloanctxentry{\oxcprov^\prime}{\oxset{\oxloans}} \in
  \oxextendctx{\oxvarctx}{\oxvarctxentry{\oxid}{\oxsitype_\oxid}}$.

  If $\forall \oxloanpkg{\oxmuta^\prime}{\oxplaceexprctx[\oxplace^\prime]} \in
  \oxset{\overline{\oxloan}}. (\oxmuta = \oxmut \vee \oxmuta^\prime = \oxmut)
  \implies \oxrelevant{\oxplace^\prime}{\oxplace}$ held in our original
  $\oxvarctx$, then it still holds in the extended typing $\oxnewframe{
    \oxextendctx{\oxvarctx}{\oxvarctxentry{\oxid}{\oxsitype_\oxid}} }{\oxframe}$
  since the loan sets are unchanged between the two stack typings.

  Otherwise, we must have used the second clause $\exists
  \oxvarctxentry{\oxplace^\prime}{\oxtref{\oxcprov^\prime}{\oxmuta^\prime}{\oxtype^\prime}}
  \in \oxexplode{\oxvarctx} \; \wedge \; \not \exists
  \oxtref{\oxcprov^\prime}{\oxmuta^\prime}{\oxtype^\prime} \in \oxkontctx \;
  \wedge (\forall
  \oxvarctxentry{\oxplace^\prime}{\oxtref{\oxcprov^\prime}{\oxmuta^\prime}{\oxtype^\prime}}
  \in \oxexplode{\oxvarctx}. \; \oxplace^\prime \in \oxset{
    \overline{\oxplace_e} })$ in the first place. To show that this is still
  true, we need to show that nothing in our newly-bound $\oxid$ shares a type
  with a reborrowed reference which would end up in our exclusion list. The
  reason for this is that the failure condition for this is that with such a
  reference now bound at (or reachable within) $\oxid$, we could violate the
  universally-quantified portion of this clause. Fortunately, we have from our
  premise that all the regions that appear in the type $\oxsitype_\oxid$ are
  distinct from the ones in the exclusion list ($\forall \oxcprov \in
  \oxfprovs{\oxsitype_\oxid}. \; \forall \oxplace_i \in
  \oxset{\overline{\oxplace}}. \; \oxvarctx(\oxplace_i) =
  \oxtref{\oxcprov^{\prime\prime}}{\oxmuta}{\oxxitype} \implies
  \oxcprov^{\prime\prime} \neq \oxcprov$). Thus, we know this cannot be the
  case.

  \vspace{1em}
  \noindent \framebox[\textwidth]{ \figuresize
    \begin{mathpar}
      \ODeref
    \end{mathpar}
  } \vspace{1em}

  We'd like to apply \oxname{O-Deref} to show $\oxmusafetyinner{\oxemptyctx}{\oxkontctx}{
    \oxnewframe{ \oxextendctx{\oxvarctx}{\oxvarctxentry{\oxid}{\oxsitype_\oxid}}
    }{\oxframe} }{\oxmuta}{ \overline{\oxplace_e} }{ \oxplace }{
    \oxset{\oxloanpkg{\oxmuta}{\oxplace}} }$. This requires us to show
  $\oxlookup{ (\oxnewframe{
      \oxextendctx{\oxvarctx}{\oxvarctxentry{\oxid}{\oxsitype_\oxid}}
    }{\oxframe}) }{\oxplace}{
    \oxtref{\oxcprov}{\oxmuta_\oxplace}{\oxtype_\oxplace} }$ and $\oxlookup{
    (\oxnewframe{
      \oxextendctx{\oxvarctx}{\oxvarctxentry{\oxid}{\oxsitype_\oxid}}
    }{\oxframe}) }{\oxcprov}{
    \oxset{\overline{\oxloanpkg{\oxmuta^\prime}{\oxplaceexpr}}^n} }$. The former
  follows from the disjointedness assumption for $\oxid$, i.e. that $\oxid$ is
  disjoint from all existing identifiers in $\oxvarctx$. The latter follows from
  the fact that no loan sets are changed between the two stack typings. Since
  the loan set is unchanged, we also have that the new extension for the
  exclusion list $\oxkey{excl}$ is the same. This leaves us with two pieces to
  show. First, that the recursive uses of ownership safety still succeed (for
  which we will use the induction hypothesis) and that our last obligation holds
  (which follows much as it did for \oxname{O-SafePlace}).

  For the inductive cases, we can very nearly just apply the induction
  hypothesis, but we first must show that our exclusion list invariant applies
  for the extensions to the exclusion list. That is, we have that $\forall
  \oxcprov \in \oxfprovs{\oxsitype_\oxid}. \; \forall \oxplace_i \in
  \oxset{\overline{\oxplace_e}}. \; \oxvarctx(\oxplace_i) =
  \oxtref{\oxcprov^{\prime\prime}}{\oxmuta}{\oxxitype} \implies
  \oxcprov^{\prime\prime} \neq \oxcprov$, and we need to show $\forall \oxcprov
  \in \oxfprovs{\oxsitype_\oxid}. \; \forall \oxplace_i \in \oxset{\textrm{excl}
    \oxcomma \oxplace}. \; \oxvarctx(\oxplace_i) =
  \oxtref{\oxcprov^{\prime\prime}}{\oxmuta}{\oxxitype} \implies
  \oxcprov^{\prime\prime} \neq \oxcprov$. The exclusion extension $\oxkey{excl}$
  is constructed by looking specifically at the reborrow loans associated with
  the region $\oxcprov$. Since we know that $\forall \oxcprov \in
  \oxfprovs{\oxsitype_\oxid}. \; \forall
  \oxvarctxentry{\oxplace}{\oxtref{\oxcprov}{\oxmuta}{\oxxitype}} \in
  \oxexplode{\oxvarctx}. \; \not\exists \oxcprov^\prime. \;
  \oxloanpkg{\oxmuta}{\oxderef{\! \oxplace}} \in \oxvarctx(\oxcprov^\prime)$
  (from inversion of \oxname{NRB-Region}), it
  follows directly that none of the places in $\oxkey{excl}$ can have a
  reference type with an region in $\oxsitype_\oxid$. If they did, that would
  mean syntactically that $\oxvarctx(\oxcprov)$ contains a loan for
  $\oxderef{\oxplace}$ which would give us a contradiction.

  For the last obligation, let
  $\oxloanctxentry{\oxcprov^\prime}{\oxset{\oxloans}} \in \oxvarctx$. Note that
  necessarily $\oxloanctxentry{\oxcprov^\prime}{\oxset{\oxloans}} \in
  \oxextendctx{\oxvarctx}{\oxvarctxentry{\oxid}{\oxsitype_\oxid}}$.

  If $\forall \oxloanpkg{\oxmuta^\prime}{\oxplaceexpr^{\prime\prime}} \in
  \oxset{\overline{\oxloan}}. (\oxmuta = \oxmut \vee \oxmuta^\prime = \oxmut)
  \implies \oxrelevant{\oxplaceexpr^{\prime\prime}}{
    \oxplaceexprctx[\oxderef{\oxplace}] }$ held in our original $\oxvarctx$,
  then it still holds in the extended typing $\oxnewframe{
    \oxextendctx{\oxvarctx}{\oxvarctxentry{\oxid}{\oxsitype_\oxid}} }{\oxframe}$
  since the loan sets are unchanged between the two stack typings.

  Otherwise, we must have used the second clause $\exists
  \oxvarctxentry{\oxplace^\prime}{\oxtref{\oxcprov^\prime}{\oxmuta^\prime}{\oxtype^\prime}}
  \in \oxexplode{\oxvarctx} \; \wedge \; \not \exists
  \oxtref{\oxcprov^\prime}{\oxmuta^\prime}{\oxtype^\prime} \in \oxkontctx \;
  \wedge (\forall
  \oxvarctxentry{\oxplace^\prime}{\oxtref{\oxcprov^\prime}{\oxmuta^\prime}{\oxtype^\prime}}
  \in \oxexplode{\oxvarctx}. \; \oxplace^\prime \in \oxset{
    \overline{\oxplace_e} \oxcomma \textrm{excl} \oxcomma \oxplace })$ in the
  first place. To show that this is still true, we need to show that nothing in
  our newly-bound $\oxid$ shares a type with a reborrowed reference which would
  end up in our exclusion list. The reason for this is that the failure
  condition for this is that with such a reference now bound at (or reachable
  within) $\oxid$, we could violate the universally-quantified portion of this
  clause. Fortunately, we have from our premise that all the regions that
  appear in the type $\oxsitype_\oxid$ are distinct from the ones in the
  exclusion list ($\forall \oxcprov \in \oxfprovs{\oxsitype_\oxid}. \; \forall
  \oxplace_i \in \oxset{\overline{\oxplace}}. \; \oxvarctx(\oxplace_i) =
  \oxtref{\oxcprov^{\prime\prime}}{\oxmuta}{\oxxitype} \implies
  \oxcprov^{\prime\prime} \neq \oxcprov$). Thus, we know this cannot be the
  case.

  \vspace{1em}
  \noindent \framebox[\textwidth]{ \figuresize
    \begin{mathpar}
      \ODerefAbs
    \end{mathpar}
  } \vspace{1em}

  Since $\oxtvarctx = \oxemptyctx$, there are no valid reference types that have
  an abstract region, meaning the first hypothesis is a contradiction.
\end{proof}

\begin{oxlemma}{Type Computation is Preserved under Well-Typed Extensions}
  {lemma:computety-well-typed-extension} If
  $\oxtypevalidity{\oxglobalctx}{\oxemptyctx}{\oxvarctx}{\oxsitype_\oxid}$ and
  $\oxcomputety{\oxemptyctx}{\oxvarctx}{\oxmuta}{\oxplaceexpr}{\oxtype}
  {\oxset{\overline{\oxprov}}}$ then $\oxcomputety{\oxemptyctx}{
    \oxextendctx{\oxvarctx}{\oxsitype_\oxid}
  }{\oxmuta}{\oxplaceexpr}{\oxtype}{\oxset{\overline{\oxprov}}}$.
\end{oxlemma}

\begin{proof}
  We proceed by induction on the type computation. This gives us three cases,
  \oxname{TC-Var}, \oxname{TC-Proj} and \oxname{TC-Deref}. In \oxname{TC-Var},
  the lookup yields the same type based on the assumption that our new binding
  is disjoint from our existing ones. \oxname{TC-Proj} and \oxname{TC-Deref}
  proceed directly from the induction hypothesis.
\end{proof}

\begin{oxlemma}{Outlives is Preserved under Well-Typed Extensions}
  {lemma:outlives-well-typed-extension} If
  $\oxtypevalidity{\oxglobalctx}{\oxemptyctx}{\oxvarctx}{\oxsitype_\oxid}$ and
  $\forall \oxcprov \in \oxfprovs{\oxsitype_\oxid}. \;
  \oxnotreborrowed{\oxvarctx}{\oxcprov}$ and
  $\oxrunify{\oxemptyctx}{\oxkontctx}{\oxvarctx}{\oxprov_1}{\oxprov_2}{\oxvarctx^\prime}$
  then $\oxrunify{\oxemptyctx}{\oxkontctx}{ \oxextendctx{\oxvarctx}{
      \oxvarctxentry{\oxid}{\oxsitype_\oxid} } }{\oxprov_1}{\oxprov_2}{
    \oxextendctx{\oxvarctx^\prime}{ \oxvarctxentry{\oxid}{\oxsitype_\oxid} } }$
\end{oxlemma}

\begin{proof}
  We proceed by induction on the outlives judgment. This gives us six cases,
  \oxname{OL-Refl}, \oxname{OL-Trans}, \oxname{OL-BothAbstract},
  \oxname{OL-CombineConcrete}, \oxname{OL-ConcreteAbstract} and
  \oxname{OL-AbstractConcrete}.

  \oxname{OL-Refl}, \oxname{OL-BothAbstract} and \oxname{OL-AbstractConcrete}
  are immediate.

  \oxname{OL-Trans} follows from the induction hypothesis.

  \oxname{OL-ConcreteAbstract} follows from the induction hypothesis and
  \Lemma{lemma:computety-well-typed-extension}.

  This leaves \oxname{OL-CombineConcrete} as the most interesting case. Here we
  use $\forall \oxcprov \in \oxfprovs{\oxsitype_\oxid}. \;
  \oxnotreborrowed{\oxvarctx}{\oxcprov}$ from our premise and note that since
  this holds for arbitrary regions $\oxcprov$, we know that it holds for both
  $\oxcprov_1$ and $\oxcprov_2$ in the premise of \oxname{OL-CombineConcrete} and
  thus we are done.
\end{proof}

\begin{oxlemma}{Region Rewriting is Preserved under Well-Typed Extensions}
  {lemma:subtyping-well-typed-extension} If
  $\oxtypevalidity{\oxglobalctx}{\oxemptyctx}{\oxvarctx}{\oxsitype_\oxid}$ and
  $\forall \oxcprov \in \oxfprovs{\oxsitype_\oxid}. \;
  \oxnotreborrowed{\oxvarctx}{\oxcprov}$ and
  $\oxtunify{\oxemptyctx}{\oxkontctx}{\oxvarctx}{\oxtype_1}{\oxtype_2}{\oxvarctx^\prime}$
  then $\oxtunify{\oxemptyctx}{\oxkontctx}{ \oxextendctx{\oxvarctx}{
      \oxvarctxentry{\oxid}{\oxsitype_\oxid} } }{\oxtype_1}{\oxtype_2}{
    \oxextendctx{\oxvarctx^\prime}{ \oxvarctxentry{\oxid}{\oxsitype_\oxid} } }$.
\end{oxlemma}

\begin{proof}
  We proceed by induction on the region rewriting judgment. This gives us seven
  cases, \oxname{RR-Refl}, \oxname{RR-Trans}, \oxname{RR-Array},
  \oxname{RR-Slice}, \oxname{RR-Reference}, \oxname{RR-Tuple}, and
  \oxname{RR-Dead}.

  \oxname{RR-Refl} is immediate.

  \oxname{RR-Trans}, \oxname{RR-Array}, \oxname{RR-Slice}, \oxname{RR-Tuple} and
  \oxname{RR-Dead} all follow directly from the induction hypothesis.

  This leaves \oxname{RR-Reference} which follows from
  \Lemma{lemma:outlives-well-typed-extension} and the induction hypothesis.
\end{proof}

\begin{oxlemma}{Closure Bodies are Well-Typed under Well-Typed Extensions}
  {lemma:exprs-well-typed-extension} If
  $\oxtypevalidity{\oxglobalctx}{\oxemptyctx}{\oxvarctx}{\oxsitype_\oxid}$ and
  $\forall \oxcprov \in \oxfprovs{\oxsitype_\oxid}. \;
  \oxnotreborrowed{\oxvarctx}{\oxcprov}$ and
  $\oxtypjudge{\oxglobalctx}{\oxemptyctx}{\oxkontctx}{
    \oxnewframe{\oxvarctx}{\oxframe} }{\oxexpr}{\oxsitype}{
    \oxnewframe{\oxvarctx^\prime}{\oxframe^\prime} }$, then
  $\oxtypjudge{\oxglobalctx}{\oxemptyctx}{\oxkontctx}{ \oxnewframe{
      \oxextendctx{\oxvarctx}{ \oxvarctxentry{\oxid}{\oxsitype_\oxid} }
    }{\oxframe} }{\oxexpr}{\oxsitype}{ \oxnewframe{
      \oxextendctx{\oxvarctx^\prime}{ \oxvarctxentry{\oxid}{\oxsitype_\oxid} }
    }{\oxframe^\prime} }$ and $\forall \oxcprov \in \oxfprovs{\oxsitype_\oxid}.
  \; \forall \oxvarctxentry{\oxplace}{\oxtref{\oxcprov}{\oxmuta}{\oxxitype}} \in
  \oxexplode{\oxvarctx^\prime}. \; \not\exists \oxcprov^\prime. \;
  \oxloanpkg{\oxmuta}{\oxderef{\! \oxplace}} \in
  \oxvarctx^\prime(\oxcprov^\prime)$.
\end{oxlemma}

\begin{proof}
  We proceed by induction on the typing derivation.

  \oxname{T-Function}, \oxname{T-Abort}, \oxname{T-Unit}, \oxname{T-u32},
  \oxname{T-True}, \oxname{T-False}, and \oxname{T-Dead} are all immediate.

  \oxname{T-LetRegion}, \oxname{T-While}, \oxname{T-ForArray},
  \oxname{T-ForSlice}, \oxname{T-Closure}, \oxname{T-Tuple}, \oxname{T-Slice},
  \oxname{T-Drop}, \oxname{T-Left}, \oxname{T-Right}, and \oxname{T-Shift},
  \oxname{T-Framed}, and \oxname{T-ClosureValue} all follow directly from the
  induction hypothesis.

  For \oxname{T-Move}, \oxname{T-Copy}, \oxname{T-Borrow},
  \oxname{T-BorrowIndex}, \oxname{T-BorrowSlice}, and \oxname{T-IndexCopy}, we
  rely on \Lemma{lemma:ownership-well-typed-extension} and the induction
  hypothesis for almost all of our obligations. For all of them except
  \oxname{T-Move}, we also have to show that the type computation for
  $\oxplaceexpr$ still works in the extended stack typing. This follows from
  \Lemma{lemma:computety-well-typed-extension}.

  For \oxname{T-Branch} and \oxname{T-Match}, we use the induction hypothesis in
  conjunction with \Lemma{lemma:subtyping-well-typed-extension} to get most of
  the premises. In the end, we also need to deal with the union between the
  output environments from the two region rewriting derivations. Fortunately, we
  know that by definition this operation unions corresponding loan sets for the
  same region and as such creates no new loans leaving our not-reborrowed
  property intact.

  \oxname{T-Seq} proceeds almost directly based on just the induction
  hypothesis, but with the required note that garbage collecting loans can only
  remove loans and thus leaves our not-reborrowed property intact.
  \oxname{T-Let} follows similarly to \oxname{T-Seq} but also requires the use
  of \Lemma{lemma:subtyping-well-typed-extension}.

  \oxname{T-Assign} and \oxname{T-AssignDeref} follow from the induction
  hypothesis combined with \Lemma{lemma:subtyping-well-typed-extension} and
  \Lemma{lemma:ownership-well-typed-extension}.

  \oxname{T-App} follows from the induction hypothesis and
  \Lemma{lemma:outlives-well-typed-extension}.

  \oxname{T-Pointer} follows from \Lemma{lemma:computety-well-typed-extension}.
\end{proof}

\begin{oxlemma}{Values are Well-Typed under Well-Typed Extensions}
  {lemma:values-well-typed-extension} If
  $\oxtypevalidity{\oxglobalctx}{\oxemptyctx}{\oxvarctx}{\oxsitype_\oxid}$ and
  $\forall \oxcprov \in \oxfprovs{\oxsitype_\oxid}. \;
  \oxnotreborrowed{\oxvarctx}{\oxcprov}$ and
  $\oxtypjudge{\oxglobalctx}{\oxemptyctx}{\oxkontctx}{\oxvarctx}
  {\oxvalue}{\oxsitype}{\oxvarctx}$, then
  $\oxtypjudge{\oxglobalctx}{\oxemptyctx}{\oxkontctx}{ \oxextendctx{\oxvarctx}{
      \oxvarctxentry{\oxid}{\oxsitype_\oxid} } }{\oxvalue}{\oxsitype}{
    \oxextendctx{\oxvarctx}{ \oxvarctxentry{\oxid}{\oxsitype_\oxid} } }$.
\end{oxlemma}

\begin{proof}
  We proceed by induction on the typing derivation. The only non-immediate case
  is \oxname{T-ClosureValue}.

  \vspace{1em}
  \noindent \framebox[\textwidth]{ \figuresize
    \begin{mathpar}
      \TClosureVal
    \end{mathpar}
  } \vspace{1em}

  First, we invert the stack frame typing hypothesis to get that $\forall \oxid
  \in \oxdomain{\oxstackframe}.\;
  \oxtypjudge{\oxglobalctx}{\oxemptyctx}{\oxkontctx}{\oxnewframe{\oxvarctx}{\oxframe_c}}{\oxstackframe(\oxid)}{\oxframe_c(\oxid)}{\oxnewframe{\oxvarctx}{\oxframe_c}}$.
  We can apply the induction hypothesis to each of these statements, and apply
  \oxname{WF-Frame} to get
  $\oxsubstorevalidity{\oxglobalctx}{\oxextendctx{\oxvarctx}{\oxvarctxentry{\oxid}{\oxsitype_x}}}{\oxstackframe_c}{\oxframe_c}$.

  Next, we need to show that
  $\oxtypjudge{\oxglobalctx}{\oxemptyctx}{\oxkontctx}{ \oxnewframe{
      \oxextendctx{\oxvarctx}{ \oxvarctxentry{\oxid}{\oxsitype_\oxid} } }{
      \oxframe_c \oxcomma \oxvarctxentry{\oxid_1}{\oxsitype_1} \oxdotsc
      \oxvarctxentry{\oxid_n}{\oxsitype_n} }
  }{\oxexpr}{\oxsitype_\oxcprov}{\oxnewframe{\oxvarctx^{\prime\prime}}{\oxframe^\prime}}$
  for some $\oxvarctx^{\prime\prime}$ and $\oxframe^\prime$. We get this by
  applying \Lemma{lemma:exprs-well-typed-extension} to
  $\oxtypjudge{\oxglobalctx}{\oxemptyctx}{\oxkontctx}{ \oxnewframe{\oxvarctx}{
      \oxframe_c \oxcomma \oxvarctxentry{\oxid_1}{\oxsitype_1} \oxdotsc
      \oxvarctxentry{\oxid_n}{\oxsitype_n} }
  }{\oxexpr}{\oxsitype_\oxcprov}{\oxnewframe{\oxvarctx^\prime}{\oxframe^\prime}}$
  (from the premise of \oxname{T-ClosureValue}).
\end{proof}

\begin{oxlemma}{Stack Validity is Preserved under Well-Typed Extensions}
  {lemma:store-validity-extension} If
  $\oxstorevalidity{\oxglobalctx}{\oxvarctx}{\oxstore}$ and
  $\forall \oxcprov \in \oxfprovs{\oxsitype_\oxid}. \;
  \oxnotreborrowed{\oxvarctx}{\oxcprov}$ and
  $\oxtypjudge{\oxglobalctx}{\oxtvarctx}{\oxkontctx}{\oxvarctx}{\oxvalue}
  {\oxsitype}{\oxvarctx}$, then $\oxstorevalidity{\oxglobalctx}{
    \oxextendctx{\oxvarctx}{ \oxvarctxentry{\oxid}{\oxsitype} } }{
    \oxextendctx{\oxstore}{ \oxstoreentry{\oxid}{\oxvalue} } }$.
\end{oxlemma}

\begin{proof}
  This proof follows directly from the definition of \oxname{WF-StackFrame}.
  
  \vspace{1em}
  \noindent \framebox[\textwidth]{ \figuresize
    \begin{mathpar}
      \VStack
    \end{mathpar}
  } \vspace{1em}

  In particular, inversion of \oxname{WF-StackFrame} on
  $\oxstorevalidity{\oxglobalctx}{\oxvarctx}{\oxstore}$ gives us well-formedness
  for the remainder of the stack, $\oxdomain{\oxstackframe} =
  \oxdomain{\oxframe}|_\oxid$ and $\forall \oxid \in \oxdomain{\oxstackframe}.
  \;
  \oxtypjudge{\oxglobalctx}{\oxemptyctx}{\oxemptyctx}{\oxnewframe{\oxvarctx}{\oxframe}}
  {(\oxnewframe{\oxstore}{\oxstackframe})(\oxid)}{(\oxnewframe{\oxvarctx}{\oxframe})(\oxid)}{\oxnewframe{\oxvarctx}{\oxframe}}$.
  We can then apply \Lemma{lemma:values-well-typed-extension} to each of these
  derivations to get $\forall \oxid \in \oxdomain{ \oxextendctx{\oxstackframe}{
      \oxstoreentry{\oxid}{\oxvalue} } }. \;
  \oxtypjudge{\oxglobalctx}{\oxemptyctx}{\oxemptyctx}{
    \oxextendctx{\oxnewframe{\oxvarctx}{\oxframe}}{
      \oxvarctxentry{\oxid}{\oxsitype} } }{ (\oxextendctx{\oxstackframe}{
      \oxstoreentry{\oxid}{\oxvalue} })(\oxid) }{
    (\oxextendctx{\oxnewframe{\oxvarctx}{\oxframe}}{
      \oxvarctxentry{\oxid}{\oxsitype} })(\oxid) }{
    \oxextendctx{\oxnewframe{\oxvarctx}{\oxframe}}{
      \oxvarctxentry{\oxid}{\oxsitype} } }$. We can then see that the
  well-formedness of the remainder of the stack is unaffected, and that the
  domains when extended with $\oxid$ remain equal. The last obligation is to
  show that the $\oxvalue$ is well-typed in the current stack typing, but we
  already have that from our premise. Thus, we can apply \oxname{WF-StackFrame}
  with the extended stack to get $\oxstorevalidity{\oxglobalctx}{
    \oxextendctx{\oxvarctx}{ \oxvarctxentry{\oxid}{\oxsitype} } }{
    \oxextendctx{\oxstore}{ \oxstoreentry{\oxid}{\oxvalue} } }$.
\end{proof}

\subsection{Preservation after Assignment Lemmas}
\label{sec:assign-lemmas}

\begin{oxlemma}{Ownership Safety is Preserved after Assignment}
  {lemma:ownership-safety-assignment}
  If $\oxlookup{\oxvarctx}{\oxplace_a}{\oxsxtype}$ and
  $\oxtunify[\oxnoop]{\oxtvarctx}{\oxkontctx}{
    \oxrsub{\oxvarctx}{\oxderef{\oxplace_a}}
  }{\oxsitype}{\oxsxtype}{\oxvarctx^\prime}$ and
  $\oxrgnuniqueto{\oxsxtype}{\oxplace_a}{\oxvarctx}$ and
  $\oxmusafetyinner{\oxemptyctx}{\oxkontctx}{
    \oxnewframe{\oxvarctx}{\oxframe}
  }{\oxmuta}{\overline{\oxplace_o}}{\oxplaceexpr}{\oxset{\oxloans}}$ and
  $\oxset{\overline{\oxplace_n}} = \oxset{\overline{\oxplace_o}} \setminus \oxplace_a$,
  then $\oxmusafetyinner{\oxemptyctx}{\oxkontctx}{
    \oxnewframe{
      \oxgcloans{\oxkontctx}{\oxtupdate{\oxvarctx^\prime}{\oxplace_a}{\oxsitype}}
    }{\oxframe}
  }{\oxmuta}{\overline{\oxplace_n}}{\oxplaceexpr}{\oxset{\oxloans^\prime}}$.
\end{oxlemma}

\begin{proof}
  We proceed by induction on the ownership safety derivation.

  \vspace{1em}
  \noindent \framebox[\textwidth]{ \figuresize
    \begin{mathpar}
      \OSafePlace
    \end{mathpar}
  } \vspace{1em}

  Consider an arbitrary region $\oxcprov$ from the domain of
  $\oxnewframe{\oxvarctx}{\oxframe}$. For this $\oxcprov$, we wish to show that
  either of the two clauses in \oxname{O-SafePlace} which were previously true
  are maintained after going through
  $\oxtunify[\oxnoop]{\oxtvarctx}{\oxkontctx}{
    \oxrsub{\oxvarctx}{\oxderef{\oxplace_a}}
  }{\oxsitype}{\oxsxtype}{\oxvarctx^\prime}$ and $\oxgcloans{\oxkontctx}{\cdot}$,
  which kills loans prefixed by $\oxderef{\oxplace}$, checks that the new type
  for $\oxplace_a$ is compatible with its old type, and clears out loans
  associated with its outermost region. So, we will consider each clause as a
  separate case.

  We will first consider the case where we have $\forall
  \oxloanpkg{\oxmuta^\prime}{\oxplaceexprctx[\oxplace^\prime]} \in \oxset{\overline{\oxloan}}.
  (\oxmuta = \oxmut \vee \oxmuta^\prime = \oxmut) \implies
  \oxrelevant{\oxplace^\prime}{\oxplace}$. In this case, we know by definition of
  $\rsub$ that $\forall \oxcprov \in \oxdomain{\oxvarctx}. \;
  (\oxrsub{\oxvarctx}{\oxderef{\oxplace_a}})(\oxcprov) \subseteq \oxvarctx(\oxcprov)$.
  We then know, again by definition (see \oxname{OL-CheckConcrete}), that
  $\forall \oxcprov \in \oxdomain{\oxvarctx}. \;
  (\oxrsub{\oxvarctx}{\oxderef{\oxplaceexpr}})(\oxcprov) =
  \oxvarctx^\prime(\oxcprov)$. As such, we know that the $\rsub$ and
  $\oxgcloans{\oxkontctx}{\cdot}$ has at most shrank the obligations in this case to
  having fewer disjointedness obligations, and it is otherwise unchanged.

  This leaves us to consider the second case
  $\exists \oxvarctxentry{\oxplace^\prime}{\oxtref{\oxcprov^\prime}{\oxmuta^\prime}{\oxtype^\prime}} \in \oxexplode{\oxvarctx} \; \wedge \; \not \exists
  \oxtref{\oxcprov^\prime}{\oxmuta^\prime}{\oxtype^\prime} \in \oxkontctx \;
  \wedge
  (\forall \oxvarctxentry{\oxplace^\prime}{\oxtref{\oxcprov^\prime}{\oxmuta^\prime}{\oxtype^\prime}} \in \oxexplode{\oxvarctx}. \;
  \oxplace^\prime \in \oxset{
    \overline{\oxplace_e}
  })$. Recall that, definitionally, neither $\rsub$ nor the region
  rewriting judgment change variable bindings and their associated types in the
  environment (instead both affect only the loan sets associated with regions,
  though the latter does not when run in the checking mode $\oxnoop$).
  Thus, we know that $\oxexplode{\oxvarctx} = \oxexplode{\oxvarctx^\prime}$. We then
  need to consider two distinct possibilities for how the exclusion list has
  changed. We know from the premise that $\oxset{\overline{\oxplace_n}} =
  \oxset{\overline{\oxplace_o}} \setminus \oxplace_a$ which means that either the two sets
  are exactly identical (when $\oxplace_a \not\in \oxset{\overline{\oxplace_o}}$) or
  smaller by $\oxplace_a$ in particular
  (when $\oxplace_a \in \oxset{\overline{\oxplace_o}}$). In the former case, the
  exclusion list is unchanged which means the whole clause is true for every
  $\oxcprov^\prime$ in $\oxvarctx^\prime$ for which it was true in $\oxvarctx$. In the
  latter case, the regions $\oxcprov^\prime$ is in the type of $\oxplace_a$ which
  \emph{has} been removed from the exclusion list $\oxset{\overline{\oxplace_n}}$.
  Thus, we need to show for the loans $\oxset{\overline{\oxloan}}$ associated
  with $\oxcprov^\prime$ that $\forall \oxloanpkg{\oxmuta^\prime}{\oxplaceexprctx[\oxplace^\prime]}
  \in \oxset{\overline{\oxloan}}. (\oxmuta = \oxmut \vee \oxmuta^\prime = \oxmut)
  \implies \oxrelevant{\oxplace^\prime}{\oxplace}$. By definition,
  $\oxrgnuniqueto{\oxsxtype}{\oxplace}{\oxvarctx}$ tells us that the outermost
  region $\oxcprov^\prime$ is unique to the type $\oxsxtype$ and place $\oxplace_a$,
  and thus when we replace it with $\oxsitype$, we ensure that $\oxcprov^\prime$ does
  not occur in any type in $\oxvarctx^\prime$. Thus, the surrounding call of
  $\oxgcloans{\oxkontctx}{\cdot}$ necessarily clears out the loan set meaning that
  the set associated with $\oxcprov^\prime$ is always empty in new environment,
  meaning the disjointness condition from \oxname{O-SafePlace} holds trivially.

  \vspace{1em}
  \noindent \framebox[\textwidth]{ \figuresize
    \begin{mathpar}
      \ODeref
    \end{mathpar}
  } \vspace{1em}

  We want to produce a new derivation using \oxname{O-Deref} for
  $\oxmusafetyinner{\oxemptyctx}{\oxkontctx}{
    \oxnewframe{\oxrsub{\oxvarctx}{\oxderef{\oxplace_a}}}{\oxframe}
  }{\oxmuta}{
    \overline{\oxplace_e}
  }{
    \oxplaceexprctx[\oxderef{\oxplace}]
  }{
    \oxset{\overline{\oxloan^{\prime\prime}}}
  }$. We have $\oxlookup{(\oxnewframe{\oxvarctx}{\oxframe})}{\oxplace}{
    \oxtref{\oxcprov}{\oxmuta_\oxplace}{\oxtype_\oxplace}
  }$ from the premise of \oxname{O-Deref}. We then know by the definition of
  $\rsub$ that $\oxlookup{
    (\oxnewframe{\oxrsub{\oxvarctx}{\oxderef{\oxplace_a}}}{\oxframe})
  }{\oxplace}{
    \oxtref{\oxcprov}{\oxmuta_\oxplace}{\oxtype_\oxplace}
  }$ since $\rsub$ only affects the loan set portion of the codomain of its
  input environment. We also know from the premise of \oxname{O-Deref} that
  $\oxlookup{
    (\oxnewframe{\oxvarctx}{\oxframe})
  }{\oxcprov}{
    \oxset{\overline{\oxloanpkg{\oxmuta^\prime}{\oxplaceexpr}}^n}
  }$. Again, by the definition of $\rsub$, we have that
  $\oxlookup{
    (\oxnewframe{\oxrsub{\oxvarctx}{\oxderef{\oxplace_a}}}{\oxframe})
  }{\oxcprov}{
    \oxset{\overline{\oxloanpkg{\oxmuta^k}{\oxplaceexpr^k}}^m}
  }$ where we know $m \leq n$ and
  $\oxset{\overline{\oxloanpkg{\oxmuta^k}{\oxplaceexpr^k}}^m} \subseteq
  \oxset{\overline{\oxloanpkg{\oxmuta^\prime}{\oxplaceexpr}}^n}$.

  We know by the definition of $\rsub$ also that every place expression
  $\oxplaceexpr_d$ in
  $\oxset{\overline{\oxloanpkg{\oxmuta^\prime}{\oxplaceexpr}}^n} \setminus \oxset{\overline{\oxloanpkg{\oxmuta^k}{\oxplaceexpr^k}}^m}$
  can be decomposed into $\oxplaceexprctx_d[\oxderef{\oxplace_a}]$. This means
  that for computing the set $\textrm{excl}$, it is either the same or has
  shrunk by precisely $\oxplace_a$. This lines up with our induction hypothesis
  which we apply to each of $\forall i \in \oxset{1 \oxdots n}. \;
  \oxmusafetyinner{\oxemptyctx}{\oxkontctx}{\oxvarctx}{\oxmuta}{
    \overline{\oxplace_e} \oxcomma \textrm{excl} \oxcomma \oxplace
  }{
    \oxplaceexprctx[\oxplaceexpr_i]
  }{
    \oxset{\overline{\oxloanpkg{\oxmuta}{\oxplaceexpr_i^\prime}}}
  }$ from the premise of \oxname{O-Deref}. This gives us
  $\forall i \in \oxset{1 \oxdots n}. \;
  \oxmusafetyinner{\oxemptyctx}{\oxkontctx}{
    \oxrsub{\oxvarctx}{\oxderef{\oxplace_a}}
  }{\oxmuta}{
    \overline{\oxplace_e} \oxcomma \textrm{excl} \oxcomma \oxplace
  }{
    \oxplaceexprctx[\oxplaceexpr_i]
  }{
    \oxset{\overline{\oxloanpkg{\oxmuta}{\oxplaceexpr_i^{\prime\prime}}}}
  }$.

  Finally, for the last premise of \oxname{O-Deref}, the proof precedes
  identically to the case for \oxname{O-SafePlace} since the obligation is
  exactly the same.

  \vspace{1em}
  \noindent \framebox[\textwidth]{ \figuresize
    \begin{mathpar}
      \ODerefAbs
    \end{mathpar}
  } \vspace{1em}

  Since $\oxtvarctx = \oxemptyctx$, there are no valid reference types that have
  an abstract region, meaning the first hypothesis is a contradiction.
\end{proof}

\begin{oxlemma}{Outlives is Preserved after Assignment}
  {lemma:outlives-assignment} If
  $\oxlookup{\oxvarctx}{\oxplace_a}{\oxsxtype}$
  and $\oxtunify[\oxnoop]{\oxtvarctx}{\oxkontctx}{
    \oxrsub{\oxvarctx}{\oxderef{\oxplace_a}}
  }{\oxsitype}{\oxsxtype}{\oxvarctx^\prime}$ and
  $\oxrunify{\oxemptyctx}{\oxkontctx}{\oxvarctx}{\oxprov_1}{\oxprov_2}{\oxvarctx_o}$, then
  $\oxrunify{\oxemptyctx}{\oxkontctx}{\oxvarctx^\prime}{\oxprov_1}{\oxprov_2}{
    \oxvarctx^\prime_o
  }$.
\end{oxlemma}

\begin{proof}
  We proceed by induction on the outlives judgment. This gives us seven cases,
  \oxname{OL-Refl}, \oxname{OL-Trans}, \oxname{OL-BothAbstract},
  \oxname{OL-CombineConcrete}, \oxname{OL-CombineConcrete},
  \oxname{OL-ConcreteAbstract} and \oxname{OL-AbstractConcrete}.

  \oxname{OL-Refl}, \oxname{OL-BothAbstract} and \oxname{OL-AbstractConcrete}
  are immediate.

  \oxname{OL-Trans} follows from the induction hypothesis.

  \oxname{OL-ConcreteAbstract} follows from the induction hypothesis and
  noting that the type computation does not depend on the contents of loan sets.

  This leaves \oxname{OL-CombineConcreteUnrestricted},
  \oxname{OL-CombineConcrete} and \oxname{OL-CheckConcrete} as the most
  interesting cases. We note that the checking mode $\oxnoop$ corresponds to
  making no changes to the environment, thus $\oxvarctx^\prime =
  \oxrsub{\oxvarctx}{\oxderef{\oxplace_a}}$. Then, for each, we note that the
  value of each associated loan set in the input environment only has an effect
  on the output environment and not whether or not the rule applies. Thus, since
  we know that $\oxvarctx^\prime$ (compared to $\oxvarctx$) has had some loans
  removed (those rooted at $\oxderef{\oxplace_a}$), then we can still produce a
  derivation, only with a different, potentially smaller output.
\end{proof}

\begin{oxlemma}{Region Rewriting is Preserved after Assignment}
  {lemma:subtyping-assignment} If
  $\oxlookup{\oxvarctx}{\oxplace_a}{\oxsxtype}$
  and $\oxtunify[\oxnoop]{\oxtvarctx}{\oxkontctx}{
    \oxrsub{\oxvarctx}{\oxderef{\oxplace_a}}
  }{\oxsitype}{\oxsxtype}{\oxvarctx^\prime}$ and
  $\oxtunify{\oxemptyctx}{\oxkontctx}{\oxvarctx}{\oxtype_1}{\oxtype_2}{\oxvarctx_o}$, then
  $\oxtunify{\oxemptyctx}{\oxkontctx}{\oxvarctx^\prime}{\oxtype_1}{\oxtype_2}{
    \oxvarctx^\prime_o
  }$.
\end{oxlemma}

\begin{proof}
  We proceed by induction on the region rewriting judgment. This gives us seven
  cases, \oxname{RR-Refl}, \oxname{RR-Trans}, \oxname{RR-Array},
  \oxname{RR-Slice}, \oxname{RR-Reference}, \oxname{RR-Tuple}, and
  \oxname{RR-Dead}.

  \oxname{RR-Refl} is immediate.

  \oxname{RR-Trans}, \oxname{RR-Array}, \oxname{RR-Slice}, \oxname{RR-Tuple} and
  \oxname{RR-Dead} all follow directly from the induction hypothesis.

  This leaves \oxname{RR-Reference} which follows from
  \Lemma{lemma:outlives-assignment} and the induction hypothesis.
\end{proof}

\begin{oxlemma}{Expressions are Well-Typed after Assignment}
  {lemma:exprs-well-typed-assignment}
  If $\oxlookup{\oxvarctx}{\oxplace_a}{\oxsxtype}$
  and $\oxtunify[\oxnoop]{\oxtvarctx}{\oxkontctx}{
    \oxrsub{\oxvarctx}{\oxderef{\oxplace_a}}
  }{\oxsitype}{\oxsxtype}{\oxvarctx^\prime}$ and
  $\oxmusafety{\oxemptyctx}{\oxkontctx}{\oxvarctx^\prime}{\oxmut}{\oxplace_a}{
    \oxset{\oxloanpkg{\oxmut}{\oxplace_a}}
  }$ and
  $\oxtypjudge{\oxglobalctx}{\oxemptyctx}{\oxkontctx}{
    \oxnewframe{\oxvarctx}{\oxframe}
  }{\oxexpr}{\oxsitype}{\oxnewframe{\oxvarctx_o}{\oxframe^\prime}}$, then
  $\oxtypjudge{\oxglobalctx}{\oxemptyctx}{\oxkontctx}{
    \oxnewframe{
      \oxgcloans{\oxkontctx}{\oxtupdate{\oxvarctx^\prime}{\oxplace_a}{\oxsitype}}
    }{\oxframe}
  }{\oxexpr}{\oxsitype}{
    \oxnewframe{
      \oxtupdate{\oxvarctx^\prime_o}{\oxplace_a}{\oxsitype}
    }{\oxframe^{\prime\prime}}
  }$.
\end{oxlemma}

\begin{proof}
  We proceed by induction on the typing derivation.

  \oxname{T-Function}, \oxname{T-Abort}, \oxname{T-Unit}, \oxname{T-u32},
  \oxname{T-True}, \oxname{T-False}, and \oxname{T-Dead} are all immediate.

  \oxname{T-Seq}, \oxname{T-LetRegion}, \oxname{T-While}, \oxname{T-ForArray},
  \oxname{T-ForSlice}, \oxname{T-Closure}, \oxname{T-Tuple}, \oxname{T-Slice},
  \oxname{T-Drop}, \oxname{T-Left}, \oxname{T-Right}, and \oxname{T-Shift},
  \oxname{T-Framed}, and \oxname{T-ClosureValue} all follow directly from the
  induction hypothesis.

  For \oxname{T-Move}, \oxname{T-Copy}, \oxname{T-Borrow},
  \oxname{T-BorrowIndex}, \oxname{T-BorrowSlice}, and \oxname{T-IndexCopy}, we
  rely on \Lemma{lemma:ownership-safety-assignment} and the induction
  hypothesis for almost all of our obligations. For all of them except
  \oxname{T-Move}, we also have to show that the type computation for
  $\oxplaceexpr$ still works in the updated environment. Since we know that
  $\oxplace_a$ is ownership safe from our premise, we know that $\oxplaceexpr$
  is disjoint from $\oxplace_a$ and thus the type update could not affect its
  type computation. Otherwise, the only difference is in loan sets associated
  with regions, and thus does not affect type computation.

  \oxname{T-Pointer} requires the same argument about type computation as in
  \oxname{T-Borrow}, but does not need the additional lemmas or the induction
  hypothesis.

  For \oxname{T-Branch} and \oxname{T-Match}, we use the induction hypothesis in
  conjunction with \Lemma{lemma:subtyping-assignment} to get most of

  the premises. In the end, we also need to deal with the union between the
  output environments from the two rewriting derivations. Fortunately, we know
  that by definition this operation unions corresponding loan sets for the same
  region and so commutes with $\oxgcloans{\oxkontctx}{\cdot}$ and the type update.

  \oxname{T-Let} follows similarly to \oxname{T-Branch} and \oxname{T-Match}
  using the induction hypothesis in conjunction with
  \Lemma{lemma:subtyping-assignment} without the need to address a combined
  environment.

  \oxname{T-Assign} and \oxname{T-AssignDeref} follow from the induction
  hypothesis combined with \Lemma{lemma:subtyping-assignment} and
  \Lemma{lemma:ownership-safety-assignment}.

  \oxname{T-App} follows from the induction hypothesis and
  \Lemma{lemma:outlives-assignment}.
\end{proof}

\begin{oxlemma}{Values are Well-Typed after Assignment}
  {lemma:values-well-typed-assignment}
  If $\oxlookup{\oxvarctx}{\oxplace_a}{\oxsxtype}$ and
  $\oxtunify[\oxnoop]{\oxtvarctx}{\oxkontctx}{
    \oxrsub{\oxvarctx}{\oxderef{\oxplace_a}}
  }{\oxsitype}{\oxsxtype}{\oxvarctx^\prime}$ and
  $\oxmusafety{\oxemptyctx}{\oxkontctx}{\oxvarctx^\prime}{\oxmut}{\oxplace_a}{
    \oxset{\oxloanpkg{\oxmut}{\oxplace_a}}
  }$ and
  $\oxtypjudge{\oxglobalctx}{\oxemptyctx}{\oxkontctx}{\oxvarctx}
  {\oxvalue}{\oxtype}{\oxvarctx}$, then
  $\oxtypjudge{\oxglobalctx}{\oxemptyctx}{\oxkontctx}{
    \oxgcloans{\oxkontctx}{\oxtupdate{\oxvarctx^\prime}{\oxplace_a}{\oxsitype}}
  }{
    \oxvalue
  }{\oxtype}{
    \oxgcloans{\oxkontctx}{\oxtupdate{\oxvarctx^\prime}{\oxplace_a}{\oxsitype}}
  }$.
\end{oxlemma}

\begin{proof}
  We proceed by induction on the value typing derivation. The only non-immediate
  cases are \oxname{T-Pointer} and \oxname{T-ClosureValue}.

  \vspace{1em}
  \noindent \framebox[\textwidth]{ \figuresize
    \begin{mathpar}
      \TPointer
    \end{mathpar}
  } \vspace{1em}

  In \oxname{T-Pointer}, we have a requirement that
  $\oxloanpkg{\oxmuta}{\oxplace} \in \oxvarctx(\oxcprov)$ which could
  potentially be affected by the kill rules. However, note that the definition
  of $\rsub$ is such that we only remove loans of the form
  $\oxderef{\oxplace_a}$ which necessarily cannot match this loan which has no
  dereference in it. Thus, we know that $\oxloanpkg{\oxmuta}{\oxplace} \in
  (\oxrsub{\oxvarctx}{\oxderef{\oxplace}})(\oxcprov)$ and thus,
  $\oxtypjudge{\oxglobalctx}{\oxemptyctx}{\oxkontctx}{
    \oxrsub{\oxvarctx}{\oxderef{\oxplace_a}}
  }{
    \oxptr{\oxreferentctx[\oxplace]}
  }{\oxsitype}{
    \oxrsub{\oxvarctx}{\oxderef{\oxplace_a}}
  }$. We then know that the checking mode $\oxnoop$ for rewriting does not
  change the output environment, and thus, $\oxvarctx^\prime =
  \oxrsub{\oxvarctx}{\oxderef{\oxplace_a}}$. Then, we know that
  $\oxplace \neq \oxplace_a$ (this would otherwise conflict with the ownership
  safety derivation for $\oxplace_a$ in our premise), so the type update for
  $\oxplace_a$ does not impact \oxname{T-Pointer}. Finally, the call to
  $\oxgcloans{\oxkontctx}{\cdot}$ clears out any unused regions, but the region
  $\oxcprov$ here is still in use and thus not changed.

  \vspace{1em}
  \noindent \framebox[\textwidth]{ \figuresize
    \begin{mathpar}
      \TClosureVal
    \end{mathpar}
  } \vspace{1em}

  
  First, we invert the stack frame typing hypothesis to get that $\forall \oxid
  \in \oxdomain{\oxstackframe}. \;
  \oxtypjudge{\oxglobalctx}{\oxemptyctx}{\oxkontctx}{
    \oxnewframe{\oxvarctx}{\oxframe_c}
  }{\oxstackframe(\oxid)}{\oxframe_c(\oxid)}
  {\oxnewframe{\oxvarctx}{\oxframe_c}}$.
  We can apply the induction hypothesis to each of these statements, and apply
  \oxname{WF-Frame} to get
  $\oxsubstorevalidity{\oxglobalctx}{
      \oxgcloans{\oxkontctx}{\oxtupdate{\oxvarctx^\prime}{\oxplace_a}{\oxsitype}}
  }{\oxstackframe_c}{\oxframe_c}$.

  Next, we need to show that
  $\oxtypjudge{\oxglobalctx}{\oxemptyctx}{\oxkontctx}{
    \oxnewframe{
      \oxgcloans{\oxkontctx}{\oxtupdate{\oxvarctx^\prime}{\oxplace_a}{\oxsitype}}
    }{
      \oxframe_c \oxcomma
      \oxvarctxentry{\oxid_1}{\oxsitype_1} \oxdotsc
      \oxvarctxentry{\oxid_n}{\oxsitype_n}
    }
  }{\oxexpr}{\oxsitype_\oxcprov}{\oxnewframe{\oxvarctx^{\prime\prime}}{\oxframe^\prime}}$
  for some $\oxvarctx^{\prime\prime}$ and $\oxframe^\prime$. We get this by
  applying \Lemma{lemma:exprs-well-typed-assignment} to
  $\oxtypjudge{\oxglobalctx}{\oxemptyctx}{\oxkontctx}{
    \oxnewframe{\oxvarctx}{
      \oxframe_c \oxcomma
      \oxvarctxentry{\oxid_1}{\oxsitype_1} \oxdotsc
      \oxvarctxentry{\oxid_n}{\oxsitype_n}
    }
  }{\oxexpr}{\oxsitype_\oxcprov}{\oxnewframe{\oxvarctx^\prime}{\oxframe^\prime}}$
  from the premise of \oxname{T-ClosureValue}.
\end{proof}

\begin{oxlemma}{Stack Validity is Preserved after Assignment}
  {lemma:stack-validity-assign}
  If $\oxstorevalidity{\oxglobalctx}{\oxvarctx}{\oxstore}$ and
  $\oxlookup{\oxvarctx}{\oxplace}{\oxsxtype}$ and
  $\oxtypjudge{\oxglobalctx}{\oxemptyctx}{\oxkontctx}{\oxvarctx}{
    \oxvalue
  }{\oxsitype}{\oxvarctx}$ and
  $\oxtunify{\oxemptyctx}{\oxkontctx}{
    \oxrsub{\oxvarctx}{\oxderef{\oxplace}}
  }{\oxsitype}{\oxsxtype}{\oxvarctx^\prime}$ and
  $\oxmusafety{\oxemptyctx}{\oxkontctx}{\oxvarctx^\prime}{\oxmut}{\oxplace}{
    \oxset{\oxloanpkg{\oxmut}{\oxplace}}
  }$ and
  $\oxnorm{\oxstore}{\oxplace}{\oxplace}{\oxvaluectx}{\_}$ and
  $\oxplace = \oxproj{\oxid}{\oxpath}$,
  then
  $\oxstorevalidity{\oxglobalctx}{
    \oxgcloans{\oxkontctx}{\oxtupdate{\oxvarctx^\prime}{\oxplace}{\oxsitype}}
  }{
    \oxvupdate{\oxstore}{\oxid}{\oxvaluectx[\oxvalue]}
  }$.
\end{oxlemma}

\begin{proof}
  The proof proceeds by induction on the stack validity judgment
  $\oxstorevalidity{\oxglobalctx}{\oxvarctx}{\oxstore}$ which has two cases,
  \oxname{WF-StackEmpty} and \oxname{WF-StackFrame}.

  \vspace{1em}
  \noindent \framebox[\textwidth]{
    \figuresize
    \begin{mathpar}
      \VStackEmpty
    \end{mathpar}
  }
  \vspace{1em}

  In this case, the stack is empty and therefore, we have a contradiction since
  our premise says that $\oxlookup{\oxvarctx}{\oxplace}{\oxsxtype}$, but
  $\oxvarctx = \oxemptyctx$ and $\oxemptyctx(\oxplace)$ necessarily fails.

  \vspace{1em}
  \noindent \framebox[\textwidth]{
    \figuresize
    \begin{mathpar}
      \VStack
    \end{mathpar}
  }
  \vspace{1em}

  In the premise of \oxname{WF-StackFrame}, we have a collection of typing
  judgments for values stored in the stack. This naturally leads us to another
  case split: either $\oxid$ (the root of $\oxplace$ from $\oxplace =
  \oxproj{\oxid}{\oxpath}$) is in the current frame or it is not.

  If $\oxid$ is not in the current frame, we apply our induction hypothesis to
  $\oxstorevalidity{\oxglobalctx}{\oxvarctx}{\oxstore}$ to get
  $\oxstorevalidity{\oxglobalctx}{
    \oxtupdate{\oxvarctx^\prime}{\oxplace}{\oxsitype}
  }{
    \oxvupdate{\oxstore}{\oxid}{\oxvaluectx[\oxvalue]}
  }$. Then, we apply \oxname{WF-StackFrame} with the same typing judgments we
  already have to reach our overall conclusion of
  $\oxstorevalidity{\oxglobalctx}{
    \oxtupdate{
      (\oxnewframe{\oxvarctx^\prime}{\oxframe})
    }{\oxplace}{\oxsitype}
  }{
    \oxvupdate{
      (\oxnewframe{\oxstore}{\oxstackframe})
    }{\oxid}{\oxvaluectx[\oxvalue]}
  }$ (noting that substituting inside or outside is definitionally equal when
  we know that $\oxid \not\in \oxdomain{\oxstore}$).

  If $\oxid$ is in the current frame, then we apply
  \Lemma{lemma:subtyping-value-typing} to
  $\oxtypjudge{\oxglobalctx}{\oxemptyctx}{\oxkontctx}{\oxvarctx}{
    \oxvalue
  }{\oxsitype}{\oxvarctx}$ and
  $\oxtunify{\oxemptyctx}{\oxkontctx}{\oxvarctx}{\oxsitype}{\oxsxtype}{\oxvarctx^\prime}$ (both
  from our premise) to get
  $\oxtypjudge{\oxglobalctx}{\oxemptyctx}{\oxkontctx}{\oxvarctx^\prime}{
    \oxvalue
  }{\oxsitype}{\oxvarctx^\prime}$. Then, we note that it would be a well-formedness
  violation for this value to depend on $\oxid$ itself (since that would mean it
  was a cyclical reference) and thus, we can get that
  $\oxtypjudge{\oxglobalctx}{\oxemptyctx}{\oxkontctx}{
    \oxtupdate{\oxvarctx^\prime}{\oxplace}{\oxsitype}
  }{
    \oxvalue
  }{\oxsitype}{
    \oxtupdate{\oxvarctx^\prime}{\oxplace}{\oxsitype}
  }$. Finally, we can garbage collect the loans from the old type to get
  $\oxtypjudge{\oxglobalctx}{\oxemptyctx}{\oxkontctx}{
    \oxgcloans{\oxkontctx}{\oxtupdate{\oxvarctx^\prime}{\oxplace}{\oxsitype}}
  }{
    \oxvalue
  }{\oxsitype}{
    \oxgcloans{\oxkontctx}{\oxtupdate{\oxvarctx^\prime}{\oxplace}{\oxsitype}}
  }$

  For the other typing judgments in this frame, we apply
  \Lemma{lemma:values-well-typed-assignment} to get
  $\forall \oxid \in \oxdomain{\oxstackframe}. \;
  \oxtypjudge{\oxglobalctx}{\oxemptyctx}{\oxemptyctx}{
    \oxgcloans{\oxkontctx}{\oxtupdate{\oxvarctx^\prime}{\oxplace}{\oxsitype}}
  }{
    (\oxnewframe{\oxstore}{\oxstackframe})(\oxid)
  }{(\oxnewframe{\oxvarctx^\prime}{\oxframe})(\oxid)}{
    \oxgcloans{\oxkontctx}{\oxtupdate{\oxvarctx^\prime}{\oxplace}{\oxsitype}}
  }$. Thus, we can apply \oxname{WF-StackFrame} to conclude
  $\oxstorevalidity{\oxglobalctx}{
    \oxgcloans{\oxkontctx}{\oxtupdate{\oxvarctx^\prime}{\oxplace}{\oxsitype}}
  }{
    \oxvupdate{\oxstore}{\oxid}{\oxvaluectx[\oxvalue]}
  }$.
\end{proof}

\subsection{Values are Well-Types at Rewritten Types Lemma}
\label{sec:values-typed-supertype-lemma}

\begin{oxlemma}{Values are Well-Typed At Rewritten Types}
  {lemma:values-typed-supertype} 
  If $\oxtypjudge{\oxglobalctx}{\oxtvarctx}{\oxkontctx}{\oxvarctx}{\oxvalue}{\oxsitype}
  {\oxvarctx_i}$ and $\oxtunify{\oxtvarctx}{\oxkontctx}{\oxvarctx_i}{\oxsitype}
  {{\oxsitype}^\prime}{\oxvarctx^\prime}$, then $\oxtypjudge{\oxglobalctx}
  {\oxtvarctx}{\oxkontctx}{\oxvarctx^\prime}{\oxvalue}{{\oxsitype}^\prime}{\oxvarctx^\prime}$.
\end{oxlemma}
\begin{proof}
  We proceed by induction on the value typing relation.

  In the case of \oxname{T-Tuple}, we need to apply the induction hypothesis for
  each entry which has a changed type, and \Lemma{lemma:subtyping-value-typing}
  for each entry which does not.

  In the case of \oxname{T-Array}, we just apply the induction hypothesis to
  each entry.

  \vspace{1em}
  \noindent \framebox[\textwidth]{
    \figuresize
    \begin{mathpar}
      \TPointer
    \end{mathpar}
  }
  \vspace{1em}

  For the \oxname{T-Pointer} case, we proceed by induction on the region
  rewriting judgement. The only interesting cases are for reference types. From
  there, we proceed by induction on the outlives relation, for which the only
  interesting case is \oxname{OL-CombineConcrete}.
 
  \vspace{1em}
  \noindent \framebox[\textwidth]{
    \figuresize
    \begin{mathpar}
      \UCombineLocalProvs
    \end{mathpar}
  }
  \vspace{1em}

  The \oxname{T-Pointer} case is immediate. We know that the referent type is
  preserved since we do not change any types in the context, and we know the
  loan is preserved since loan sets only grow.

  In all other cases, we know the types cannot change, which means $\oxvarctx =
  \oxvarctx^\prime$, so we are done.
\end{proof}

\subsection{Function Definitions are Self-Contained Lemma}
\label{sec:function-defns-lemma}

\begin{oxlemma}{Function Definitions are Self-Contained}
  {lemma:function-defns-self-contained}\hfill\\
  If $\oxctxswellformed{\oxglobalctx}{\oxemptyctx}{\oxvarctx}{\oxkontctx}$ and
  $\oxlookup{\oxglobalctx}{\oxfnname}{
    \oxfuncdef{\oxfnname}{
      \overline{\oxenvvar} \oxcomma \overline{\oxabsprov} \oxcomma \overline{\oxtvar}
    }{
      \oxascribe{\oxid_1}{\oxsitype_1} \oxdotsc
      \oxascribe{\oxid_n}{\oxsitype_n}
    }{\oxsitype_r}{\overline{\oxabsprov_1: \oxabsprov_2}}{\oxexpr}
  }$, then
  $\oxtypjudge{\oxglobalctx}{
    \overline{\oxtvarctxentry{\oxenvvar}{\oxkenv}} \oxcomma
    \overline{\oxtvarctxentry{\oxabsprov}{\oxkprov}} \oxcomma
    \overline{\oxabsprov_1 \oxoutlives \oxabsprov_2} \oxcomma
    \overline{\oxtvarctxentry{\oxtvar}{\oxktype}}
  }{\oxkontctx}{
    \oxnewframe{\oxvarctx}{
      \oxvarctxentry{\oxid_1}{\oxsitype_1} \oxdotsc
      \oxvarctxentry{\oxid_n}{\oxsitype_n}
    }
  }{
    \oxframed{\oxexpr}
  }{\oxsitype_f}{\oxvarctx}$.
\end{oxlemma}

\begin{proof}
  Begin by noting that \oxname{WF-FunctionDefinition} gives us that
  \\ $\oxtypjudge{\oxglobalctx}{\overline{\oxtvarctxentry{\oxenvvar}{\oxkenv}}
    \oxcomma \overline{\oxtvarctxentry{\oxabsprov}{\oxkprov}} \oxcomma
    \overline{\oxabsprov_1 \oxoutlives \oxabsprov_2} \oxcomma
    \overline{\oxtvarctxentry{\oxtvar}{\oxktype}}}{\oxemptyctx}{\oxnewframe{\oxemptyctx}
    {\oxvarctxentry{\oxid_1}{\oxsitype_1} \oxdotsc
      \oxvarctxentry{\oxid_n}{\oxsitype_n}}}{\oxexpr}{\oxsitype_f}{\oxvarctx^\prime}$.
  We also have by inspection of the typing rules that $\oxvarctx^\prime =
  \oxnewframe{\oxemptyctx}{\oxframe^\prime}$ for some frame $\oxframe^\prime$.
  Then by \oxname{T-Framed}, it suffices to show that
  $\oxtypjudge{\oxglobalctx}{\overline{\oxtvarctxentry{\oxenvvar}{\oxkenv}}
    \oxcomma \overline{\oxtvarctxentry{\oxabsprov}{\oxkprov}} \oxcomma
    \overline{\oxabsprov_1 \oxoutlives \oxabsprov_2} \oxcomma
    \overline{\oxtvarctxentry{\oxtvar}{\oxktype}}}{\oxkontctx}{\oxnewframe{\oxvarctx}
    {\oxvarctxentry{\oxid_1}{\oxsitype_1} \oxdotsc
      \oxvarctxentry{\oxid_n}{\oxsitype_n}}}{\oxexpr}{\oxsitype_f}{\oxnewframe{\oxvarctx}{\oxframe^\prime}}$.
  But note that this is immediate. The typing derivation with $\oxemptyctx$ and
  the current frame means that there's absolutely no reliance on context outside
  $\oxid_1, \ldots \oxid_n$, and these places are necessarily completely
  disjoint from places in $\oxvarctx$ since any regions in their types must
  be abstract.
\end{proof}

\subsection{Subset Related Environments Lemma}
\label{sec:subset-related-lemmas}

\begin{oxlemma}{Outlives Produces Subset-Related Environments}
  {lemma:outlives-subset-related-envs}
  If $\oxrunify{\oxtvarctx}{\oxkontctx}{\oxvarctx}{\oxcprov_1}{\oxcprov_2}{\oxvarctx^\prime}$, then
  $\forall \oxcprov. \; \oxvarctx(\oxcprov) \subseteq \oxvarctx^\prime(\oxcprov)$.
\end{oxlemma}

\begin{proof}
  The proof proceeds by induction on the outlives relation $\oxrunify{\oxtvarctx}{\oxkontctx}
  {\oxvarctx}{\oxcprov_1}{\oxcprov_2}{\oxvarctx^\prime}$. We will consider each case.

  \vspace{1em}
  \noindent \framebox[\textwidth]{
    \figuresize
    \begin{mathpar}
      \UReflProv \and \UAbsProvs \and \OLocalAbsProvs \and \UCheckLocalProvs
    \end{mathpar}
  }
  \vspace{1em}

  Each of \oxname{OL-Refl}, \oxname{OL-BothAbstract},
  \oxname{OL-AbstractConcrete}, and \oxname{OL-CheckConcrete} are immediate
  since $\oxvarctx^\prime = \oxvarctx$.

  \vspace{1em}
  \noindent \framebox[\textwidth]{
    \figuresize
    \begin{mathpar}
      \UTransProv \and \OAbsLocalProvs
    \end{mathpar}
  }
  \vspace{1em}

  Both \oxname{OL-Trans} and \oxname{OL-ConcreteAbstract} follow by applying the induction
  hypothesis to all instances of the outlives judgment in their premise and then relying
  on transitivity of subset.

  \vspace{1em}
  \noindent \framebox[\textwidth]{
    \figuresize
    \begin{mathpar}
      \UCombineLocalProvs \and \UCombineLocalProvsUnrest
    \end{mathpar}
  }
  \vspace{1em}

  For \oxname{OL-CombineConcrete} and \oxname{OL-CombineConcreteUnrestricted},
  the conclusion is almost immediate since $\oxvarctx^\prime$ is very nearly
  $\oxvarctx$. However, it differs in the loan set for one particular region
  $\oxcprov_2$. Fortunately, its new loan set in $\oxvarctx^\prime$ is the union
  of its loan set with the loan set for $\oxcprov_1$ and thus we immediately
  have $\oxvarctx(\oxcprov_2) \subseteq \oxvarctx^\prime(\oxcprov_2)$.
\end{proof}

\begin{oxlemma}{Region Rewriting Produces Subset-Related Environments}
  {lemma:subtyping-subset-related-envs}
  If $\oxtunify{\oxtvarctx}{\oxkontctx}{\oxvarctx}{\oxtype_1}{\oxtype_2}{\oxvarctx^\prime}$, then
  $\forall \oxcprov. \; \oxvarctx(\oxcprov) \subseteq \oxvarctx^\prime(\oxcprov)$.
\end{oxlemma}

\begin{proof}
  This proof proceeds by induction on the region rewriting relation
  $\oxtunify{\oxtvarctx}{\oxkontctx}{\oxvarctx}{\oxtype_1}{\oxtype_2}{\oxvarctx^\prime}$. We will
  consider each case. For \oxname{RR-Refl} and \oxname{RR-Uninit}, the output
  environment $\oxvarctx^\prime$ is precisely $\oxvarctx$ and thus the result is
  immediate. For \oxname{RR-Trans}, \oxname{RR-Array}, \oxname{RR-Slice} and
  \oxname{RR-Tuple}, the result follows from applying the induction hypothesis
  to every region rewriting derivation in their premise and combining the
  results by transitivity of subset. This leaves us with one more interesting
  case, \oxname{RR-Reference}. For this case, apply
  \Lemma{lemma:outlives-subset-related-envs} to the outlives derivation in the
  premise. Then, apply our induction hypothesis to the region rewriting
  derivation in their premise. Finally, combine the two by transitivity of
  subset.
\end{proof}

\begin{oxlemma}{Frame Typing Union Produces Subset-Related Environments}
  {lemma:union-subset-frames}
  If $\oxframe = \oxframe_1 \oxintersect \oxframe_2$, then
  $\forall \oxcprov. \; \oxframe_1(\oxcprov) \subseteq \oxframe(\oxcprov)$ and
  $\forall \oxcprov. \; \oxframe_2(\oxcprov) \subseteq \oxframe(\oxcprov)$.
\end{oxlemma}

\begin{proof}
  First, note that the definition of $\oxintersect$ is symmetric and thus we will only
  prove the first conclusion, the second proceeding immediately the same in all cases.

  We proceed by induction over the frame typing. For $\oxemptyctx \oxintersect
  \oxemptyctx$, the case follows immediately. For
  $\oxextendctx{\oxframe_1}{\oxvarctxentry{\oxid}{\oxtype}} \oxintersect
  \oxextendctx{\oxframe_2}{\oxvarctxentry{\oxid}{\oxtype}}$, the result follows
  directly from applying the induction hypothesis to $\oxframe_1 \oxintersect
  \oxframe_2$. The last case is the interesting one,
  $\oxextendctx{\oxframe_1}{\oxloanctxentry{\oxcprov}{\overline{\oxloan}}}
  \oxintersect
  \oxextendctx{\oxframe_2}{\oxloanctxentry{\oxcprov}{\overline{\oxloan^\prime}}}$.
  In this case, we can apply our induction hypothesis to get $\forall
  \oxcprov^\prime \in \oxdomain{\oxframe_1}$, $\oxframe_1(\oxcprov^\prime)
  \subseteq \oxframe_2(\oxcprov^\prime)$. Now we just need that
  $\oxset{\overline{\oxloan}} \subseteq \oxframe(\oxcprov)$. But this holds immediately
  since $\oxframe(\oxcprov) = \oxset{\overline{\oxloan}} \cup
  \oxset{\overline{\oxloan^\prime}}$, so we're done.
\end{proof}

\begin{oxlemma}{Stack Typing Union Produces Subset-Related Environments}
  {lemma:union-subset-related-envs}
  If $\oxvarctx = \oxvarctx_1 \oxintersect \oxvarctx_2$, then
  $\forall \oxcprov. \; \oxvarctx_1(\oxcprov) \subseteq \oxvarctx(\oxcprov)$ and
  $\forall \oxcprov. \; \oxvarctx_2(\oxcprov) \subseteq \oxvarctx(\oxcprov)$.
\end{oxlemma}

\begin{proof}
  First, note that the definition of $\oxintersect$ is symmetric and thus we will only
  prove the first conclusion, the second proceeding immediately the same in all cases.

  We proceed by induction over the Stack Typing. 

  For $\oxemptyctx \oxintersect \oxemptyctx$, the result is trivial and thus
  immediate. For $\oxnewframe{\oxvarctx_1}{\oxframe} \oxintersect
  \oxnewframe{\oxvarctx_2}{\oxframe}$, we apply the induction hypothesis and
  \Lemma{lemma:union-subset-frames}.
\end{proof}

\begin{oxlemma}{Subset-Related Frames are also Frame Typing Union Related}
  {lemma:subset-related-frame-related} If $\forall \oxcprov \in \oxframe.$
  $\oxframe^\prime(\oxcprov) \subseteq \oxframe(\oxcprov)$ and
  $\oxdomain{\oxframe} = \oxdomain{\oxframe^\prime}$, then $\exists \oxframe_o$ such
  that $\oxframe = \oxframe^\prime \oxintersect \oxframe_o$.
\end{oxlemma}

\begin{proof}
  We proceed by induction over the frame typing. For the $\oxemptyctx$ case, the
  proof follows immediately.

  For $\oxframe = \oxframe_i, \oxid : \oxtype$ and $\oxframe^\prime =
  \oxframe_i^\prime, \oxid : \oxtype$, we just apply the induction hypothesis
  on $\oxframe_i$ and $\oxframe_i^\prime$, and add $\oxid : \oxtype$ to
  $\oxframe_o$.

  For $\oxframe = \oxframe_i, \oxcprov \mapsto \oxset{\oxloans}$ and
  $\oxframe^\prime = \oxframe_i^\prime, \oxcprov \mapsto
  \oxset{\oxloans^\prime}$, we apply the induction hypothesis, and add on
  $\oxcprov \mapsto \oxset{\oxloans} \setminus \oxset{\oxloans^\prime}$, which
  is well defined because from our premise we have $\oxset{\oxloans^\prime}
  \subseteq \oxset{\oxloans}$.
\end{proof}
    

\begin{oxlemma}{Subset-Related Environments are also Stack Typing Union Related}
  {lemma:subset-related-stack-related} If $\forall \oxcprov \in \oxvarctx.$
  $\oxvarctx^\prime(\oxcprov) \subseteq \oxvarctx(\oxcprov)$ and
  $\oxdomain{\oxvarctx} = \oxdomain{\oxvarctx^\prime}$, then $\exists \oxvarctx_o$ such
  that $\oxvarctx = \oxvarctx^\prime \oxintersect \oxvarctx_o$.
\end{oxlemma}

\begin{proof}
  Proceed by induction on the stack typing. In the $\oxemptyctx$ case, the proof
  is immediate. In the $\oxvarctx = \oxnewframe{\oxvarctx_i}{\oxframe}$ case, we
  apply \Lemma{lemma:subset-related-frame-related} and the induction hypothesis.
\end{proof}

\subsection{Preservation in More Precise Environments Lemmas}
\label{sec:more-precise-lemmas}

\begin{oxlemma}{Type Computation Preserved in More Precise Environments}
  {lemma:computety-more-precise}\hfill\\ If $\oxdomain{\oxvarctx} =
  \oxdomain{\oxvarctx^\prime}$, and $\forall \oxid.$ $\oxvarctx(\oxid) = \oxvarctx^\prime(\oxid)$ and $\forall \oxcprov.$ $\oxvarctx^\prime(\oxcprov)
  \subseteq \oxvarctx(\oxcprov)$ and
  $\oxcomputetynoprov{\oxtvarctx}{\oxvarctx}{\oxmuta}{\oxplaceexpr}{\oxxitype}$
  then
  $\oxcomputetynoprov{\oxtvarctx}{\oxvarctx^\prime}{\oxmuta}{\oxplaceexpr}{\oxxitype}$.
\end{oxlemma}

\begin{proof}
  Proceed by induction over the type computation judgement with $\oxvarctx$. The
  only non-immediate case is \oxname{TC-Deref}.

  \vspace{1em}
  \noindent \framebox[\textwidth]{
    \figuresize
    \begin{mathpar}
      \TCDeref
    \end{mathpar}
  }
  \vspace{1em}

  This follows from the induction hypothesis and
  \Lemma{lemma:outlives-subset-related-envs}.

\end{proof}

\begin{oxlemma}{Type Well-Formedness Preserved in More Precise Environments}
  {lemma:wf-type-more-precise}\hfill\\
  If $\oxdomain{\oxvarctx} =
  \oxdomain{\oxvarctx^\prime}$, and $\forall \oxcprov.$ $\oxvarctx^\prime(\oxcprov)
  \subseteq \oxvarctx(\oxcprov)$ and
  $\oxtypevalidity{\oxglobalctx}{\oxtvarctx}{\oxvarctx}{\oxsitype}$
  then
  $\oxtypevalidity{\oxglobalctx}{\oxtvarctx}{\oxvarctx^\prime}{\oxsitype}$.
\end{oxlemma}

\begin{proof}
  Proceed by induction over the type well formedness under $\oxvarctx$. The only
  case that isn't immediate or doesn't proceed directly from the induction
  hypothesis is \oxname{WF-Ref}.

  \vspace{1em}
  \noindent \framebox[\textwidth]{
    \figuresize
    \begin{mathpar}
      \VRef
    \end{mathpar}
  }
  \vspace{1em}

  We can apply the induction hypothesis to get the premise that $\oxxitype$ is
  well formed. For our other premise, we need to show that our region is still
  well-formed. We can look at this by cases. If the region $\oxprov$ is local,
  we can apply \oxname{WF-LocalRegion} since we know $\oxdomain{\oxvarctx} =
  \oxdomain{\oxvarctx^\prime}$. If the region $\oxprov$ is abstract, we can apply
  \oxname{WF-AbstractRegion} since we know that $\oxtvarctx$ is unchanged.
\end{proof}


\begin{oxlemma}{Stack Typing Validity Preserved in More Precise Environments}
  {lemma:wf-varctx-more-precise}\hfill\\
  If $\oxdomain{\oxvarctx} =
  \oxdomain{\oxvarctx^\prime}$, and $\forall \oxcprov.$ $\oxvarctx^\prime(\oxcprov)
  \subseteq \oxvarctx(\oxcprov)$ and
  $\oxvarctxwellformed{\oxglobalctx}{\oxtvarctx}{\oxvarctx}$
  then
  $\oxvarctxwellformed{\oxglobalctx}{\oxtvarctx}{\oxvarctx^\prime}$
\end{oxlemma}

\begin{proof}
  Proceed by induction over the stack typing validity judgement for $\oxvarctx$.
  The empty case is trivial, so the interesting case is \oxname{WF-StackTyping}.

  \vspace{1em}
  \noindent \framebox[\textwidth]{
    \figuresize
    \begin{mathpar}
      \WFVarCtx
    \end{mathpar}
  }
  \vspace{1em}

  We get $\oxvarctxwellformed{\oxglobalctx}{\oxtvarctx}{\oxvarctx^\prime}$ by
  induction. We get the type well formedness from
  \Lemma{lemma:wf-type-more-precise}. We get the type computation from
  \Lemma{lemma:computety-more-precise}, and that's all we needed to show.

\end{proof}

\begin{oxlemma}{Ownership Safety Preserved in More Precise Environments}
  {lemma:ownership-safety-more-precise}\hfill\\ If $\oxdomain{\oxvarctx} =
  \oxdomain{\oxvarctx^\prime}$, and $\forall \oxcprov.$ $\oxvarctx^\prime(\oxcprov)
  \subseteq \oxvarctx(\oxcprov)$ and
  $\oxmusafety{\oxtvarctx}{\oxkontctx}{\oxvarctx}{\oxmuta}{\oxplaceexpr}{
    \oxset{\overline{\oxloan}}
  }$
  then
  $\oxmusafety{\oxtvarctx}{\oxkontctx}{\oxvarctx^\prime}{\oxmuta}{\oxplaceexpr}{
    \oxset{\overline{\oxloan^\prime}}
  }$ and
  $\oxset{\overline{\oxloan^\prime}} \subseteq \oxset{\overline{\oxloan}}$.
\end{oxlemma}

\begin{proof}
  The proof proceeds by induction on the ownership safety judgment
  $\oxmusafety{\oxtvarctx}{\oxkontctx}{\oxvarctx}{\oxmuta}{\oxplaceexpr}{
    \oxset{\overline{\oxloan}}
  }$. This gives us three cases: \oxname{O-SafePlace}, \oxname{O-Deref}, and
  \oxname{O-DerefAbs}.

  \vspace{1em}
  \noindent \framebox[\textwidth]{
    \figuresize
    \begin{mathpar}
      \OSafePlace
    \end{mathpar}
  }
  \vspace{1em}

  We need to show $\oxmusafety{\oxtvarctx}{\oxkontctx}{\oxvarctx^\prime}{\oxmuta}{\oxplaceexpr}{
    \oxset{\oxloanpkg{\oxmuta}{\oxplace}}
  }$ and that
  $\oxset{\oxloanpkg{\oxmuta}{\oxplace}} \subseteq \oxset{\oxloanpkg{\oxmuta}{\oxplace}}$.
  The latter is immediate from the definition of subset which leaves us with the former.
  For the former, we'll correspondingly wish to apply \oxname{O-SafePlace} but using
  $\oxvarctx^\prime$ as our context. This means we need to show that
  $\forall \oxloanctxentry{\oxcprov^\prime}{\oxset{\overline{\oxloan}}} \in \oxvarctx^\prime.
  \; (\forall \oxloanpkg{\oxmuta^\prime}{\oxplaceexpr} \in \oxset{\overline{\oxloan}}.
  (\oxmuta = \oxmut \vee \oxmuta^\prime = \oxmut)
  \implies
  \oxrelevant{\oxplace^\prime}{\oxplace}) \vee \;
  (\exists \oxvarctxentry{\oxplace^\prime}{\oxtref{\oxcprov^\prime}{\oxmuta^\prime}{\oxtype^\prime}} \in \oxvarctx \;
  \wedge
  (\forall \oxvarctxentry{\oxplace^\prime}{\oxtref{\oxcprov^\prime}{\oxmuta^\prime}{\oxtype^\prime}} \in \oxvarctx. \;
  \oxplace^\prime \in \oxset{
    \overline{\oxplace_e}
  }))$. Fortunately, from $\oxdomain{\oxvarctx} = \oxdomain{\oxvarctx^\prime}$, we know that
  for every $\oxcprov^\prime \in \oxdomain{\oxvarctx}$, $\oxcprov^\prime \in
  \oxdomain{\oxvarctx^\prime}$, and further that $\oxvarctx^\prime(\oxcprov^\prime)
  \subseteq \oxvarctx(\oxcprov^\prime)$. Thus, for each $\oxcprov^\prime$, we know there are only
  potentially fewer loans to show if the obligation was met using the clause of
  $\forall \oxloanpkg{\oxmuta^\prime}{\oxplaceexpr} \in \oxset{\overline{\oxloan}}.
  (\oxmuta = \oxmut \vee \oxmuta^\prime = \oxmut)
  \implies \oxrelevant{\oxplace^\prime}{\oxplace}$. If the obligation was met using the other
  clause, note that $\oxvarctx$ and $\oxvarctx^\prime$ can only differ in the loan sets they
  associate with any given region and so the exact fact must still be true for
  $\oxvarctx^\prime$.

  \vspace{1em}
  \noindent \framebox[\textwidth]{
    \figuresize
    \begin{mathpar}
      \ODeref
    \end{mathpar}
  }
  \vspace{1em}

  This case proceeds much like the \oxname{O-SafePlace} case in terms of meeting the direct
  ownership safety criterion (the last premise of \oxname{O-Deref}) for the new derivation
  using \oxname{O-Deref} with $\oxvarctx^\prime$. It differs only in that we need also apply
  our induction hypothesis to each of the $n$ derivations of ownership safety used in
  $\forall i \in \oxset{1 \oxdots n}. \;
  \oxmusafetyinner{\oxtvarctx}{\oxkontctx}{\oxvarctx}{\oxmuta}{
    \overline{\oxplace_e} \oxcomma \overline{\oxplace_i} \oxcomma \oxplace
  }{
    \oxplaceexprctx[\oxplaceexpr_i]
  }{
    \oxset{\overline{\oxloanpkg{\oxmuta}{\oxplaceexpr_i^\prime}}}
  }$. This gives us
  $\forall i \in \oxset{1 \oxdots n}. \;
  \oxmusafetyinner{\oxtvarctx}{\oxkontctx}{\oxvarctx^\prime}{\oxmuta}{
    \overline{\oxplace_e} \oxcomma \overline{\oxplace_i} \oxcomma \oxplace
  }{
    \oxplaceexprctx[\oxplaceexpr_i]
  }{
    \oxset{\overline{\oxloanpkg{\oxmuta}{\oxplaceexpr_i^{\prime\prime}}}}
  }$ and $\forall i \in \oxset{1 \oxdots n}. \;
  \oxset{\overline{\oxloanpkg{\oxmuta}{\oxplaceexpr_i^{\prime\prime}}}} \subseteq
  \oxset{\overline{\oxloanpkg{\oxmuta}{\oxplaceexpr_i^\prime}}}$. The former combined with
  the same reasoning from the \oxname{O-SafePlace} case gives us
  $\oxmusafety{\oxtvarctx}{\oxkontctx}{\oxvarctx^\prime}{\oxmuta}{
    \oxplaceexprctx[\oxderef{\oxplace}]
  }{
    \oxset{
      \overline{\oxloanpkg{\oxmuta}{\oxplaceexpr^{\prime\prime}_1}} \oxcomma
      \oxdots
      \overline{\oxloanpkg{\oxmuta}{\oxplaceexpr^{\prime\prime}_n}} \oxcomma
      \oxloanpkg{\oxmuta}{\oxplaceexprctx[\oxderef{\oxplace}]}
    }
  }$ and the latter allows us to conclude
  $\oxset{
    \overline{\oxloanpkg{\oxmuta}{\oxplaceexpr^{\prime\prime}_1}} \oxcomma
    \oxdots
    \overline{\oxloanpkg{\oxmuta}{\oxplaceexpr^{\prime\prime}_n}} \oxcomma
    \oxloanpkg{\oxmuta}{\oxplaceexprctx[\oxderef{\oxplace}]}
  } \subseteq \oxset{
    \overline{\oxloanpkg{\oxmuta}{\oxplaceexpr^\prime_1}} \oxcomma
    \oxdots
    \overline{\oxloanpkg{\oxmuta}{\oxplaceexpr^\prime_n}} \oxcomma
    \oxloanpkg{\oxmuta}{\oxplaceexprctx[\oxderef{\oxplace}]}
  }$ since we know that each individual collection of loans has the subset relation from
  above and the whole set is simply their union.

  \vspace{1em}
  \noindent \framebox[\textwidth]{
    \figuresize
    \begin{mathpar}
      \ODerefAbs
    \end{mathpar}
  }
  \vspace{1em}

  This case proceeds identically to the case for \oxname{O-SafePlace}.
  We need to show $\oxmusafety{\oxtvarctx}{\oxkontctx}{\oxvarctx^\prime}{\oxmuta}{\oxplaceexpr}{
    \oxset{\oxloanpkg{\oxmuta}{\oxplaceexprctx[\oxderef{\oxplace}]}}
  }$ and that
  $\oxset{
    \oxset{\oxloanpkg{\oxmuta}{\oxplaceexprctx[\oxderef{\oxplace}]}}
  } \subseteq \oxset{
    \oxset{\oxloanpkg{\oxmuta}{\oxplaceexprctx[\oxderef{\oxplace}]}}
  }$.
  The latter is immediate from the definition of subset which leaves us with the former.
  For the former, we'll correspondingly wish to apply \oxname{O-DerefAbs} but using
  $\oxvarctx^\prime$ as our context. This means we need to show that
  $\forall \oxloanctxentry{\oxcprov^\prime}{\oxset{\overline{\oxloan}}} \in \oxvarctx^\prime.
  \; (\forall \oxloanpkg{\oxmuta^\prime}{\oxplaceexpr} \in \oxset{\overline{\oxloan}}.
  (\oxmuta = \oxmut \vee \oxmuta^\prime = \oxmut)
  \implies
  \oxrelevant{\oxplace^\prime}{\oxplaceexprctx[\oxderef{\oxplace}]}) \vee \;
  (\exists \oxvarctxentry{\oxplace^\prime}{\oxtref{\oxcprov^\prime}{\oxmuta^\prime}{\oxtype^\prime}} \in \oxvarctx \;
  \wedge
  (\forall \oxvarctxentry{\oxplace^\prime}{\oxtref{\oxcprov^\prime}{\oxmuta^\prime}{\oxtype^\prime}} \in \oxvarctx. \;
  \oxplace^\prime \in \oxset{
    \overline{\oxplace_e}
  }))$. Fortunately, from $\oxdomain{\oxvarctx} = \oxdomain{\oxvarctx^\prime}$, we know that
  for every $\oxcprov^\prime \in \oxdomain{\oxvarctx}$, $\oxcprov^\prime \in
  \oxdomain{\oxvarctx^\prime}$, and further that $\oxvarctx^\prime(\oxcprov^\prime)
  \subseteq \oxvarctx(\oxcprov^\prime)$. Thus, for each $\oxcprov^\prime$, we know there are only
  potentially fewer loans to show if the obligation was met using the clause of
  $\forall \oxloanpkg{\oxmuta^\prime}{\oxplaceexpr} \in \oxset{\overline{\oxloan}}.
  (\oxmuta = \oxmut \vee \oxmuta^\prime = \oxmut)
  \implies \oxrelevant{\oxplace^\prime}{\oxplace}$. If the obligation was met using the other
  clause, note that $\oxvarctx$ and $\oxvarctx^\prime$ can only differ in the loan sets they
  associate with any given region and so the exact fact must still be true for
  $\oxvarctx^\prime$.
\end{proof}

\begin{oxlemma}{Expressions Remain Well-Typed in More Precise Environments}
  {lemma:expressions-typing-more-precise}\hfill\\ If $\oxdomain{\oxvarctx} =
  \oxdomain{\oxvarctx^\prime}$, and $\forall \oxcprov.$ $\oxvarctx^\prime(\oxcprov)
  \subseteq \oxvarctx(\oxcprov)$ and 
  $\oxtypjudge{\oxglobalctx}{\oxemptyctx}{\oxkontctx}{\oxvarctx}
  {\oxexpr}{\oxtype}{\oxvarctx_f}$ then
  $\oxtypjudge{\oxglobalctx}{\oxemptyctx}{\oxkontctx}{\oxvarctx^\prime}
  {\oxexpr}{\oxtype}{\oxvarctx_f^\prime}$ and
  $\oxdomain{\oxvarctx_f} = \oxdomain{\oxvarctx_f^\prime}$, and
  $\forall \oxcprov.$ $\oxvarctx_f^\prime(\oxcprov) \subseteq \oxvarctx_f(\oxcprov)$
\end{oxlemma}

\begin{proof}
  We proceed by induction over the expression typing.

  \oxname{T-Move}, \oxname{T-Copy}, and \oxname{T-Borrow} all follow from the
  \Lemma{lemma:ownership-safety-more-precise} and the \Lemma{lemma:computety-more-precise}.
  
  \oxname{T-BorrowIndex}, \oxname{T-BorrowSlice}, and \oxname{T-IndexCopy} all follow from the
  induction hypothesis, \Lemma{lemma:ownership-safety-more-precise}, and the
  \Lemma{lemma:computety-more-precise}.

  \oxname{T-Seq} follows from the observation that garbage collecting loans will
  preserve subsets and clear in exactly both or neither, and from the induction
  hypothesis.

  \oxname{T-Branch} and \oxname{T-Match} follow from applying the induction hypothesis,
  \Lemma{lemma:subtyping-subset-related-envs},
  \Lemma{lemma:union-subset-related-envs}.

  \oxname{T-Let} follows from the same observation about garbage collection in
  the \oxname{T-Seq} case, the induction hypothesis, and
  \Lemma{lemma:subtyping-subset-related-envs}.

  \oxname{T-Assign} follows from the induction hypothesis,
  \Lemma{lemma:ownership-safety-more-precise},
  \Lemma{lemma:union-subset-related-envs}.

  \oxname{T-AssignDeref} both follow from the induction hypothesis, follows from
  the induction hypothesis, \Lemma{lemma:ownership-safety-more-precise},
  \Lemma{lemma:union-subset-related-envs}, and
  \Lemma{lemma:computety-more-precise}.

  \oxname{T-AppFunction} follows from the induction hypothesis, and a few pieces
  about the well-formedness of the instantiations happening in
  \oxname{T-AppFunction} (namely, frame expressions, regions, and types). For
  the frame expression validity, we consider each case and note that
  \oxname{WF-EnvVar} depends only on $\oxtvarctx$ which is unchanged and that
  \oxname{WF-Env} appeals to stack typing validity and so it suffices to show
  that that still holds in our more precise environment which we do by appealing
  to \Lemma{lemma:wf-varctx-more-precise}. For the region validity, there are
  again two cases to consider \oxname{WF-LocalProv} which applies if the domain
  of $\oxvarctx$ is the same as the domain of $\oxvarctx^\prime$ which we have
  directly from our premise and \oxname{WF-AbstractProv} which depends only on
  $\oxtvarctx$ which is unchanged. For the type validity, we appeal to
  \Lemma{lemma:wf-type-more-precise}.

  \oxname{T-AppClosure} follows from the induction hypothesis and
  \Lemma{lemma:subtyping-subset-related-envs}.

  \oxname{T-LetRegion}, \oxname{T-While}, \oxname{T-ForArray},
  \oxname{T-ForSlice}, \oxname{T-Closure}, \oxname{T-Tuple}, \oxname{T-Array},
  \oxname{T-Slice}, \oxname{T-Drop}, \oxname{T-Left}, \oxname{T-Right} follow
  immediately from the induction hypothesis.

  \oxname{T-Function}, \oxname{T-Abort}, \oxname{T-Unit}, \oxname{T-U32},
  \oxname{T-True}, and \oxname{T-False} are immediate.

\end{proof}

\subsection{Preservation under Safe Loan Updates Lemmas}
\label{sec:safe-loan-lemmas}

\begin{oxlemma}{Ownership Safety Produces Non Conflicting Loans}
  {lemma:ownership-safety-no-conflicts}
  If
  \begin{enumerate}
  \item $\oxplaceexprctx_{\oxcprov_p}[\oxplaceexpr_{\oxcprov_p}] \in \oxset{\oxloans}$,
    where $\oxcprov_p \mapsto \oxset{\oxloans} \in \oxloanmappings{\oxvarctx}{\oxkontctx}$
  \item $\oxvarctx(\oxcprov_b) = \emptyset$
  \item $\oxnotreborrowed{\oxvarctx[\oxcprov_b \mapsto \oxset{\oxloans_b}]}{\oxcprov_p}$
  \item $\oxmusafetyinner{\oxemptyctx}{\oxkontctx}{\oxvarctx}{\oxmuta}{\overline{\oxplace_b}}{\oxplaceexpr_b}{\oxset{\oxloans_b}}$
  \end{enumerate}
  then $\forall \oxloanpkg{\oxmuta}{\oxplaceexpr^\prime} \in \oxset{\oxloans_b}$.
  $\oxrelevant{\oxplaceexpr_{\oxcprov_p}}{\oxplaceexpr^\prime}$.
\end{oxlemma}
\begin{proof}
  Proceed by induction on the ownership safety judgement for $\oxplaceexpr$ in
  the premise.
 
  \vspace{1em}
  \noindent \framebox[\textwidth]{ \figuresize
    \begin{mathpar}
      \OSafePlace
    \end{mathpar}
  } \vspace{1em}

  We need to show that $\oxrelevant{\oxplaceexpr_{\oxcprov_p}}{\oxplace}$.
  Since $\oxcprov_p \in \oxdomain{\oxloanmappings{\oxvarctx}{\oxkontctx}}$, we know from the hypothesis of
  \oxname{O-SafePlace} that either $\oxcprov_p$ is excluded, or all loans in it
  are disjoint from $\oxplace$.

  $\oxcprov_p$ cannot have been exluded because
  $\oxnotreborrowed{\oxvarctx}{\oxcprov_p}$.

  This just leaves the case where all loans in $\oxvarctx(\oxcprov_p)$ are
  disjoint from $\oxplace$. Let $\oxplace_{\oxcprov_p}$ be the inner place of
  $\oxplaceexpr_{\oxcprov_p}$. More formally, $\oxplaceexpr_{\oxcprov_p} =
  \oxplaceexprctx[\oxplace_{\oxcprov_p}]$. By this disjointness, we know
  $\oxrelevant{\oxplaceexprctx_{\oxcprov_p}[\oxplaceexprctx[\oxplace_{\oxcprov_p}]]}{\oxplace}$,
  which directly implies
  $\oxrelevant{\oxplaceexprctx[\oxplace_{\oxcprov_p}]}{\oxplace}$, which is what
  we wanted to show.
  
  \vspace{1em}
  \noindent \framebox[\textwidth]{ \figuresize
    \begin{mathpar}
      \ODeref
    \end{mathpar}
  } \vspace{1em}

  We need to show that $\oxrelevant{\oxplaceexprctx[\oxderef
      \oxplace]}{\oxplaceexpr_{\oxcprov_p}}$, and that $\forall i,
  \ \oxloanpkg{\oxmuta}{\oxplaceexpr} \in
  \oxset{\overline{\oxloanpkg{\oxmuta}{\oxplaceexpr^\prime_i}}}$. $\oxrelevant{\oxplaceexpr_{\oxcprov_p}}{\oxplaceexpr}$.

  The latter we get from applying the induction hypothesis.

  The former follows from the same reasoning as in the previous case. 

  \vspace{1em}
  \noindent \framebox[\textwidth]{ \figuresize
    \begin{mathpar}
      \ODerefAbs
    \end{mathpar}
  } \vspace{1em}

  Since $\oxtvarctx = \oxemptyctx$, there are no valid reference types that have
  an abstract region, meaning the first hypothesis is a contradiction.
\end{proof}

\begin{oxlemma}{Ownership Safety is Preserved under Safe Loan Updates}
  {lemma:ownership-safety-loan-update} If
  \begin{enumerate}
  \item $\oxmusafety{\oxemptyctx}{\oxkontctx}{\oxvarctx}{\oxmuta_b}{\oxplaceexpr_b}{\oxset{\oxloans_b}}$
  \item $\oxmusafetyinner{\oxemptyctx}{\oxkontctx}{\oxnewframe{\oxvarctx}{\oxframe}}{\oxmuta}
     {\overline{\oxplace_1}}{\oxplaceexpr}{\oxset{\oxloans}}$
   \item and $\oxlookup{\oxvarctx}{\oxcprov_b}{\{\}}$
   \item $\oxnotreborrowed{\oxvarctx[\oxcprov_b \mapsto \oxset{\oxloans_b}]}{\oxcprov_p}$
  \item $\oxplace_1 = \oxplace_2$ or $\oxplace_1 = \oxplace_2 \cup \oxset{\oxplace
    \ | \ \oxplaceexprctx[\oxderef\oxplace] \in
    \oxset{\oxloans_b}}$
  \item either $\oxrootof{\oxplaceexpr} \in \oxdomain{\oxframe}$ or
  $\exists \oxcprov_{p} \in \oxdomain{\oxnewframe{\oxvarctx}{\oxframe}}, \oxplaceexprctx, \oxplaceexpr^\prime.$
  \item $\oxplaceexpr = \oxplaceexprctx[\oxplaceexpr^\prime]$
  \item $\oxplaceexpr^\prime \in \oxnewframe{\oxvarctx}{\oxframe}(\oxcprov_p)$
  \item $\forall \oxcprov \in \oxdomain{\oxframe},
    \oxloanpkg{\oxmuta}{\oxplaceexpr} \in \oxframe(\oxcprov).$ either
    $\oxrootof{\oxplaceexpr} \in \oxdomain{\oxframe}$, or $\exists \oxcprov^\prime \mapsto \oxset{\oxloans} \in
    \oxloanmappings{\oxvarctx}{\oxkontctx}.$ $\oxloanpkg{\oxmuta}{\oxplaceexpr} \in \oxset{\oxloans}$.

      (In english, loans in the closure's frame come from the current frame or a closure's captured frame in $\oxvarctx$ or $\oxkontctx$)
  \end{enumerate}
  then $\oxmusafetyinner{\oxemptyctx}{\oxkontctx}{\oxnewframe{\oxvarctx[\oxcprov_b \mapsto
        \oxloans_b]}{\oxframe}}{\oxmuta}{\overline{\oxplace_2}}{\oxplaceexpr}{\oxset{\oxloans^\prime}}$.
\end{oxlemma}

\begin{proof}
  Proceed by induction on the ownership safety judgement for $\oxplaceexpr$ in
  the premise.
 
  \vspace{1em}
  \noindent \framebox[\textwidth]{ \figuresize
    \begin{mathpar}
      \OSafePlace
    \end{mathpar}
  } \vspace{1em}

  Let $\oxcprov^\prime$ be an arbitrary region. If the disjunction was proven
  using the right part, which talks about the exclusion list, then we can prove
  it again the same way, since none of the types are changed and the exclusion
  list $\overline{\oxplace_2}$ includes all of $\overline{\oxplace_1}$ in either case.

  If the disjunction was proven using the left part, the only interesting case
  is when $\oxcprov^\prime = \oxcprov_b$. If $\oxrootof{\oxplace} \in \oxdomain{\oxframe}$,
  then we're done, because all loans in $\oxset{\oxloans_b}$ are disjoint just
  by the well formedness of the environment $\oxvarctx$.

  Otherwise, $\oxrootof{\oxplace} \in \oxdomain{\oxvarctx}$ and $\exists
  \oxcprov_p$ such that $\oxplace \in
  \oxnewframe{\oxvarctx}{\oxframe}(\oxcprov_p)$, and we want to show that
  $\forall \oxloanpkg{\oxmuta^\prime}{\oxplaceexprctx[\oxplace^\prime]} \in
  \oxset{\oxloans_b}$, $\oxrelevant{\oxplace}{\oxplace^\prime}$.

  If $\oxcprov_p \in \oxdomain{\oxframe}$, then we're done because each loan in
  $\oxcprov_p$ either comes from $\oxdomain{\oxframe}$, in which case
  disjointness is immediate, or it comes from a loan mapping in $\oxvarctx$ or
  $\oxkontctx$, in which case we can apply
  \Lemma{lemma:ownership-safety-no-conflicts} to finish the proof.

  Otherwise, $\oxcprov_p \in \oxdomain{\oxvarctx}$. By the well formedness 
  of $\oxnewframe{\oxvarctx[\oxcprov_b \mapsto \oxset{\oxloans_b}]}{\oxframe}$,
  $\oxnotreborrowed{\oxvarctx[\oxcprov_b \mapsto
      \oxset{\oxloans_b}]}{\oxcprov^\prime}$. Given all of this we can apply
  \Lemma{lemma:ownership-safety-no-conflicts} to finish the proof.

  \vspace{1em}
  \noindent \framebox[\textwidth]{ \figuresize
    \begin{mathpar}
      \ODeref
    \end{mathpar}
  } \vspace{1em}

  Firstly, we would like to apply our induction hypothesis. In order to do so,
  we need to show that the new \texttt{excl} is either the same, or only has
  places from $\oxset{\oxloans_b}$ added. If $\oxcprov \neq \oxcprov_b$, then
  $\oxnewframe{\oxvarctx}{\oxframe}(\oxcprov) = \oxnewframe{\oxvarctx[\oxcprov_b
      \mapsto \oxset{\oxloans_b}]}{\oxframe}$, so the exclusion list is the
  same. If $\oxcprov = \oxcprov_b$, then \texttt{excl} was empty, and now
  includes the places $\oxset{\oxplace \ | \ \oxplaceexprctx[\oxderef\oxplace]
    \in \oxset{\oxloans_b}}$. We also need to show that either that either the
  place expression is in the domain of $\oxframe$ or that it has a sub place
  expression in a loan set, but this is immediate from our hypotheses. So in
  either case we satisfy the necessary hypothesis and can apply the induction
  hypothesis.

  Whats left to show is the disjointness or exclusion condition, which follows
  identically to the reasoning in the previous case.

  \vspace{1em}
  \noindent \framebox[\textwidth]{ \figuresize
    \begin{mathpar}
      \ODerefAbs
    \end{mathpar}
  } \vspace{1em}

  Since $\oxtvarctx = \oxemptyctx$, there are no valid reference types that have
  an abstract region, meaning the first hypothesis is a contradiction.
\end{proof}

\begin{oxlemma}{Ownership Safety is Preserved after Environment Union}
  {lemma:ownership-safety-environment-union} If
  $\oxmusafetyinner{\oxemptyctx}{\oxkontctx}{\oxvarctx_1}{\oxmuta}{\overline{\oxplace_1}}{\oxplaceexpr}{\oxset{\oxloans}}$
  and
  $\oxmusafetyinner{\oxemptyctx}{\oxkontctx}{\oxvarctx_2}{\oxmuta}{\overline{\oxplace_2}}{\oxplaceexpr}{\oxset{\oxloans^\prime}}$
  then $\oxmusafetyinner{\oxemptyctx}{\oxkontctx}{\oxvarctx_1 \oxintersect
    \oxvarctx_2}{\oxmuta}{\overline{\oxplace_1},\overline{\oxplace_2}}{\oxplaceexpr}{\oxset{\oxloans^{\prime\prime}}}$
  and $\oxplace_1 = \oxplace_2$ or $\oxplace_1 = \oxplace_2 \cup \oxset{\oxplace
    \ | \ \oxplaceexprctx[\oxderef\oxplace] \in
    \oxset{\overline{\oxloans_b}}}$.
\end{oxlemma}
\begin{proof}
  Proceed by induction on the ownership safety judgements in the premise. Note
  that they both have the same sequence of proof rule applications, because the
  judgement is inductive over $\oxplaceexpr$.
 
  \vspace{1em}
  \noindent \framebox[\textwidth]{ \figuresize
    \begin{mathpar}
      \OSafePlace
    \end{mathpar}
  } \vspace{1em}

  Let $\oxcprov$ be an arbitrary region.

  If either the derivation with
  $\oxvarctx_1$ or $\oxvarctx_2$ used the second part of the disjunction, we can
  proceed by the second part of the disjunction.

  Otherwise, both derivations used the first part of the disjunction. Since
  $\oxvarctx_1 \oxintersect \oxvarctx_2(\oxcprov) = \oxvarctx_1(\oxcprov) \cup
  \oxvarctx_2(\oxcprov)$, we can just combine these two facts and proceed by the
  first part of the disjunction.

  \vspace{1em}
  \noindent \framebox[\textwidth]{ \figuresize
    \begin{mathpar}
      \ODeref
    \end{mathpar}
  } \vspace{1em}

  In order to apply the induction hypothesis, we need show that our new
  \texttt{excl} is the union of the two from the derivations. This is immediate
  because $\oxvarctx_1 \oxintersect \oxvarctx_2(\oxcprov) =
  \oxvarctx_1(\oxcprov) \cup \oxvarctx_2(\oxcprov)$.

  To finish the case, follow the same reasoning from the previous case for the
  disjunction.


  \vspace{1em}
  \noindent \framebox[\textwidth]{ \figuresize
    \begin{mathpar}
      \ODerefAbs
    \end{mathpar}
  } \vspace{1em}

  Since $\oxtvarctx = \oxemptyctx$, there are no valid reference types that have
  an abstract region, meaning the first hypothesis is a contradiction.

\end{proof}

\begin{oxlemma}{Ownership Safety is Preserved after Type Checking a Closure Body}
  {lemma:ownership-safety-closure-body} If
  \begin{enumerate}
  \item $\oxmusafety{\oxemptyctx}{\oxkontctx}{\oxvarctx}{\oxmuta}{\oxplaceexpr_b}{\oxset{\oxloans_b}}$
  \item $\oxlookup{\oxvarctx}{\oxcprov}{\{\}}$ 
  \item $\forall \oxid \in \oxfvars{\oxexpr}.$ $\oxid \in \oxdomain{\oxframe}$ 
  \item $\forall \oxcprov \in \oxfprovs{\oxexpr}.$ $\oxcprov \in \oxdomain{\oxframe}$
  \item $\forall \oxcprov \in \oxdomain{\oxframe},
    \oxloanpkg{\oxmuta}{\oxplaceexpr} \in \oxframe(\oxcprov).
    \oxrootof{\oxplaceexpr} \in \oxdomain{\oxframe} \vee
    \exists \oxcprov^\prime \mapsto \oxset{\oxloans} \in
    \oxloanmappings{\oxvarctx}{\oxkontctx}.$ $\oxloanpkg{\oxmuta}{\oxplaceexpr} \in \oxset{\oxloans}$.
    
    (In English, loans in the closure's frame come from the current frame or a closure's captured
    frame in $\oxvarctx$ or $\oxkontctx$)
  \item $\oxtypjudge{\oxglobalctx}{\oxemptyctx}{\oxkontctx}{\oxnewframe{\oxvarctx}{\oxframe}}
    {\oxexpr}{\oxsitype}{\oxnewframe{\oxvarctx_o}{\oxframe_o}}$
  \end{enumerate}
  then
  \begin{enumerate}
  \item $\oxmusafety{\oxemptyctx}{\oxkontctx}{\oxvarctx_o}{\oxmuta}{\oxplaceexpr_b}{\oxset{\oxloans_b^\prime}}$
  \item $\forall \oxcprov \in \oxdomain{\oxframe_o},
    \oxloanpkg{\oxmuta}{\oxplaceexpr} \in \oxframe_o(\oxcprov). \oxrootof{\oxplaceexpr} \in \oxdomain{\oxframe_o} \vee \exists \oxcprov^\prime \mapsto \oxset{\oxloans} \in
    \oxloanmappings{\oxvarctx_o}{\oxkontctx}.$ $\oxloanpkg{\oxmuta}{\oxplaceexpr} \in \oxset{\oxloans}$.
    
    (In English, loans in the closure's frame come from the current frame or a closure's captured
    frame in $\oxvarctx_o$ or $\oxkontctx$)
  \end{enumerate}
\end{oxlemma}
\begin{proof}
  In the \oxname{T-Move} case, the context is updated, but since $\oxplace \in
  \oxdomain{\oxframe}$, $\oxvarctx_o = \oxvarctx$, so the conclusions follow
  from the premises.

  In the \oxname{T-Copy}, \oxname{T-Function}, \oxname{T-Abort},
  \oxname{T-Unit}, \oxname{T-u32}, \oxname{T-True}, and \oxname{T-False} cases,
  $\oxvarctx_o = \oxvarctx$ so the conclusions follow from the premises.

  In the \oxname{T-IndexCopy}, \oxname{T-LetRegion}, \oxname{T-While},
  \oxname{T-ForArray}, \oxname{T-ForSlice}, \oxname{T-Closure},
  \oxname{T-Tuple}, \oxname{T-Array}, \oxname{T-Slice}, \oxname{T-Drop},
  \oxname{T-Left}, and \oxname{T-Right} cases, the proof is immediate from the
  induction hypothesis.

  In the \oxname{T-Borrow}, \oxname{T-BorrowIndex}, and \oxname{T-BorrowSlice}
  cases, the proof follows from the induction hypothesis and
  \Lemma{lemma:ownership-safety-loan-update}. The last condition is the
  restriction on loans in $\oxframe$, but this is immediate by inspection of the
  ownership safety judgement. The only way to create new loans is to directly
  borrow a place, in which case we'd have that the loan is in the domain of
  $\oxframe$, and otherwise the loans originated from the loan set in
  $\oxnewframe{\oxvarctx}{\oxframe}$ of a reference being reborrowed, and we
  already know the property for $\oxnewframe{\oxvarctx}{\oxframe}$.

  In the \oxname{T-Seq} case, the proof follows from the induction hypothesis
  and \Lemma{lemma:ownership-safety-related-envs} (note that
  \texttt{gc-loans} does not change any types in the environment and produces
  related environments).

  In the \oxname{T-Branch} and \oxname{T-Match} cases, the proof follows from
  the induction hypothesis, \Lemma{lemma:subtyping-ownership-safety}, and
  \Lemma{lemma:ownership-safety-environment-union}. We also need to show that
  rewriting preserves the restriction on loans in $\oxframe$, but this is
  immediate because the most that rewriting can do is union together loan sets.
  The rest of the cases that use rewriting also use this same reasoning.

  In the \oxname{T-Let} case, the proof follows from the induction hypothesis,
  \Lemma{lemma:subtyping-ownership-safety}, and
  \Lemma{lemma:ownership-safety-related-envs} (note that
  \texttt{gc-loans} does not change any types in the environment and produces
  related environments).

  In the \oxname{T-AssignDeref} case the proof follows from the induction
  hypothesis, and \Lemma{lemma:subtyping-ownership-safety}.

  In the \oxname{T-Assign} case, the proof follows from the induction
  hypothesis, \Lemma{lemma:subtyping-ownership-safety}, and
  \Lemma{lemma:ownership-safety-assignment}.

  In the \oxname{T-App} case, the proof follows from the induction hypothesis
  and \Lemma{lemma:subtyping-ownership-safety}.
\end{proof}

\begin{oxlemma}{Outlives is Preserved under Safe Loan Updates}
  {lemma:outlives-loan-update} If
  \begin{enumerate}
  \item $\oxmusafety{\oxemptyctx}{\oxkontctx}{\oxvarctx}{\oxmuta}{\oxplaceexpr}{
    \oxset{\overline{\oxloan}} }$
  \item $\oxlookup{\oxvarctx}{\oxcprov}{\{\}}$
  \item $\oxcprov \not \in \oxdomain{\oxframe}$
  \item $\oxcprov \not = \oxcprov_1$ and $\oxcprov \not = \oxcprov_2$
  \item $\forall \oxtype^\prime \in \oxdomain{\oxvarctx}.$ $\oxcprov_1$ and
    $\oxcprov_2$ does not occur outside of a closure in $\oxtype^\prime$
  \item $\oxrunify{\oxemptyctx}{\oxkontctx}{
    \oxnewframe{\oxvarctx}{\oxframe}
  }{\oxcprov_1}{\oxcprov_2}{
    \oxnewframe{\oxvarctx^\prime}{\oxframe^\prime}
  }$
  \end{enumerate}
  then
  $\oxrunify{\oxemptyctx}{\oxkontctx}{
    \oxnewframe{\oxvarctx[\oxcprov \mapsto \oxset{\oxloans_b}]}{\oxframe}
  }{\oxcprov_1}{\oxcprov_2}{
    \oxnewframe{\oxvarctx^{\prime}[\oxcprov \mapsto \oxset{\oxloans_b}]}{\oxframe^{\prime}}
  }$
\end{oxlemma}

\begin{proof}
  Proceed by induction on the outlives relation. The only interesting cases are
  \oxname{OL-CombineConcrete}, \oxname{OL-CombineConcreteUnrestricted}, and
  \oxname{OL-CheckConcrete}. They all proceed similarly. The closure restriction
  is immediate from the premise because no types are changed. The region not
  reborrowed restriction follows from the fact that $\oxcprov_1$ and
  $\oxcprov_2$ don't occur in any types in $\oxvarctx$ outside of a closure,
  which means there's no place in the domain of $\oxvarctx$ for there to be a
  reborrow of in $\oxloans$.
\end{proof}

\begin{oxlemma}{Rewriting is Preserved under Safe Loan Updates}
  {lemma:rewriting-loan-update} If
  \begin{enumerate}
  \item $\oxmusafety{\oxemptyctx}{\oxkontctx}{\oxvarctx}{\oxmuta}{\oxplaceexpr}{
    \oxset{\overline{\oxloan}} }$
  \item $\oxlookup{\oxvarctx}{\oxcprov}{\{\}}$
  \item $\oxcprov \not \in \oxdomain{\oxframe}$
  \item $\forall \oxtype \in \oxframe$. $\oxcprov$ does not occur in $\oxtype$ or $\oxtype_1$ or $\oxtype_2$
  \item $\forall \oxcprov \in \oxfprovs{\oxtype_1} \cup \oxfprovs{\oxtype_2}$.
    $\forall \oxtype^\prime \in \oxdomain{\oxvarctx}.$ $\oxcprov$ does not occur
    outside of a closure in $\oxtype^\prime$.
  \item $\oxtunify{\oxemptyctx}{\oxkontctx}{
    \oxnewframe{\oxvarctx}{\oxframe}
  }{\oxtype_1}{\oxtype_2}{
    \oxnewframe{\oxvarctx^\prime}{\oxframe^\prime}
  }$
  \end{enumerate}
  then
  $\oxtunify{\oxemptyctx}{\oxkontctx}{
    \oxnewframe{\oxvarctx[\oxcprov \mapsto \oxset{\oxloans_b}]}{\oxframe}
  }{\oxtype_1}{\oxtype_2}{
    \oxnewframe{\oxvarctx^{\prime}[\oxcprov \mapsto \oxset{\oxloans_b}]}{\oxframe^{\prime}}
  }$.

\end{oxlemma}
\begin{proof}
  Proceed by induction on the rewriting derivation. The only interesting case is
  \oxname{RR-Reference}, in which case we apply
  \Lemma{lemma:outlives-loan-update} and the induction hypothesis. The other
  cases all follow immediately or from the induction hypothesis.
\end{proof}

\begin{oxlemma}{Closure Bodies are Well-Typed under Safe Loan Updates}
  {lemma:exprs-well-typed-loan-update}
  If
  \begin{enumerate}
  \item $\oxmusafety{\oxtvarctx}{\oxkontctx}{\oxvarctx}{\oxmuta}{\oxplaceexpr}{
    \oxset{\overline{\oxloan}}
  }$
  \item $\oxlookup{\oxvarctx}{\oxcprov}{\{\}}$
  \item $\oxcprov \not \in \oxdomain{\oxframe}$
  \item $\forall \oxtype \in \oxframe$. $\oxcprov$ does not occur in $\oxtype$
  \item $\forall \oxid \in \oxfvars{\oxexpr}.$ $\oxid \in \oxdomain{\oxframe}$
  \item $\forall \oxcprov \in \oxfprovs{\oxexpr}.$ $\oxcprov \in \oxdomain{\oxframe}$
  \item $\forall \oxcprov \in \oxdomain{\oxframe},
    \oxloanpkg{\oxmuta}{\oxplaceexpr} \in \oxframe(\oxcprov).
    \oxrootof{\oxplaceexpr} \in \oxdomain{\oxframe} \vee
    \exists \oxcprov^\prime \mapsto \oxset{\oxloans} \in
    \oxloanmappings{\oxvarctx}{\oxkontctx}.$ $\oxloanpkg{\oxmuta}{\oxplaceexpr} \in \oxset{\oxloans}$.
    
    (In English, loans in the closure's frame come from the current frame or a closure's captured
    frame in $\oxvarctx$ or $\oxkontctx$)
  \item $\forall \oxcprov \in \oxfprovs{\oxsitype}.\ \forall \oxtype^\prime \in
    \oxdomain{\oxvarctx}.\ \oxcprov\text{ does not occur in }\oxtype^\prime$
  \item $\oxtypjudge{\oxglobalctx}{\oxemptyctx}{\oxkontctx}{\oxnewframe{\oxvarctx}{\oxframe}}
    {\oxexpr}{\oxsitype}{\oxnewframe{\oxvarctx_o}{\oxframe_o}}$
  \item $\oxvarctxwellformed{\oxglobalctx}{\oxemptyctx}{\oxnewframe{
      \oxvarctx[\oxcprov \mapsto \oxset{\oxloans}]
    }{\oxframe}}$
  \end{enumerate}
  then
  \begin{enumerate}
  \item $\oxtypjudge{\oxglobalctx}{\oxemptyctx}{\oxkontctx}{
    \oxnewframe{
      \oxupdate{\oxvarctx}{
        \oxloanctxentry{\oxcprov}{ \oxset{\overline{\oxloan}} }
      }
    }{\oxframe}
  }{\oxexpr}{\oxsitype}{
    \oxnewframe{
      \oxupdate{\oxvarctx_o}{
        \oxloanctxentry{\oxcprov}{\oxset{\overline{\oxloan}}
        }
      }
    }{\oxframe_o}
  }$.
  \item $\forall \oxcprov \in \oxdomain{\oxframe_o},
    \oxloanpkg{\oxmuta}{\oxplaceexpr} \in \oxframe_o(\oxcprov).
    \oxrootof{\oxplaceexpr} \in \oxdomain{\oxframe_o} \vee
    \exists \oxcprov^\prime \mapsto \oxset{\oxloans} \in
    \oxloanmappings{\oxvarctx_o}{\oxkontctx}.$ $\oxloanpkg{\oxmuta}{\oxplaceexpr} \in \oxset{\oxloans}$.
    
    (In English, loans in the closure's frame come from the current frame or a closure's captured
    frame in $\oxvarctx_o$ or $\oxkontctx$)
  \end{enumerate}
\end{oxlemma}
\begin{proof}
  Proceed by induction over the typing derivation for $\oxexpr$.

  In \oxname{T-Move} and \oxname{T-Copy}, and \oxname{T-Borrow} cases, we can apply
  \Lemma{lemma:ownership-safety-loan-update} and note that type computation is
  unaffected by changes in loan sets.

  In the \oxname{T-Borrow} case, we can apply
  \Lemma{lemma:ownership-safety-loan-update} and note that type computation is
  unaffected by changes in loan sets. The last condition is the restriction on
  loans in $\oxframe$, but this is immediate by inspection of the ownership
  safety judgement. The only way to create new loans is to directly borrow a
  place, in which case we'd have that the loan is in the domain of $\oxframe$,
  and otherwise the loans originated from the loan set in
  $\oxnewframe{\oxvarctx}{\oxframe}$ of a reference being reborrowed, and we
  already know the loan set restriction for $\oxnewframe{\oxvarctx}{\oxframe}$.
  The rest of the cases that involve borrowing use similar reasoning.

  In the \oxname{T-BorrowIndex} and \oxname{T-IndexCopy} cases, we can apply
  \Lemma{lemma:ownership-safety-loan-update}, the induction hypothesis, and the
  note about type computation.

  In the \oxname{T-BorrowSlice} case, we can apply
  \Lemma{lemma:ownership-safety-loan-update}, the induction hypothesis, the note
  about type computation, and \Lemma{lemma:ownership-safety-closure-body}.

  In the \oxname{T-Seq} case, we can apply the induction hypothesis, for which
  we need to apply \Lemma{lemma:ownership-safety-closure-body} and use the fact
  that garbage collection produces related environments with
  \Lemma{lemma:ownership-safety-related-envs}.

  In the \oxname{T-Branch} case we can apply the induction hypothesis,
  \Lemma{lemma:ownership-safety-closure-body}, and
  \Lemma{lemma:rewriting-loan-update}. We also need to show that
  rewriting preserves the restriction on loans in $\oxframe$, but this is
  immediate because the most that rewriting can do is union together loan sets.
  The rest of the cases that use rewriting also use this same reasoning.

  In the \oxname{T-Let} case, we can apply the induction hypothesis,
  \Lemma{lemma:ownership-safety-closure-body},
  \Lemma{lemma:rewriting-loan-update}, and the fact that garbage collection
  produces related environments with
  \Lemma{lemma:ownership-safety-related-envs}. The last obligation is the region
  not reborrowed condition. Note that by environment well formedness and our
  hypothesis, any free regions in the types are not in any non closure types.
  Therefore, there are no places in $\oxvarctx$ with the region for the new
  loans in $\oxloans$ to even contain a reborrow of, meaning the regions are not
  reborrowed. The other cases which require region not reborrowed proceed by the
  same reasoning.

  In the \oxname{T-LetRegion}, \oxname{T-While}, \oxname{T-Closure},
  \oxname{T-Tuple}, \oxname{T-Array}, \oxname{T-Slice}, \oxname{T-Left}, and
  \oxname{T-Right} cases, the proofs follows from the induction hypothesis and
  \Lemma{lemma:ownership-safety-closure-body}.

  In the \oxname{T-AssignDeref} case, we can apply the induction hypothesis,
  \Lemma{lemma:ownership-safety-closure-body}, the fact that type computation is
  unaffected by loan updates, \Lemma{lemma:rewriting-loan-update}, and
  \Lemma{lemma:ownership-safety-loan-update}.

  In the \oxname{T-Assign} case, we can apply the induction hypothesis,
  \Lemma{lemma:ownership-safety-closure-body}, the fact that type computation is
  unaffected by loan updates, \Lemma{lemma:rewriting-loan-update}, and
  \Lemma{lemma:ownership-safety-loan-update}. The last obligation is the unique
  to judgement, which is unaffected by loan updates.

  In the \oxname{T-ForArray} and \oxname{T-ForSlice} cases, we can apply the
  induction hypothesis and \Lemma{lemma:ownership-safety-closure-body}. The
  remaining region not reborrowed obligation follows from the same reasoning in
  the \oxname{T-Let} case.

  In the \oxname{T-Function}, \oxname{T-Abort}, \oxname{T-Unit}, \oxname{T-u32},
  \oxname{T-True}, and \oxname{T-False} cases, the proof is immediate.

  In the \oxname{T-App} case, we can apply the induction hypothesis,
  \Lemma{lemma:ownership-safety-closure-body} and
  \Lemma{lemma:rewriting-loan-update}. Note that the well formedness judgements
  are unaffected by loan updates. The last obligation is the region not
  reborrowed judgement, follows from the same reasoning in
  the \oxname{T-Let} case.

  In the \oxname{T-Drop} case, we can apply the induction hypothesis,
  \Lemma{lemma:ownership-safety-closure-body}, and
  \Lemma{lemma:ownership-safety-related-envs}, noting that making the type of a
  place dead produces a related environment.

  In the \oxname{T-Match} case we can apply the induction hypothesis,
  \Lemma{lemma:ownership-safety-closure-body}, and
  \Lemma{lemma:rewriting-loan-update}. The last obligation is the region not
  reborrowed judgement, follows from the same reasoning in
  the \oxname{T-Let} case.
\end{proof}

\begin{oxlemma}{Values are Well-Typed under Safe Loan Updates}
  {lemma:values-well-typed-loan-update} If
  $\oxmusafety{\oxemptyctx}{\oxkontctx}{\oxvarctx}{\oxmuta}{\oxplaceexpr}{
    \oxset{\overline{\oxloan}} }$ and
  $\oxnotinclosure{\oxkontctx}{\oxvarctx}{\oxcprov}$ and
  $\oxlookup{\oxvarctx}{\oxcprov}{\emptyset}$ and
  $\oxtypjudge{\oxglobalctx}{\oxemptyctx}{\oxkontctx}{\oxvarctx}
  {\oxvalue}{\oxsitype}{\oxvarctx}$, then
  $\oxtypjudge{\oxglobalctx}{\oxemptyctx}{\oxkontctx}{ \oxupdate{\oxvarctx}{
      \oxloanctxentry{\oxcprov}{ \oxset{\overline{\oxloan}} } }
  }{\oxvalue}{\oxsitype}{ \oxupdate{\oxvarctx}{ \oxloanctxentry{\oxcprov}{
        \oxset{\overline{\oxloan}} } } }$.
\end{oxlemma}

\begin{proof}
  We proceed by induction on the value typing relation.

  For \oxname{T-u32}, \oxname{T-True}, \oxname{T-False}, the result is
  immediate.

  For \oxname{T-Tuple} and \oxname{T-Array}, we apply the induction
  hypothesis to each entry.

  \vspace{1em}
  \noindent \framebox[\textwidth]{
    \figuresize
    \begin{mathpar}
      \TPointer
    \end{mathpar}
  }
  \vspace{1em}

  In the \oxname{T-Pointer} case, both judgements in the premise are unaffected
  by taking an empty loan set and adding loans to it, so the case is immediate.

  \vspace{1em}
  \noindent \framebox[\textwidth]{
    \figuresize
    \begin{mathpar}
      \TClosureVal
    \end{mathpar}
  }
  \vspace{1em}

  In the \oxname{T-ClosureValue} case, first we invert the stack frame typing
  hypothesis to get that $\forall x \in \oxdomain{\oxstackframe}.$
  $\oxtypjudge{\oxglobalctx}{\oxemptyctx}{\oxkontctx}{
    \oxnewframe{\oxvarctx}{\oxframe_c} }{\oxstackframe(x)}{\oxframe_c(x)}{
    \oxnewframe{\oxvarctx}{\oxframe_c} }$. We can apply the induction hypothesis
  to each of these statements, and apply \oxname{WF-Frame} to get
  $\oxsubstorevalidity{\oxglobalctx}{
    \oxvarctx[\oxloanctxentry{\oxcprov}{\oxset{\overline{\oxloan}}}]
  }{\oxstackframe_c}{\oxframe_c}$.

  Next we need to show that the body remains well typed. This follows from
  \Lemma{lemma:exprs-well-typed-loan-update}.

  In all other value cases, the typing judgement holds immediately.
\end{proof}

\begin{oxlemma}{Stack Validity is Preserved under Safe Loan Updates}
  {lemma:store-validity-loan-update} If
  $\oxmusafety{\oxemptyctx}{\oxkontctx}{\oxvarctx}{\oxmuta}{\oxplaceexpr}{
    \oxset{\overline{\oxloan}} }$ and
  $\oxnotinclosure{\oxkontctx}{\oxvarctx}{\oxcprov}$ and
  and $\oxvarctx(\oxcprov) = \emptyset$ and
  $\oxstorevalidity{\oxglobalctx}{\oxvarctx}{\oxstore}$, then
  $\oxstorevalidity{\oxglobalctx}{ \oxupdate{\oxvarctx}{
      \oxloanctxentry{\oxcprov}{\oxset{\overline{\oxloan}}} } }{\oxstore}$.
\end{oxlemma}

\begin{proof}
  We proceed by induction on the stack validity. There are two cases,
  \oxname{WF-StackEmpty}, and \oxname{WF-StackFrame}. \oxname{WF-StackEmpty} is
  impossible, since we already know that $\oxcprov$ is in $\oxvarctx$.

  \vspace{1em}
  \noindent \framebox[\textwidth]{ \figuresize
    \begin{mathpar}
      \VStack\and \VStackEmpty
    \end{mathpar}
  } \vspace{1em}

  In the case of \oxname{WF-StackFrame}, we have to show that the
  values remain well-typed in the updated environment. For the remaining
  $\oxvarctx^\prime$, if $\oxcprov \in \oxvarctx^\prime$, then we apply the
  induction hypothesis, otherwise we just use the derivation from the premise.

  To show that the values in the stack are still well typed in
  $\oxvarctx[\oxcprov \mapsto \{\overline{\oxloan}\}]$, we apply
  \Lemma{lemma:values-well-typed-loan-update}.
\end{proof}

\subsection{Preservation of Rewriting under Parallel Type Checking Lemmas}
\label{sec:parallel-lemmas}

\begin{oxlemma}{Region Rewriting is Preserved by Parallel Loan Updates}
  {lemma:rewriting-parallel-ownership}
  If
  \begin{enumerate}
  \item $\oxtunify{\oxemptyctx}{\oxkontctx}{\oxvarctx^\prime}{\oxtype_1}{\oxtype_2}{\oxvarctx_s}$
  \item $\oxmusafety{\oxemptyctx}{\oxkontctx, \oxtype_1, \oxkontctx^\prime}{\oxvarctx}{\oxmuta}{\oxplaceexpr}{\oxset{\oxloans}}$
  \item $\oxmusafety{\oxemptyctx}{\oxkontctx, \oxtype_2, \oxkontctx^\prime}{\oxvarctx^\prime}{\oxmuta}{\oxplaceexpr}{\oxset{\oxloans^\prime}}$
  \item $\oxdomain{\oxvarctx} = \oxdomain{\oxvarctx^\prime}$ and $\oxdomain{\oxvarctx_o} = \oxdomain{\oxvarctx_o^\prime}$
  \item $\forall \oxcprov \in \oxdomain{\oxvarctx}.$ $\oxvarctx^\prime(\oxcprov) \subseteq \oxvarctx(\oxcprov)$
  \item $\forall \oxcprov \in \oxdomain{\oxvarctx_o}.$ $\oxvarctx^\prime_o(\oxcprov) \subseteq \oxvarctx_o(\oxcprov)$
  \end{enumerate}
  then
  $\oxtunify{\oxemptyctx}{\oxkontctx}{\oxvarctx^\prime[\oxcprov \mapsto \oxset{\oxloans^\prime}]}{\oxtype_1}{\oxtype_2}{\oxvarctx_s^\prime}$.
\end{oxlemma}

\begin{proof}
  Proceed by induction on the rewriting derivation. The only interesting case is
  \oxname{RR-Reference}, in which case we proceed by induction on the outlives
  derivation.

  The interesting cases are \oxname{OL-CombineConcrete},
  \oxname{OL-CombineConcreteUnrestricted}, and \oxname{OL-CheckConcrete} since
  the type variable environment is guaranteed to be empty and the other cases
  all involve abstract regions. The only interesting obligations are the region
  not reborrowed ones, since the types in the environments are unchanged. This
  amounts to showing that for all loans in $\oxset{\oxloans^\prime}$, none are
  reborrows of references that have either $\oxcprov_1$ or $\oxcprov_2$ as their
  region.

  Assume one such loan $\oxloanpkg{\oxmuta^\prime}{\oxderef \oxplace} \in
  \oxloans^\prime$ exists. Assume without loss of generality that $\oxplace:
  \oxtref{\oxcprov_1}{\oxmuta^{\prime\prime}}{\oxtype} \in \oxvarctx$. By the
  well formedness of $\oxkontctx, \oxtype_1, \oxkontctx^\prime$, since
  $\oxtype_1$ contains $\oxcprov_1$, it must be the case that
  $\oxvarctx(\oxcprov_1) \not = \emptyset$. When checking ownership safety for
  $\oxderef \oxplace$, we'll need to show that excluding $\oxplace$, there are
  no conflicts in $\oxcprov^\prime$ with any lons in $\oxcprov^\prime$. But
  since the exclusion clause doesn't exclude when we have a reference in theta,
  and $\oxtype_1$ contains $\oxcprov_1$, it must be the case that $\oxcprov_1$
  will not be excluded, and we will then find loan conflicts, which is a
  contradiction.
\end{proof}

\begin{oxlemma}{Region Rewriting is Preserved by Type Checking Parallel Expressions}
  {lemma:rewriting-parallel-expression} \; \\
  If
  \begin{enumerate}
  \item $\oxtunify{\oxemptyctx}{\oxkontctx}{\oxvarctx^\prime}{\oxtype_1}{\oxtype_2}{\oxvarctx_s}$
  \item $\oxtypjudge{\oxglobalctx}{\oxemptyctx}{\oxkontctx, \oxtype_1, \oxkontctx^\prime}{\oxvarctx}{\oxexpr}{\oxtype}{\oxvarctx_o}$
  \item $\oxtypjudge{\oxglobalctx}{\oxemptyctx}{\oxkontctx, \oxtype_2, \oxkontctx^\prime}{\oxvarctx^\prime}{\oxexpr}{\oxtype}{\oxvarctx_o^\prime}$
  \item $\oxdomain{\oxvarctx} = \oxdomain{\oxvarctx^\prime}$ and $\oxdomain{\oxvarctx_o} = \oxdomain{\oxvarctx_o^\prime}$
  \item $\forall \oxcprov \in \oxdomain{\oxvarctx}.$ $\oxvarctx^\prime(\oxcprov) \subseteq \oxvarctx(\oxcprov)$
  \item $\forall \oxcprov \in \oxdomain{\oxvarctx_o}.$ $\oxvarctx^\prime_o(\oxcprov) \subseteq \oxvarctx_o(\oxcprov)$
  \end{enumerate}
  then
  $\oxtunify{\oxemptyctx}{\oxkontctx}{\oxvarctx_o^\prime}{\oxtype_1}{\oxtype_2}{\oxvarctx_s^\prime}$.
\end{oxlemma}

\begin{proof}
  Proceed by induction on the typing derivation using $\oxvarctx^\prime$ (note
  that since the typing derivation is by the structure of $\oxexpr$, we can
  simultaneously induct on the typing derivation using $\oxvarctx$).

  In the \oxname{T-Copy}, \oxname{T-Function}, \oxname{T-Abort},
  \oxname{T-Unit}, \oxname{T-u32}, \oxname{T-True}, and \oxname{T-False} cases,
  the proof is immediate, with $\oxvarctx_o = \oxvarctx$ and
  $\oxvarctx^{\prime\prime} = \oxvarctx^\prime$.

  In the \oxname{T-IndexCopy}, \oxname{T-LetRegion}, \oxname{T-While},
  \oxname{T-ForArray}, \oxname{T-ForSlice}, \oxname{T-Closure},
  \oxname{T-Tuple}, \oxname{T-Array}, \oxname{T-Slice}, \oxname{T-Left}, and
  \oxname{T-Right} cases, the proof follows immediately from the induction
  hypothesis.

  In the \oxname{T-Move} case, the proof is immediate because making a type a
  dead does not add additional obligations in the rewriting judgement.

  In the \oxname{T-Borrow}, \oxname{T-BorrowIndex}, and \oxname{T-BorrowSlice}
  cases, we proceed by the induction hypothesis and apply
  \Lemma{lemma:rewriting-parallel-ownership}.

  In the \oxname{T-Seq} case, we just need to show that garbage collecting loans
  preserves rewriting. But this is immediate, because garbage collection loans
  can only clear loan sets, which makes the requirements in rewriting strictly
  easier since it could only \emph{remove} reborrows, not add them.

  In the \oxname{T-AssignDeref}, \oxname{T-AppFunction}, and
  \oxname{T-AppClosure} cases, we first apply the induction hypothesis. Note
  that the output of the region rewriting at most only combines loan sets. As
  such, the region rewriting is preserved, because the main condition, the
  region not reborrowed requirement on the regions in the types, will still
  consider the same set of loans. For all region rewriting cases below, use this
  same reasoning.

  In the \oxname{T-Assign} case, we apply the induction hypothesis and reason
  about the rewriting as above, but we additionally need to know that the type
  update maintains the region not reborrowed and closure restrictions. For both,
  it's immediate because the rewriting in the hypothesis of the typing rule will
  check the same restrictions.

  In the \oxname{T-Branch} and \oxname{T-Match} cases, we apply the induction
  hypothesis, the reasoning above about rewriting, and the fact that
  $\oxintersect$ only unions together the loan sets from $\oxvarctx_2$ and
  $\oxvarctx_3$, both of which had the rewriting restrictions true by the
  induction hypothesis.

  In the \oxname{T-Let} case, we again apply the induction hypothesis and the
  reasoning above about rewriting, but additionally use the same reasoning as
  the garbage collection case as well.

  In the \oxname{T-Drop} case, we just need to show that rewriting is preserved
  by making a place dead in order to apply the induction hypothesis. This is
  immediate though, because all of the obligations in rewriting are either the
  same difficulty or made easier by making a place dead.
\end{proof}

\begin{oxlemma}{Outlives Still Holds with Smaller Continuation Contexts}
  {lemma:outlives-smaller-kont}
  If $\oxrunify{\oxemptyctx}{\oxkontctx, \oxtype}{\oxvarctx}{\oxcprov_1}{\oxcprov_2}{\oxvarctx^\prime}$
  then $\oxrunify{\oxemptyctx}{\oxkontctx}{\oxvarctx}{\oxcprov_1}{\oxcprov_2}{\oxvarctx^\prime}$.
\end{oxlemma}
\begin{proof}
  Proceed by induction on the outlives judgement. The only interesting cases are
  the \oxname{OL-CombineConcrete}, \oxname{OL-CombineConcreteUnrestricted}, and
  \oxname{OL-CheckConcrete} cases which all proceed similarly. The main
  obligations, the region not reborrowed and closure restriction judgements, are
  immediate since they are either unaffected by or have strictly fewer
  obligations in the smaller temporary typing $\oxkontctx$.
\end{proof}

\begin{oxlemma}{Region Rewriting Still Holds with Smaller Continuation Contexts}
  {lemma:rewriting-smaller-kont}
  If $\oxtunify{\oxemptyctx}{\oxkontctx, \oxtype}{\oxvarctx}{\oxtype_1}{\oxtype_2}{\oxvarctx^\prime}$
  then $\oxtunify{\oxemptyctx}{\oxkontctx}{\oxvarctx}{\oxtype_1}{\oxtype_2}{\oxvarctx^\prime}$.
\end{oxlemma}
\begin{proof}
  Proceed by induction on the rewriting judgement. The only interesting case is
  \oxname{RR-Reference}, in which case we just apply
  \Lemma{lemma:outlives-smaller-kont}.
\end{proof}

\newpage
\subsection{Progress}
\label{sec:progress}

\begin{oxlemma}{Progress}{lemma:progress-app}
  If
  \oxtypjudge{\oxglobalctx}{\oxemptyctx}{\oxkontctx}{\oxvarctx}{\oxexpr}
  {\oxsitype}{\oxvarctx^\prime} and
  \oxstorevalidity{\oxglobalctx}{\oxvarctx}{\oxstore}, then either
  $\oxexpr$ is a value, $\oxexpr$ is an \oxabort{\,\dots\,}, or $\oxexists
  \oxstore^\prime, \oxexpr^\prime. \;
  \oxreduce{\oxglobalctx}{\oxstore}{\oxexpr}{\oxstore^\prime}{\oxexpr^\prime}$.
\end{oxlemma}

\paragraph{Proof} We proceed by induction on the derivation
\oxtypjudge{\oxglobalctx}{\oxemptyctx}{\oxkontctx}{\oxvarctx}{\oxexpr}
{\oxtype}{\oxvarctx^\prime}. \\

\oxprogresscase{T-Move}{\TMove}{\EMove}{
  Applying \Lemma{lemma:place-exprs-reduce} to
  $\oxmusafety{\oxtvarctx}{\oxvarctx}{\oxmut}{\oxplace}{\oxset{
      \oxloanpkg{\oxmut}{\oxplace}
  }}$, $\oxcomputetynoprov{\oxtvarctx}{\oxvarctx}{\oxmut}{\oxplace}{\oxsitype}$
  (from $\oxlookup{\oxvarctx}{\oxplace}{\oxsitype}$ by \oxname{TC-Place}), and
  $\oxstorevalidity{\oxglobalctx}{\oxvarctx}{\oxstore}$ to conclude that
  $\oxnorm{\oxstore}{\oxplace}{\_}{\_}{\oxvalue}$. Thus, we can step
  with \oxname{E-Move}.
}

\oxprogresscase{T-Copy}{\TCopy}{\ECopy}{
  Applying \Lemma{lemma:place-exprs-reduce} to
  $\oxmusafety{\oxtvarctx}{\oxvarctx}{\oximm}{\oxplaceexpr}{\oxset{
      \overline{\oxloan}
  }}$, $\oxcomputetynoprov{\oxtvarctx}{\oxvarctx}{\oximm}{\oxplaceexpr}
  {\oxsitype}$, and $\oxstorevalidity{\oxglobalctx}{\oxvarctx}{\oxstore}$ to
  conclude that $\oxnorm{\oxstore}{\oxplaceexpr}{\_}{\_}{\oxvalue}$. Thus, we can
  step with \oxname{E-Copy}.
}

\oxprogresscasewide{T-Borrow}{\TBorrow}{\EBorrow}{
  Applying \Lemma{lemma:place-exprs-reduce} to
  $\oxmusafety{\oxtvarctx}{\oxvarctx}{\oxmuta}{\oxplaceexpr}{\oxset{
      \overline{\oxloan}
    }}$, $\oxcomputetynoprov{\oxtvarctx}{\oxvarctx}{\oxmuta}{\oxplaceexpr}
  {\oxxitype}$, and $\oxstorevalidity{\oxglobalctx}{\oxvarctx}{\oxstore}$ to
  conclude that $\oxnorm{\oxstore}{\oxplaceexpr}{\oxreferent}{\_}{\_}$. Thus, we can
  step with \oxname{E-Borrow}.
}

\oxprogresscaseinductive{T-BorrowIndex}{\TBorrowIndex}{
  We proceed based on  whether or not $\oxexpr$ is a value. If it is not, we can
  decompose our expression into the evaluation context
  $\oxref{\oxprov}{\oxmuta}{\oxindex{\oxplaceexpr}{\oxhole}}$ and redex
  $\oxexpr$. Then, by applying our induction hypothesis to the typing derivation
  for $\oxexpr$, we know either that $\oxexpr$ is an $\oxkey{abort!}$ expression
  or it $\oxexpr$ steps to some $\oxexpr^\prime$.  In the former case, we can
  step with \oxname{E-EvalCtxAbort}.  In the latter case, we can plug
  $\oxexpr^\prime$ back into our evaluation context and step with
  \oxname{E-EvalCtx}.
}{
  If $\oxexpr$ is a value, we would like to step with one of:
}{\EBorrowIndex \and \EBorrowIndexOOB}{
  Since $\oxexpr$ is a value, we can apply \Lemma{lemma:values-dont-change} to
  get $\oxsubtypectx{\oxglobalctx}{\oxemptyctx}{\oxvarctx}{\oxvarctx^\prime}$.
  Applying \Lemma{lemma:stack-validity-related-envs}, then gives us
  $\oxstorevalidity{\oxglobalctx}{\oxvarctx^\prime}{\oxstore}$.

  Then, we can apply \Lemma{lemma:place-exprs-reduce} to
  $\oxmusafety{\oxtvarctx}{\oxvarctx^\prime}{\oxmuta}{\oxplaceexpr}{
    \oxset{\overline{\oxloan}}
  }$, $\oxcomputetynoprov{\oxtvarctx}{\oxvarctx^\prime}{\oxmuta}{\oxplaceexpr}
  {\oxxitype}$, and $\oxstorevalidity{\oxglobalctx}{\oxvarctx^\prime}{\oxstore}$
  to get $\oxnorm{\oxstore}{\oxplaceexpr}{\oxreferent}{\_}{\oxvalue}$. By
  \Lemma{lemma:canonical-forms}, we know that
  $\oxvalue =  \oxarr{\oxvalue_0 \oxdotsc \oxvalue_n}$ since the type tells us
  the shape of the resultant value.

  Since we wish to step with one of \oxname{E-BorrowIndex} and
  \oxname{E-BorrowIndexOOB}, we should observe that we now have
  their shared requirement: $\oxnorm{\oxstore}{\oxplaceexpr}{\oxreferent}{\_}{
    \oxarr{\oxvalue_0 \oxdotsc \oxvalue_n}
  }$. Their other obligations are a bounds check which together are a tautology
  (i.e. one of them must hold). Thus, we can step with the appropriate rule
  based on whether or not the bounds check succeeds.
}

\oxprogresscaseinductive{T-BorrowSlice}{\TBorrowSlice}{
  The proof proceeds along similar lines as for \oxname{T-BorrowIndex}. We
  proceed based on whether or not $\oxexpr_1$ and $\oxexpr_2$ are values.

  If $\oxexpr_1$ is not a value, then we can decompose our whole expression into
  the evaluation context
  $\oxref{\oxprov}{\oxmuta}{\oxslice{\oxplaceexpr}{\oxhole}{\oxexpr_2}}$ and
  redex $\oxexpr_1$. Then, by applying our induction hypothesis to $\oxexpr_1$,
  we know either that $\oxexpr_1$ steps to some $\oxexpr_1^\prime$ or is an
  $\oxkey{abort!}$ expression. In the former case, this satisfies our
  requirement since we can plug $\oxexpr^\prime$ back into our evaluation
  context. In the latter case, we can step with \oxname{E-EvalCtxAbort}.


  If $\oxexpr_1$ is a value and $\oxexpr_2$ is not a value, then we can
  decompose our whole expression into the evaluation context
  $\oxref{\oxprov}{\oxmuta}{\oxslice{\oxplaceexpr}{\oxvalue_1}{\oxhole}}$ and
  redex $\oxexpr_2$. Then, by applying our induction hypothesis to $\oxexpr_2$,
  we know either that $\oxexpr_2$ steps to some $\oxexpr_2^\prime$ or is an
  $\oxkey{abort!}$ expression. In the former case, this satisfies our
  requirement since we can plug $\oxexpr^\prime$ back into our evaluation
  context. In the latter case, we can step with \oxname{E-EvalCtxAbort}.
}{
  If $\oxexpr_1$ and $\oxexpr_2$ are values, we would like to step with one of:
}{\EBorrowSlice \and \EBorrowSliceOOB}{
  Since $\oxexpr_1$ is a value, we can apply \Lemma{lemma:values-dont-change} to
  get $\oxsubtypectx{\oxglobalctx}{\oxemptyctx}{\oxvarctx}{\oxvarctx_1}$.
  Then, since $\oxexpr_2$ is also a value, we can apply
  \Lemma{lemma:values-dont-change} to get
  $\oxsubtypectx{\oxglobalctx}{\oxemptyctx}{\oxvarctx_1}{\oxvarctx_2}$. Then, by
  transitivity, we get
  $\oxsubtypectx{\oxglobalctx}{\oxemptyctx}{\oxvarctx}{\oxvarctx_2}$.
  Then, applying \Lemma{lemma:stack-validity-related-envs} gives us
  $\oxstorevalidity{\oxglobalctx}{\oxvarctx_2}{\oxstore}$.

  Then, we can apply \Lemma{lemma:place-exprs-reduce} to
  $\oxmusafety{\oxtvarctx}{\oxvarctx_2}{\oxmuta}{\oxplaceexpr}{
    \oxset{\overline{\oxloan}}
  }$, $\oxcomputetynoprov{\oxtvarctx}{\oxvarctx_2}{\oxmuta}{\oxplaceexpr}
  {\oxtslice{\oxsitype}}$, and $\oxstorevalidity{\oxglobalctx}{\oxvarctx_2}{
    \oxstore
  }$ to get $\oxnorm{\oxstore}{\oxplaceexpr}{\oxreferent}{\_}{\oxvalue}$. By
  \Lemma{lemma:canonical-forms}, we know that
  $\oxvalue =  \oxarr{\oxvalue_0 \oxdotsc \oxvalue_n}$ since the type tells us
  the shape of the resultant value.

  Since we wish to step with one of \oxname{E-BorrowSlice} and
  \oxname{E-BorrowSliceOOB}, we should observe that we now have
  their shared requirement: $\oxnorm{\oxstore}{\oxplaceexpr}{\oxreferent}{\_}{
    \oxarr{\oxvalue_0 \oxdotsc \oxvalue_n}
  }$. Their other obligations are a bounds check which together are a tautology
  (i.e. one of them must hold). Thus, we can step with the appropriate rule
  based on whether or not the bounds check succeeds.
}

\oxprogresscaseinductive{T-IndexCopy}{\TIndexCopy}{
  We proceed based on whether or not $\oxexpr$ is a value. If it is not, we can
  decompose our expression into the evaluation context
  $\oxindex{\oxplaceexpr}{\oxhole}$ and redex $\oxexpr$. Then, by applying our
  induction hypothesis to $\oxexpr$, we know either that $\oxexpr$ steps to
  some $\oxexpr^\prime$ or is an $\oxkey{abort!}$ expression. In the former
  case, this satisfies our requirement since we can plug $\oxexpr^\prime$ back
  into our evaluation context. In the latter case, we can step with
  \oxname{E-EvalCtxAbort}.
}{If $\oxexpr$ is a value, we would like to step with one of:}{
  \EIndexCopy \and \EIndexCopyOOB
}{
  Since $\oxexpr$ is a value, we can apply \Lemma{lemma:values-dont-change} to
  get $\oxsubtypectx{\oxglobalctx}{\oxemptyctx}{\oxvarctx}{\oxvarctx^\prime}$.
  Applying \Lemma{lemma:stack-validity-related-envs}, then gives us
  $\oxstorevalidity{\oxglobalctx}{\oxvarctx^\prime}{\oxstore}$.

  Then, we can apply \Lemma{lemma:place-exprs-reduce} to
  $\oxmusafety{\oxtvarctx}{\oxvarctx^\prime}{\oxmuta}{\oxplaceexpr}{
    \oxset{\overline{\oxloan}}
  }$, $\oxcomputetynoprov{\oxtvarctx}{\oxvarctx^\prime}{\oxmuta}{\oxplaceexpr}
  {\oxxitype}$, and $\oxstorevalidity{\oxglobalctx}{\oxvarctx^\prime}{\oxstore}$
  to get $\oxnorm{\oxstore}{\oxplaceexpr}{\_}{\_}{\oxvalue}$. By
  \Lemma{lemma:canonical-forms}, we know that
  $\oxvalue =  \oxarr{\oxvalue_0 \oxdotsc \oxvalue_n}$ since the type tells us
  the shape of the resultant value.

  Since we wish to step with one of \oxname{E-IndexCopy} and
  \oxname{E-IndexCopyOOB}, we should observe that we now have
  their shared requirement: $\oxnorm{\oxstore}{\oxplaceexpr}{\_}{\_}{
    \oxarr{\oxvalue_0 \oxdotsc \oxvalue_n}
  }$. Their other obligations are a bounds check which together are a tautology
  (i.e. one of them must hold). Thus, we can step with the appropriate rule
  based on whether or not the bounds check succeeds.
}

\oxprogresscaseinductive{T-Seq}{\TSeq}{
  We proceed based on whether or not $\oxexpr_1$ is a value. If it is not, we can
  decompose our expression into the evaluation context
  $\oxseq{\oxhole}{\oxexpr_2}$ and redex $\oxexpr_1$. Then, by applying our
  induction hypothesis to $\oxexpr_1$, we know either that $\oxexpr_1$ steps to
  some $\oxexpr_1^\prime$ or is an $\oxkey{abort!}$ expression. In the former
  case, this satisfies our requirement since we can plug $\oxexpr_1^\prime$ back
  into our evaluation context. In the latter case, we can step with
  \oxname{E-EvalCtxAbort}.
}{If $\oxexpr_1$ is a value, we can step with:}{\ESeq}{}

\oxprogresscaseinductive{T-Branch}{\TBranch}{
  We proceed based on whether or not $\oxexpr_1$ is a value. If it is not, we can
  decompose our expression into the evaluation context
  $\oxbranch{\oxhole}{\oxexpr_2}{\oxexpr_3}$ and redex $\oxexpr_1$. Then, by
  applying our induction hypothesis to $\oxexpr_1$, we know either that
  $\oxexpr_1$ steps to some $\oxexpr_1^\prime$ or is an $\oxkey{abort!}$
  expression. In the former case, this satisfies our requirement since we can
  plug $\oxexpr_1^\prime$ back into our evaluation context. In the latter case,
  we can step with \oxname{E-EvalCtxAbort}.
}{If $\oxexpr_1$ is a value, we would like to step with one of:}{
  \EIfTrue \and \EIfFalse
}{
  Since $\oxexpr_1$ is a value, applying \Lemma{lemma:canonical-forms} tells us
  that $\oxexpr_1$ is either $\oxtrue$ or $\oxfalse$. In the former case, we can
  step with \oxname{E-IfTrue} and in the latter case, we can step with
  \oxname{E-IfFalse}
}

\oxprogresscaseinductive{T-Let}{\TLet}{
  We proceed based on whether or not $\oxexpr_1$ is a value. If it is not, we can
  decompose our expression into the evaluation context
  $\oxlet{\oxid}{\oxsitype_a}{\oxhole}{\oxexpr_2}$ and redex $\oxexpr_1$. Then,
  by applying our induction hypothesis to $\oxexpr_1$, we know either that
  $\oxexpr_1$ steps to some $\oxexpr_1^\prime$ or is an $\oxkey{abort!}$
  expression. In the former case, this satisfies our requirement since we can
  plug $\oxexpr_1^\prime$ back into our evaluation context. In the latter case,
  we can step with \oxname{E-EvalCtxAbort}.
}{If $\oxexpr_1$ is a value, we can step with:}{\ELet}{}

\oxprogresscaseinductive{T-LetRegion}{\TLetProv}{
  We proceed based on whether or not $\oxexpr$ is a value. If it is not, we can
  decompose our expression into the evaluation context
  $\oxletrgn{\oxcprov}{\oxhole}$ and redex $\oxexpr$. Then, by applying our
  induction hypothesis to $\oxexpr$, we know either that $\oxexpr$ steps to some
  $\oxexpr^\prime$ or is an $\oxkey{abort!}$ expression. In the former case,
  this satisfies our requirement since we can plug $\oxexpr^\prime$ back into
  our evaluation context. In the latter case, we can step with
  \oxname{E-EvalCtxAbort}.
}{If $\oxexpr$ is a value, we can step with:}{\ELetProv}{}

\oxprogresscaseinductive{T-Assign}{\TAssign}{
  We proceed based on whether or not $\oxexpr$ is a value. If it is not, we can
  decompose our expression into the evaluation context
  $\oxassign{\oxplaceexpr}{\oxhole}$ and redex $\oxexpr$. Then, by applying our
  induction hypothesis to $\oxexpr$, we know either that $\oxexpr$ steps to some
  $\oxexpr^\prime$ or is an $\oxkey{abort!}$ expression. In the former case,
  this satisfies our requirement since we can plug $\oxexpr^\prime$ back into
  our evaluation context. In the latter case, we can step with
  \oxname{E-EvalCtxAbort}.
}{
  If $\oxexpr$ is a value, we would like to step with:
}{\EAssign}{
  Since $\oxexpr$ is a value, we can apply \Lemma{lemma:values-dont-change} to
  get $\oxsubtypectx{\oxglobalctx}{\oxemptyctx}{\oxvarctx}{\oxvarctx^\prime}$.
  Applying \Lemma{lemma:stack-validity-related-envs}, then gives us
  $\oxstorevalidity{\oxglobalctx}{\oxvarctx^\prime}{\oxstore}$.

  Then, we can apply \Lemma{lemma:place-exprs-reduce} to
  $\oxmusafety{\oxtvarctx}{\oxvarctx^\prime}{\oxmuta}{\oxplaceexpr}{
    \oxset{\overline{\oxloan}}
  }$, $\oxcomputetynoprov{\oxtvarctx}{\oxvarctx^\prime}{\oxmuta}{\oxplaceexpr}
  {\oxxitype}$, and $\oxstorevalidity{\oxglobalctx}{\oxvarctx^\prime}{\oxstore}$
  to get $\oxnorm{\oxstore}{\oxplaceexpr}{\oxreferent}{\oxvaluectx}{\_}$.
}

\oxprogresscaseinductive{T-ForArray}{\TForArray}{
  We proceed based on whether or not $\oxexpr_1$ is a value. If it is not, we can
  decompose our expression into the evaluation context
  $\oxfor{\oxid}{\oxhole}{\oxexpr_2}$ and redex $\oxexpr_1$. Then, by applying our
  induction hypothesis to $\oxexpr_1$, we know either that $\oxexpr_1$ steps to
  some $\oxexpr_1^\prime$ or is an $\oxkey{abort!}$ expression. In the former
  case, this satisfies our requirement since we can plug $\oxexpr_1^\prime$ back
  into our evaluation context. In the latter case, we can step with
  \oxname{E-EvalCtxAbort}.
}{If $\oxexpr_1$ is a value, we would like to step with one of:}{
  \EForArray \and \EForEmptyArray
}{
  Since $\oxexpr_1$ is a value, then by \Lemma{lemma:canonical-forms}, we know
  that $\oxexpr_1$ is of the form $\oxarr{\oxvalue_1 \oxdotsc \oxvalue_n}$. If
  $n > 0$, then we can step with \oxname{E-ForArray}, and if $n = 0$, then we
  can step with \oxname{E-ForEmptyArray}.
}

\oxprogresscaseinductive{T-ForSlice}{\TForSlice}{
  We proceed based on whether or not $\oxexpr_1$ is a value. If it is not, we can
  decompose our expression into the evaluation context
  $\oxfor{\oxid}{\oxhole}{\oxexpr_2}$ and redex $\oxexpr_1$. Then, by applying our
  induction hypothesis to $\oxexpr_1$, we know either that $\oxexpr_1$ steps to
  some $\oxexpr_1^\prime$ or is an $\oxkey{abort!}$ expression. In the former
  case, this satisfies our requirement since we can plug $\oxexpr_1^\prime$ back
  into our evaluation context. In the latter case, we can step with
  \oxname{E-EvalCtxAbort}.
}{If $\oxexpr_1$ is a value, we would like to step with one of:}{
  \EForSlice \and \EForEmptySlice
}{
  If $\oxexpr_1$ is a value, then by \Lemma{lemma:canonical-forms}, we know that
  $\oxexpr_1$ is of the form $\oxsliceptr{\oxreferent}{\oxnum_1}{\oxnum_2}$.
  Further, by inversion of \oxname{T-Pointer} for the typing derivation of
  $\oxexpr_1$, we get $\oxreferentvalidity{\oxglobalctx}{\oxvarctx}{
    \oxslice{\oxreferent}{i}{j}
  }{\oxtslice{\oxsitype}}$. By inversion of \oxname{WF-RefSliceArray} or
  \oxname{WF-RefSliceSlice} (one of which must apply since the referent ends in
  a slice), we know that $i \leq j$. If $i < j$, we step with
  \oxname{E-ForSlice} and if $i = j$, we step with \oxname{E-ForEmptySlice}.
}

\oxprogresscasewide{T-Closure}{\TClosure}{\EClosure}{
  Since $\oxfvars{\cdot}$ and $\oxfncvars{\oxstore}{\cdot}$ are total, we can
  always step with \oxname{E-Closure}.
}

\oxprogresscaseinductive{T-AppClosure}{\TAppClosure}{
  We proceed based on whether or not $\oxexpr_f$ is a value. If it is not, we can
  decompose our expression into the evaluation context
  $\oxapply{\oxhole}{\oxexpr_1 \oxdotsc \oxexpr_n}$ and redex
  $\oxexpr_f$. Then, by applying our induction hypothesis to $\oxexpr_f$, we
  know either that $\oxexpr_f$ steps to some $\oxexpr_f^\prime$ or is an
  $\oxkey{abort!}$ expression. In the former case, this satisfies our
  requirement since we can plug $\oxexpr_f^\prime$ back into our evaluation
  context. In the latter case, we can step with \oxname{E-EvalCtxAbort}.

  Next, we'll proceed based on whether or not each expression $\oxexpr_i$ is a
  value. If any of them are not, we can decompose our expression into the
  evaluation context $\oxapply{\oxvalue_f}{\oxvalue_1 \oxdotsc \oxvalue_m
    \oxcomma \oxhole \oxcomma \oxexpr_1 \oxdotsc \oxexpr_{n^\prime}}$ and redex
  $\oxexpr_i$. Then, by applying our induction hypothesis to $\oxexpr_i$, we
  know either that $\oxexpr_i$ steps to some $\oxexpr_i^\prime$ or is an
  $\oxkey{abort!}$ expression. In the former case, this satisfies our
  requirement since we can plug $\oxexpr_i^\prime$ back into our evaluation
  context. In the latter case, we can step with \oxname{E-EvalCtxAbort}.
}{
  If $\oxexpr_f$ is a value and every $\oxexpr_i$ is a value, we would like to
  step with one of:
}{\EApp}{
  Since $\oxexpr_f$ is a value, then by \Lemma{lemma:canonical-forms}, we know
  that it has the form
  $\oxclosureval{\oxstore_c}{ \oxascribe{\oxid_1}{\oxsitype_1} \oxdotsc
    \oxascribe{\oxid_n}{\oxsitype_n} }{\oxsitype_r}{\oxexpr}$.
  Then, since all of the $e_i$ are values, then we can step using
  \oxname{E-AppClosure}.
}

\oxprogresscaseinductive{T-AppFunction}{\TApp}{
  We proceed based on whether or not $\oxexpr_f$ is a value. If it is not, we can
  decompose our expression into the evaluation context
  $\oxapp{\oxhole}{\overline{\oxenv} \oxcomma \overline{\oxprov} \oxcomma
    \overline{\oxsitype}}{\oxexpr_1 \oxdotsc \oxexpr_n}$ and redex
  $\oxexpr_f$. Then, by applying our induction hypothesis to $\oxexpr_f$, we
  know either that $\oxexpr_f$ steps to some $\oxexpr_f^\prime$ or is an
  $\oxkey{abort!}$ expression. In the former case, this satisfies our
  requirement since we can plug $\oxexpr_f^\prime$ back into our evaluation
  context. In the latter case, we can step with \oxname{E-EvalCtxAbort}.

  Next, we'll proceed based on whether or not each expression $\oxexpr_i$ is a
  value. If any of them are not, we can decompose our expression into the
  evaluation context
  $\oxapp{\oxvalue_f}{\overline{\oxenv} \oxcomma \overline{\oxprov} \oxcomma \overline{\oxsitype} }{\oxvalue_1 \oxdotsc \oxvalue_m \oxcomma \oxhole \oxcomma \oxexpr_1 \oxdotsc \oxexpr_{n^\prime}}$
  and redex $\oxexpr_i$. Then, by applying our induction hypothesis to
  $\oxexpr_i$, we know either that $\oxexpr_i$ steps to some $\oxexpr_i^\prime$
  or is an $\oxkey{abort!}$ expression. In the former case, this satisfies our
  requirement since we can plug $\oxexpr_i^\prime$ back into our evaluation
  context. In the latter case, we can step with \oxname{E-EvalCtxAbort}. }{ If
  $\oxexpr_f$ is a value and every $\oxexpr_i$ is a value, we would like to step
  with one of: }{\EAppFun}{ Since $\oxexpr_f$ is a value, then by
  \Lemma{lemma:canonical-forms}, we know that it has the form
  $\oxfnname$.
  Then, since all of the $e_i$ are values, then we can step using
  \oxname{E-AppFunction}.
}

\oxprogresscasevalue{T-Unit}{\TUnit}{
  By inspection of the value grammar, we know that $\oxunit$ is already a value.
}

\oxprogresscasevalue{T-u32}{\TuThreeTwo}{
  By inspection of the value grammar, we know that $\oxnum$ is already a value.
}

\oxprogresscasevalue{T-True}{\TTrue}{
  By inspection of the value grammar, we know that $\oxtrue$ is already a value.
}

\oxprogresscasevalue{T-False}{\TFalse}{
  By inspection of the value grammar, we know that $\oxfalse$ is already a
  value.
}

\oxprogresscasevalue{T-Tuple}{\TTuple}{
  We'll proceed based on whether or not each expression $\oxexpr_i$ is a
  value. If any of them are not, we can decompose our expression into the
  evaluation context $\oxprod{\oxvalue_1 \oxdotsc \oxvalue_m \oxcomma \oxhole
    \oxcomma \oxexpr_1 \oxdotsc \oxexpr_{n^\prime}}$ and redex
  $\oxexpr_i$. Then, by applying our induction hypothesis to $\oxexpr_i$, we
  know either that $\oxexpr_i$ steps to some $\oxexpr_i^\prime$ or to an
  $\oxkey{abort!}$ expression. In either case, this satisfies our requirement,
  since we can plug $\oxexpr_i^\prime$ back into our evaluation context.

  If every expression $e_i$ is a value, then the whole expression is a value by
  the definition of values.
}

\oxprogresscasevalue{T-Array}{\TArray}{
  We'll proceed based on whether or not each expression $\oxexpr_i$ is a
  value. If any of them are not, we can decompose our expression into the
  evaluation context $\oxarr{\oxvalue_1 \oxdotsc \oxvalue_m \oxcomma \oxhole
    \oxcomma \oxexpr_1 \oxdotsc \oxexpr_{n^\prime}}$ and redex
  $\oxexpr_i$. Then, by applying our induction hypothesis to $\oxexpr_i$, we
  know either that $\oxexpr_i$ steps to some $\oxexpr_i^\prime$ or to an
  $\oxkey{abort!}$ expression. In either case, this satisfies our requirement,
  since we can plug $\oxexpr_i^\prime$ back into our evaluation context.

  If every expression $e_i$ is a value, then the whole expression is a value by
  the definition of values.
}

\oxprogresscasevalue{T-Abort}{\TAbort}{
  By definition, $\oxabort{\oxdots}$ is an $\oxkey{abort!}$ expression.
}

% \oxprogresscasevalue{T-TupleStruct}{\TTupleStruct}{
%   We'll proceed based on whether or not each expression $\oxexpr_i$ is a
%   value. If any of them are not, we can decompose our expression into the
%   evaluation context $\oxtupstruct{
%     \oxturbofish{\oxdatatype}{\overline{\oxprov} \oxcomma \overline{\oxsitype}}
%   }{\oxvalue_1 \oxdotsc \oxvalue_m \oxcomma \oxhole
%     \oxcomma \oxexpr_1 \oxdotsc \oxexpr_{n^\prime}}$ and redex
%   $\oxexpr_i$. Then, by applying our induction hypothesis to $\oxexpr_i$, we
%   know either that $\oxexpr_i$ steps to some $\oxexpr_i^\prime$ or to an
%   $\oxkey{abort!}$ expression. In either case, this satisfies our requirement,
%   since we can plug $\oxexpr_i^\prime$ back into our evaluation context.

%   If every expression $e_i$ is a value, then the whole expression is a value by
%   the definition of values.
% }

% \oxprogresscasevalue{T-RecordStruct}{\TRecordStruct}{
%   We'll proceed based on whether or not each expression $\oxexpr_i$ is a
%   value. If any of them are not, we can decompose our expression into the
%   following evaluation context:
%   \begin{center}
%     $\oxrecstruct{
%       \oxturbofish{\oxdatatype}{\overline{\oxprov} \oxcomma \overline{\oxsitype}}
%     }{\oxid_1 \oxcolon \; \oxvalue_1 \oxdotsc \oxid_m \oxcolon \; \oxvalue_m \oxcomma
%       \oxid_{m+1} \oxcolon \; \oxhole
%       \oxcomma \oxid_{m+2} \oxcolon \; \oxexpr_1 \oxdotsc \oxid_n \oxcolon \;
%       \oxexpr_{n^\prime}}$
%   \end{center}
%   and the redex $\oxexpr_i$. Then, by applying our induction hypothesis to
%   $\oxexpr_i$, we know either that $\oxexpr_i$ steps to some $\oxexpr_i^\prime$
%   or to an $\oxkey{abort!}$ expression. In either case, this satisfies our
%   requirement, since we can plug $\oxexpr_i^\prime$ back into our evaluation
%   context.

%   If every expression $e_i$ is a value, then the whole expression is a value by
%   the definition of values.
% }

\oxprogresscaseinductive{T-Framed}{\TFramed}{
  We proceed based on whether or not $\oxexpr$ is a value. If it is not, we can
  decompose our expression into the evaluation context
  $\oxframed{\oxhole}$ and redex $\oxexpr$. Then, by applying our
  induction hypothesis to $\oxexpr$, we know either that $\oxexpr$ steps to some
  $\oxexpr^\prime$ or to an $\oxkey{abort!}$ expression. In either case, this
  satisfies our requirement, since we can plug $\oxexpr^\prime$ back into our
  evaluation context.
}{If $\oxexpr$ is a value, then we would like to step with:}{\EFramed}{
  In order to do so, we need to know $\oxid \in \oxdomain{\oxstore}$.
  Fortunately, we know from our assumption that
  \oxstorevalidity{\oxglobalctx}{\oxvarctx}{\oxstore} (via
  \oxname{WF-Stack}). The premise of \oxname{WF-Stack} tells us that
  $\oxdomain{\oxstore} = \oxdomain{\oxvarctx}$, and thus the $\oxid \in
  \oxdomain{\oxvarctx}$ from the premise of \oxname{T-Framed} is sufficient to
  tell us that $\oxid \in \oxdomain{\oxstore}$. Thus, we can step with
  \oxname{E-Framed}.
}

\oxprogresscasevalue{T-Pointer}{\TPointer}{
  By inspection of the value grammar, we know that
  $\oxptr{\oxplace}$ is already a value.
}

\oxprogresscasevalue{T-ClosureValue}{\TClosureVal}{
  By inspection of the value grammar, we know that
  $\oxclosureval{\oxstore}{ \oxascribe{\oxid_1}{\oxsitype_1} \oxdotsc
    \oxascribe{\oxid_n}{\oxsitype_n} }{\oxsitype_r}{\oxexpr}$ is already a
  value.
}

\oxprogresscasevalue{T-Dead}{\TUninit}{
  The type $\oxsidtype$ is not in the grammar of $\oxsitype$. Thus, we have a
  contradiction.
}

\oxprogresscasevalue{T-Drop}{\TDrop}{
  By \oxname{R-Env}, we have that $\oxsubtypectx{\oxglobalctx}{\oxemptyctx}
  {\oxvarctx}{\oxtupdate{\oxvarctx}{\oxplace}{\oxsidtype_\oxplace}}$. Then,
  applying \Lemma{lemma:stack-validity-related-envs} with
  $\oxstorevalidity{\oxglobalctx}{\oxvarctx}{\oxstore}$ (from our premise) gives
  us $\oxstorevalidity{\oxglobalctx}{
    \oxtupdate{\oxvarctx}{\oxplace}{\oxsidtype_\oxplace}
  }{\oxstore}$. We can then apply our induction hypothesis to this and
  $\oxtypjudge{\oxglobalctx}{\oxemptyctx}{\oxkontctx}{
    \oxtupdate{\oxvarctx}{\oxplace}{\oxsidtype_\oxplace}
  }{\oxexpr}{\oxsxtype}{\oxvarctx_f}$ to reach our goal.
}

\oxprogresscasevalue{T-Left}{\TInl}{
  We proceed based on whether or not $\oxexpr$ is a value. If it is not, we can
  decompose our expression into the evaluation context
  $\oxinl{\oxsitype_1}{\oxsitype_2}{\oxhole}$ and redex $\oxexpr$. Then, by
  applying our induction hypothesis to $\oxexpr$, we know either that $\oxexpr$
  steps to some $\oxexpr^\prime$ or to an $\oxkey{abort!}$ expression. In either
  case, this satisfies our requirement, since we can plug $\oxexpr^\prime$ back
  into our evaluation context.

  If $\oxexpr$ is a value, then we know that the whole expression
  $\oxinl{\oxsitype_1}{\oxsitype_2}{\oxvalue}$ is a value and thus we are done.
}

\oxprogresscasevalue{T-Right}{\TInr}{
  We proceed based on whether or not $\oxexpr$ is a value. If it is not, we can
  decompose our expression into the evaluation context
  $\oxinr{\oxsitype_1}{\oxsitype_2}{\oxhole}$ and redex $\oxexpr$. Then, by
  applying our induction hypothesis to $\oxexpr$, we know either that $\oxexpr$
  steps to some $\oxexpr^\prime$ or to an $\oxkey{abort!}$ expression. In either
  case, this satisfies our requirement, since we can plug $\oxexpr^\prime$ back
  into our evaluation context.

  If $\oxexpr$ is a value, then we know that the whole expression
  $\oxinr{\oxsitype_1}{\oxsitype_2}{\oxvalue}$ is a value and thus we are done.
}

\oxprogresscaseinductive{T-Match}{\TMatch}{
  We proceed based on whether or not $\oxexpr$ is a value. If it is not, we can
  decompose our expression into the evaluation context
  $\oxmatch{\oxhole}{\oxid_1}{\oxexpr_1}{\oxid_2}{\oxexpr_2}$ and redex
  $\oxexpr$. Then, by applying our induction hypothesis to $\oxexpr$, we know
  either that $\oxexpr$ steps to some $\oxexpr^\prime$ or to an $\oxkey{abort!}$
  expression. In either case, this satisfies our requirement, since we can plug
  $\oxexpr^\prime$ back into our evaluation context.
}{If $\oxexpr$ is a value, then we would like to step with either:}{
  \EMatchLeft \and \EMatchRight
}{
  Since $\oxexpr$ is a value, applying \Lemma{lemma:canonical-forms} tells us
  that $\oxexpr$ is either of the form $\oxinl{\oxsitype_1}{\oxsitype_2}{\oxvalue}$ or
  $\oxinr{\oxsitype_1}{\oxsitype_2}{\oxvalue}$. In the former case, we can
  step with \oxname{E-MatchLeft} and in the latter case, we can step with
  \oxname{E-MatchRight}.
}

%%% Local Variables:
%%% mode: latex
%%% fill-column: 80
%%% End:


\newpage
\subsection{Preservation}
\label{sec:preservation}

\begin{oxlemma}{Preservation}{lemma:preservation-app}
  If $\oxtypjudge{\oxglobalctx}{\oxemptyctx}{\oxkontctx}{\oxvarctx}{\oxexpr}{\oxsitype_1}{\oxvarctx_f} and
  \oxstorevalidity{\oxglobalctx}{\oxvarctx}{\oxstore}$ and
  $\oxwfkontctx{\oxglobalctx}{\oxvarctx}{\overline{\oxvalue}}{\oxkontctx}$ and
  $\oxreduce{\oxglobalctx}{\oxstore}{\oxexpr}{\oxstore^\prime}{\oxexpr^\prime}$, then there exists
  $\oxvarctx_i$ such that
  $\oxstorevalidity{\oxglobalctx}{\oxvarctx_i}{\oxstore^\prime}$ and
  $\oxwfkontctx{\oxglobalctx}{\oxvarctx_i}{\overline{\oxvalue}}{\oxkontctx}$ and
  $\oxtypjudge{\oxglobalctx}{\oxemptyctx}{\oxkontctx}{\oxvarctx_i}{\oxexpr^\prime}
  {\oxsitype_2}{\oxvarctx_f^\prime}$ and
  $\oxtunify[\oxcombine]{\oxemptyctx}{\oxkontctx}{\oxvarctx_f^\prime}
  {\oxsitype_2}{\oxsitype_1}{\oxvarctx_s}$ and there exists $\oxvarctx_o$ such
  that $\oxvarctx_f = \oxvarctx_s \oxintersect \oxvarctx_o$.
\end{oxlemma}

\paragraph{Proof} We proceed by induction on the derivation
\oxtypjudge{\oxglobalctx}{\oxemptyctx}{\oxkontctx}{\oxvarctx}{\oxexpr}{\oxtype}{\oxvarctx_f} \\

\oxpreservationcase{T-Move}{\TMove}{\oxplace}{\EMove}{
  \oxtupdate{\oxvarctx}{\oxplace}{\oxsidtype}
}

\begin{preservation}
  %% new store is valid
  \oxpresline{
    \oxstorevalidity{\oxglobalctx}{
      \oxtupdate{\oxvarctx}{\oxplace}{\oxsidtype}
    }{
      \oxvupdate{\oxstore}{\oxplace}{\oxuninit}
    }
  }{
    Applying \Lemma{lemma:stack-validity-related-envs} to
    $\oxsubtypectx{\oxglobalctx}{\oxemptyctx}{\oxvarctx}{
      \oxtupdate{\oxvarctx}{\oxplace}{\oxsidtype}
    }$ (immediate by \oxname{R-Env}) and
    $\oxstorevalidity{\oxglobalctx}{\oxvarctx}{\oxstore}$ (from premise) gives us
    $\oxstorevalidity{\oxglobalctx}{
      \oxtupdate{\oxvarctx}{\oxplace}{\oxsidtype}
    }{\oxstore}$. Then, since we know
    $\oxtypjudge{\oxglobalctx}{\oxemptyctx}{\oxkontctx}{
      \oxtupdate{\oxvarctx}{\oxplace}{\oxsidtype}
    }{\oxuninit}{\oxsidtype}{
      \oxtupdate{\oxvarctx}{\oxplace}{\oxsidtype}
    }$ (by \oxname{T-Dead}), we can conclude
    $\oxstorevalidity{\oxglobalctx}{
      \oxtupdate{\oxvarctx}{\oxplace}{\oxsidtype}
    }{
      \oxvupdate{\oxstore}{\oxplace}{\oxuninit}
    }$.
  }

  %% kontctx values are still valid
  \oxpresline{
    \oxwfkontctx{\oxglobalctx}{
      \oxtupdate{\oxvarctx}{\oxplace}{\oxsidtype}
    }{\overline{\oxvalue}}{\oxkontctx}
  }{
    Inverting \oxname{WF-Temporaries} gives us $\forall i \in 1 \oxdots n. \;
    \oxtypjudge{\oxglobalctx}{\oxemptyctx}{\oxsitype_1 \oxdotsc \oxsitype_{i-1}}{\oxvarctx}{
      \oxvalue_i
    }{\oxsitype_i}{\oxvarctx}$. By \oxname{R-Env}, we have that
    \oxsubtypectx{\oxglobalctx}{\oxemptyctx}{\oxvarctx}{
      \oxtupdate{\oxvarctx}{\oxplace}{\oxsidtype}
    }. Thus, we can apply \Lemma{lemma:value-typing-related-envs} to each of the
    typing judgments and then apply \oxname{WF-Temporaries} to get
    $\oxwfkontctx{\oxglobalctx}{
      \oxtupdate{\oxvarctx}{\oxplace}{\oxsidtype}
    }{\overline{\oxvalue}}{\oxkontctx}$.
  }

  %% new expression is well-typed
  \oxpresline{
    \oxtypjudge{\oxglobalctx}{\oxemptyctx}{\oxkontctx}{
      \oxtupdate{\oxvarctx}{\oxplace}{\oxsidtype}
    }{\oxvalue}{\oxsitype}{
      \oxtupdate{\oxvarctx}{\oxplace}{\oxsidtype}
    }
  }{
    Applying \Lemma{lemma:place-exprs-reduce} to
    $\oxmusafety{\oxemptyctx}{\oxvarctx}{\oxmut}{\oxplace}{
      \oxset{\oxloanpkg{\oxmut}{\oxplace}}
    }$, $\oxcomputetynoprov{\oxemptyctx}{\oxvarctx}{\oxmut}{\oxplace}{\oxsitype}$
    (immediate by \oxname{TC-Place} with $\oxlookup{\oxvarctx}{\oxplace}{
      \oxsitype
    }$), and $\oxstorevalidity{\oxglobalctx}{\oxvarctx}{\oxstore}$ gives us
    $\oxtypjudge{\oxglobalctx}{\oxemptyctx}{\oxkontctx}{\oxvarctx}{\oxvalue}
    {\oxsitype}{\oxvarctx}$. Then, by applying
    \Lemma{lemma:value-typing-related-envs} with
    $\oxsubtypectx{\oxglobalctx}{\oxemptyctx}
    {\oxvarctx}{\oxtupdate{\oxvarctx}{\oxplace}{\oxsidtype}}$ (immediate by
    \oxname{R-Env}), we can conclude $\oxtypjudge{\oxglobalctx}{\oxemptyctx}{\oxkontctx}{
      \oxtupdate{\oxvarctx}{\oxplace}{\oxsidtype}
    }{\oxvalue}{\oxsitype}{
      \oxtupdate{\oxvarctx}{\oxplace}{\oxsidtype}
    }$.
  }

  %% new type is a subtype of the old type
  \oxpresline{
    \oxtunify[\oxcombine]{\oxemptyctx}{\oxkontctx}{
      \oxtupdate{\oxvarctx}{\oxplace}{\oxsidtype}
    }{\oxsitype}{\oxsitype}{
      \oxtupdate{\oxvarctx}{\oxplace}{\oxsidtype}
    }
  }{
    Immediate by \oxname{RR-Refl}.
  }

  %% new loan context is a sub-context of old loan context
  \oxpresline{
    \exists \oxvarctx_o. \oxtupdate{\oxvarctx}{\oxplace}{\oxsidtype} \oxintersect \oxvarctx_o
    = \oxtupdate{\oxvarctx}{\oxplace}{\oxsidtype}
  }{
    $\oxvarctx_o = \oxtupdate{\oxvarctx}{\oxplace}{\oxsidtype}$
  }
\end{preservation}

\oxpreservationcase{T-Copy}{\TCopy}{\oxplaceexpr}{\ECopy}{
  \oxvarctx
}

\begin{preservation}
  %% new store is valid
  \oxpresline{
    \oxstorevalidity{\oxglobalctx}{\oxvarctx}{\oxstore}
  }{
    Immediate from our premise.
  }

  %% kontctx values are still valid
  \oxpresline{
    \oxwfkontctx{\oxglobalctx}{\oxvarctx}{\overline{\oxvalue}}{\oxkontctx}
  }{
    Immediate from our premise.
  }

  %% new expression is well-typed
  \oxpresline{
    \oxtypjudge{\oxglobalctx}{\oxemptyctx}{\oxkontctx}{\oxvarctx}{\oxvalue}
    {\oxsitype}{\oxvarctx}
  }{
    Applying \Lemma{lemma:place-exprs-reduce} to
    $\oxmusafety{\oxemptyctx}{\oxvarctx}{\oximm}{\oxplaceexpr}{
      \oxset{\overline{\oxloan}}
    }$, $\oxcomputetynoprov{\oxemptyctx}{\oxvarctx}{\oximm}{\oxplaceexpr}
    {\oxsitype}$, and $\oxstorevalidity{\oxglobalctx}{\oxvarctx}{\oxstore}$
    gives us $\oxtypjudge{\oxglobalctx}{\oxemptyctx}{\oxkontctx}{\oxvarctx}{\oxvalue}
    {\oxsitype}{\oxvarctx}$.
  }

  %% new type is a subtype of the old type
  \oxpresline{
    \oxtunify[\oxcombine]{\oxemptyctx}{\oxkontctx}{\oxvarctx}
    {\oxsitype}{\oxsitype}{\oxvarctx}
  }{
    Immediate by \oxname{RR-Refl}.
  }

  %% new loan context is a sub-context of the old loan context
  \oxpresline{
    \exists \oxvarctx_o. \oxvarctx \oxintersect \oxvarctx_o = \oxvarctx
  }{
    $\oxvarctx_o = \oxvarctx$
  }
\end{preservation}

\oxpreservationcase{T-Borrow}{\TBorrow}{
  \oxref{\oxprov}{\oxmuta}{\oxplaceexpr}
}{\EBorrow}{
  \oxupdate{\oxvarctx}{
    \oxloanctxentry{\oxcprov}{
      \oxset{\overline{\oxloan}}
    }
  }
}

\begin{preservation}
  %% new store is valid
  \oxpresline{
    \oxstorevalidity{\oxglobalctx}{
      \oxupdate{\oxvarctx}{
        \oxloanctxentry{\oxcprov}{
          \oxset{\overline{\oxloan}}
        }
      }
    }{\oxstore}
  }{
    Apply \Lemma{lemma:store-validity-loan-update} to
    $\oxstorevalidity{\oxglobalctx}{\oxvarctx}{\oxstore}$ (from our premise)
    and $\oxctxswellformed{\oxglobalctx}{\oxemptyctx}{
      \oxupdate{\oxvarctx}{
        \oxloanctxentry{\oxcprov}{
          \oxset{\overline{\oxloan}}
        }
      }
    }{\oxkontctx}$ (from our premise)
    gives us $\oxstorevalidity{\oxglobalctx}{
      \oxupdate{\oxvarctx}{
        \oxloanctxentry{\oxcprov}{
          \oxset{\overline{\oxloan}}
        }
      }
    }{\oxstore}$.
  }

  %% kontctx values are still valid
  \oxpresline{
    \oxwfkontctx{\oxglobalctx}{
      \oxupdate{\oxvarctx}{
        \oxloanctxentry{\oxcprov}{
          \oxset{\overline{\oxloan}}
        }
      }
    }{\overline{\oxvalue}}{\oxkontctx}
  }{
    Inverting \oxname{WF-Temporaries} gives us $\forall i \in 1 \oxdots n. \;
    \oxtypjudge{\oxglobalctx}{\oxemptyctx}{\oxsitype_1 \oxdotsc \oxsitype_{i-1}}{\oxvarctx}{
      \oxvalue_i
    }{\oxsitype_i}{\oxvarctx}$. We can then apply
    \Lemma{lemma:values-well-typed-loan-update} to each of the typing judgments
    along with
    $\oxmusafety{\oxemptyctx}{\oxkontctx}{\oxvarctx}{\oxmuta}{\oxplaceexpr}{
      \oxset{\overline{\oxloan}}
    }$ and $\oxnotinclosure{\oxkontctx}{\oxvarctx}{\oxcprov}$
    and $\oxlookup{\oxvarctx}{\oxcprov}{\emptyset}$ (all from the premise of
    \oxname{T-Borrow}) to get $\forall i \in 1 \oxdots n. \;
    \oxtypjudge{\oxglobalctx}{\oxemptyctx}{\oxsitype_1 \oxdotsc \oxsitype_{i-1}}{
      \oxupdate{\oxvarctx}{
        \oxloanctxentry{\oxcprov}{
          \oxset{\overline{\oxloan}}
        }
      }
    }{\oxvalue_i}{\oxsitype_i}{
      \oxupdate{\oxvarctx}{
        \oxloanctxentry{\oxcprov}{
          \oxset{\overline{\oxloan}}
        }
      }
    }$. We can then apply \oxname{WF-Temporaries} to get
    $\oxwfkontctx{\oxglobalctx}{
      \oxupdate{\oxvarctx}{
        \oxloanctxentry{\oxcprov}{
          \oxset{\overline{\oxloan}}
        }
      }
    }{\overline{\oxvalue}}{\oxkontctx}$.
  }

  %% new expression is well-typed
  \oxpresline{
    \oxtypjudge{\oxglobalctx}{\oxemptyctx}{\oxkontctx}{\oxvarctx_i}{
      \oxptr{\oxreferent}
    }{\oxtref{\oxcprov}{\oxmuta}{\oxxitype}}{\oxvarctx_i}
  }{
    Applying \Lemma{lemma:norm-for-valid-referents} to
    $\oxstorevalidity{\oxglobalctx}{\oxvarctx}{\oxstore}$, and
    $\oxnorm{\oxstore}{\oxplaceexpr}{\oxreferent}{\_}{\_}$ gives us
     $\oxreferentvalidity{\oxglobalctx}{\oxvarctx}{
      \oxreferentctx[\oxplace]
    }{\oxxitype}$. Then, note that referent well-formedness does not
    depend on the contents of loan sets. This means we can also conclude
    $\oxreferentvalidity{\oxglobalctx}{
      \oxupdate{\oxvarctx}{
        \oxloanctxentry{\oxcprov}{
          \oxset{\overline{\oxloan}}
        }
      }
    }{
      \oxreferentctx[\oxplace]
    }{\oxxitype}$.

    Applying \Lemma{lemma:norm-place-prefix-in-loans} to
    $\oxstorevalidity{\oxglobalctx}{\oxvarctx}{\oxstore}$,
    $\oxnorm{\oxstore}{\oxplaceexpr}{\oxreferent}{\_}{\_}$, and
    $\oxmusafety{\oxemptyctx}{\oxvarctx}{\oxmuta}{\oxplaceexpr}{
      \oxset{\overline{\oxloan}}
    }$ gives us $\oxreferent = \oxreferentctx[\oxplace]$ and
    $\oxloanpkg{\oxmuta}{\oxplace} \in \oxset{\overline{\oxloan}}$.

    Finally, we can apply \oxname{T-Pointer} to the two facts above to get
    $\oxtypjudge{\oxglobalctx}{\oxemptyctx}{\oxkontctx}{
      \oxupdate{\oxvarctx}{
        \oxloanctxentry{\oxcprov}{
          \oxset{\overline{\oxloan}}
        }
      }
    }{
      \oxptr{\oxreferent}
    }{\oxtref{\oxcprov}{\oxmuta}{\oxxitype}}{
      \oxupdate{\oxvarctx}{
        \oxloanctxentry{\oxcprov}{
          \oxset{\overline{\oxloan}}
        }
      }
    }$.
  }

  %% new type is a subtype of the old type
  \oxpresline{
    \oxtunify[\oxcombine]{\oxemptyctx}{\oxkontctx}{\oxvarctx_i}{
      \oxtref{\oxcprov}{\oxmuta}{\oxxitype}
    }{
      \oxtref{\oxcprov}{\oxmuta}{\oxxitype}
    }{\oxvarctx_i}
  }{
    Immediate by \oxname{RR-Refl}.
  }

  %% new loan context is a sub-context of the old loan context
  \oxpresline{
    \exists \oxvarctx_o. \oxvarctx_i \oxintersect \oxvarctx_o = \oxvarctx_i
  }{
    $\oxvarctx_o = \oxvarctx_i$
  }
\end{preservation}

\oxpreservationcasemulti{T-BorrowIndex}{\TBorrowIndex}{
  \oxref{\oxprov}{\oxmuta}{\oxindex{\oxplaceexpr}{\oxexpr_i}}
}{
  \EBorrowIndex \and \EBorrowIndexOOB
}

\oxpreservationsubcaseheading{E-BorrowIndex}{
  \oxupdate{\oxvarctx^\prime}{
    \oxloanctxentry{\oxcprov}{
      \oxset{\overline{\oxloan}}
    }
  }
}

\begin{preservation}
  %% new store is valid
  \oxpresline{
    \oxstorevalidity{\oxglobalctx}{
      \oxupdate{\oxvarctx^\prime}{
        \oxloanctxentry{\oxcprov}{
          \oxset{\overline{\oxloan}}
        }
      }
    }{\oxstore}
  }{
    Applying \Lemma{lemma:values-dont-change} to the typing derivation (from
    \oxname{T-BorrowIndex}) for
    $\oxexpr$ (which we know is a value from \oxname{E-BorrowIndex}) gives us
    $\oxsubtypectx{\oxglobalctx}{\oxemptyctx}{\oxvarctx}{\oxvarctx^\prime}$.

    Then, applying \Lemma{lemma:stack-validity-related-envs} to
    $\oxsubtypectx{\oxglobalctx}{\oxemptyctx}{\oxvarctx}{\oxvarctx^\prime}$ and
    $\oxstorevalidity{\oxglobalctx}{\oxvarctx}{\oxstore}$ (from premise) gives us
    $\oxstorevalidity{\oxglobalctx}{\oxvarctx^\prime}{\oxstore}$.

    Finally, applying \Lemma{lemma:store-validity-loan-update} to
    $\oxstorevalidity{\oxglobalctx}{\oxvarctx^\prime}{\oxstore}$ gives us
    $\oxstorevalidity{\oxglobalctx}{
      \oxupdate{\oxvarctx^\prime}{
        \oxloanctxentry{\oxcprov}{
          \oxset{\overline{\oxloan}}
        }
      }
    }{\oxstore}$.
  }

  %% kontctx values are still valid
  \oxpresline{
    \oxwfkontctx{\oxglobalctx}{
      \oxupdate{\oxvarctx^\prime}{
        \oxloanctxentry{\oxcprov}{
          \oxset{\overline{\oxloan}}
        }
      }
    }{\overline{\oxvalue}}{\oxkontctx}
  }{
    Inverting \oxname{WF-Temporaries} gives us $\forall i \in 1 \oxdots n. \;
    \oxtypjudge{\oxglobalctx}{\oxemptyctx}{\oxsitype_1 \oxdotsc \oxsitype_{i-1}}{\oxvarctx}{
      \oxvalue_i
    }{\oxsitype_i}{\oxvarctx}$.

    Applying \Lemma{lemma:values-dont-change} to the typing derivation (from
    \oxname{T-BorrowIndex}) for
    $\oxexpr$ (which we know is a value from \oxname{E-BorrowIndex}) gives us
    $\oxsubtypectx{\oxglobalctx}{\oxemptyctx}{\oxvarctx}{\oxvarctx^\prime}$.

    Thus, we can apply \Lemma{lemma:value-typing-related-envs} to each of the
    typing judgments to get $\forall i \in 1 \oxdots n. \;
    \oxtypjudge{\oxglobalctx}{\oxemptyctx}{\oxemptyctx}{\oxvarctx^\prime}{
      \oxvalue_i
    }{\oxsitype_i}{\oxvarctx^\prime}$.

    We can then apply \Lemma{lemma:values-well-typed-loan-update} to each of the
    typing judgments along with
    $\oxmusafety{\oxemptyctx}{\oxkontctx}{\oxvarctx^\prime}{\oxmuta}{\oxplaceexpr}{
      \oxset{\overline{\oxloan}}
    }$ and $\oxnotinclosure{\oxkontctx}{\oxvarctx^\prime}{\oxcprov}$
    and $\oxlookup{\oxvarctx^\prime}{\oxcprov}{\emptyset}$ (all from the premise
    of \oxname{T-BorrowIndex}) to get $\forall i \in 1 \oxdots n. \;
    \oxtypjudge{\oxglobalctx}{\oxemptyctx}{\oxemptyctx}{
      \oxupdate{\oxvarctx^\prime}{
        \oxloanctxentry{\oxcprov}{
          \oxset{\overline{\oxloan}}
        }
      }
    }{\oxvalue_i}{\oxsitype_i}{
      \oxupdate{\oxvarctx^\prime}{
        \oxloanctxentry{\oxcprov}{
          \oxset{\overline{\oxloan}}
        }
      }
    }$. We can then apply \oxname{WF-Temporaries} to get
    $\oxwfkontctx{\oxglobalctx}{
      \oxupdate{\oxvarctx^\prime}{
        \oxloanctxentry{\oxcprov}{
          \oxset{\overline{\oxloan}}
        }
      }
    }{\overline{\oxvalue}}{\oxkontctx}$.
  }

  %% new expression is well-typed
  \oxpresline{
    \oxtypjudge{\oxglobalctx}{\oxemptyctx}{\oxkontctx}{\oxvarctx_i}{
      \oxindexptr{\oxreferent}{\oxnum_i}
    }{\oxtref{\oxcprov}{\oxmuta}{\oxsitype}}{\oxvarctx_i}
  }{
    Applying \Lemma{lemma:values-dont-change} to the typing derivation (from
    \oxname{T-BorrowIndex}) for
    $\oxexpr$ (which we know is a value from \oxname{E-BorrowIndex}) gives us
    $\oxsubtypectx{\oxglobalctx}{\oxemptyctx}{\oxvarctx}{\oxvarctx^\prime}$.

    Then, applying \Lemma{lemma:stack-validity-related-envs} to
    $\oxsubtypectx{\oxglobalctx}{\oxemptyctx}{\oxvarctx}{\oxvarctx^\prime}$ and
    $\oxstorevalidity{\oxglobalctx}{\oxvarctx}{\oxstore}$ (from premise) gives us
    $\oxstorevalidity{\oxglobalctx}{\oxvarctx^\prime}{\oxstore}$.

    Applying \Lemma{lemma:norm-for-valid-referents} to
    $\oxstorevalidity{\oxglobalctx}{\oxvarctx^\prime}{\oxstore}$, and
    $\oxnorm{\oxstore}{\oxplaceexpr}{\oxreferent}{\_}{
      \oxarr{\oxvalue_0 \oxdotsc \oxvalue_n}
    }$ gives us $\oxreferentvalidity{\oxglobalctx}{\oxvarctx^\prime}{
      \oxreferentctx[\oxplace]
    }{\oxxitype}$. Then, note that referent well-formedness does not
    depend on the contents of loan sets. This means we can also conclude
    $\oxreferentvalidity{\oxglobalctx}{
      \oxupdate{\oxvarctx^\prime}{
        \oxloanctxentry{\oxcprov}{
          \oxset{\overline{\oxloan}}
        }
      }
    }{
      \oxreferentctx[\oxplace]
    }{\oxxitype}$. We can then apply \oxname{WF-RefIndexArray} or
    \oxname{WF-RefIndexSlice} to get $\oxreferentvalidity{\oxglobalctx}{
      \oxupdate{\oxvarctx^\prime}{
        \oxloanctxentry{\oxcprov}{
          \oxset{\overline{\oxloan}}
        }
      }
    }{
      \oxindex{\oxreferentctx[\oxplace]}{\oxnum_i}
    }{\oxxitype}$.

    Applying \Lemma{lemma:norm-place-prefix-in-loans} to
    $\oxstorevalidity{\oxglobalctx}{\oxvarctx^\prime}{\oxstore}$,
    $\oxnorm{\oxstore}{\oxplaceexpr}{\oxreferent}{\_}{
      \oxarr{\oxvalue_0 \oxdotsc \oxvalue_n}
    }$, and
    $\oxmusafety{\oxemptyctx}{\oxvarctx^\prime}{\oxmuta}{\oxplaceexpr}{
      \oxset{\overline{\oxloan}}
    }$ gives us $\oxreferent = \oxreferent[\oxplace]$ and
    $\oxloanpkg{\oxmuta}{\oxplace} \in \oxset{\overline{\oxloan}}$.

    Finally, we can apply \oxname{T-Pointer} to the two facts above to get
    $\oxtypjudge{\oxglobalctx}{\oxemptyctx}{\oxkontctx}{
      \oxupdate{\oxvarctx}{
        \oxloanctxentry{\oxcprov}{
          \oxset{\overline{\oxloan}}
        }
      }
    }{
      \oxindexptr{\oxreferent}{\oxnum_i}
    }{\oxtref{\oxcprov}{\oxmuta}{\oxxitype}}{
      \oxupdate{\oxvarctx}{
        \oxloanctxentry{\oxcprov}{
          \oxset{\overline{\oxloan}}
        }
      }
    }$.
  }

  %% new type is a subtype of the old type
  \oxpresline{
    \oxtunify[\oxcombine]{\oxemptyctx}{\oxkontctx}{\oxvarctx_i}{
      \oxtref{\oxcprov}{\oxmuta}{\oxsitype}
    }{
      \oxtref{\oxcprov}{\oxmuta}{\oxsitype}
    }{\oxvarctx_i}
  }{
    Immediate by \oxname{RR-Refl}.
  }

  %% new loan context is a sub-context of the old loan context
  \oxpresline{
    \exists \oxvarctx_o. \oxvarctx_i \oxintersect \oxvarctx_o = \oxvarctx_i
  }{
    $\oxvarctx_o = \oxvarctx_i$
  }
\end{preservation}

\oxpreservationsubcaseheading{E-BorrowIndexOOB}{
  \oxvarctx
}

\begin{preservation}
  %% new store is valid
  \oxpresline{
    \oxstorevalidity{\oxglobalctx}{\oxvarctx}{\oxstore}
  }{
    Immediate from our premise.
  }

  %% kontctx values are still valid
  \oxpresline{
    \oxwfkontctx{\oxglobalctx}{\oxvarctx}{\overline{\oxvalue}}{\oxkontctx}
  }{
    Immediate from our premise.
  }

  %% new expression is well-typed
  \oxpresline{
    \oxtypjudge{\oxglobalctx}{\oxemptyctx}{\oxkontctx}{\oxvarctx}
    {\oxkey{abort!}(\oxdots)}{\oxtref{\oxcprov}{\oxmuta}{\oxsitype}}{\oxvarctx}
  }{
    An \oxkey{abort!} expression is well-typed (at any type) via the rule
    \oxname{T-Abort}.
  }

  %% new type is a subtype of the old type
  \oxpresline{
    \oxtunify[\oxcombine]{\oxemptyctx}{\oxkontctx}{\oxvarctx}{
      \oxtref{\oxcprov}{\oxmuta}{\oxsitype}
    }{
      \oxtref{\oxcprov}{\oxmuta}{\oxsitype}
    }{\oxvarctx}
  }{
    Immediate by \oxname{RR-Refl}.
  }

  %% new loan context is a sub-context of the old loan context
  \oxpresline{
    \exists \oxvarctx_o. \oxvarctx_i \oxintersect \oxvarctx_o = \oxvarctx_i
  }{
    $\oxvarctx_o = \oxvarctx_i$
  }
\end{preservation}

\oxpreservationcasemulti{T-BorrowSlice}{\TBorrowSlice}{
  \oxref{\oxcprov}{\oxmuta}{\oxslice{\oxplaceexpr}{\oxexpr_1}{\oxexpr_2}}
}{
  \EBorrowSlice \and \EBorrowSliceOOB
}

\oxpreservationsubcaseheading{E-BorrowSlice}{
  \oxupdate{\oxvarctx_2}{
    \oxloanctxentry{\oxcprov}{
      \oxset{\overline{\oxloan}}
    }
  }
}

\begin{preservation}
  %% new store is valid
  \oxpresline{
    \oxstorevalidity{\oxglobalctx}{
      \oxupdate{\oxvarctx_2}{
        \oxloanctxentry{\oxcprov}{
          \oxset{\overline{\oxloan}}
        }
      }
    }{\oxstore}
  }{
    Applying \Lemma{lemma:values-dont-change} to the typing derivation (from
    \oxname{T-BorrowSlice}) for $\oxexpr_1$ (which we know is a value from
    \oxname{E-BorrowSlice}) gives us
    $\oxsubtypectx{\oxglobalctx}{\oxemptyctx}{\oxvarctx}{\oxvarctx_1}$. Then,
    applying \Lemma{lemma:values-dont-change} to the typing derivation (from
    \oxname{T-BorrowSlice}) for  $\oxexpr_2$ (which we know is a value from
    \oxname{E-BorrowSlice}) gives us
    $\oxsubtypectx{\oxglobalctx}{\oxemptyctx}{\oxvarctx_1}{\oxvarctx_2}$. Then,
    by transitivity, we have
    $\oxsubtypectx{\oxglobalctx}{\oxemptyctx}{\oxvarctx}{\oxvarctx_2}$.

    Then, applying \Lemma{lemma:stack-validity-related-envs} to
    $\oxsubtypectx{\oxglobalctx}{\oxemptyctx}{\oxvarctx}{\oxvarctx_2}$ and
    $\oxstorevalidity{\oxglobalctx}{\oxvarctx}{\oxstore}$ (from premise) gives us
    $\oxstorevalidity{\oxglobalctx}{\oxvarctx_2}{\oxstore}$. Finally,
    applying \Lemma{lemma:store-validity-loan-update} to
    $\oxstorevalidity{\oxglobalctx}{\oxvarctx_2}{\oxstore}$ gives us
    $\oxstorevalidity{\oxglobalctx}{
      \oxupdate{\oxvarctx_2}{
        \oxloanctxentry{\oxcprov}{
          \oxset{\overline{\oxloan}}
        }
      }
    }{\oxstore}$.
  }

  %% kontctx values are still valid
  \oxpresline{
    \oxwfkontctx{\oxglobalctx}{
      \oxupdate{\oxvarctx_2}{
        \oxloanctxentry{\oxcprov}{
          \oxset{\overline{\oxloan}}
        }
      }
    }{\overline{\oxvalue}}{\oxkontctx}
  }{
    Inverting \oxname{WF-Temporaries} gives us $\forall i \in 1 \oxdots n. \;
    \oxtypjudge{\oxglobalctx}{\oxemptyctx}{\oxsitype_1 \oxdotsc \oxsitype_{i-1}}{\oxvarctx}{
      \oxvalue_i
    }{\oxsitype_i}{\oxvarctx}$.

    Applying \Lemma{lemma:values-dont-change} to the typing derivation (from
    \oxname{T-BorrowSlice}) for $\oxexpr_1$ (which we know is a value from
    \oxname{E-BorrowSlice}) gives us
    $\oxsubtypectx{\oxglobalctx}{\oxemptyctx}{\oxvarctx}{\oxvarctx_1}$. Then,
    applying \Lemma{lemma:values-dont-change} to the typing derivation (from
    \oxname{T-BorrowSlice}) for  $\oxexpr_2$ (which we know is a value from
    \oxname{E-BorrowSlice}) gives us
    $\oxsubtypectx{\oxglobalctx}{\oxemptyctx}{\oxvarctx_1}{\oxvarctx_2}$. Then,
    by transitivity, we have
    $\oxsubtypectx{\oxglobalctx}{\oxemptyctx}{\oxvarctx}{\oxvarctx_2}$.

    Then, we apply \Lemma{lemma:value-typing-related-envs} to each of the
    typing judgments to get $\forall i \in 1 \oxdots n. \;
    \oxtypjudge{\oxglobalctx}{\oxemptyctx}{\oxsitype_1 \oxdotsc \oxsitype_{i-1}}{\oxvarctx_2}{
      \oxvalue_i
    }{\oxsitype_i}{\oxvarctx_2}$.

    We can then apply \Lemma{lemma:values-well-typed-loan-update} to each of the
    typing judgments along with
    $\oxmusafety{\oxemptyctx}{\oxkontctx}{\oxvarctx_2}{\oxmuta}{\oxplaceexpr}{
      \oxset{\overline{\oxloan}}
    }$ and $\oxnotinclosure{\oxkontctx}{\oxvarctx_2}{\oxcprov}$
    and $\oxlookup{\oxvarctx_2}{\oxcprov}{\emptyset}$ (all from the premise
    of \oxname{T-BorrowSlice}) to get $\forall i \in 1 \oxdots n. \;
    \oxtypjudge{\oxglobalctx}{\oxemptyctx}{\oxsitype_1 \oxdotsc \oxsitype_{i-1}}{
      \oxupdate{\oxvarctx_2}{
        \oxloanctxentry{\oxcprov}{
          \oxset{\overline{\oxloan}}
        }
      }
    }{\oxvalue_i}{\oxsitype_i}{
      \oxupdate{\oxvarctx_2}{
        \oxloanctxentry{\oxcprov}{
          \oxset{\overline{\oxloan}}
        }
      }
    }$. We can then apply \oxname{WF-Temporaries} to get
    $\oxwfkontctx{\oxglobalctx}{
      \oxupdate{\oxvarctx_2}{
        \oxloanctxentry{\oxcprov}{
          \oxset{\overline{\oxloan}}
        }
      }
    }{\overline{\oxvalue}}{\oxkontctx}$.
  }

  %% new expression is well-typed
  \oxpresline{
    \oxtypjudge{\oxglobalctx}{\oxemptyctx}{\oxkontctx}{\oxvarctx_i}{
      \oxsliceptr{\oxreferent}{\oxnum_1}{\oxnum_2}
    }{
      \oxtref{\oxcprov}{\oxmuta}{\oxtslice{\oxsitype}}
    }{\oxvarctx_i}
  }{
    Applying \Lemma{lemma:values-dont-change} to the typing derivation (from
    \oxname{T-BorrowSlice}) for $\oxexpr_1$ (which we know is a value from
    \oxname{E-BorrowSlice}) gives us
    $\oxsubtypectx{\oxglobalctx}{\oxemptyctx}{\oxvarctx}{\oxvarctx_1}$. Then,
    applying \Lemma{lemma:values-dont-change} to the typing derivation (from
    \oxname{T-BorrowSlice}) for  $\oxexpr_2$ (which we know is a value from
    \oxname{E-BorrowSlice}) gives us
    $\oxsubtypectx{\oxglobalctx}{\oxemptyctx}{\oxvarctx_1}{\oxvarctx_2}$. Then,
    by transitivity, we have
    $\oxsubtypectx{\oxglobalctx}{\oxemptyctx}{\oxvarctx}{\oxvarctx_2}$.

    Then, applying \Lemma{lemma:stack-validity-related-envs} to
    $\oxsubtypectx{\oxglobalctx}{\oxemptyctx}{\oxvarctx}{\oxvarctx_2}$ and
    $\oxstorevalidity{\oxglobalctx}{\oxvarctx}{\oxstore}$ (from premise) gives us
    $\oxstorevalidity{\oxglobalctx}{\oxvarctx_2}{\oxstore}$.

    Applying \Lemma{lemma:norm-for-valid-referents} to
    $\oxstorevalidity{\oxglobalctx}{\oxvarctx_2}{\oxstore}$, and
    $\oxnorm{\oxstore}{\oxplaceexpr}{\oxreferent}{\_}{
      \oxarr{\oxvalue_0 \oxdotsc \oxvalue_n}
    }$ gives us $\oxreferentvalidity{\oxglobalctx}{\oxvarctx_2}{
      \oxreferentctx[\oxplace]
    }{\oxxitype}$. Then, note that referent well-formedness does not
    depend on the contents of loan sets. This means we can also conclude
    $\oxreferentvalidity{\oxglobalctx}{
      \oxupdate{\oxvarctx_2}{
        \oxloanctxentry{\oxcprov}{
          \oxset{\overline{\oxloan}}
        }
      }
    }{
      \oxreferentctx[\oxplace]
    }{\oxxitype}$. We can then apply \oxname{WF-RefSliceArray} or
    \oxname{WF-RefSliceSlice} to get $\oxreferentvalidity{\oxglobalctx}{
      \oxupdate{\oxvarctx^\prime}{
        \oxloanctxentry{\oxcprov}{
          \oxset{\overline{\oxloan}}
        }
      }
    }{
      \oxslice{\oxreferentctx[\oxplace]}{\oxnum_1}{\oxnum_2}
    }{\oxxitype}$.

    Applying \Lemma{lemma:norm-place-prefix-in-loans} to
    $\oxstorevalidity{\oxglobalctx}{\oxvarctx_2}{\oxstore}$,
    $\oxnorm{\oxstore}{\oxplaceexpr}{\oxreferent}{\_}{
      \oxarr{\oxvalue_0 \oxdotsc \oxvalue_n}
    }$, and
    $\oxmusafety{\oxemptyctx}{\oxvarctx_2}{\oxmuta}{\oxplaceexpr}{
      \oxset{\overline{\oxloan}}
    }$ gives us $\oxloanpkg{\oxmuta}{\oxplace} \in \oxset{\overline{\oxloan}}$.

    Finally, we can apply \oxname{T-Pointer} to the two facts above to get
    $\oxtypjudge{\oxglobalctx}{\oxemptyctx}{\oxkontctx}{
      \oxupdate{\oxvarctx}{
        \oxloanctxentry{\oxcprov}{
          \oxset{\overline{\oxloan}}
        }
      }
    }{
      \oxsliceptr{\oxreferent[\oxplace]}{\oxnum_1}{\oxnum_2}
    }{\oxtref{\oxcprov}{\oxmuta}{\oxxitype}}{
      \oxupdate{\oxvarctx}{
        \oxloanctxentry{\oxcprov}{
          \oxset{\overline{\oxloan}}
        }
      }
    }$.
  }

  %% new type is a subtype of the old type
  \oxpresline{
    \oxtunify[\oxcombine]{\oxemptyctx}{\oxkontctx}{\oxvarctx_i}{
      \oxtref{\oxcprov}{\oxmuta}{\oxtslice{\oxsitype}}
    }{
      \oxtref{\oxcprov}{\oxmuta}{\oxtslice{\oxsitype}}
    }{\oxvarctx_i}
  }{
    Immediate by \oxname{RR-Refl}.
  }

  %% new loan context is a sub-context of the old loan context
  \oxpresline{
    \exists \oxvarctx_o. \oxvarctx_i \oxintersect \oxvarctx_o = \oxvarctx_i
  }{
    $\oxvarctx_o = \oxvarctx_i$
  }
\end{preservation}

\oxpreservationsubcaseheading{E-BorrowSliceOOB}{
  \oxvarctx
}

\begin{preservation}
  %% new store is valid
  \oxpresline{
    \oxstorevalidity{\oxglobalctx}{\oxvarctx}{\oxstore}
  }{
    Immediate from our premise.
  }

  %% kontctx values are still valid
  \oxpresline{
    \oxwfkontctx{\oxglobalctx}{\oxvarctx}{\overline{\oxvalue}}{\oxkontctx}
  }{
    Immediate from our premise.
  }

  %% new expression is well-typed
  \oxpresline{
    \oxtypjudge{\oxglobalctx}{\oxemptyctx}{\oxkontctx}{\oxvarctx}
    {\oxabort{\oxdots}}{\oxtref{\oxcprov}{\oxmuta}{\oxtslice{\oxsitype}}}
    {\oxvarctx^\prime}
  }{
    An \oxkey{abort!} expression is well-typed (at any type) via the rule \oxname{T-Abort}.
  }

  %% new type is a subtype of the old type
  \oxpresline{
    \oxtunify[\oxcombine]{\oxemptyctx}{\oxkontctx}{\oxvarctx}{
      \oxtref{\oxcprov}{\oxmuta}{\oxtslice{\oxsitype}}
    }{
      \oxtref{\oxcprov}{\oxmuta}{\oxtslice{\oxsitype}}
    }{\oxvarctx}
  }{
    Immediate by \oxname{RR-Refl}.
  }

  %% new loan context is a sub-context of the old loan context
  \oxpresline{
    \exists \oxvarctx_o. \oxvarctx_i \oxintersect \oxvarctx_o = \oxvarctx_i
  }{
    $\oxvarctx_o = \oxvarctx_i$
  }
\end{preservation}

\oxpreservationcase{T-IndexCopy}{\TIndexCopy}{
  \oxindex{\oxplaceexpr}{\oxexpr_i}
}{\EIndexCopy}{
  \oxvarctx^\prime
}

\begin{preservation}
  %% new store is valid
  \oxpresline{
    \oxstorevalidity{\oxglobalctx}{\oxvarctx^\prime}{\oxstore}
  }{
    Applying \Lemma{lemma:values-dont-change} to the typing derivation (from
    \oxname{T-IndexCopy}) for $\oxexpr$ (which we know is a value from
    \oxname{E-IndexCopy}) gives us
    $\oxsubtypectx{\oxglobalctx}{\oxemptyctx}{\oxvarctx}{\oxvarctx^\prime}$.
    Then, applying \Lemma{lemma:stack-validity-related-envs} to
    $\oxsubtypectx{\oxglobalctx}{\oxemptyctx}{\oxvarctx}{\oxvarctx^\prime}$ and
    $\oxstorevalidity{\oxglobalctx}{\oxvarctx}{\oxstore}$ (from premise) gives us
    $\oxstorevalidity{\oxglobalctx}{\oxvarctx^\prime}{\oxstore}$.
  }

  %% kontctx values are still valid
  \oxpresline{
    \oxwfkontctx{\oxglobalctx}{\oxvarctx^\prime}{\overline{\oxvalue}}{\oxkontctx}
  }{
    Inverting \oxname{WF-Temporaries} gives us $\forall i \in 1 \oxdots n. \;
    \oxtypjudge{\oxglobalctx}{\oxemptyctx}{\oxsitype_1 \oxdotsc \oxsitype_{i-1}}{\oxvarctx}{
      \oxvalue_i
    }{\oxsitype_i}{\oxvarctx}$.

    Applying \Lemma{lemma:values-dont-change} to the typing derivation (from
    \oxname{T-IndexCopy}) for
    $\oxexpr$ (which we know is a value from \oxname{E-IndexCopy}) gives us
    $\oxsubtypectx{\oxglobalctx}{\oxemptyctx}{\oxvarctx}{\oxvarctx^\prime}$.

    Thus, we can apply \Lemma{lemma:value-typing-related-envs} to each of the
    typing judgments to get $\forall i \in 1 \oxdots n. \;
    \oxtypjudge{\oxglobalctx}{\oxemptyctx}{\oxsitype_1 \oxdotsc \oxsitype_{i-1}}{\oxvarctx^\prime}{
      \oxvalue_i
    }{\oxsitype_i}{\oxvarctx^\prime}$.

    Finally, applying \oxname{WF-Temporaries} gives us
    $\oxwfkontctx{\oxglobalctx}{\oxvarctx^\prime}{\overline{\oxvalue}}{\oxkontctx}$.
  }

  %% new expression is well-typed
  \oxpresline{
    \oxtypjudge{\oxglobalctx}{\oxemptyctx}{\oxkontctx}{\oxvarctx^\prime}{
      \oxvalue_{n_i}
    }{\oxsitype}{\oxvarctx^\prime}
  }{
    Applying \Lemma{lemma:values-dont-change} to the typing derivation (from
    \oxname{T-IndexCopy}) for $\oxexpr$ (which we know is a value from
    \oxname{E-IndexCopy}) gives us
    $\oxsubtypectx{\oxglobalctx}{\oxemptyctx}{\oxvarctx}{\oxvarctx^\prime}$.
    Then, applying \Lemma{lemma:stack-validity-related-envs} to
    $\oxsubtypectx{\oxglobalctx}{\oxemptyctx}{\oxvarctx}{\oxvarctx^\prime}$ and
    $\oxstorevalidity{\oxglobalctx}{\oxvarctx}{\oxstore}$ (from premise) gives us
    $\oxstorevalidity{\oxglobalctx}{\oxvarctx^\prime}{\oxstore}$.

    Applying \Lemma{lemma:place-exprs-reduce} to
    $\oxmusafety{\oxemptyctx}{\oxvarctx^\prime}{\oximm}{\oxplaceexpr}{
      \oxset{\overline{\oxloan}}
    }$, $\oxcomputetynoprov{\oxemptyctx}{\oxvarctx^\prime}{\oximm}{\oxplaceexpr}
    {\oxsitype}$, and $\oxstorevalidity{\oxglobalctx}{\oxvarctx^\prime}{\oxstore}$
    % gives us $\oxtypjudge{\oxglobalctx}{\oxemptyctx}{\oxvarctx^\prime}{
    %   \oxarr{\oxvalue_0 \oxdotsc \oxvalue_{\oxnum_i} \oxdotsc \oxvalue_n}
    % }{\oxxitype}{\oxvarctx^\prime}$. Then, by inversion of either \oxname{T-Array} or
    \oxname{T-Slice} (based on whether $\oxxitype = \oxtarr{\oxsitype}{\oxnum}$ or
    $\oxtslice{\oxsitype}$ respectively), we get
    $\forall i \in \oxset{1 \oxdots n}. \;
    \oxtypjudge{\oxglobalctx}{\oxemptyctx}{\oxkontctx}{\oxvarctx^\prime}{
      \oxvalue_i
    }{\oxsitype}{\oxvarctx^\prime}$ (after accounting for the fact that the
    constituent expressions are values and the output environment matches the
    input environment). Thus, we can pick out specifically that
    $\oxtypjudge{\oxglobalctx}{\oxemptyctx}{\oxkontctx}{\oxvarctx^\prime}{
      \oxvalue_i
    }{\oxsitype}{\oxvarctx^\prime}$.
  }

  %% new type is a subtype of the old type
  \oxpresline{
    \oxtunify[\oxcombine]{\oxemptyctx}{\oxkontctx}{\oxvarctx^\prime}{\oxsitype}{\oxsitype}{\oxvarctx^\prime}
  }{
    Immediate by \oxname{RR-Refl}.
  }

  %% new loan context is a sub-context of the old loan context
  \oxpresline{
    \exists \oxvarctx_o. \oxvarctx^\prime \oxintersect \oxvarctx_o = \oxvarctx^\prime
  }{
    $\oxvarctx_o = \oxvarctx^\prime$
  }
\end{preservation}

\oxpreservationcase{T-Seq}{\TSeq}{\oxseq{\oxexpr_1}{\oxexpr_2}}{\ESeq}{
  \oxgcloans{\oxkontctx}{\oxvarctx_1}
}

\begin{preservation}
  %% new store is valid
  \oxpresline{
    \oxstorevalidity{\oxglobalctx}{
      \oxgcloans{\oxkontctx}{\oxvarctx_1}
    }{\oxstore}
  }{
    Applying \Lemma{lemma:values-dont-change} to the typing derivation (from
    \oxname{T-Seq}) for $\oxexpr_1$ (which we know is a value from
    \oxname{E-Seq}) gives us
    $\oxsubtypectx{\oxglobalctx}{\oxemptyctx}{\oxvarctx}{\oxvarctx_1}$.
    Applying \Lemma{lemma:stack-validity-related-envs} with this and
    $\oxstorevalidity{\oxglobalctx}{\oxvarctx}{\oxstore}$ (from
    premise) gives us
    $\oxstorevalidity{\oxglobalctx}{\oxvarctx_1}{\oxstore}$. By definition of
    \oxname{R-Env}, we know that $\oxsubtypectx{\oxglobalctx}{\oxemptyctx}
    {\oxvarctx_1}{\oxgcloans{\oxkontctx}{\oxvarctx_1}}$. Then, we can apply
    \Lemma{lemma:stack-validity-related-envs} to get
    $\oxstorevalidity{\oxglobalctx}{
      \oxgcloans{\oxkontctx}{\oxvarctx_1}
    }{\oxstore}$.
  }

  %% kontctx values are still valid
  \oxpresline{
    \oxwfkontctx{\oxglobalctx}{
      \oxgcloans{\oxkontctx}{\oxvarctx_1}
    }{\overline{\oxvalue}}{\oxkontctx}
  }{
    Inverting \oxname{WF-Temporaries} gives us $\forall i \in 1 \oxdots n. \;
    \oxtypjudge{\oxglobalctx}{\oxemptyctx}{\oxsitype_1 \oxdotsc \oxsitype_{i-1}}{\oxvarctx}{
      \oxvalue_i
    }{\oxsitype_i}{\oxvarctx}$.

    Applying \Lemma{lemma:values-dont-change} to the typing derivation (from
    \oxname{T-Seq}) for
    $\oxexpr_1$ (which we know is a value from \oxname{E-Seq}) gives us
    $\oxsubtypectx{\oxglobalctx}{\oxemptyctx}{\oxvarctx}{\oxvarctx_1}$. Then,
    we note that definitionally $\oxsubtypectx{\oxglobalctx}{\oxemptyctx}{\oxvarctx_1}{
      \oxgcloans{\oxkontctx}{\oxvarctx_1}
    }$ since we can only clear loans that are not present in any of the types.
    Then, by transitivity, we have
    $\oxsubtypectx{\oxglobalctx}{\oxemptyctx}{\oxvarctx}{
      \oxgcloans{\oxkontctx}{\oxvarctx_1}
    }$.

    Thus, we can apply \Lemma{lemma:value-typing-related-envs} to each of the
    typing judgments to get $\forall i \in 1 \oxdots n. \;
    \oxtypjudge{\oxglobalctx}{\oxemptyctx}{\oxsitype_1 \oxdotsc \oxsitype_{i-1}}{
      \oxgcloans{\oxkontctx}{\oxvarctx_1}
    }{
      \oxvalue_i
    }{\oxsitype_i}{
      \oxgcloans{\oxkontctx}{\oxvarctx_1}
    }$.

    Finally, applying \oxname{WF-Temporaries} gives us
    $\oxwfkontctx{\oxglobalctx}{
      \oxgcloans{\oxkontctx}{\oxvarctx_1}
    }{\overline{\oxvalue}}{\oxkontctx}$.
  }

  %% new expression is well-typed
  \oxpresline{
    \oxtypjudge{\oxglobalctx}{\oxemptyctx}{\oxkontctx}{\oxgcloans{\oxkontctx}{\oxvarctx_1}}
    {\oxexpr_2}{\oxsitype_2}{\oxvarctx_2}
  }{
    Immediate from the premise of \oxname{T-Seq}.
  }

  %% new type is a subtype of the old type
  \oxpresline{
    \oxtunify[\oxcombine]{\oxemptyctx}{\oxkontctx}{\oxvarctx_2}
    {\oxsitype_2}{\oxsitype_2}{\oxvarctx_2}
  }{
    Immediate by \oxname{RR-Refl}.
  }

  %% new loan context is a sub-context of the old loan context
  \oxpresline{
    \exists \oxvarctx_o. \oxvarctx_2 \oxintersect \oxvarctx_o = \oxvarctx_2
  }{
    $\oxvarctx_o = \oxvarctx_2$
  }
\end{preservation}

\oxpreservationcasemulti{T-Branch}{\TBranch}{
  \oxbranch{\oxexpr_1}{\oxexpr_2}{\oxexpr_3}
}{\EIfTrue \and \EIfFalse}

\oxpreservationsubcaseheading{E-IfTrue}{
  \oxvarctx_1
}

\begin{preservation}
  %% new store is valid
  \oxpresline{
    \oxstorevalidity{\oxglobalctx}{\oxvarctx_1}{\oxstore}
  }{
    Applying \Lemma{lemma:values-dont-change} to the typing derivation (from
    \oxname{T-Branch}) for $\oxexpr_1$ (which we know is a value from
    \oxname{E-IfTrue}) gives us
    $\oxsubtypectx{\oxglobalctx}{\oxemptyctx}{\oxvarctx}{\oxvarctx_1}$. Then,
    applying \Lemma{lemma:stack-validity-related-envs} with this and
    $\oxstorevalidity{\oxglobalctx}{\oxvarctx}{\oxstore}$ (from
    premise) gives us
    $\oxstorevalidity{\oxglobalctx}{\oxvarctx_1}{\oxstore}$.
  }

  %% kontctx values are still valid
  \oxpresline{
    \oxwfkontctx{\oxglobalctx}{\oxvarctx_1}{\overline{\oxvalue}}{\oxkontctx}
  }{
    Inverting \oxname{WF-Temporaries} gives us $\forall i \in 1 \oxdots n. \;
    \oxtypjudge{\oxglobalctx}{\oxemptyctx}{\oxsitype_1 \oxdotsc \oxsitype_{i-1}}{\oxvarctx}{
      \oxvalue_i
    }{\oxsitype_i}{\oxvarctx}$.

    Applying \Lemma{lemma:values-dont-change} to the typing derivation (from
    \oxname{T-Branch}) for
    $\oxexpr_1$ (which we know is a value from \oxname{E-IfTrue}) gives us
    $\oxsubtypectx{\oxglobalctx}{\oxemptyctx}{\oxvarctx}{\oxvarctx_1}$.

    Thus, we can apply \Lemma{lemma:value-typing-related-envs} to each of the
    typing judgments to get $\forall i \in 1 \oxdots n. \;
    \oxtypjudge{\oxglobalctx}{\oxemptyctx}{\oxsitype_1 \oxdotsc \oxsitype_{i-1}}{\oxvarctx_1}{
      \oxvalue_i
    }{\oxsitype_i}{\oxvarctx_1}$.

    Finally, applying \oxname{WF-Temporaries} gives us
    $\oxwfkontctx{\oxglobalctx}{\oxvarctx_1}{\overline{\oxvalue}}{\oxkontctx}$.
  }

  %% new expression is well-typed
  \oxpresline{
    \oxtypjudge{\oxglobalctx}{\oxemptyctx}{\oxkontctx}{\oxvarctx_1}{\oxexpr_2}
    {\oxsitype}{\oxvarctx_2}
  }{
    Immediate from premise of \oxname{T-Branch}.
  }

  %% new type is a subtype of the old type
  \oxpresline{
    \oxtunify[\oxcombine]{\oxemptyctx}{\oxkontctx}{\oxvarctx_2}
    {\oxsitype_2}{\oxsitype}{\oxvarctx_2^\prime}
  }{
    Immediate from premise of \oxname{T-Branch}.
  }

  %% new loan context is a sub-context of the old loan context
  \oxpresline{
    \exists \oxvarctx_o. \oxvarctx_2^\prime \oxintersect \oxvarctx_o = \oxvarctx^\prime
  }{
    $\oxvarctx_o = \oxvarctx_3^\prime$
  }
\end{preservation}

\oxpreservationsubcaseheading{E-IfFalse}{
  \oxvarctx_1
}

\begin{preservation}
  %% new store is valid
  \oxpresline{
    \oxstorevalidity{\oxglobalctx}{\oxvarctx_1}{\oxstore}
  }{
    Applying \Lemma{lemma:values-dont-change} to the typing derivation (from
    \oxname{T-Branch}) for $\oxexpr_1$ (which we know is a value from
    \oxname{E-IfFalse}) gives us
    $\oxsubtypectx{\oxglobalctx}{\oxemptyctx}{\oxvarctx}{\oxvarctx_1}$. Then,
    applying \Lemma{lemma:stack-validity-related-envs} with this and
    $\oxstorevalidity{\oxglobalctx}{\oxvarctx}{\oxstore}$ (from
    premise) gives us
    $\oxstorevalidity{\oxglobalctx}{\oxvarctx_1}{\oxstore}$.
  }

  %% kontctx values are still valid
  \oxpresline{
    \oxwfkontctx{\oxglobalctx}{\oxvarctx_1}{\overline{\oxvalue}}{\oxkontctx}
  }{
    Inverting \oxname{WF-Temporaries} gives us $\forall i \in 1 \oxdots n. \;
    \oxtypjudge{\oxglobalctx}{\oxemptyctx}{\oxsitype_1 \oxdotsc \oxsitype_{i-1}}{\oxvarctx}{
      \oxvalue_i
    }{\oxsitype_i}{\oxvarctx}$.

    Applying \Lemma{lemma:values-dont-change} to the typing derivation (from
    \oxname{T-Branch}) for
    $\oxexpr_1$ (which we know is a value from \oxname{E-IfFalse}) gives us
    $\oxsubtypectx{\oxglobalctx}{\oxemptyctx}{\oxvarctx}{\oxvarctx_1}$.

    Thus, we can apply \Lemma{lemma:value-typing-related-envs} to each of the
    typing judgments to get $\forall i \in 1 \oxdots n. \;
    \oxtypjudge{\oxglobalctx}{\oxemptyctx}{\oxsitype_1 \oxdotsc \oxsitype_{i-1}}{\oxvarctx_1}{
      \oxvalue_i
    }{\oxsitype_i}{\oxvarctx_1}$.

    Finally, applying \oxname{WF-Temporaries} gives us
    $\oxwfkontctx{\oxglobalctx}{\oxvarctx_1}{\overline{\oxvalue}}{\oxkontctx}$.
  }

  %% new expression is well-typed
  \oxpresline{
    \oxtypjudge{\oxglobalctx}{\oxemptyctx}{\oxkontctx}{\oxvarctx_1}
    {\oxexpr_3}{\oxsitype_3}{\oxvarctx_3}
  }{
    Immediate from premise of \oxname{T-Branch}.
  }

  %% new type is a subtype of the old type
  \oxpresline{
    \oxtunify[\oxcombine]{\oxemptyctx}{\oxkontctx}{\oxvarctx_3}
    {\oxsitype_3}{\oxsitype}{\oxvarctx_3^\prime}
  }{
    Immediate from premise of \oxname{T-Branch}.
  }

  %% new loan context is a sub-context of the old loan context
  \oxpresline{
    \exists \oxvarctx_o. \oxvarctx_3^\prime \oxintersect \oxvarctx_o = \oxvarctx^\prime
  }{
    $\oxvarctx_o = \oxvarctx_2^\prime$ (note that $\oxintersect$ commutes)
  }
\end{preservation}

\oxpreservationcasemulti{T-Match}{\TMatch}{
  \oxmatch{\oxexpr}{\oxid_1}{\oxexpr_1}{\oxid_2}{\oxexpr_2}
}{\EMatchLeft \and \EMatchRight}

\oxpreservationsubcaseheading{E-MatchLeft}{
  \oxextendctx{\oxvarctx^\prime}{
    \oxvarctxentry{\oxid_1}{\oxsitype_l}
  }
}

\begin{preservation}
  %% new store is valid
  \oxpresline{
    \oxstorevalidity{\oxglobalctx}{
      \oxextendctx{\oxvarctx^\prime}{
        \oxvarctxentry{\oxid_1}{\oxsitype_l}
      }
    }{
      \oxextendctx{\oxstore}{
        \oxstoreentry{\oxid_1}{\oxvalue}
      }
    }
  }{
    Applying \Lemma{lemma:values-dont-change} to the typing derivation (from
    \oxname{T-Match}) for $\oxexpr$ (which we know is a value from
    \oxname{E-MatchLeft}) gives us
    $\oxsubtypectx{\oxglobalctx}{\oxemptyctx}{\oxvarctx}{\oxvarctx^\prime}$. Then,
    applying \Lemma{lemma:stack-validity-related-envs} with this and
    $\oxstorevalidity{\oxglobalctx}{\oxvarctx}{\oxstore}$ (from
    premise) gives us
    $\oxstorevalidity{\oxglobalctx}{\oxvarctx^\prime}{\oxstore}$. Then, using
    $\oxtypjudge{\oxglobalctx}{\oxemptyctx}{\oxkontctx}{\oxvarctx}{\oxvalue}{\oxsitype_l}
    {\oxvarctx^\prime}$ (which we get by inversion of \oxname{T-Inl} in the
    derivation for $\oxexpr$) with \oxname{WF-StackFrame} allows us to finally
    conclude $\oxstorevalidity{\oxglobalctx}{
      \oxextendctx{\oxvarctx^\prime}{
        \oxvarctxentry{\oxid_1}{\oxsitype_l}
      }
    }{
      \oxextendctx{\oxstore}{
        \oxstoreentry{\oxid_1}{\oxvalue}
      }
    }$.
  }

  %% kontctx values are still valid
  \oxpresline{
    \oxwfkontctx{\oxglobalctx}{
      \oxextendctx{\oxvarctx^\prime}{
        \oxvarctxentry{\oxid_1}{\oxsitype_l}
      }
    }{\overline{\oxvalue}}{\oxkontctx}
  }{
    Inverting \oxname{WF-Temporaries} gives us $\forall i \in 1 \oxdots n. \;
    \oxtypjudge{\oxglobalctx}{\oxemptyctx}{\oxsitype_1 \oxdotsc \oxsitype_{i-1}}{\oxvarctx}{
      \oxvalue_i
    }{\oxsitype_i}{\oxvarctx}$.

    Applying \Lemma{lemma:values-dont-change} to the typing derivation (from
    \oxname{T-Match}) for
    $\oxexpr$ (which we know is a value from \oxname{E-MatchLeft}) gives us
    $\oxsubtypectx{\oxglobalctx}{\oxemptyctx}{\oxvarctx}{\oxvarctx^\prime}$.

    Thus, we can apply \Lemma{lemma:value-typing-related-envs} to each of the
    typing judgments to get $\forall i \in 1 \oxdots n. \;
    \oxtypjudge{\oxglobalctx}{\oxemptyctx}{\oxsitype_1 \oxdotsc \oxsitype_{i-1}}{\oxvarctx^\prime}{
      \oxvalue_i
    }{\oxsitype_i}{\oxvarctx^\prime}$.

    Applying \Lemma{lemma:values-well-typed-extension} to each of the value
    typing judgments and $\forall \oxcprov \in
    \oxfprovs{\oxtsum{\oxsitype_l}{\oxsitype_r}}. \;
    \oxnotreborrowed{\oxvarctx^\prime}{\oxcprov}$ (from the premise of
    \oxname{T-Match}) gives us
    $\forall i \in 1 \oxdots n. \;
    \oxtypjudge{\oxglobalctx}{\oxemptyctx}{\oxsitype_1 \oxdotsc \oxsitype_{i-1}}{
      \oxextendctx{\oxvarctx^\prime}{
        \oxvarctxentry{\oxid_1}{\oxsitype_l}
      }
    }{
      \oxvalue_i
    }{\oxsitype_i}{
      \oxextendctx{\oxvarctx^\prime}{
        \oxvarctxentry{\oxid_1}{\oxsitype_l}
      }
    }$.

    Finally, applying \oxname{WF-Temporaries} gives us
    $\oxwfkontctx{\oxglobalctx}{
      \oxextendctx{\oxvarctx^\prime}{
        \oxvarctxentry{\oxid_1}{\oxsitype_l}
      }
    }{\overline{\oxvalue}}{\oxkontctx}$.
  }

  %% new expression is well-typed
  \oxpresline{
    \oxtypjudge{\oxglobalctx}{\oxemptyctx}{\oxkontctx}{
      \oxextendctx{\oxvarctx^\prime}{
        \oxvarctxentry{\oxid_1}{\oxsitype_l}
      }
    }{\oxexpr_1}
    {\oxsitype_1}{\oxvarctx_1}
  }{
    Immediate from applying \oxname{T-Shift} to
    $\oxtypjudge{\oxglobalctx}{\oxemptyctx}{\oxkontctx}{
      \oxextendctx{\oxvarctx^\prime}{
        \oxvarctxentry{\oxid_1}{\oxsitype_l}
      }
    }{
      \oxexpr_1
    }{\oxsitype_1}{
      \oxextendctx{\oxvarctx_1}{
        \oxvarctxentry{\oxid_1}{\oxsdtype_l}
      }
    }$ from the premise of \oxname{T-Match}.
  }

  %% new type is a subtype of the old type
  \oxpresline{
    \oxtunify[\oxcombine]{\oxemptyctx}{\oxkontctx}{\oxvarctx_1}
    {\oxsitype_1}{\oxsitype}{\oxvarctx_1^\prime}
  }{
    Immediate from premise of \oxname{T-Match}.
  }

  %% new loan context is a sub-context of the old loan context
  \oxpresline{
    \exists \oxvarctx_o. \oxvarctx_1^\prime \oxintersect \oxvarctx_o = \oxvarctx^\prime
  }{
    $\oxvarctx_o = \oxvarctx_2^\prime$
  }
\end{preservation}

\oxpreservationsubcaseheading{E-MatchRight}{
  \oxextendctx{\oxvarctx^\prime}{
    \oxvarctxentry{\oxid_2}{\oxsitype_r}
  }
}

\begin{preservation}
  %% new store is valid
  \oxpresline{
    \oxstorevalidity{\oxglobalctx}{
      \oxextendctx{\oxvarctx^\prime}{
        \oxvarctxentry{\oxid_2}{\oxsitype_r}
      }
    }{
      \oxextendctx{\oxstore}{
        \oxstoreentry{\oxid_2}{\oxvalue}
      }
    }
  }{
    Applying \Lemma{lemma:values-dont-change} to the typing derivation (from
    \oxname{T-Match}) for $\oxexpr$ (which we know is a value from
    \oxname{E-MatchRight}) gives us
    $\oxsubtypectx{\oxglobalctx}{\oxemptyctx}{\oxvarctx}{\oxvarctx^\prime}$. Then,
    applying \Lemma{lemma:stack-validity-related-envs} with this and
    $\oxstorevalidity{\oxglobalctx}{\oxvarctx}{\oxstore}$ (from
    premise) gives us
    $\oxstorevalidity{\oxglobalctx}{\oxvarctx^\prime}{\oxstore}$. Then, using
    $\oxtypjudge{\oxglobalctx}{\oxemptyctx}{\oxkontctx}{\oxvarctx}{\oxvalue}{\oxsitype_r}
    {\oxvarctx^\prime}$ (which we get by inversion of \oxname{T-Inr} in the
    derivation for $\oxexpr$) with \oxname{WF-StackFrame} allows us to finally
    conclude $\oxstorevalidity{\oxglobalctx}{
      \oxextendctx{\oxvarctx^\prime}{
        \oxvarctxentry{\oxid_2}{\oxsitype_r}
      }
    }{
      \oxextendctx{\oxstore}{
        \oxstoreentry{\oxid_2}{\oxvalue}
      }
    }$.
  }

  %% kontctx values are still valid
  \oxpresline{
    \oxwfkontctx{\oxglobalctx}{
      \oxextendctx{\oxvarctx^\prime}{
        \oxvarctxentry{\oxid_2}{\oxsitype_r}
      }
    }{\overline{\oxvalue}}{\oxkontctx}
  }{
    Inverting \oxname{WF-Temporaries} gives us $\forall i \in 1 \oxdots n. \;
    \oxtypjudge{\oxglobalctx}{\oxemptyctx}{\oxsitype_1 \oxdotsc \oxsitype_{i-1}}{\oxvarctx}{
      \oxvalue_i
    }{\oxsitype_i}{\oxvarctx}$.

    Applying \Lemma{lemma:values-dont-change} to the typing derivation (from
    \oxname{T-Match}) for
    $\oxexpr$ (which we know is a value from \oxname{E-MatchLeft}) gives us
    $\oxsubtypectx{\oxglobalctx}{\oxemptyctx}{\oxvarctx}{\oxvarctx^\prime}$.

    Thus, we can apply \Lemma{lemma:value-typing-related-envs} to each of the
    typing judgments to get $\forall i \in 1 \oxdots n. \;
    \oxtypjudge{\oxglobalctx}{\oxemptyctx}{\oxsitype_1 \oxdotsc \oxsitype_{i-1}}{\oxvarctx^\prime}{
      \oxvalue_i
    }{\oxsitype_i}{\oxvarctx^\prime}$.

    Applying \Lemma{lemma:values-well-typed-extension} to each of the value
    typing judgments and $\forall \oxcprov \in
    \oxfprovs{\oxtsum{\oxsitype_l}{\oxsitype_r}}. \;
    \oxnotreborrowed{\oxvarctx^\prime}{\oxcprov}$ (from the premise of
    \oxname{T-Match}) gives us
    $\forall i \in 1 \oxdots n. \;
    \oxtypjudge{\oxglobalctx}{\oxemptyctx}{\oxsitype_1 \oxdotsc \oxsitype_{i-1}}{
      \oxextendctx{\oxvarctx^\prime}{
        \oxvarctxentry{\oxid_2}{\oxsitype_r}
      }
    }{
      \oxvalue_i
    }{\oxsitype_i}{
      \oxextendctx{\oxvarctx^\prime}{
        \oxvarctxentry{\oxid_2}{\oxsitype_r}
      }
    }$.

    Finally, applying \oxname{WF-Temporaries} gives us
    $\oxwfkontctx{\oxglobalctx}{
      \oxextendctx{\oxvarctx^\prime}{
        \oxvarctxentry{\oxid_2}{\oxsitype_r}
      }
    }{\overline{\oxvalue}}{\oxkontctx}$.
  }

  %% new expression is well-typed
  \oxpresline{
    \oxtypjudge{\oxglobalctx}{\oxemptyctx}{\oxkontctx}{
      \oxextendctx{\oxvarctx^\prime}{
        \oxvarctxentry{\oxid_2}{\oxsitype_r}
      }
    }{\oxexpr_2}
    {\oxsitype_2}{\oxvarctx_2}
  }{
    Immediate from applying \oxname{T-Shift} to
    $\oxtypjudge{\oxglobalctx}{\oxemptyctx}{\oxkontctx}{
      \oxextendctx{\oxvarctx^\prime}{
        \oxvarctxentry{\oxid_2}{\oxsitype_r}
      }
    }{
      \oxexpr_2
    }{\oxsitype_2}{
      \oxextendctx{\oxvarctx_2}{
        \oxvarctxentry{\oxid_2}{\oxsdtype_r}
      }
    }$ from the premise of \oxname{T-Match}.
  }

  %% new type is a subtype of the old type
  \oxpresline{
    \oxtunify[\oxcombine]{\oxemptyctx}{\oxkontctx}{\oxvarctx_2}
    {\oxsitype_2}{\oxsitype}{\oxvarctx_2^\prime}
  }{
    Immediate from premise of \oxname{T-Match}.
  }

  %% new loan context is a sub-context of the old loan context
  \oxpresline{
    \exists \oxvarctx_o. \oxvarctx_2^\prime \oxintersect \oxvarctx_o = \oxvarctx^\prime
  }{
    $\oxvarctx_o = \oxvarctx_1^\prime$ (note that $\oxintersect$ commutes)
  }
\end{preservation}

\oxpreservationcase{T-Assign}{\TAssign}{
  \oxassign{\oxplace}{\oxexpr_a}
}{\EAssign}{
  \oxgcloans{\oxkontctx}{\oxtupdate{\oxvarctx^\prime}{\oxplace}{\oxsitype}}
}

\begin{preservation}
  %% new store is valid
  \oxpresline{
    \oxstorevalidity{\oxglobalctx}{
      \oxgcloans{\oxkontctx}{\oxtupdate{\oxvarctx^\prime}{\oxplace}{\oxsitype}}
    }{
      \oxvupdate{\oxstore}{\oxplace}{\oxvaluectx[\oxvalue]}
    }
  }{
    From our premise, we have
    $\oxstorevalidity{\oxglobalctx}{\oxvarctx}{\oxstore}$. Then, applying
    \Lemma{lemma:values-dont-change} to the typing derivation
    $\oxtypjudge{\oxglobalctx}{\oxemptyctx}{\oxkontctx}{\oxvarctx}{
      \oxvalue
    }{\oxsitype}{\oxvarctx_1}$ (from the premise of \oxname{T-Assign} combined
    with the fact that $\oxexpr$ is a value from \oxname{E-Assign}) gives us
    $\oxsubtypectx{\oxglobalctx}{\oxemptyctx}{\oxvarctx}{\oxvarctx_1}$. Then,
    applying \Lemma{lemma:stack-validity-related-envs} to these two facts gives
    us $\oxstorevalidity{\oxglobalctx}{\oxvarctx_1}{\oxstore}$.

    Finally, we apply \Lemma{lemma:stack-validity-assign} to
    $\oxstorevalidity{\oxglobalctx}{\oxvarctx_1}{\oxstore}$ and
    $\oxlookup{{\oxvarctx_1}}{\oxplace}{\oxsxtype}$ (from premise) and
    $\oxtunify[\oxnoop]{\oxemptyctx}{\oxkontctx}{
      \oxrsub{\oxvarctx_1}{\oxplace}
    }{\oxsitype}{\oxsxtype}{\oxvarctx^\prime}$ (from the premise of
    \oxname{T-Assign}) and
    $\oxmusafety{\oxemptyctx}{\oxvarctx^\prime}{\oxmut}{\oxplace}{
      \oxset{\oxloanpkg{\oxmut}{\oxplace}}
    }$ (from premise of \oxname{T-Assign} and
    $\oxnorm{\oxstore}{\oxplace}{\oxplace}{\oxvaluectx}{\_}$ and $\oxplace =
    \oxproj{\oxid}{\oxpath}$ (both from the premise of \oxname{E-Assign}),
    $\oxtypjudge{\oxglobalctx}{\oxemptyctx}{\oxkontctx}{
      \oxrsub{\oxvarctx_1}{\oxplaceexpr} }{\oxvalue}{\oxsitype}{
      \oxrsub{\oxvarctx_1}{\oxplaceexpr} }$ (from above) and
    $\oxstorevalidity{\oxglobalctx}{
      \oxgcloans{\oxkontctx}{\oxtupdate{\oxvarctx^\prime}{\oxplace}{\oxsitype}}
    }{
      \oxvupdate{\oxstore}{\oxid}{\oxvaluectx[\oxvalue]}
    }$.
  }

  %% kontctx values are still valid
  \oxpresline{
    \oxwfkontctx{\oxglobalctx}{
      \oxgcloans{\oxkontctx}{\oxtupdate{\oxvarctx^\prime}{\oxplace}{\oxsitype}}
    }{\overline{\oxvalue}}{\oxkontctx}
  }{
    Inverting \oxname{WF-Temporaries} gives us $\forall i \in 1 \oxdots n. \;
    \oxtypjudge{\oxglobalctx}{\oxemptyctx}{\oxsitype_1 \oxdotsc \oxsitype_{i-1}}{\oxvarctx}{
      \oxvalue_i
    }{\oxsitype_i}{\oxvarctx}$.

    Applying \Lemma{lemma:values-dont-change} to the typing derivation (from
    \oxname{T-Let}) for
    $\oxexpr_1$ (which we know is a value from \oxname{E-Let}) gives us
    $\oxsubtypectx{\oxglobalctx}{\oxemptyctx}{\oxvarctx}{\oxvarctx_1}$.
    Then, we can apply \Lemma{lemma:value-typing-related-envs} to each of the
    typing judgments to get $\forall i \in 1 \oxdots n. \;
    \oxtypjudge{\oxglobalctx}{\oxemptyctx}{\oxsitype_1 \oxdotsc \oxsitype_{i-1}}{\oxvarctx_1}{
      \oxvalue_i
    }{\oxsitype_i}{\oxvarctx_1}$.

    Then, we apply \Lemma{lemma:values-well-typed-assignment} to each of these
    typing judgments to get $\forall i \in 1 \oxdots n. \;
    \oxtypjudge{\oxglobalctx}{\oxemptyctx}{\oxsitype_1 \oxdotsc \oxsitype_{i-1}}{
      \oxgcloans{\oxkontctx}{\oxtupdate{\oxvarctx^\prime}{\oxplace}{\oxsitype}}
    }{
      \oxvalue_i
    }{\oxsitype_i}{
      \oxgcloans{\oxkontctx}{\oxtupdate{\oxvarctx^\prime}{\oxplace}{\oxsitype}}
    }$. After which we can apply \oxname{WF-Temporaries} to conclude
    $\oxwfkontctx{\oxglobalctx}{
      \oxgcloans{\oxkontctx}{\oxtupdate{\oxvarctx^\prime}{\oxplace}{\oxsitype}}
    }{\overline{\oxvalue}}{\oxkontctx}$.
  }

  %% new expression is well-typed
  \oxpresline{
    \oxtypjudge{\oxglobalctx}{\oxemptyctx}{\oxkontctx}{
      \oxvarctx_i
    }{\oxunit}{\oxtunit}{
      \oxvarctx_i
    }
  }{
    Immediate by \oxname{T-Unit}.
  }

  %% new type is a subtype of the old type
  \oxpresline{
    \oxtunify[\oxcombine]{\oxemptyctx}{\oxkontctx}{
      \oxvarctx_i
    }{\oxtunit}{\oxtunit}{
      \oxvarctx_i
    }
  }{
    Immediate by \oxname{RR-Refl}.
  }

  %% new loan context is a sub-context of the old loan context
  \oxpresline{
    \exists \oxvarctx_o.
    \oxvarctx_i
    \oxintersect \oxvarctx_o =
    \oxvarctx_i
  }{
    $\oxvarctx_o = \oxvarctx_i =
    \oxgcloans{\oxkontctx}{\oxtupdate{\oxvarctx^\prime}{\oxplace}{\oxsitype}}$
  }
\end{preservation}

\oxpreservationcase{T-AssignDeref}{\TAssignDeref}{
  \oxassign{\oxplaceexpr}{\oxexpr_a}
}{\EAssign}{
  \oxvarctx^\prime
}

\begin{preservation}
  %% new store is valid
  \oxpresline{
    \oxstorevalidity{\oxglobalctx}{\oxvarctx^\prime}{
      \oxvupdate{\oxstore}{\oxplace}{\oxvaluectx[\oxvalue]}
    }
  }{
    From our premise, we have
    $\oxstorevalidity{\oxglobalctx}{\oxvarctx}{\oxstore}$. Then, applying
    \Lemma{lemma:values-dont-change} to the typing derivation
    $\oxtypjudge{\oxglobalctx}{\oxemptyctx}{\oxkontctx}{\oxvarctx}{
      \oxvalue
    }{\oxsitype}{\oxvarctx_1}$ (from the premise of \oxname{T-AssignDeref} combined
    with the fact that $\oxexpr$ is a value from \oxname{E-Assign}) gives us
    $\oxsubtypectx{\oxglobalctx}{\oxemptyctx}{\oxvarctx}{\oxvarctx_1}$. Then,
    applying \Lemma{lemma:stack-validity-related-envs} to these two facts gives
    us $\oxstorevalidity{\oxglobalctx}{\oxvarctx_1}{\oxstore}$.

    Then, applying \Lemma{lemma:stack-validity-region-rewriting} to
    $\oxstorevalidity{\oxglobalctx}{\oxvarctx_1}{\oxstore}$ and
    $\oxtunify[\oxcombine]{\oxemptyctx}{\oxkontctx}{\oxvarctx_1}
    {\oxsitype_n}{\oxsitype_o}{\oxvarctx^\prime}$ (from the premise of
    \oxname{T-AssignDeref}) gives us
    $\oxstorevalidity{\oxglobalctx}{\oxvarctx^\prime}{\oxstore}$.
  }

  %% kontctx values are still valid
  \oxpresline{
    \oxwfkontctx{\oxglobalctx}{
      \oxvarctx^\prime
    }{\overline{\oxvalue}}{\oxkontctx}
  }{
    Inverting \oxname{WF-Temporaries} gives us $\forall i \in 1 \oxdots n. \;
    \oxtypjudge{\oxglobalctx}{\oxemptyctx}{\oxsitype_1 \oxdotsc \oxsitype_{i-1}}{\oxvarctx}{
      \oxvalue_i
    }{\oxsitype_i}{\oxvarctx}$.

    Applying \Lemma{lemma:values-dont-change} to the typing derivation (from
    \oxname{T-AssignDeref}) for
    $\oxexpr$ (which we know is a value from \oxname{E-Assign}) gives us
    $\oxsubtypectx{\oxglobalctx}{\oxemptyctx}{\oxvarctx}{\oxvarctx_1}$.
    Then, we can apply \Lemma{lemma:value-typing-related-envs} to each of the
    typing judgments to get $\forall i \in 1 \oxdots n. \;
    \oxtypjudge{\oxglobalctx}{\oxemptyctx}{\oxsitype_1 \oxdotsc \oxsitype_{i-1}}{\oxvarctx_1}{
      \oxvalue_i
    }{\oxsitype_i}{\oxvarctx_1}$.

    Then, applying \Lemma{lemma:subtyping-value-typing} to each of these
    derivations along with
    $\oxtunify[\oxcombine]{\oxemptyctx}{\oxkontctx}{\oxvarctx_1}
    {\oxsitype_n}{\oxsitype_o}{\oxvarctx^\prime}$ (from the premise of
    \oxname{T-AssignDeref}) gives us $\forall i \in 1 \oxdots n. \;
    \oxtypjudge{\oxglobalctx}{\oxemptyctx}{\oxsitype_1 \oxdotsc \oxsitype_{i-1}}{
      \oxvarctx^\prime
    }{\oxvalue_i}{\oxsitype_i}{\oxvarctx^\prime}$. We can then apply
    \oxname{WF-Temporaries} to conclude
    $\oxwfkontctx{\oxglobalctx}{
      \oxvarctx^\prime
    }{\overline{\oxvalue}}{\oxkontctx}$.
  }

  %% new expression is well-typed
  \oxpresline{
    \oxtypjudge{\oxglobalctx}{\oxemptyctx}{\oxkontctx}{
      \oxvarctx^\prime
    }{\oxunit}{\oxtunit}{
      \oxvarctx^\prime
    }
  }{
    Immediate by \oxname{T-Unit}.
  }

  %% new type is a subtype of the old type
  \oxpresline{
    \oxtunify[\oxcombine]{\oxemptyctx}{\oxkontctx}{
      \oxvarctx^\prime
    }{\oxtunit}{\oxtunit}{
      \oxvarctx^\prime
    }
  }{
    Immediate by \oxname{RR-Refl}.
  }

  %% new loan context is a sub-context of the old loan context
  \oxpresline{
    \exists \oxvarctx_o. \oxvarctx^\prime \oxintersect
    \oxvarctx_o = \oxvarctx^\prime
  }{
    $\oxvarctx_o = \oxvarctx^\prime$
  }
\end{preservation}

\oxpreservationcase{T-Let}{\TLet}{
  \oxlet{\oxid}{\oxsitype}{\oxexpr_1}{\oxexpr_2}
}{\ELet}{
  \oxgcloans{\oxkontctx}{
    \oxextendctx{\oxvarctx_1^\prime}{
      \oxvarctxentry{\oxid}{\oxsitype_a}
    }
  }
}

\begin{preservation}
  %% new store is valid
  \oxpresline{
    \oxstorevalidity{\oxglobalctx}{
      \oxgcloans{\oxkontctx}{
        \oxextendctx{\oxvarctx_1^\prime}{
          \oxvarctxentry{\oxid}{\oxsitype_a}
        }
      }
    }{
      \oxextendctx{\oxstore}{\oxstoreentry{\oxid}{\oxvalue}}
    }
  }{
    Applying \Lemma{lemma:values-dont-change} to the typing derivation (from
    \oxname{T-Let}) for $\oxexpr_1$ (which we know is a value from
    \oxname{E-Let}) gives us
    $\oxsubtypectx{\oxglobalctx}{\oxemptyctx}{\oxvarctx}{\oxvarctx_1}$.
    Then, we can apply \Lemma{lemma:stack-validity-related-envs} to get
    $\oxstorevalidity{\oxglobalctx}{\oxvarctx_1}{\oxstore}$.
    Then,
    applying \Lemma{lemma:stack-validity-region-rewriting} to
    $\oxtunify[\oxcombine]{\oxemptyctx}{\oxkontctx}{\oxvarctx_1}{\oxsitype_1}{\oxsitype_a}
    {\oxvarctx_1^\prime}$ (from premise of \oxname{T-Let}) gives us
    $\oxstorevalidity{\oxglobalctx}{\oxvarctx_1^\prime}{\oxstore}$. We can also
    apply \Lemma{lemma:values-typed-supertype} to $\oxtypjudge{\oxglobalctx}
    {\oxemptyctx}{\oxkontctx}{\oxvarctx}{\oxvalue}{\oxsitype_1}{\oxvarctx_1}$ (from premise
    of \oxname{T-Let}) and $\oxtunify[\oxcombine]{\oxemptyctx}{\oxkontctx}{\oxvarctx_1}
    {\oxsitype_1}{\oxsitype_a}{\oxvarctx_1^\prime}$ (from premise of
    \oxname{T-Let}) gives us $\oxtypjudge{\oxglobalctx}{\oxemptyctx}{\oxkontctx}
    {\oxvarctx^\prime}{\oxexpr_1}{\oxsitype_a}{\oxvarctx^\prime}$

    Then, apply \Lemma{lemma:store-validity-extension} to
    $\oxstorevalidity{\oxglobalctx}{\oxvarctx_1^\prime}{\oxstore}$ and
    $\oxtypjudge{\oxglobalctx}{\oxemptyctx}{\oxkontctx}{\oxvarctx^\prime}
    {\oxexpr_1}{\oxsitype_a}{\oxvarctx^\prime}$ gives us
    $\oxstorevalidity{\oxglobalctx}{
      \oxextendctx{\oxvarctx_1^\prime}{
        \oxvarctxentry{\oxid}{\oxsitype_a}
      }
    }{
      \oxextendctx{\oxstore}{\oxstoreentry{\oxid}{\oxvalue}}
    }$. By definition of
    \oxname{R-Env}, we know that $\oxsubtypectx{\oxglobalctx}{\oxemptyctx}{
      \oxextendctx{\oxvarctx_1^\prime}{
        \oxvarctxentry{\oxid}{\oxsitype_a}
      }
    }{
      \oxgcloans{\oxkontctx}{
        \oxextendctx{\oxvarctx_1^\prime}{
          \oxvarctxentry{\oxid}{\oxsitype_a}
        }
      }
    }$. Then, we can apply
    \Lemma{lemma:stack-validity-related-envs} to conclude
    $\oxstorevalidity{\oxglobalctx}{
      \oxgcloans{\oxkontctx}{
        \oxextendctx{\oxvarctx_1^\prime}{
          \oxvarctxentry{\oxid}{\oxsitype_a}
        }
      }
    }{
      \oxextendctx{\oxstore}{\oxstoreentry{\oxid}{\oxvalue}}
    }$.
  }

  %% kontctx values are still valid
  \oxpresline{
    \oxwfkontctx{\oxglobalctx}{
      \oxgcloans{\oxkontctx}{
        \oxextendctx{\oxvarctx_1^\prime}{
          \oxvarctxentry{\oxid}{\oxsitype_a}
        }
      }
    }{\overline{\oxvalue}}{\oxkontctx}
  }{
    Inverting \oxname{WF-Temporaries} gives us $\forall i \in 1 \oxdots n. \;
    \oxtypjudge{\oxglobalctx}{\oxemptyctx}{\oxsitype_1 \oxdotsc \oxsitype_{i-1}}{\oxvarctx}{
      \oxvalue_i
    }{\oxsitype_i}{\oxvarctx}$.

    Applying \Lemma{lemma:values-dont-change} to the typing derivation (from
    \oxname{T-Let}) for
    $\oxexpr_1$ (which we know is a value from \oxname{E-Let}) gives us
    $\oxsubtypectx{\oxglobalctx}{\oxemptyctx}{\oxvarctx}{\oxvarctx_1}$.
    Then, we can apply \Lemma{lemma:value-typing-related-envs} to each of the
    typing judgments to get $\forall i \in 1 \oxdots n. \;
    \oxtypjudge{\oxglobalctx}{\oxemptyctx}{\oxsitype_1 \oxdotsc \oxsitype_{i-1}}{\oxvarctx_1}{
      \oxvalue_i
    }{\oxsitype_i}{\oxvarctx_1}$.

    Applying \Lemma{lemma:values-well-typed-extension} to each of the value
    typing judgments and $\forall \oxcprov \in \oxfprovs{\oxsitype_a}. \;
    \oxnotreborrowed{\oxvarctx_1^\prime}{\oxcprov}$ (from the premise of
    \oxname{T-Let}) gives us
    $\forall i \in 1 \oxdots n. \;
    \oxtypjudge{\oxglobalctx}{\oxemptyctx}{\oxsitype_1 \oxdotsc \oxsitype_{i-1}}{
      \oxextendctx{\oxvarctx_1}{
        \oxvarctxentry{\oxid}{\oxsitype_a}
      }
    }{
      \oxvalue_i
    }{\oxsitype_i}{
      \oxextendctx{\oxvarctx_1}{
        \oxvarctxentry{\oxid}{\oxsitype_a}
      }
    }$.

    Then, we note that definitionally
    $\oxsubtypectx{\oxglobalctx}{\oxemptyctx}{
      \oxextendctx{\oxvarctx_1}{
        \oxvarctxentry{\oxid}{\oxsitype_a}
      }
    }{
      \oxgcloans{\oxkontctx}{
        \oxextendctx{\oxvarctx_1}{
          \oxvarctxentry{\oxid}{\oxsitype_a}
        }
      }
    }$ since we can only clear loans that are not present in any of the types.
    We can use this with each typing derivation in \Lemma{lemma:values-dont-change},
    and then finally apply \oxname{WF-Temporaries} to get
    $\oxwfkontctx{\oxglobalctx}{
      \oxgcloans{\oxkontctx}{
        \oxextendctx{\oxvarctx_1}{
          \oxvarctxentry{\oxid}{\oxsitype_a}
        }
      }
    }{\overline{\oxvalue}}{\oxkontctx}$.
  }

  %% new expression is well-typed
  \oxpresline{
    \oxtypjudge{\oxglobalctx}{\oxemptyctx}{\oxkontctx}{
      \oxgcloans{\oxkontctx}{
        \oxextendctx{\oxvarctx_1^\prime}{
          \oxvarctxentry{\oxid}{\oxsitype_a}
        }
      }
    }{\oxshift{\oxexpr}}{\oxsitype_2}{\oxvarctx_2}
  }{
    Immediate by applying \oxname{T-Shift} to the derivation
    \oxtypjudge{\oxglobalctx}{\oxemptyctx}{\oxkontctx}{
      \oxgcloans{\oxkontctx}{
        \oxextendctx{\oxvarctx_1^\prime}{
          \oxvarctxentry{\oxid}{\oxsitype_a}
        }
      }
    }{\oxexpr}{\oxsitype_2}{
      \oxextendctx{\oxvarctx_2}{
        \oxvarctxentry{\oxid}{\oxsdtype}
      }
    } (from premise of \oxname{T-Let}).
  }

  %% new type is a subtype of the old type
  \oxpresline{
    \oxtunify[\oxcombine]{\oxemptyctx}{\oxkontctx}{\oxvarctx_2}
    {\oxsitype_2}{\oxsitype_2}{\oxvarctx_2}
  }{
    Immediate by \oxname{RR-Refl}.
  }

  %% new loan context is a sub-context of the old loan context
  \oxpresline{
    \exists \oxvarctx_o. \oxvarctx_2 \oxintersect \oxvarctx_o = \oxvarctx_2
  }{
    $\oxvarctx_o = \oxvarctx_2$
  }
\end{preservation}

\oxpreservationcase{T-LetRegion}{\TLetProv}{
  \oxletrgn{\oxcprov}{\oxexpr}
}{\ELetProv}{
  \oxvarctx^\prime
}

\begin{preservation}
  %% new store is valid
  \oxpresline{
    \oxstorevalidity{\oxglobalctx}{\oxvarctx^\prime}{\oxstore}
  }{
    Applying \Lemma{lemma:values-dont-change} to the typing derivation (from
    \oxname{T-LetRegion}) for $\oxexpr$ (which we know is a value from
    \oxname{E-LetRegion}) gives us
    $\oxsubtypectx{\oxglobalctx}{\oxemptyctx}{\oxvarctx}{\oxvarctx^\prime}$. Then,
    applying \Lemma{lemma:stack-validity-related-envs} with this and
    $\oxstorevalidity{\oxglobalctx}{\oxvarctx}{\oxstore}$ (from
    premise) gives us
    $\oxstorevalidity{\oxglobalctx}{\oxvarctx^\prime}{\oxstore}$.
  }

  %% kontctx values are still valid
  \oxpresline{
    \oxwfkontctx{\oxglobalctx}{\oxvarctx^\prime}{\overline{\oxvalue}}{\oxkontctx}
  }{
    Inverting \oxname{WF-Temporaries} gives us $\forall i \in 1 \oxdots n. \;
    \oxtypjudge{\oxglobalctx}{\oxemptyctx}{\oxsitype_1 \oxdotsc \oxsitype_{i-1}}{\oxvarctx}{
      \oxvalue_i
    }{\oxsitype_i}{\oxvarctx}$.

    Applying \Lemma{lemma:values-dont-change} to the typing derivation (from
    \oxname{T-LetRegion}) for
    $\oxexpr$ (which we know is a value from \oxname{E-LetRegion}) gives us
    $\oxsubtypectx{\oxglobalctx}{\oxemptyctx}{\oxvarctx}{\oxvarctx^\prime}$.

    Thus, we can apply \Lemma{lemma:value-typing-related-envs} to each of the
    typing judgments to get $\forall i \in 1 \oxdots n. \;
    \oxtypjudge{\oxglobalctx}{\oxemptyctx}{\oxsitype_1 \oxdotsc \oxsitype_{i-1}}{\oxvarctx^\prime}{
      \oxvalue_i
    }{\oxsitype_i}{\oxvarctx^\prime}$.

    Finally, applying \oxname{WF-Temporaries} gives us
    $\oxwfkontctx{\oxglobalctx}{\oxvarctx^\prime}{\overline{\oxvalue}}{\oxkontctx}$.
  }

  %% new expression is well-typed
  \oxpresline{
    \oxtypjudge{\oxglobalctx}{\oxemptyctx}{\oxkontctx}{\oxvarctx^\prime}{\oxvalue}
    {\oxsitype}{\oxvarctx^\prime}
  }{
    We know from \oxname{E-LetRegion} that $\oxexpr$ is a value $\oxvalue$. Thus,
    we can apply \Lemma{lemma:value-typing-output} to
    $\oxtypjudge{\oxglobalctx}{\oxemptyctx}{\oxkontctx}{\oxvarctx}{\oxvalue}{\oxsitype}{
      \oxextendctx{\oxvarctx^\prime}{
        \oxloanctxentry{\oxcprov}{\oxset{\overline{\oxloan}}}
      }
    }$ to get $\oxtypjudge{\oxglobalctx}{\oxemptyctx}{\oxkontctx}{
      \oxextendctx{\oxvarctx^\prime}{
        \oxloanctxentry{\oxcprov}{\oxset{\overline{\oxloan}}}
      }
    }{\oxvalue}{\oxsitype}{
      \oxextendctx{\oxvarctx^\prime}{
        \oxloanctxentry{\oxcprov}{\oxset{\overline{\oxloan}}}
      }
    }$.

    We now wish to show that
    $\oxtypjudge{\oxglobalctx}{\oxemptyctx}{\oxkontctx}{\oxvarctx^\prime}{\oxvalue}
    {\oxsitype}{\oxvarctx^\prime}$. By inspecting the grammar of values and
    their typing rules, we know that the only values who depend on the
    context are pointers and closure values. But by inversion on
    $\oxtypjudge{\oxglobalctx}{\oxemptyctx}{\oxkontctx}{\oxvarctx}{
      \oxletrgn{\oxcprov}{\oxvalue}
    }{\oxsitype}{\oxvarctx^\prime}$, we know that
    $\oxtypevalidity{\oxglobalctx}{\oxemptyctx}{\oxvarctx^\prime}{\oxsitype}$.
    Since the type is valid without the frame, we know that the values
    cannot depend on that frame. Thus, we can conclude
    $\oxtypjudge{\oxglobalctx}{\oxemptyctx}{\oxkontctx}{\oxvarctx^\prime}{\oxvalue}
    {\oxsitype}{\oxvarctx^\prime}$.
  }

  %% new type is a subtype of the old type
  \oxpresline{
    \oxtunify[\oxcombine]{\oxemptyctx}{\oxkontctx}{\oxvarctx^\prime}
    {\oxsitype}{\oxsitype}{\oxvarctx^\prime}
  }{
    Immediate by \oxname{RR-Refl}.
  }

  %% new loan context is a sub-context of the old loan context
  \oxpresline{
    \exists \oxvarctx_o. \oxvarctx^\prime \oxintersect \oxvarctx_o = \oxvarctx^\prime
  }{
    $\oxvarctx_o = \oxvarctx^\prime$
  }
\end{preservation}

\oxpreservationcase{T-While}{\TWhile}{
  \oxwhile{\oxexpr_1}{\oxexpr_2}
}{\EWhile}{
  \oxvarctx
}

\begin{preservation}
  %% new store is valid
  \oxpresline{
    \oxstorevalidity{\oxglobalctx}{\oxvarctx}{\oxstore}
  }{
    Immediate from our premise.
  }

  %% kontctx values are still valid
  \oxpresline{
    \oxwfkontctx{\oxglobalctx}{\oxvarctx}{\overline{\oxvalue}}{\oxkontctx}
  }{
    Immediate from our premise.
  }

  %% new expression is well-typed
  \oxpresline{
    \oxtypjudge{\oxglobalctx}{\oxemptyctx}{\oxkontctx}{\oxvarctx}{
      \oxexpr^\prime
      % \oxbranch{\oxexpr_1}{
      %   \oxseq{\oxexpr_2}{
      %     \oxwhile{\oxexpr_1}{\oxexpr_2}
      %   }
      % }{\oxunit}
    }{\oxtunit}{\oxvarctx_2}
  }{
    We would like to build a derivation to show that the expression
    $\oxbranch{\oxexpr_1}{
      \oxseq{\oxexpr_2}{
        \oxwhile{\oxexpr_1}{\oxexpr_2}
      }
    }{\oxunit}$ is well-typed. We thus start by applying \oxname{T-Branch}.

    This requires us to show three things. First, $\oxtypjudge{\oxglobalctx}
    {\oxemptyctx}{\oxkontctx}{\oxvarctx}{\oxexpr_1}{\oxtbool}{\oxvarctx_1}$ which we have
    from the premise of \oxname{T-While}. Second, $\oxtypjudge{\oxglobalctx}
    {\oxemptyctx}{\oxkontctx}{\oxvarctx_1}{
      \oxseq{\oxexpr_2}{
        \oxwhile{\oxexpr_1}{\oxexpr_2}
      }
    }{\oxtunit}{\oxvarctx_2}$. We build this by applying \oxname{T-Seq} to
    $\oxtypjudge{\oxglobalctx}{\oxemptyctx}{\oxkontctx}{\oxvarctx_1}{\oxexpr_2}
    {\oxtunit}{\oxvarctx_2}$ and $\oxtypjudge{\oxglobalctx}{\oxemptyctx}{\oxkontctx}
    {\oxvarctx_2}{\oxwhile{\oxexpr_1}{\oxexpr_2}}{\oxvarctx_2}$. The former is
    directly in the premise of \oxname{T-While} and the latter can be built by
    applying \oxname{T-While} to
    $\oxtypjudge{\oxglobalctx}{\oxemptyctx}{\oxkontctx}{\oxvarctx_2}{
      \oxexpr_1
    }{\oxtbool}{\oxvarctx_2}$ and
    $\oxtypjudge{\oxglobalctx}{\oxemptyctx}{\oxkontctx}{\oxvarctx_2}{
      \oxexpr_2
    }{\oxtunit}{\oxvarctx_2}$, both from the premise of our original
    \oxname{T-While}. Finally, we need to show
    $\oxtypjudge{\oxglobalctx}{\oxemptyctx}{\oxkontctx}{\oxvarctx_2}{
      \oxunit
    }{\oxtunit}{\oxvarctx_2}$, which is immediate from \oxname{T-Unit}.
  }

  %% new type is a subtype of the old type
  \oxpresline{
    \oxtunify[\oxcombine]{\oxemptyctx}{\oxkontctx}{\oxvarctx_2}
    {\oxtunit}{\oxtunit}{\oxvarctx_2}
  }{
    Immediate by \oxname{RR-Refl}.
  }

  %% new loan context is a sub-context of the old loan context
  \oxpresline{
    \exists \oxvarctx_o. \oxvarctx_2 \oxintersect \oxvarctx_o = \oxvarctx_2
  }{
    $\oxvarctx_o = \oxvarctx_2$
  }
\end{preservation}

\oxpreservationcasemulti{T-ForArray}{\TForArray}{
  \oxfor{\oxid}{\oxexpr_1}{\oxexpr_2}
}{
  \EForArray \and \EForEmptyArray
}

\oxpreservationsubcaseheading{E-ForArray}{
  \oxextendctx{\oxvarctx_1}{\oxvarctxentry{\oxid}{\oxsitype}}
}

\begin{preservation}
  %% new store is valid
  \oxpresline{
    \oxstorevalidity{\oxglobalctx}{
      \oxextendctx{\oxvarctx_1}{\oxvarctxentry{\oxid}{\oxsitype}}
    }{
      \oxextendctx{\oxstore}{
        \oxstoreentry{\oxid}{\oxvalue_0}
      }
    }
  }{
    Applying \Lemma{lemma:values-dont-change} to the typing derivation (from
    \oxname{T-ForArray}) for $\oxexpr_1$ (which we know is a value from
    \oxname{E-ForArray}) gives us
    $\oxsubtypectx{\oxglobalctx}{\oxemptyctx}{\oxvarctx}{\oxvarctx_1}$. Then,
    we can apply \Lemma{lemma:stack-validity-related-envs} to get
    $\oxstorevalidity{\oxglobalctx}{\oxvarctx_1}{\oxstore}$. Applying
    \Lemma{lemma:value-typing-output} to the derivation
    $\oxtypjudge{\oxglobalctx}{\oxemptyctx}{\oxkontctx}{\oxvarctx}
    {\oxarr{\oxvalue_0 \oxdotsc \oxvalue_n}}{\oxtarr{\oxsitype}{\oxnum}}{\oxvarctx_1}$
    gives us $\oxtypjudge{\oxglobalctx}{\oxemptyctx}{\oxkontctx}{\oxvarctx_1}
    {\oxarr{\oxvalue_0 \oxdotsc \oxvalue_n}}{\oxtarr{\oxsitype}{\oxnum}}{\oxvarctx_1}$.
    Then, using inversion (of \oxname{T-Array}), we get
    $\oxtypjudge{\oxglobalctx}{\oxemptyctx}{\oxkontctx}{\oxvarctx_1}
    {\oxvalue_0}{\oxsitype}{\oxvarctx_1}$. Finally, applying
    \Lemma{lemma:store-validity-extension} to
    $\oxstorevalidity{\oxglobalctx}{\oxvarctx_1}{\oxstore}$ and
    $\oxtypjudge{\oxglobalctx}{\oxemptyctx}{\oxkontctx}{\oxvarctx_1}
    {\oxvalue_0}{\oxsitype}{\oxvarctx_1}$ gives us
    $\oxstorevalidity{\oxglobalctx}{
      \oxextendctx{\oxvarctx_1}{
        \oxvarctxentry{\oxid}{\oxsitype}
      }
    }{
      \oxextendctx{\oxstore}{\oxstoreentry{\oxid}{\oxvalue_0}}
    }$.
  }

  %% kontctx values are still valid
  \oxpresline{
    \oxwfkontctx{\oxglobalctx}{
      \oxextendctx{\oxvarctx_1}{
        \oxvarctxentry{\oxid}{\oxsitype}
      }
    }{\overline{\oxvalue}}{\oxkontctx}
  }{
    Inverting \oxname{WF-Temporaries} gives us $\forall i \in 1 \oxdots n. \;
    \oxtypjudge{\oxglobalctx}{\oxemptyctx}{\oxsitype_1 \oxdotsc \oxsitype_{i-1}}{\oxvarctx}{
      \oxvalue_i
    }{\oxsitype_i}{\oxvarctx}$.

    Applying \Lemma{lemma:values-dont-change} to the typing derivation (from
    \oxname{T-ForArray}) for
    $\oxexpr_1$ (which we know is a value from \oxname{E-ForArray}) gives us
    $\oxsubtypectx{\oxglobalctx}{\oxemptyctx}{\oxvarctx}{\oxvarctx_1}$.

    Thus, we can apply \Lemma{lemma:value-typing-related-envs} to each of the
    typing judgments to get $\forall i \in 1 \oxdots n. \;
    \oxtypjudge{\oxglobalctx}{\oxemptyctx}{\oxemptyctx}{\oxvarctx_1}{
      \oxvalue_i
    }{\oxsitype_i}{\oxvarctx_1}$.

    Applying \Lemma{lemma:values-well-typed-extension} to each of the value
    typing judgments and $\forall \oxcprov \in \oxfprovs{\oxsitype}. \;
    \oxnotreborrowed{\oxvarctx_1}{\oxcprov}$ (from the premise of
    \oxname{T-ForArray}) gives us
    $\forall i \in 1 \oxdots n. \;
    \oxtypjudge{\oxglobalctx}{\oxemptyctx}{\oxemptyctx}{
      \oxextendctx{\oxvarctx_1}{
        \oxvarctxentry{\oxid}{\oxsitype}
      }
    }{
      \oxvalue_i
    }{\oxsitype_i}{
      \oxextendctx{\oxvarctx_1}{
        \oxvarctxentry{\oxid}{\oxsitype}
      }
    }$.
  }

  %% new expression is well-typed
  \oxpresline{
    \oxtypjudge{\oxglobalctx}{\oxemptyctx}{\oxkontctx}{
      \oxextendctx{\oxvarctx_1}{\oxvarctxentry{\oxid}{\oxsitype}}
    }{\oxexpr^\prime}{\oxtunit}{\oxvarctx_1}
  }{
    We need to build a derivation for the expression $\oxseq{\oxshift{\oxexpr}}{
      \oxfor{\oxid}{\oxarr{\oxvalue_1 \oxdotsc \oxvalue_n}}{\oxexpr}
    }$. The bottom of this derivation will be \oxname{T-Seq} which requires us
    to show that $\oxtypjudge{\oxglobalctx}{\oxemptyctx}{\oxkontctx}{
      \oxextendctx{\oxvarctx_1}{
        \oxvarctxentry{\oxid}{\oxsitype}
      }
    }{\oxshift{\oxexpr}}{\oxtunit}{\oxvarctx_1}$ and that
    $\oxtypjudge{\oxglobalctx}{\oxemptyctx}{\oxkontctx}{\oxvarctx_1}{
      \oxfor{\oxid}{\oxarr{\oxvalue_1 \oxdotsc \oxvalue_n}}{\oxexpr}
    }{\oxtunit}{\oxvarctx_1}$.

    To show $\oxtypjudge{\oxglobalctx}{\oxemptyctx}{\oxkontctx}{
      \oxextendctx{\oxvarctx_1}{
        \oxvarctxentry{\oxid}{\oxsitype}
      }
    }{\oxshift{\oxexpr}}{\oxtunit}{\oxvarctx_1}$, we apply \oxname{T-Shift} to
    $\oxtypjudge{\oxglobalctx}{\oxemptyctx}{\oxkontctx}{
      \oxextendctx{\oxvarctx_1}{
        \oxvarctxentry{\oxid}{\oxsitype}
      }
    }{\oxexpr}{\oxtunit}{
      \oxextendctx{\oxvarctx_1}{
        \oxvarctxentry{\oxid}{\oxsdtype}
      }
    }$ (from the premise of \oxname{T-ForArray}).

    To show $\oxtypjudge{\oxglobalctx}{\oxemptyctx}{\oxkontctx}{\oxvarctx_1}{
      \oxfor{\oxid}{\oxarr{\oxvalue_1 \oxdotsc \oxvalue_n}}{\oxexpr}
    }{\oxtunit}{\oxvarctx_1}$, we apply \Lemma{lemma:value-typing-output} to the
    derivation $\oxtypjudge{\oxglobalctx}{\oxemptyctx}{\oxkontctx}{\oxvarctx}
    {\oxarr{\oxvalue_0 \oxdotsc
        \oxvalue_n}}{\oxtarr{\oxsitype}{\oxnum}}{\oxvarctx_1}$ to get
    $\oxtypjudge{\oxglobalctx}{\oxemptyctx}{\oxkontctx}{\oxvarctx_1} {\oxarr{\oxvalue_0
        \oxdotsc \oxvalue_n}}{\oxtarr{\oxsitype}{\oxnum}}{\oxvarctx_1}$. Then,
    we rewrite the derivation (inverting and then reapply \oxname{T-Array}) to
    exclude $\oxvalue_0$ giving us
    $\oxtypjudge{\oxglobalctx}{\oxemptyctx}{\oxkontctx}{\oxvarctx_1} {\oxarr{\oxvalue_1
        \oxdotsc \oxvalue_n}}{\oxtarr{\oxsitype}{\oxnum - 1}}{\oxvarctx_1}$.
    Finally, we apply \oxname{T-ForArray} using this combined with
    $\oxtypjudge{\oxglobalctx}{\oxemptyctx}{\oxkontctx}{\oxvarctx_1}
    {\oxexpr}{\oxtunit}{\oxvarctx_1}$ (from the premise of \oxname{T-ForArray}).
  }

  %% new type is a subtype of the old type
  \oxpresline{
    \oxtunify[\oxcombine]{\oxemptyctx}{\oxkontctx}{\oxvarctx_1}
    {\oxtunit}{\oxtunit}{\oxvarctx_1}
  }{
    Immediate by \oxname{RR-Refl}.
  }

  %% new loan context is a sub-context of the old loan context
  \oxpresline{
    \exists \oxvarctx_o. \oxvarctx_1 \oxintersect \oxvarctx_o = \oxvarctx_1
  }{
    $\oxvarctx_o = \oxvarctx_1$
  }
\end{preservation}

\oxpreservationsubcaseheading{E-ForEmptyArray}{
  \oxvarctx
}

\begin{preservation}
  %% new store is valid
  \oxpresline{
    \oxstorevalidity{\oxglobalctx}{\oxvarctx}{\oxstore}
  }{
    Immediate from our premise.
  }

  %% kontctx values are still valid
  \oxpresline{
    \oxwfkontctx{\oxglobalctx}{\oxvarctx}{\overline{\oxvalue}}{\oxkontctx}
  }{
    Immediate from our premise.
  }

  %% new expression is well-typed
  \oxpresline{
    \oxtypjudge{\oxglobalctx}{\oxemptyctx}{\oxkontctx}{\oxvarctx}
    {\oxunit}{\oxtunit}{\oxvarctx}
  }{
    Immediate by \oxname{T-Unit}.
  }

  %% new type is a subtype of the old type
  \oxpresline{
    \oxtunify[\oxcombine]{\oxemptyctx}{\oxkontctx}{\oxvarctx}
    {\oxtunit}{\oxtunit}{\oxvarctx}
  }{
    Immediate by \oxname{RR-Refl}.
  }

  %% new loan context is a sub-context of the old loan context
  \oxpresline{
    \exists \oxvarctx_o. \oxvarctx \oxintersect \oxvarctx_o = \oxvarctx
  }{
    $\oxvarctx_o = \oxvarctx$
  }
\end{preservation}

\oxpreservationcasemulti{T-ForSlice}{\TForSlice}{
  \oxfor{\oxid}{\oxexpr_1}{\oxexpr_2}
}{
  \EForSlice \and \EForEmptySlice
}

\oxpreservationsubcaseheading{E-ForSlice}{
  \oxextendctx{\oxvarctx_1}{
    \oxvarctxentry{\oxid}{
      \oxtref{\oxcprov}{\oxmuta}{\oxsitype}
    }
  }
}

\begin{preservation}
  %% new store is valid
  \oxpresline{
    \oxstorevalidity{\oxglobalctx}{
      \oxextendctx{\oxvarctx_1}{
        \oxvarctxentry{\oxid}{
          \oxtref{\oxcprov}{\oxmuta}{\oxsitype}
        }
      }
    }{
      \oxextendctx{\oxstore}{
        \oxstoreentry{\oxid}{
          \oxindexptr{\oxreferent}{\oxnum_1}
        }
      }
    }
  }{
    Applying \Lemma{lemma:values-dont-change} to the typing derivation (from
    \oxname{T-ForSlice}) for $\oxexpr_1$ (which we know is a value from
    \oxname{E-ForSlice}) gives us
    $\oxsubtypectx{\oxglobalctx}{\oxemptyctx}{\oxvarctx}{\oxvarctx_1}$. Then,
    we can apply \Lemma{lemma:stack-validity-related-envs} to get
    $\oxstorevalidity{\oxglobalctx}{\oxvarctx_1}{\oxstore}$. Applying
    \Lemma{lemma:value-typing-output} to the derivation
    $\oxtypjudge{\oxglobalctx}{\oxemptyctx}{\oxkontctx}{\oxvarctx}
    {\oxsliceptr{\oxreferent}{i}{j}}{
      \oxtref{\oxcprov}{\oxmuta}{\oxtslice{\oxsitype}}
    }{\oxvarctx_1}$
    gives us $\oxtypjudge{\oxglobalctx}{\oxemptyctx}{\oxkontctx}{\oxvarctx_1}
    {\oxsliceptr{\oxreferent}{i}{j}}{
      \oxtref{\oxcprov}{\oxmuta}{\oxtslice{\oxsitype}}
    }{\oxvarctx_1}$.
    Then, using inversion (on \oxname{T-Pointer}), we get
    $\oxreferentvalidity{\oxglobalctx}{\oxvarctx_1}{
      \oxslice{\oxreferent}{i}{j}
    }{\oxxitype}$ (where $\oxxitype = \oxtarr{\oxsitype}{\oxnum}$ or
    $\oxtslice{\oxsitype}$) and
    $\oxloanpkg{\oxmuta}{\oxplace} \in \oxvarctx_1(\oxcprov)$ (where
    $\oxreferent = \oxreferentctx[\oxplace]$). We can invert
    \oxname{WF-RefSliceArray} or \oxname{WF-RefSliceSlice} (based on
    $\oxxitype$) for $\oxreferentvalidity{\oxglobalctx}{\oxvarctx_1}{
      \oxslice{\oxreferent}{i}{j}
    }{\oxxitype}$ and then apply \oxname{WF-RefIndexArray} or
    \oxname{WF-RefIndexSlice} appropriately to get
    $\oxreferentvalidity{\oxglobalctx}{\oxvarctx_1}{
      \oxindex{\oxreferent}{i}
    }{\oxxitype}$. We can then use
    \oxname{T-Pointer} to get
    $\oxtypjudge{\oxglobalctx}{\oxemptyctx}{\oxkontctx}{\oxvarctx_1}
    {\oxindexptr{\oxreferent}{\oxnum_1}}{
      \oxtref{\oxcprov}{\oxmuta}{\oxsitype}
    }{\oxvarctx_1}$. Finally, applying
    \Lemma{lemma:store-validity-extension} to
    $\oxstorevalidity{\oxglobalctx}{\oxvarctx_1}{\oxstore}$ and
    $\oxtypjudge{\oxglobalctx}{\oxemptyctx}{\oxkontctx}{\oxvarctx_1}
    {\oxindexptr{\oxreferent}{\oxnum_1}}{
      \oxtref{\oxcprov}{\oxmuta}{\oxsitype}
    }{\oxvarctx_1}$ gives us
    $\oxstorevalidity{\oxglobalctx}{
      \oxextendctx{\oxvarctx_1}{
        \oxvarctxentry{\oxid}{\oxtref{\oxcprov}{\oxmuta}{\oxsitype}}
      }
    }{
      \oxextendctx{\oxstore}{
        \oxstoreentry{\oxid}{\oxindexptr{\oxreferent}{\oxnum_1}}
      }
    }$.
  }

  %% kontctx values are still valid
  \oxpresline{
    \oxwfkontctx{\oxglobalctx}{
      \oxextendctx{\oxvarctx_1}{
        \oxvarctxentry{\oxid}{
          \oxtref{\oxcprov}{\oxmuta}{\oxsitype}
        }
      }
    }{\overline{\oxvalue}}{\oxkontctx}
  }{
    Inverting \oxname{WF-Temporaries} gives us $\forall i \in 1 \oxdots n. \;
    \oxtypjudge{\oxglobalctx}{\oxemptyctx}{\oxsitype_1 \oxdotsc \oxsitype_{i-1}}{\oxvarctx}{
      \oxvalue_i
    }{\oxsitype_i}{\oxvarctx}$.

    Applying \Lemma{lemma:values-dont-change} to the typing derivation (from
    \oxname{T-ForSlice}) for
    $\oxexpr_1$ (which we know is a value from \oxname{E-ForSlice}) gives us
    $\oxsubtypectx{\oxglobalctx}{\oxemptyctx}{\oxvarctx}{\oxvarctx_1}$.

    Thus, we can apply \Lemma{lemma:value-typing-related-envs} to each of the
    typing judgments to get $\forall i \in 1 \oxdots n. \;
    \oxtypjudge{\oxglobalctx}{\oxemptyctx}{\oxemptyctx}{\oxvarctx_1}{
      \oxvalue_i
    }{\oxsitype_i}{\oxvarctx_1}$.

    Applying \Lemma{lemma:values-well-typed-extension} to each of the value
    typing judgments and $\forall \oxcprov \in
    \oxfprovs{\oxtref{\oxprov}{\oxmuta}{\oxsitype}}. \;
    \oxnotreborrowed{\oxvarctx_1}{\oxcprov}$ (from the premise of
    \oxname{T-ForSlice}) gives us
    $\forall i \in 1 \oxdots n. \;
    \oxtypjudge{\oxglobalctx}{\oxemptyctx}{\oxemptyctx}{
      \oxextendctx{\oxvarctx_1}{
        \oxvarctxentry{\oxid}{
          \oxtref{\oxcprov}{\oxmuta}{\oxsitype}
        }
      }
    }{
      \oxvalue_i
    }{\oxsitype_i}{
      \oxextendctx{\oxvarctx_1}{
        \oxvarctxentry{\oxid}{
          \oxtref{\oxcprov}{\oxmuta}{\oxsitype}
        }
      }
    }$.
  }

  %% new expression is well-typed
  \oxpresline{
    \oxtypjudge{\oxglobalctx}{\oxemptyctx}{\oxkontctx}{
      \oxextendctx{\oxvarctx_1}{
        \oxvarctxentry{\oxid}{
          \oxtref{\oxcprov}{\oxmuta}{\oxsitype}
        }
      }
    }{\oxexpr^\prime}{\oxtunit}{\oxvarctx_1}
  }{
    We need to build a derivation for the expression $\oxseq{\oxshift{\oxexpr}}{
      \oxfor{\oxid}{
        \oxsliceptr{\oxreferent}{\oxnum_1^\prime}{\oxnum_2}
      }{\oxexpr}
    }$. The bottom of this derivation will be \oxname{T-Seq} which requires us
    to show that $\oxtypjudge{\oxglobalctx}{\oxemptyctx}{\oxkontctx}{
      \oxextendctx{\oxvarctx_1}{
        \oxvarctxentry{\oxid}{\oxsitype}
      }
    }{\oxshift{\oxexpr}}{\oxtunit}{\oxvarctx_1}$ and that
    $\oxtypjudge{\oxglobalctx}{\oxemptyctx}{\oxkontctx}{\oxvarctx_1}{
      \oxfor{\oxid}{
        \oxsliceptr{\oxreferent}{\oxnum_1^\prime}{\oxnum_2}
      }{\oxexpr}
    }{\oxtunit}{\oxvarctx_1}$.

    To show $\oxtypjudge{\oxglobalctx}{\oxemptyctx}{\oxkontctx}{
      \oxextendctx{\oxvarctx_1}{
        \oxvarctxentry{\oxid}{\oxsitype}
      }
    }{\oxshift{\oxexpr}}{\oxtunit}{\oxvarctx_1}$, we apply \oxname{T-Shift} to
    $\oxtypjudge{\oxglobalctx}{\oxemptyctx}{\oxkontctx}{
      \oxextendctx{\oxvarctx_1}{
        \oxvarctxentry{\oxid}{\oxsitype}
      }
    }{\oxexpr}{\oxtunit}{
      \oxextendctx{\oxvarctx_1}{
        \oxvarctxentry{\oxid}{\oxsdtype}
      }
    }$ (from the premise of \oxname{T-ForSlice}).

    To show $\oxtypjudge{\oxglobalctx}{\oxemptyctx}{\oxkontctx}{\oxvarctx_1}{
      \oxfor{\oxid}{\oxsliceptr{\oxreferent}{\oxnum_1^\prime}{\oxnum_2}}{\oxexpr}
    }{\oxtunit}{\oxvarctx_1}$, we apply
    \Lemma{lemma:value-typing-output} to the derivation
    $\oxtypjudge{\oxglobalctx}{\oxemptyctx}{\oxkontctx}{\oxvarctx}
    {\oxsliceptr{\oxreferent}{\oxnum_1}{\oxnum_2}}{
      \oxtref{\oxcprov}{\oxmuta}{\oxtslice{\oxsitype}}
    }{\oxvarctx_1}$
    to get $\oxtypjudge{\oxglobalctx}{\oxemptyctx}{\oxkontctx}{\oxvarctx_1}
    {\oxsliceptr{\oxreferent}{\oxnum_1}{\oxnum_2}}{
      \oxtref{\oxcprov}{\oxmuta}{\oxtslice{\oxsitype}}
    }{\oxvarctx_1}$. Then, we rewrite the derivation (inverting and then reapply
    \oxname{T-Pointer}) to increment $\oxnum_1$ to $\oxnum_1^\prime$ giving us
    $\oxtypjudge{\oxglobalctx}{\oxemptyctx}{\oxkontctx}{\oxvarctx_1}
    {\oxsliceptr{\oxreferent}{\oxnum_1^\prime}{\oxnum_2}}{
      \oxtref{\oxcprov}{\oxmuta}{\oxtslice{\oxsitype}}
    }{\oxvarctx_1}$.
    Finally, we apply \oxname{T-ForSlice} using this combined with
    $\oxtypjudge{\oxglobalctx}{\oxemptyctx}{\oxkontctx}{\oxvarctx_1}
    {\oxexpr}{\oxtunit}{\oxvarctx_1}$ (from the premise of \oxname{T-ForSlice}).
  }

  %% new type is a subtype of the old type
  \oxpresline{
    \oxtunify[\oxcombine]{\oxemptyctx}{\oxkontctx}{\oxvarctx_1}
    {\oxtunit}{\oxtunit}{\oxvarctx_1}
  }{
    Immediate by \oxname{RR-Refl}.
  }

  %% new loan context is a sub-context of the old loan context
  \oxpresline{
    \exists \oxvarctx_o. \oxvarctx_1 \oxintersect \oxvarctx_o = \oxvarctx_1
  }{
    $\oxvarctx_o = \oxvarctx_1$
  }
\end{preservation}

\oxpreservationsubcaseheading{E-ForEmptySlice}{
  \oxvarctx
}

\begin{preservation}
  %% new store is valid
  \oxpresline{
    \oxstorevalidity{\oxglobalctx}{\oxvarctx}{\oxstore}
  }{
    Immediate from our premise.
  }

  %% kontctx values are still valid
  \oxpresline{
    \oxwfkontctx{\oxglobalctx}{\oxvarctx}{\overline{\oxvalue}}{\oxkontctx}
  }{
    Immediate from our premise.
  }

  %% new expression is well-typed
  \oxpresline{
    \oxtypjudge{\oxglobalctx}{\oxemptyctx}{\oxkontctx}{\oxvarctx}
    {\oxunit}{\oxtunit}{\oxvarctx}
  }{
    Immediate by \oxname{T-Unit}.
  }

  %% new type is a subtype of the old type
  \oxpresline{
    \oxtunify[\oxcombine]{\oxemptyctx}{\oxkontctx}{\oxvarctx}
    {\oxtunit}{\oxtunit}{\oxvarctx}
  }{
    Immediate by \oxname{RR-Refl}.
  }

  %% new loan context is a sub-context of the old loan context
  \oxpresline{
    \exists \oxvarctx_o. \oxvarctx \oxintersect \oxvarctx_o = \oxvarctx
  }{
    $\oxvarctx_o = \oxvarctx$
  }
\end{preservation}

\oxpreservationcase[
% Let $\oxsitype_{\textsc{s}1} =
% \oxsubst{\oxsubst{\oxsitype_1}{\overline{\oxprov^\prime}}{\overline{\oxprov}}}
% {\overline{\oxsitype}}{\overline{\oxtvar}} \oxdotsc \oxsitype_{\textsc{s}n} =
% \oxsubst{\oxsubst{\oxsitype_n}{\overline{\oxprov^\prime}}{\overline{\oxprov}}}
% {\overline{\oxsitype}}{\overline{\oxtvar}}$.
]{T-AppFunction}{\TApp}{
  \oxapp{\oxexpr_f}{
    \overline{\oxenv} \oxcomma
    \overline{\oxprov^\prime} \oxcomma
    \overline{\oxsitype}
  }{
    \oxexpr_1 \oxdotsc \oxexpr_n
  }
}{\EAppFun}{
  \oxnewframe{\oxvarctx_b}{
    \oxvarctxentry{\oxid_1}{\delta(\oxsitype_1)} \oxdotsc
    \oxvarctxentry{\oxid_n}{\delta(\oxsitype_n)}
  }
}

\begin{preservation}
  %% new store is valid
  \oxpresline{
    \oxstorevalidity{\oxglobalctx}{
      \oxvarctx_i
    }{
      \oxstore^\prime
    }
  }{
    Applying \Lemma{lemma:values-dont-change} to the derivation for $\oxvalue_f$
    gives us $\oxsubtypectx{\oxglobalctx}{\oxemptyctx}{\oxvarctx}{\oxvarctx_0}$.
    Then, applying \Lemma{lemma:values-dont-change} to the derivations for
    every $\oxvalue_i$ gives us $\forall i \in \oxset{1 \oxdots n}. \;
    \oxsubtypectx{\oxglobalctx}{\oxemptyctx}{\oxvarctx_{i-1}}{\oxvarctx_i}$.

    Inversion on $\oxrunifymany{\oxemptyctx}{\oxkontctx}{\oxvarctx_n}{
      \oxsubstmany{\oxabsprov_2}{\oxprov}{\oxabsprov}
    }{
      \oxsubstmany{\oxabsprov_1}{\oxprov}{\oxabsprov}
    }{\oxvarctx_b}$ gives us a sequence of outlives relations with
    intermediate contexts. Applying  \Lemma{lemma:outlives-related} to each of
    them and then combining the result by transitivity gives us
    $\oxsubtypectx{\oxglobalctx}{\oxemptyctx}{\oxvarctx_n}{\oxvarctx_b}$.
    Combining both by transitivity, we have
    $\oxsubtypectx{\oxglobalctx}{\oxemptyctx}{\oxvarctx}{\oxvarctx_b}$.

    We can then apply \Lemma{lemma:stack-validity-related-envs} with
    $\oxstorevalidity{\oxglobalctx}{\oxvarctx}{\oxstore}$ to get
    $\oxstorevalidity{\oxglobalctx}{\oxvarctx_b}{\oxstore}$.

    We can apply \oxname{WF-StackFrame} to get
    $\oxstorevalidity{\oxglobalctx}{
      \oxnewframe{\oxvarctx_b}{\oxemptyctx}
    }{
      \oxnewframe{\oxstore}{\oxemptyctx}
    }$. Then, we repeatedly apply \Lemma{lemma:store-validity-extension} to
    the derivations for the arguments ($\oxvalue_1 \oxdots \oxvalue_n$) to get
    $\oxstorevalidity{\oxglobalctx}{
      \oxnewframe{\oxvarctx}{
        \oxvarctxentry{\oxid_1}{\delta(\oxsitype_1)} \oxdotsc
        \oxvarctxentry{\oxid_n}{\delta(\oxsitype_n)}
      }
    }{
      \oxnewframe{\oxstore}{
        \oxstoreentry{\oxid_1}{\oxvalue_1} \oxdotsc
        \oxstoreentry{\oxid_n}{\oxvalue_n}
      }
    }$.
  }

  %% kontctx values are still valid
  \oxpresline{
    \oxwfkontctx{\oxglobalctx}{\oxvarctx_i}{\overline{\oxvalue}}{\oxkontctx}
  }{
    Inverting \oxname{WF-Temporaries} gives us $\forall i \in 1 \oxdots n. \;
    \oxtypjudge{\oxglobalctx}{\oxemptyctx}{\oxsitype_1 \oxdotsc \oxsitype_{i-1}}{\oxvarctx}{
      \oxvalue_i
    }{\oxsitype_i}{\oxvarctx}$.

    Applying \Lemma{lemma:values-dont-change} to the derivation for $\oxvalue_f$
    gives us $\oxsubtypectx{\oxglobalctx}{\oxemptyctx}{\oxvarctx}{\oxvarctx_0}$.
    Then, applying \Lemma{lemma:values-dont-change} to the derivations for
    every $\oxvalue_i$ gives us $\forall i \in \oxset{1 \oxdots n}. \;
    \oxsubtypectx{\oxglobalctx}{\oxemptyctx}{\oxvarctx_{i-1}}{\oxvarctx_i}$.

    Inversion on $\oxrunifymany{\oxemptyctx}{\oxkontctx}{\oxvarctx_n}{
      \oxsubstmany{\oxabsprov_2}{\oxprov}{\oxabsprov}
    }{
      \oxsubstmany{\oxabsprov_1}{\oxprov}{\oxabsprov}
    }{\oxvarctx_b}$ gives us a sequence of outlives relations with
    intermediate contexts. Applying  \Lemma{lemma:outlives-related} to each of
    them and then combining the result by transitivity gives us
    $\oxsubtypectx{\oxglobalctx}{\oxemptyctx}{\oxvarctx_n}{\oxvarctx_b}$.
    Combining both by transitivity, we have
    $\oxsubtypectx{\oxglobalctx}{\oxemptyctx}{\oxvarctx}{\oxvarctx_b}$.

    Thus, we can apply \Lemma{lemma:value-typing-related-envs} to each of the
    typing judgments to get $\forall i \in 1 \oxdots n. \;
    \oxtypjudge{\oxglobalctx}{\oxemptyctx}{\oxemptyctx}{\oxvarctx_b}{
      \oxvalue_i
    }{\oxsitype_i}{\oxvarctx_b}$.

    We also immediately have that adding a new empty frame leaves all the typing
    derivations good since it doesn't actually bind any new names on its own.
    Thus, we have $\forall i \in 1 \oxdots n. \;
    \oxtypjudge{\oxglobalctx}{\oxemptyctx}{\oxemptyctx}{
      \oxnewframe{\oxvarctx_b}{\oxemptyctx}
    }{
      \oxvalue_i
    }{\oxsitype_i}{
      \oxnewframe{\oxvarctx_b}{\oxemptyctx}
    }$.

    Repeatedly applying \Lemma{lemma:values-well-typed-extension} to each of the value
    typing judgments gives us
    $\forall i \in 1 \oxdots n. \;
    \oxtypjudge{\oxglobalctx}{\oxemptyctx}{\oxemptyctx}{
      \oxnewframe{\oxvarctx_b}{
        \oxvarctxentry{\oxid_1}{\delta(\oxsitype_1)} \oxdotsc
        \oxvarctxentry{\oxid_n}{\delta(\oxsitype_n)}
      }
    }{
      \oxvalue_i
    }{\oxsitype_i}{
      \oxnewframe{\oxvarctx_b}{
        \oxvarctxentry{\oxid_1}{\delta(\oxsitype_1)} \oxdotsc
        \oxvarctxentry{\oxid_n}{\delta(\oxsitype_n)}
      }
    }$.
  }

  %% new expression is well-typed
  \oxpresline{
    \oxtypjudge{\oxglobalctx}{\oxemptyctx}{\oxkontctx}{\oxvarctx_i}{
      \oxexpr^\prime
    }{\delta(\oxsitype_f)}{\oxvarctx_b}
  }{
    Applying \Lemma{lemma:function-defns-self-contained} with
    $\oxctxswellformed{\oxglobalctx}{\oxemptyctx}{\oxvarctx_b}{\oxkontctx}$ and
    $\oxlookup{\oxglobalctx}{\oxfnname}{
      \oxfuncdef{\oxfnname}{
        \overline{\oxenvvar} \oxcomma \overline{\oxabsprov} \oxcomma \overline{\oxtvar}
      }{
        \oxascribe{\oxid_1}{\oxsitype_1} \oxdotsc
        \oxascribe{\oxid_n}{\oxsitype_n}
      }{\oxsitype_r}{\overline{\oxabsprov_1: \oxabsprov_2}}{\oxexpr}
    }$ (from the premise of on \oxname{E-AppFunction}) gives us
    $\oxtypjudge{\oxglobalctx}{
      \overline{\oxtvarctxentry{\oxenvvar}{\oxkenv}} \oxcomma
      \overline{\oxtvarctxentry{\oxabsprov}{\oxkprov}} \oxcomma
      \overline{\oxabsprov_1 \oxoutlives \oxabsprov_2} \oxcomma
      \overline{\oxtvarctxentry{\oxtvar}{\oxktype}}
    }{\oxkontctx}{
      \oxnewframe{\oxvarctx_b}{
        \oxvarctxentry{\oxid_1}{\oxsitype_1} \oxdotsc
        \oxvarctxentry{\oxid_n}{\oxsitype_n}
      }
    }{
      \oxframed{\oxexpr}
    }{
      \oxsitype_f
    }{\oxvarctx_b}$.

    We then repeatedly apply \Lemma{lemma:substitution} for all of our type,
    region, and environment variables to get
    $\oxtypjudge{\oxglobalctx}{\oxemptyctx}{\oxkontctx}{
      \oxnewframe{\oxvarctx_b}{
        \oxvarctxentry{\oxid_1}{\delta(\oxsitype_1)} \oxdotsc
        \oxvarctxentry{\oxid_n}{\delta(\oxsitype_n)}
      }
    }{
      \oxframed{\oxexpr}
    }{\delta(\oxsitype_f)}{\oxvarctx_b}$.
  }

  %% new type is a subtype of the old type
  \oxpresline{
    \oxtunify[\oxcombine]{\oxemptyctx}{\oxkontctx}{\oxvarctx_b}{
      \oxtype^\prime
    }{
      \oxtype^\prime
    }{\oxvarctx_b}
  }{
    Immediate by \oxname{RR-Refl}.
  }

  %% new loan context is a sub-context of the old loan context
  \oxpresline{
    \exists \oxvarctx_o. \oxvarctx_b \oxintersect \oxvarctx_o = \oxvarctx_b
  }{
    $\oxvarctx_o = \oxvarctx_b$
  }
\end{preservation}

\oxpreservationcase{T-AppClosure}{\TAppClosure}{
  \oxapply{\oxexpr_f}{
    \oxexpr_1 \oxdotsc \oxexpr_n
  }
}{\EApp}{
  \oxnewframe{\oxvarctx^\prime_n}{
    \oxextendctx{\oxframe_c}{
      \oxvarctxentry{\oxid_1}{\oxsitype_1} \oxdotsc
      \oxvarctxentry{\oxid_n}{\oxsitype_n}
    }
  }
}

\begin{preservation}
  %% new store is valid
  \oxpresline{
    \oxstorevalidity{\oxglobalctx}{
      \oxvarctx_i
    }{
      \oxstore^\prime
    }
  }{
    Applying \Lemma{lemma:values-dont-change} to the derivation for $\oxvalue_f$
    gives us $\oxsubtypectx{\oxglobalctx}{\oxemptyctx}{\oxvarctx}{\oxvarctx_0}$.
    We can apply \Lemma{lemma:stack-validity-related-envs} with
    $\oxstorevalidity{\oxglobalctx}{\oxvarctx}{\oxstore}$ to get
    $\oxstorevalidity{\oxglobalctx}{\oxvarctx_0}{\oxstore}$.

    Then, for each pair of value typing judgment and region rewriting, we follow
    the same pattern we will describe in terms of an arbitrary index $i$. First,
    we apply \Lemma{lemma:values-dont-change} to the derivation for $\oxvalue_i$
    to get $\oxsubtypectx{\oxglobalctx}{\oxemptyctx}
    {\oxvarctx_{i-1}}{\oxvarctx_i}$. Then, applying
    \Lemma{lemma:stack-validity-related-envs} gives us
    $\oxstorevalidity{\oxglobalctx}{\oxvarctx_i}{\oxstore}$. Then, we can apply
    \Lemma{lemma:stack-validity-region-rewriting} to
    $\oxtunify[\oxcombineunrest]{\oxemptyctx}{\oxkontctx}{\oxvarctx_i}
    {\oxsitype_{i^\prime}}{\oxsitype_i}{\oxvarctx_i^\prime}$ to get
    $\oxstorevalidity{\oxglobalctx}{\oxvarctx_i^\prime}{\oxstore}$. Going through
    this for all the values in sequence gives us
    $\oxstorevalidity{\oxglobalctx}{\oxvarctx_n^\prime}{\oxstore}$

    Then, inversion of \oxname{T-ClosureValue} for the typing derivation for
    $\oxvalue_f$ gives us
    $\oxsubstorevalidity{\oxglobalctx}{\oxvarctx}{\oxstackframe_c}{\oxframe_c}$.
    We can then invert \oxname{WF-Frame} here to get
    $\oxdomain{\oxstackframe} = \oxdomain{\oxframe_c}|_\oxid$
    $\forall \oxid \in \oxdomain{\oxstackframe}. \;
    \oxtypjudge{\oxglobalctx}{\oxemptyctx}{\oxkontctx}{\oxnewframe{\oxvarctx}{\oxframe_c}}
    {\oxstackframe(\oxid)}{\oxframe_c(\oxid)}{\oxnewframe{\oxvarctx}{\oxframe_c}}$
    which we can then use with $\oxstorevalidity{\oxglobalctx}{\oxvarctx_n}{\oxstore}$
    in \oxname{WF-StackFrame} to get
    $\oxstorevalidity{\oxglobalctx}{
      \oxnewframe{\oxvarctx_n^{\prime}}{\oxframe_c}
    }{
      \oxnewframe{\oxstore}{\oxstackframe_c}
    }$.

    Finally, we repeatedly apply \Lemma{lemma:store-validity-extension} to
    the derivations for the arguments ($\oxvalue_1 \oxdots \oxvalue_n$) to get
    $\oxstorevalidity{\oxglobalctx}{
      \oxvarctx_i
    }{
      \oxstore^\prime
    }$.
  }

  %% kontctx values are still valid
  \oxpresline{
    \oxwfkontctx{\oxglobalctx}{\oxvarctx_i}{\overline{\oxvalue^\prime}}{\oxkontctx}
  }{
    Inverting \oxname{WF-Temporaries} gives us $\forall j \in 1 \oxdots n. \;
    \oxtypjudge{\oxglobalctx}{\oxemptyctx}{\oxsitype_1 \oxdotsc \oxsitype_{i-1}}{\oxvarctx}{
      \oxvalue^\prime_j
    }{\oxsitype_j}{\oxvarctx}$.

    Applying \Lemma{lemma:values-dont-change} to the derivation for $\oxvalue_f$
    gives us $\oxsubtypectx{\oxglobalctx}{\oxemptyctx}{\oxvarctx}{\oxvarctx_0}$.
    For each of $\forall j \in 1 \oxdots n. \;
    \oxtypjudge{\oxglobalctx}{\oxemptyctx}{\oxemptyctx}{\oxvarctx}{
      \oxvalue^\prime_j
    }{\oxsitype_i}{\oxvarctx}$, we apply \Lemma{lemma:value-typing-related-envs}
    to get $\forall j \in 1 \oxdots n. \;
    \oxtypjudge{\oxglobalctx}{\oxemptyctx}{\oxemptyctx}{\oxvarctx_0}{
      \oxvalue^\prime_j
    }{\oxsitype_j}{\oxvarctx_0}$

    From here, we have an alternating pattern of lemmas to apply to deal with
    the interleaved value typing and then region rewriting in the premise of
    \oxname{T-AppClosure}. For each value $\oxvalue_i$ in \oxname{T-AppClosure},
    we'll first apply \Lemma{lemma:values-dont-change} to get
    $\oxsubtypectx{\oxglobalctx}{\oxemptyctx}{\oxvarctx_{i-1}}{\oxvarctx_i}$.
    Then, we'll apply \Lemma{lemma:value-typing-related-envs} to
    $\forall j \in 1 \oxdots n. \;
    \oxtypjudge{\oxglobalctx}{\oxemptyctx}{\oxemptyctx}{\oxvarctx_{i-1}}{
      \oxvalue^\prime_j
    }{\oxsitype_j}{\oxvarctx_{i-1}}$ to get $\forall j \in 1 \oxdots n. \;
    \oxtypjudge{\oxglobalctx}{\oxemptyctx}{\oxemptyctx}{\oxvarctx_i}{
      \oxvalue^\prime_j
    }{\oxsitype_j}{\oxvarctx_i}$. Then, we apply
    \Lemma{lemma:subtyping-value-typing} to each of these along with the
    rewriting $\oxtunify[\oxcombineunrest]{\oxemptyctx}{\oxkontctx}{\oxvarctx_i}
    {\oxsitype_{i^\prime}}{\oxsitype_i}{\oxvarctx_i^\prime}$ to get
    $\forall j \in 1 \oxdots n. \;
    \oxtypjudge{\oxglobalctx}{\oxemptyctx}{\oxemptyctx}{\oxvarctx_i^\prime}{
      \oxvalue^\prime_j
    }{\oxsitype_j}{\oxvarctx_i^\prime}$. After going through this for each value
    $\oxvalue_i$, we have $\forall j \in 1 \oxdots n. \;
    \oxtypjudge{\oxglobalctx}{\oxemptyctx}{\oxemptyctx}{\oxvarctx_n^\prime}{
      \oxvalue^\prime_j
    }{\oxsitype_j}{\oxvarctx_n^\prime}$.

    We also immediately have that adding a new empty frame leaves all the typing
    derivations good since it doesn't actually bind any new names on its own.
    Thus, we have $\forall i \in 1 \oxdots n. \;
    \oxtypjudge{\oxglobalctx}{\oxemptyctx}{\oxemptyctx}{
      \oxnewframe{\oxvarctx_n^\prime}{\oxemptyctx}
    }{
      \oxvalue_i
    }{\oxsitype_i}{
      \oxnewframe{\oxvarctx_n^\prime}{\oxemptyctx}
    }$.

    Repeatedly applying \Lemma{lemma:values-well-typed-extension} to each of the
    value typing judgments and each of $\forall i \in \oxset{ 1 \oxdots n }. \;
    \forall \oxcprov \in \oxfprovs{\oxsitype_i}. \;
    \oxnotreborrowed{\oxvarctx_1}{\oxcprov}$ (from the premise of
    \oxname{T-App}) gives us
    $\forall i \in 1 \oxdots n. \;
    \oxtypjudge{\oxglobalctx}{\oxemptyctx}{\oxemptyctx}{
      \oxnewframe{\oxvarctx_n^\prime}{
        \oxextendctx{\oxframe_c}{
          \oxvarctxentry{\oxid_1}{\oxsitype_1} \oxdotsc
          \oxvarctxentry{\oxid_n}{\oxsitype_n}
        }
      }
    }{
      \oxvalue_i
    }{\oxsitype_i}{
      \oxnewframe{\oxvarctx_n^\prime}{
        \oxextendctx{\oxframe_c}{
          \oxvarctxentry{\oxid_1}{\oxsitype_1} \oxdotsc
          \oxvarctxentry{\oxid_n}{\oxsitype_n}
        }
      }
    }$.
  }
\end{preservation} %% breaking for pagebreak
\begin{preservation}
  %% new expression is well-typed
  \oxpresline{
    \oxtypjudge{\oxglobalctx}{\oxemptyctx}{\oxkontctx}{\oxvarctx_i}{
      \oxexpr^\prime
    }{\oxsitype_f}{\oxvarctx^\prime_n}
  }{
    Consider first the typing derivation for $\oxvalue_f$ of
    $\oxtypjudge{\oxglobalctx}{\oxemptyctx}{\oxkontctx}{\oxvarctx}{
      \oxvalue_f
    }{
      \oxtfun{}{
        \oxsitype_1 \oxdotsc \oxsitype_n
      }{\oxsitype_f}{\oxenv_c}
    }{\oxvarctx_0^\prime}$.

    Applying \Lemma{lemma:value-typing-output} gives us
    $\oxtypjudge{\oxglobalctx}{\oxemptyctx}{\oxkontctx}{\oxvarctx_0^\prime}{
      \oxvalue_f
    }{
      \oxtfun{}{
        \oxsitype_1 \oxdotsc \oxsitype_n
      }{\oxsitype_f}{\oxenv_c}
    }{\oxvarctx_0^\prime}$.

    Then, as in the previous cases, we can apply
    \Lemma{lemma:values-dont-change}, \Lemma{lemma:value-typing-related-envs},
    and \Lemma{lemma:subtyping-value-typing} in sequence starting with the above
    typing derivation with each of the typing derivations for the arguments
    $\oxvalue_i$ and their subsequent corresponding region rewriting derivation.
    At the end, this gives us
    $\oxtypjudge{\oxglobalctx}{\oxemptyctx}{\oxkontctx}{\oxvarctx_n^\prime}{
      \oxvalue_f
    }{
      \oxtfun{}{
        \oxsitype_1 \oxdotsc \oxsitype_n
      }{\oxsitype_f}{\oxenv_c}
    }{\oxvarctx_n^\prime}$

    Then, inversion on \oxname{T-ClosureValue} on this typing derivation for
    $\oxvalue_f$ gives us $\oxfvars{\oxexpr} \setminus \overline{\oxid} =
    \overline{\oxid_f} = \oxdomain{\oxframe_c}|_\oxid$,
    $\overline{\oxcprov} = \overline{\oxfprovs{\oxvarctx(\oxid_f)}} \oxcomma
    \oxfprovs{\oxexpr} = \oxdomain{\oxframe_c}|_\oxcprov$, and
    $\oxtypjudge{\oxglobalctx}{\oxemptyctx}{\oxkontctx}{
      \oxnewframe{\oxvarctx_n^\prime}{
        \oxframe_c \oxcomma
        \oxvarctxentry{\oxid_1}{\oxsitype_1} \oxdotsc
        \oxvarctxentry{\oxid_n}{\oxsitype_n}
      }
    }{\oxexpr}{\oxsitype_r}{\oxnewframe{\oxvarctx_n^\prime}{\oxframe}}$. We can
    then apply \oxname{T-Framed} to get
    $\oxtypjudge{\oxglobalctx}{\oxemptyctx}{\oxkontctx}{
      \oxnewframe{\oxvarctx_n}{
        \oxextendctx{\oxframe_c}{
          \oxvarctxentry{\oxid_1}{\oxsitype_1} \oxdotsc
          \oxvarctxentry{\oxid_n}{\oxsitype_n}
        }
      }
    }{
      \oxframed{\oxexpr}
    }{\oxsitype_f}{\oxvarctx_n}$.
  }

  %% new type is a subtype of the old type
  \oxpresline{
    \oxtunify[\oxcombine]{\oxemptyctx}{\oxkontctx}{\oxvarctx^\prime_n}{
      \oxsitype_f
    }{
      \oxsitype_f
    }{\oxvarctx^\prime_n}
  }{
    Immediate by \oxname{RR-Refl}.
  }

  %% new loan context is a sub-context of the old loan context
  \oxpresline{
    \exists \oxvarctx_o. \oxvarctx^\prime_n \oxintersect \oxvarctx_o =
    \oxvarctx^\prime_n
  }{
    $\oxvarctx_o = \oxvarctx^\prime_n$
  }
\end{preservation}

\oxpreservationcase{T-Function}{\TFunction}{
  \oxtfun{\overline{\oxprov} \oxcomma \overline{\oxtvar}}{
    \oxsitype_1 \oxdotsc \oxsitype_n
  }{\oxsitype_r}{}
}{\EFunction}{
  \oxvarctx
}

\begin{preservation}
  %% new store is valid
  \oxpresline{\oxstorevalidity{\oxglobalctx}{\oxvarctx}{\oxstore}}{
    Immediate from our premise.
  }

  %% kontctx values are still valid
  \oxpresline{
    \oxwfkontctx{\oxglobalctx}{\oxvarctx}{\overline{\oxvalue}}{\oxkontctx}
  }{
    Immediate from our premise.
  }

  %% new expression is well-typed
  \oxpresline{
    \oxtypjudge{\oxglobalctx}{\oxemptyctx}{\oxkontctx}{\oxvarctx}
    {\oxexpr^\prime}{\oxtype^\prime}{\oxvarctx}
  }{
    By \oxname{T-ClosureValue} since $\oxstore_c$ is empty, and we know that the
    body itself is well-typed as a consequence of inversion on
    \oxname{WF-FunctionDefinition} for \oxfnname.
  }

  %% new type is a subtype of the old type
  \oxpresline{
    \oxtunify[\oxcombine]{\oxemptyctx}{\oxkontctx}{\oxvarctx}
    {\oxtype^\prime}{\oxtype}{\oxvarctx}
  }{
    Immediate by \oxname{RR-Refl}.
  }

  %% new loan context is a sub-context of the old loan context
  \oxpresline{
    \exists \oxvarctx_o. \oxvarctx \oxintersect \oxvarctx_o = \oxvarctx
  }{
    $\oxvarctx_o = \oxvarctx$
  }
\end{preservation}


\oxpreservationcase{T-Closure}{\TClosure}{
  \oxfunction{
    \overline{\oxenvvar} \oxcomma \overline{\oxprov} \oxcomma \overline{\oxtvar}
  }{
    \oxascribe{\oxid_1}{\oxsitype_1} \oxdotsc
    \oxascribe{\oxid_n}{\oxsitype_n}
  }{\oxsitype_r}{
    \oxexpr
  }
}{\EClosure}{
  \oxtupdatemany{\oxvarctx}{\oxid_{nc}}{\oxtuninit{\oxvarctx(\oxid_{nc})}}
}

\begin{preservation}
  %% new store is valid
  \oxpresline{
    \oxstorevalidity{\oxglobalctx}{
      \oxtupdatemany{\oxvarctx}{\oxid_{nc}}{\oxtuninit{\oxvarctx(\oxid_{nc})}}
    }{
      \oxupdate{\oxstore}{
        \overline{\oxstoreentry{\oxid_{nc}}{\oxuninit}}
      }
    }
  }{
    Compared to $\oxvarctx$, we know that
    $\oxtupdatemany{\oxvarctx}{\oxid_{nc}}{\oxtuninit{\oxvarctx(\oxid_{nc})}}$
    has more things marked dead and no other changes. Thus, we can apply
    \oxname{R-Env} to get $\oxsubtypectx{\oxglobalctx}{\oxemptyctx}{\oxvarctx}{
      \oxtupdatemany{\oxvarctx}{\oxid_{nc}}{\oxtuninit{\oxvarctx(\oxid_{nc})}}
    }$. Then, we can apply \Lemma{lemma:stack-validity-related-envs} to get
    $\oxstorevalidity{\oxglobalctx}{
      \oxtupdatemany{\oxvarctx}{\oxid_{nc}}{\oxtuninit{\oxvarctx(\oxid_{nc})}}
    }{\oxstore}$. Since $\oxuninit$ is good at every dead type $\oxsdtype$ by
    \oxname{T-Dead}, we can then build a new derivation using that rule instead
    for every $\oxid_{nc}$ that is now at a dead type. This gives us
    $\oxstorevalidity{\oxglobalctx}{
      \oxtupdatemany{\oxvarctx}{\oxid_{nc}}{\oxtuninit{\oxvarctx(\oxid_{nc})}}
    }{
      \oxupdate{\oxstore}{
        \overline{\oxstoreentry{\oxid_{nc}}{\oxuninit}}
      }
    }$.
  }

  %% kontctx values are still valid
  \oxpresline{
    \oxwfkontctx{\oxglobalctx}{
      \oxtupdatemany{\oxvarctx}{\oxid_{nc}}{\oxtuninit{\oxvarctx(\oxid_{nc})}}
    }{\overline{\oxvalue}}{\oxkontctx}
  }{
    Inverting \oxname{WF-Temporaries} gives us $\forall i \in 1 \oxdots n. \;
    \oxtypjudge{\oxglobalctx}{\oxemptyctx}{\oxsitype_1 \oxdotsc \oxsitype_{i-1}}{\oxvarctx}{
      \oxvalue_i
    }{\oxsitype_i}{\oxvarctx}$. By \oxname{R-Env}, we have that
    \oxsubtypectx{\oxglobalctx}{\oxemptyctx}{\oxvarctx}{
      \oxtupdatemany{\oxvarctx}{\oxid_{nc}}{\oxtuninit{\oxvarctx(\oxid_{nc})}}
    }. Thus, we can apply \Lemma{lemma:value-typing-related-envs} to each of the
    typing judgments and then apply \oxname{WF-Temporaries} to get
    $\oxwfkontctx{\oxglobalctx}{
      \oxtupdatemany{\oxvarctx}{\oxid_{nc}}{\oxtuninit{\oxvarctx(\oxid_{nc})}}
    }{\overline{\oxvalue}}{\oxkontctx}$.
  }

  %% new expression is well-typed
  \oxpresline{
    \oxtypjudge{\oxglobalctx}{\oxemptyctx}{\oxkontctx}{\oxvarctx_i}
    {\oxexpr^\prime}{\oxtype^\prime}{\oxvarctx_i}
  }{
    Immediate by inversion of \oxname{T-Closure} and application of
    \oxname{T-ClosureValue}. Note that they have identical premises.
  }

  %% new type is a subtype of the old type
  \oxpresline{
    \oxtunify[\oxcombine]{\oxemptyctx}{\oxkontctx}{\oxvarctx_i}
    {\oxtype^\prime}{\oxtype^\prime}{\oxvarctx_i}
  }{
    Immediate by \oxname{RR-Refl}.
  }

  %% new loan context is a sub-context of the old loan context
  \oxpresline{
    \exists \oxvarctx_o. \oxvarctx_i \oxintersect \oxvarctx_o = \oxvarctx_i
  }{
    $\oxvarctx_o = \oxvarctx_i$
  }
\end{preservation}

\oxpreservationcase{T-Shift}{\TShift}{
  \oxshift{\oxexpr_i}
}{\EShift}{\oxvarctx^\prime}

\begin{preservation}
  %% new store is valid
  \oxpresline{
    \oxstorevalidity{\oxglobalctx}{\oxvarctx^\prime}{\oxstore}
  }{
    By inversion of \oxname{WF-StackFrame} on $\oxstorevalidity{\oxglobalctx}{
      \oxextendctx{\oxvarctx^\prime}{\oxvarctxentry{\oxid}{\oxsdtype}}
    }{\oxextendctx{\oxstore}{\oxstoreentry{\oxid}{\oxvalue^\prime}}}$, we get
    $\oxstorevalidity{\oxglobalctx}{\oxvarctx_i}{\oxstore}$,
    $\oxdomain{
      \oxextendctx{\oxstackframe}{
        \oxstoreentry{\oxid}{\oxvalue^\prime}
      }
    } = \oxdomain{
      \oxextendctx{\oxframe}{
        \oxvarctxentry{\oxid}{\oxsdtype}
      }
    }|_\oxid$, and
    $\forall \oxid \in \oxdomain{
      \oxnewframe{\oxstore}{
        \oxextendctx{\oxstackframe}{
          \oxstoreentry{\oxid}{\oxvalue^\prime}
        }
      }
    }. \;
    \oxtypjudge{\oxglobalctx}{\oxemptyctx}{\oxkontctx}{
      \oxnewframe{\oxvarctx_i}{
        \oxextendctx{\oxframe}{
          \oxvarctxentry{\oxid}{\oxsdtype}
        }
      }
    }{
      (\oxnewframe{\oxstore}{
        \oxextendctx{\oxstackframe}{
          \oxstoreentry{\oxid}{\oxvalue^\prime}
        }
      })(\oxid)
    }{(\oxnewframe{\oxvarctx_i}{
        \oxextendctx{\oxframe}{
          \oxvarctxentry{\oxid}{\oxsdtype}
        }
      })(\oxid)}{
      \oxnewframe{\oxvarctx_i}{
        \oxextendctx{\oxframe}{
          \oxvarctxentry{\oxid}{\oxsdtype}
        }
      }
    }$. Note that $\oxnewframe{\oxvarctx_i}{\oxframe} = \oxvarctx^\prime$.
    We can then immediately see that the above implies
    $\oxdomain{\oxstackframe} = \oxdomain{\oxframe}|_\oxid$ and
    $\forall \oxid \in \oxdomain{
      \oxnewframe{\oxstore}{
        \oxextendctx{\oxstackframe}{
          \oxstoreentry{\oxid}{\oxvalue^\prime}
        }
      }
    }. \;
    \oxtypjudge{\oxglobalctx}{\oxemptyctx}{\oxkontctx}{
      \oxnewframe{\oxvarctx_i}{\oxframe}
    }{
      (\oxnewframe{\oxstore}{\oxstackframe})(\oxid)
    }{(\oxnewframe{\oxvarctx_i}{\oxframe})(\oxid)}{
      \oxnewframe{\oxvarctx_i}{\oxframe}
    }$. Thus, we can apply \oxname{WF-StackFrame} to get
    $\oxstorevalidity{\oxglobalctx}{
      \oxnewframe{\oxvarctx_i}{\oxstackframe}
    }{\oxstore}$ which can be rewritten as
    $\oxstorevalidity{\oxglobalctx}{\oxvarctx^\prime}{\oxstore}$.
  }

  %% kontctx values are still valid
  \oxpresline{
    \oxwfkontctx{\oxglobalctx}{\oxvarctx^\prime}{\overline{\oxvalue}}{\oxkontctx}
  }{
    Inverting \oxname{WF-Temporaries} gives us $\forall i \in 1 \oxdots n. \;
    \oxtypjudge{\oxglobalctx}{\oxemptyctx}{\oxsitype_1 \oxdotsc \oxsitype_{i-1}}{\oxvarctx}{
      \oxvalue_i
    }{\oxsitype_i}{\oxvarctx}$.

    Applying \Lemma{lemma:values-dont-change} to the typing derivation (from
    \oxname{T-Shift}) for
    $\oxexpr$ (which we know is a value from \oxname{E-Shift}) gives us
    $\oxsubtypectx{\oxglobalctx}{\oxemptyctx}{\oxvarctx}{\oxvarctx^\prime}$.

    Thus, we can apply \Lemma{lemma:value-typing-related-envs} to each of the
    typing judgments to get $\forall i \in 1 \oxdots n. \;
    \oxtypjudge{\oxglobalctx}{\oxemptyctx}{\oxsitype_1 \oxdotsc \oxsitype_{i-1}}{\oxvarctx^\prime}{
      \oxvalue_i
    }{\oxsitype_i}{\oxvarctx^\prime}$.

    Finally, applying \oxname{WF-Temporaries} gives us
    $\oxwfkontctx{\oxglobalctx}{\oxvarctx^\prime}{\overline{\oxvalue}}{\oxkontctx}$.
  }

  %% new expression is well-typed
  \oxpresline{
    \oxtypjudge{\oxglobalctx}{\oxemptyctx}{\oxkontctx}{\oxvarctx^\prime}{\oxvalue}
    {\oxsitype}{\oxvarctx^\prime}
  }{
    We know from \oxname{E-Shift} that $\oxexpr$ is a value $\oxvalue$. Thus,
    we can apply \Lemma{lemma:value-typing-output} to
    $\oxtypjudge{\oxglobalctx}{\oxemptyctx}{\oxkontctx}{\oxvarctx}{\oxvalue}{\oxsitype}{
      \oxextendctx{\oxvarctx^\prime}{
        \oxvarctxentry{\oxid}{\oxsdtype}
      }
    }$ to get $\oxtypjudge{\oxglobalctx}{\oxemptyctx}{\oxkontctx}{
      \oxextendctx{\oxvarctx^\prime}{
        \oxvarctxentry{\oxid}{\oxsdtype}
      }
    }{\oxvalue}{\oxsitype}{
      \oxextendctx{\oxvarctx^\prime}{
        \oxvarctxentry{\oxid}{\oxsdtype}
      }
    }$.

    We now wish to show that
    $\oxtypjudge{\oxglobalctx}{\oxemptyctx}{\oxkontctx}{\oxvarctx^\prime}{\oxvalue}
    {\oxsitype}{\oxvarctx^\prime}$. By inspecting the grammar of values and
    their typing rules, we know that the only values who depend on the
    context are pointers and closure values. But by inversion on
    $\oxtypjudge{\oxglobalctx}{\oxemptyctx}{\oxkontctx}{\oxvarctx}{
      \oxshift{\oxvalue}
    }{\oxsitype}{\oxvarctx^\prime}$, we know that
    $\oxtypevalidity{\oxglobalctx}{\oxemptyctx}{\oxvarctx^\prime}{\oxsitype}$.
    Since the type is valid without the frame, we know that the values
    cannot depend on that frame. Thus, we can conclude
    $\oxtypjudge{\oxglobalctx}{\oxemptyctx}{\oxkontctx}{\oxvarctx^\prime}{\oxvalue}
    {\oxsitype}{\oxvarctx^\prime}$.
  }

  %% new type is a subtype of the old type
  \oxpresline{
    \oxtunify[\oxcombine]{\oxemptyctx}{\oxkontctx}{\oxvarctx^\prime}
    {\oxsitype}{\oxsitype}{\oxvarctx^\prime}
  }{
    Immediate by \oxname{RR-Refl}.
  }

  %% new loan context is a sub-context of the old loan context
  \oxpresline{
    \exists \oxvarctx_o. \oxvarctx^\prime \oxintersect \oxvarctx_o = \oxvarctx^\prime
  }{
    $\oxvarctx_o = \oxvarctx^\prime$
  }
\end{preservation}

\oxpreservationcase{T-Framed}{\TFramed}{
  \oxframed{\oxexpr_i}
}{\EFramed}{\oxvarctx^\prime}

\begin{preservation}
  %% new store is valid
  \oxpresline{
    \oxstorevalidity{\oxglobalctx}{\oxvarctx^\prime}{\oxstore}
  }{
    Applying \Lemma{lemma:stack-envminus} to $\oxstorevalidity{\oxglobalctx}{
      \oxnewframe{\oxvarctx^\prime}{\oxframe}
    }{\oxnewframe{\oxstore}{\oxstackframe}}$ gives us
    $\oxstorevalidity{\oxglobalctx}{\oxvarctx^\prime}{\oxstore}$.
  }

  %% kontctx values are still valid
  \oxpresline{
    \oxwfkontctx{\oxglobalctx}{\oxvarctx^\prime}{\overline{\oxvalue}}{\oxkontctx}
  }{
    Inverting \oxname{WF-Temporaries} gives us $\forall i \in 1 \oxdots n. \;
    \oxtypjudge{\oxglobalctx}{\oxemptyctx}{\oxsitype_1 \oxdotsc \oxsitype_{i-1}}{\oxvarctx}{
      \oxvalue_i
    }{\oxsitype_i}{\oxvarctx}$.

    Applying \Lemma{lemma:values-dont-change} to the typing derivation (from
    \oxname{T-Framed}) for
    $\oxexpr$ (which we know is a value from \oxname{E-Framed}) gives us
    $\oxsubtypectx{\oxglobalctx}{\oxemptyctx}{\oxvarctx}{\oxvarctx^\prime}$.

    Thus, we can apply \Lemma{lemma:value-typing-related-envs} to each of the
    typing judgments to get $\forall i \in 1 \oxdots n. \;
    \oxtypjudge{\oxglobalctx}{\oxemptyctx}{\oxsitype_1 \oxdotsc \oxsitype_{i-1}}{\oxvarctx^\prime}{
      \oxvalue_i
    }{\oxsitype_i}{\oxvarctx^\prime}$.

    Finally, applying \oxname{WF-Temporaries} gives us
    $\oxwfkontctx{\oxglobalctx}{\oxvarctx^\prime}{\overline{\oxvalue}}{\oxkontctx}$.
  }

  %% new expression is well-typed
  \oxpresline{
    \oxtypjudge{\oxglobalctx}{\oxemptyctx}{\oxkontctx}{\oxvarctx^\prime}{\oxvalue}
    {\oxsitype}{\oxvarctx^\prime}
  }{
    We know from \oxname{E-Framed} that $\oxexpr$ is a value $\oxvalue$. Thus,
    we can apply \Lemma{lemma:value-typing-output} to
    $\oxtypjudge{\oxglobalctx}{\oxemptyctx}{\oxkontctx}{
      \oxnewframe{\oxvarctx}{\oxframe}
    }{\oxvalue}{\oxsitype}{
      \oxnewframe{\oxvarctx^\prime}{\oxframe^\prime}
    }$ to get $\oxtypjudge{\oxglobalctx}{\oxemptyctx}{\oxkontctx}{
      \oxnewframe{\oxvarctx^\prime}{\oxframe^\prime}
    }{\oxvalue}{\oxsitype}{
      \oxnewframe{\oxvarctx^\prime}{\oxframe^\prime}
    }$.

    We now wish to show that
    $\oxtypjudge{\oxglobalctx}{\oxemptyctx}{\oxkontctx}{\oxvarctx^\prime}{\oxvalue}
    {\oxsitype}{\oxvarctx^\prime}$. By inspecting the grammar of values and
    their typing rules, we know that the only values who depend on the
    context are pointers and closure values. But by inversion on
    $\oxtypjudge{\oxglobalctx}{\oxemptyctx}{\oxkontctx}{\oxvarctx}{
      \oxframed{\oxvalue}
    }{\oxsitype}{\oxvarctx^\prime}$, we know that
    $\oxtypevalidity{\oxglobalctx}{\oxemptyctx}{\oxvarctx^\prime}{\oxsitype}$.
    Since the type is valid without the frame, we know that the values
    cannot depend on that frame. Thus, we can conclude
    $\oxtypjudge{\oxglobalctx}{\oxemptyctx}{\oxkontctx}{\oxvarctx^\prime}{\oxvalue}
    {\oxsitype}{\oxvarctx^\prime}$.
  }

  %% new type is a subtype of the old type
  \oxpresline{
    \oxtunify[\oxcombine]{\oxemptyctx}{\oxkontctx}{\oxvarctx^\prime}
    {\oxsitype}{\oxsitype}{\oxvarctx^\prime}
  }{
    Immediate by \oxname{RR-Refl}.
  }

  %% new loan context is a sub-context of the old loan context
  \oxpresline{
    \exists \oxvarctx_o. \oxvarctx^\prime \oxintersect \oxvarctx_o = \oxvarctx^\prime
  }{
    $\oxvarctx_o = \oxvarctx^\prime$
  }
\end{preservation}

\noindent Case \oxname{T-BorrowIndex}: \\[0.5em]
\begin{tabular}{l}
  \pbox{\textwidth}{\vspace*{0.5em} From premise: \vspace*{0.5em}} \\

  \framebox[\textwidth]{
    \figuresize
    \begin{mathpar}
      \TBorrowIndex
    \end{mathpar}
  } \\

  \pbox{\textwidth}{
    \vspace*{0.5em}
    Since $\oxexpr =
    \oxref{\oxcprov}{\oxmuta}{\oxindex{\oxplaceexpr}{\oxexpr}}$, by inspection of
    the reduction rules, we know that $\oxexpr$ steps with the following rule:
    \vspace*{0.5em}
  } \\

  \framebox[\textwidth]{
    \figuresize
    \begin{mathpar}
      \inferrule[E-EvalCtx]{
        %% premises
        \oxreduce{\oxglobalctx}{\oxstore}{\oxexpr}
                 {\oxstore^\prime}{\oxexpr^\prime}
      }{
        \oxreduce{\oxglobalctx}{\oxstore}{\oxref{\oxcprov}{\oxmuta}{\oxindex{\oxplaceexpr}{\oxexpr}}}
                 {\oxstore^\prime}{\oxref{\oxcprov}{\oxmuta}{\oxindex{\oxplaceexpr}{\oxexpr^\prime}}}
      }
    \end{mathpar}
  } \\[1.5em]

\end{tabular}

We begin by applying the induction hypothesis on $\oxexpr$. We get that $\exists
\oxvarctx_i$ such that
$\oxstorevalidity{\oxglobalctx}{\oxvarctx_i}{\oxstore^\prime}$ and
$\oxwfkontctx{\oxglobalctx}{\oxvarctx_i}{\overline{\oxvalue}}{\oxkontctx}$ and
$\oxtypjudge{\oxglobalctx}{\oxemptyctx}{\oxkontctx}{\oxvarctx_i}{\oxexpr^\prime}{\oxtnum}{\oxvarctx_f^\prime}$
and
$\oxtunify[\oxcombine]{\oxemptyctx}{\oxkontctx}{\oxvarctx_f^\prime}{\oxtnum}{\oxtnum}{\oxvarctx_s}$
and there exists $\oxvarctx_o$ such that $\oxvarctx^\prime = \oxvarctx_s
\oxintersect \oxvarctx_o$. Note that since rewriting does nothing with
$\oxtnum$, $\oxvarctx_f^\prime = \oxvarctx_s$. We chose $\oxvarctx_i$ in the
lemma statement to be \framebox{$\oxvarctx_i$} and need to show:

\begin{preservation}
  %% new store is valid
  \oxpresline{
    \oxstorevalidity{\oxglobalctx}{\oxvarctx_i}{\oxstore^\prime}
  }{
    Immediate.
  }
  \oxpresline{
    \oxwfkontctx{\oxglobalctx}{\oxvarctx_i}{\overline{\oxvalue}}{\oxkontctx}
  }{
    Immediate
  }

  %% new expression is well-typed
  \oxpresline{
    \oxtypjudge{\oxglobalctx}{\oxemptyctx}{\oxkontctx}{\oxvarctx_i}{\oxref{\oxcprov}{\oxmuta}
      {\oxplaceexpr[\oxexpr^\prime]}}{\oxtref{\oxcprov}{\oxmuta}{\oxsitype}}
               {\oxvarctx_f^\prime[\oxcprov \mapsto \oxset{\oxloans^\prime}]}
  }{
    We'd like to apply \oxname{T-BorrowIndex}. In order to do so, we still need
    to show $\oxvarctx^\prime_f(\oxcprov) = \emptyset$,
    $\oxnotinclosure{\oxkontctx}{\oxvarctx^\prime_f}{\oxcprov}$,
    $\oxmusafety{\oxglobalctx}{\oxemptyctx}{\oxvarctx_f^\prime}{\oxmuta}{\oxplaceexpr}{\oxset{\oxloans^\prime}}$,
    and
    $\oxcomputetynoprov{\oxemptyctx}{\oxvarctx^\prime}{\oxmuta}{\oxplaceexpr}{\oxxitype}$.

    $\oxvarctx^\prime_f(\oxcprov) = \emptyset$ is immediate because $\oxvarctx^\prime_f(\oxcprov) \subseteq \oxvarctx^\prime(\oxcprov)$.

    Both $\oxnotinclosure{\oxkontctx}{\oxvarctx^\prime_f}{\oxcprov}$ and
    $\oxcomputetynoprov{\oxemptyctx}{\oxvarctx^\prime}{\oxmuta}{\oxplaceexpr}{\oxxitype}$
    follow from the fact that types are unchanged between $\oxvarctx^\prime$ and
    $\oxvarctx^\prime_f$.

    The last obligation is
    $\oxmusafety{\oxglobalctx}{\oxemptyctx}{\oxvarctx_f^\prime}{\oxmuta}{\oxplaceexpr}{\oxset{\oxloans^\prime}}$,
    which follows from \Lemma{lemma:ownership-safety-more-precise}
  }

  %% new type is a subtype of the old type
  \oxpresline{
    \oxtunify[\oxcombine]{\oxemptyctx}{\oxkontctx}{\oxvarctx_f^{\prime}[\oxcprov \mapsto \oxset{\oxloans^\prime}]}
    {\oxtnum}{\oxtnum}{\oxvarctx_s}
  }{
    Immediate, with $\oxvarctx_s = \oxvarctx_f^\prime[\oxcprov \mapsto \oxset{\oxloans^\prime}]$.
  }

  %% new loan context is a sub-context of the old loan context
  \oxpresline{
    \exists \oxvarctx_o. \oxvarctx_s \oxintersect \oxvarctx_o =
    \oxvarctx^\prime[\oxcprov \mapsto \oxset{\oxloans}]
  }{
    From our application of the induction hypothesis above, we have that
    $\oxvarctx^\prime = \oxvarctx_f^\prime \oxintersect \oxvarctx_o$. As we
    found in the typing judgement section above, $\oxset{\oxloans^\prime}
    \subseteq \oxset{\oxloans}$. Therefore, for this judgement, we can use
    $\oxvarctx_o[\oxcprov \mapsto \oxset{\oxloans} \setminus
      \oxset{\oxloans^\prime}]$.
  }
\end{preservation}



\noindent Case \oxname{T-Tuple}: \\[0.5em]
\begin{tabular}{l}
  \pbox{\textwidth}{\vspace*{0.5em} From premise: \vspace*{0.5em}} \\

  \framebox[\textwidth]{
    \figuresize
    \begin{mathpar}
      \TTuple
    \end{mathpar}
  } \\
  \pbox{\textwidth}{
    \vspace*{0.5em}
    Since $\oxexpr = \oxprod{\oxvalue_1 \oxdotsc \oxvalue_{i-1}, \oxexpr_i \oxdotsc \oxexpr_n}$, by inspection of the
    reduction rules, we know that $\oxexpr$ steps with the following rule:
    \vspace*{0.5em}
  } \\

  \framebox[\textwidth]{
    \figuresize
    \begin{mathpar}
      \inferrule[E-EvalCtx]{
        %% premises
        \oxreduce{\oxglobalctx}{\oxstore}{\oxexpr_i}
                 {\oxstore^\prime}{\oxexpr_i^\prime}
      }{
        \oxreduce{\oxglobalctx}{\oxstore}{\oxprod{\oxvalue_1 \oxdotsc \oxvalue_{i-1}, \oxexpr_i \oxdotsc \oxexpr_n}}
                 {\oxstore^\prime}{\oxprod{\oxvalue_1 \oxdotsc \oxvalue_{i-1}, \oxexpr_i^\prime, \oxexpr_{i+1} \oxdotsc \oxexpr_n}}
      }
    \end{mathpar}
  } \\[1.5em]
\end{tabular} \\[1em]

We want to apply our induction hypothesis to the premise of \oxname{E-EvalCtx}
with the typing judgement $\oxtypjudge{\oxglobalctx}{\oxemptyctx}{\oxkontctx,
  \oxsitype_1 \oxdotsc \oxsitype_{i-1}}{\oxvarctx_{i-1}}
{\oxexpr_i}{\oxsitype_i}{\oxvarctx_i}$. In order to do so we need to show that
$\forall \ \overline{\oxvalue}$.
$\oxwfkontctx{\oxglobalctx}{\oxvarctx_{i-1}}{\overline{\oxvalue}, \oxvalue_1
  \oxdotsc \oxvalue_{i-1}}{\oxkontctx, \oxsitype_1 \oxdotsc \oxsitype_{i-1}}$
given that
$\oxwfkontctx{\oxglobalctx}{\oxvarctx_0}{\overline{\oxvalue}}{\oxkontctx}$. So
let $\oxvalue \in \overline{\oxvalue}$. For each $\oxvalue_{j}$, where $0 < j <
i$, we can repeatedly apply \Lemma{lemma:values-dont-change} to get
$\oxsubtypectx{\oxglobalctx}{\oxemptyctx}{\oxvarctx_{j-1}}{\oxvarctx_{i-1}}$, and
then apply \Lemma{lemma:value-typing-related-envs} to get that $\forall \oxvalue_k \in \overline{\oxvalue}$. 
$\oxtypjudge{\oxglobalctx}{\oxemptyctx}{\oxkontctx, \oxsitype_1 \oxdotsc
  \oxsitype_{k-1}}{\oxvarctx_{i-1}}{\oxvalue_k}{\oxkontctx[k]}{\oxvarctx_{i-1}}$.


Now we can apply the induction hypothesis to get there is some
$\oxvarctx_{i-1}^\prime$ such that $\forall\, \overline{\oxvalue}$.
$\oxwfkontctx{\oxglobalctx}{\oxvarctx_{i-1}^\prime}{\overline{\oxvalue}}
{\oxkontctx, \oxsitype_1 \oxdotsc \oxsitype_{i-1}}$ and
$\oxtypjudge{\oxglobalctx}{\oxemptyctx}{\oxkontctx, \oxsitype_1 \oxdotsc \oxsitype_{i-1}}{\oxvarctx_{i-1}^{\prime}}
{\oxexpr_i^{\prime}}{\oxtype_i^{\textsc{si} \prime}}{\oxvarctx_i^{\prime}}$,
with $\oxstorevalidity{\oxglobalctx}{\oxvarctx_{i-1}^\prime}{\oxstore^\prime}$
and $\oxtunify[\oxcombine]{\oxemptyctx}{\oxkontctx, \oxsitype_1 \oxdotsc \oxsitype_{i-1}}{\oxvarctx_i^\prime}
{\oxtype_i^{\textsc{si}\prime}}{\oxsitype_i}{\oxvarctx_{si}}$ and there is some
$\oxvarctx_{oi}$ such that $\oxvarctx_i = \oxvarctx_{si} \oxintersect \oxvarctx_{oi}$.

We then pick $\oxvarctx_i$ in the lemma statement to be
\framebox{$\oxvarctx_{i-1}^{\prime}$} and need to show:


\begin{preservation}
  %% new store is valid
  \oxpresline{
    \oxstorevalidity{\oxglobalctx}{\oxvarctx_{i-1}^{\prime}}{\oxstore^\prime}
  }{
    Immediate.
  }
  \oxpresline{
    \oxwfkontctx{\oxglobalctx}{\oxvarctx_{i-1}^\prime}{\overline{\oxvalue}}{\oxkontctx}
  }{
    Immediate
  }

  %% new expression is well-typed
  \oxpresline{
    \begin{array}{r}
    \oxglobalctx; \, \oxemptyctx; \, \oxkontctx; \, \oxvarctx_{i-1}^{ \prime} \vdash
    \framebox{$\oxprod{\oxvalue_1 \oxdotsc \oxvalue_{i-1}, \oxexpr_i^\prime,
        \oxexpr_{i+1} \oxdotsc \oxexpr_n}$}\\
    : \oxprod{\oxsitype_1 \oxdotsc
        \oxsitype_{i-1}, \oxtype_i^{\textsc{si} \prime}, \oxsitype_{i+1}
        \oxdotsc \oxsitype_{n}} \Rightarrow \oxvarctx_n^{\prime}
    \end{array}
  }{ We want to apply \oxname{T-Tuple} in order to typecheck this tuple
    expression. We firstly want to show that we can typecheck $\oxvalue_1$
    through $\oxvalue_{i-1}$ with initial input context
    $\oxvarctx_{i-1}^{\prime}$. For each $\oxvalue_{j}$, where $0 < j < i$, we
    can repeatedly apply \Lemma{lemma:values-dont-change} to get
    $\oxsubtypectx{\oxglobalctx}{\oxemptyctx}{\oxvarctx_{j-1}}{\oxvarctx_{i-1}}$,
    and then apply \Lemma{lemma:value-typing-related-envs} to get that
    $\oxtypjudge{\oxglobalctx}{\oxemptyctx}{\oxkontctx, \oxsitype_1 \oxdotsc \oxsitype_{j-1}}{\oxvarctx_{i-1}}{\oxvalue_j}{\oxtype_j}{\oxvarctx_{i-1}}$,
    and finally applying our continuation typing hypothesis gives us
    $\oxtypjudge{\oxglobalctx}{\oxemptyctx}{\oxkontctx, \oxsitype_1 \oxdotsc \oxsitype_{j-1}}{\oxvarctx_{i-1}^\prime}{\oxvalue_j}{\oxtype_j}{\oxvarctx_{i-1}^\prime}$.

    Then we can use the result of our induction hypothesis,
    $\oxtypjudge{\oxglobalctx}{\oxemptyctx}{\oxkontctx, \oxsitype_1 \oxdotsc
      \oxsitype_{i-1}}{\oxvarctx_{i-1}^{\prime}}{\oxexpr_i^{\prime}}{\oxtype_i^{\textsc{si}
        \prime}}{\oxvarctx_i^{\prime}}$.

    And finally for $\oxexpr_{i+1} \oxdotsc \oxexpr_n$, we can repeatedly apply
    \Lemma{lemma:expressions-typing-more-precise}. }

  %% new type is a subtype of the old type
  \oxpresline{
    \oxtunify[\oxcombine]{\oxemptyctx}{\oxkontctx}{\oxvarctx_n^{\prime}}
             {\oxprod{\oxsitype_1 \oxdotsc \oxsitype_{i-1},
                 \oxtype_i^{\textsc{si} \prime}, \oxsitype_{i+1} \oxdotsc
                 \oxsitype_{n}}}{\oxprod{\oxsitype_1 \oxdotsc
                 \oxsitype_n}}{\oxvarctx_s}
  }{
    Apply \oxname{RR-Tuple} to
    \oxname{RR-Refl} for every pair of types that are identical, leaving just
    $\oxtype_i^{\textsc{si}\prime}$ and $\oxsitype_i$ for which we can use
    $\oxtunify[\oxcombine]{\oxemptyctx}
    {\oxkontctx, \oxsitype_1 \oxdotsc \oxsitype_{i-1}}{\oxvarctx_i^\prime}
    {\oxtype_i^{\textsc{si}\prime}}{\oxsitype_i}{\oxvarctx_{si}}$ from
    our induction hypothesis with
    \Lemma{lemma:rewriting-parallel-expression} and
    \Lemma{lemma:rewriting-smaller-kont} to get that
    $\oxtunify[\oxcombine]{\oxemptyctx}{\oxkontctx}{\oxvarctx_n^\prime}
    {\oxtype_i^{\textsc{si}\prime}}{\oxsitype_i}{\oxvarctx_s}$.
  }

  %% new loan context is a sub-context of the old loan context
  \oxpresline{
    \exists \oxvarctx_o. \oxvarctx_s \oxintersect \oxvarctx_o = \oxvarctx_n
  }{
    Note from the induction hypothesis above, we have $\oxvarctx_{i} =
    \oxvarctx_{si} \oxintersect \oxvarctx_{oi}$. This means that for whatever
    loan sets were unioned between $\oxvarctx_i^\prime$ and $\oxvarctx_{si}$,
    those loan sets will be unioned between $\oxvarctx^\prime_n$ and
    $\oxvarctx_s$. And since $\oxvarctx_i$ contained $\oxvarctx_{si}$, it will
    also be true that $\oxvarctx_n$ will contain $\oxvarctx_s$, which means we
    can use $\oxvarctx_o = \oxvarctx_n \setminus \oxvarctx_s$.
   }
\end{preservation}

\oxpreservationcaseheader{T-Drop}{\TDrop}

Since \oxname{T-Drop} applies to \emph{any} expression $\oxexpr$, we cannot
determine anything about what rule we stepped with. So, we will instead try to
apply our induction hypothesis to the typing derivation in the premise of
\oxname{T-Drop}
\oxtypjudge{(\oxglobalctx}{\oxemptyctx}{\oxkontctx}{
  \oxtupdate{\oxvarctx}{\oxplace}{\oxsidtype_\oxplace}
}{\oxexpr}{\oxsxtype}{\oxvarctx_f)}. To do this, we need
to establish the premises of our inductive hypothesis.

Namely, we need to show:
\begin{enumerate}
\item \oxtypjudge{\oxglobalctx}{\oxemptyctx}{\oxkontctx}{
    \oxtupdate{\oxvarctx}{\oxplace}{\oxsidtype_\oxplace}
  }{\oxexpr}{\oxsxtype}{\oxvarctx_f},
\item $\oxstorevalidity{\oxglobalctx}{
    \oxtupdate{\oxvarctx}{\oxplace}{\oxsidtype_\oxplace}
  }{\oxstore}$,
\item $\oxreduce{\oxglobalctx}{\oxstore}{\oxexpr}
  {\oxstore^\prime}{\oxexpr^\prime}$.
\end{enumerate}

\vspace{0.5em}

\begin{enumerate}
\item follows immediately from our premise.
\item follows almost directly from our premise, which tells us that
  $\oxstorevalidity{\oxglobalctx}{\oxvarctx}{\oxstore}$.
  We just need to show that the value at $\oxid$ (where $\oxid$ is the root of
  $\oxplace$) is still valid at its new type. Fortunately, it's old derivation
  works almost perfectly except for typing the part that corresponds directly to
  $\oxplace$. In this case, we can use \oxname{T-Dead} on the value to get the
  new derivation with that part of the aggregate structure at the uninitialized
  type $\oxsidtype_\oxplace$.
\item follows immediately from our premise.
\end{enumerate}

\vspace{0.5em}

This allows us to use our induction hypothesis to conclude that there exists
$\oxvarctx^\prime_i$ such that:
\begin{enumerate}
  \setcounter{enumi}{4}
\item $\oxstorevalidity{\oxglobalctx}{\oxvarctx^\prime_i}{\oxstore^\prime}$,
\item $\oxwfkontctx{\oxglobalctx}{\oxvarctx^\prime_i}{\overline{\oxvalue}}{\oxkontctx}$,
\item $\oxtypjudge{\oxglobalctx}{\oxemptyctx}{\oxkontctx}{\oxvarctx^\prime_i}
  {\oxexpr^\prime}{\oxtype^\prime}{\oxvarctx_f^\prime}$,
\item $\oxtunify[\oxcombine]{\oxemptyctx}{\oxkontctx}{\oxvarctx_f^\prime}
  {\oxtype^\prime}{\oxsxtype}{\oxvarctx_s}$, and
\item $\exists \oxvarctx_o. \; \oxvarctx_s \oxintersect \oxvarctx_o = \oxvarctx_f$.
\end{enumerate}
\hfill$\square$

%%% Local Variables:
%%% mode: latex
%%% fill-column: 80
%%% End:


\subsection{Type Safety}
\label{sec:type-safety}

\begin{oxtheorem}{Type Safety}{theorem:type-safety-app}
  If $\oxtypjudge{\oxglobalctx}{\oxemptyctx}{\oxemptyctx}{\oxemptyctx}{\oxexpr}
  {\oxsitype}{\oxvarctx}$ and $\oxglobalctxvalidity{\oxglobalctx}$, then
  $\oxreducemany{\oxglobalctx}{\oxemptyctx}{\oxexpr}{\oxstore^\prime}{\oxvalue}$ or
  evaluation of $\oxexpr$ aborts or diverges.
\end{oxtheorem}

\begin{proof}
  By the interleaved use of \textbf{Progress} and \textbf{Preservation}.
\end{proof}

%%% Local Variables:
%%% mode: latex
%%% fill-column: 80
%%% End:
